\documentclass[10pt,a4paper]{article}

% ----------------- Paquetes -------------------%
\usepackage[T1]{fontenc}
\usepackage[utf8]{inputenc}
\usepackage[a4paper,margin=2.5cm]{geometry}
\usepackage{mathptmx}             
\usepackage{changepage} 
\usepackage{setspace}
\usepackage[spanish]{babel}
\usepackage[labelfont=bf,labelsep=space]{caption}
\usepackage{amssymb}
\usepackage{amsmath}
\usepackage{ebgaramond}
\usepackage{placeins}
\usepackage{etoolbox}
\usepackage{setspace}
\usepackage{parskip} 
\usepackage{tcolorbox}
\usepackage{needspace}
\usepackage{xparse}
\usepackage{ocgx2}
\usepackage{hyperref}
\usepackage{hypcap}


%%%%%%%%%%%%%%%%%%%%%%%%%%%%%%%%%%%%%%%%%%%%%%%%%%%%%%%%%%%%%%%%
% --- Fija primero tus valores globales buenos ---
\setstretch{1.2}                 % Interlineado global
\setlength{\parskip}{0.7\baselineskip}
\setlength{\parindent}{0pt}
\newlength{\GlobalParskip}
\setlength{\GlobalParskip}{\parskip}

\newcounter{eqnum}[subsection]
\renewcommand{\theeqnum}{\arabic{eqnum}}
\newcounter{alphaitem}
\newcounter{numitem}

%% ------------------- Recuadros----------------------%%
\newcommand{\puntosclave}[1]{%
  \begin{tcolorbox}[
    colframe=black,
    title={Puntos clave},
    before upper={\setlength{\parindent}{0.5cm}\setlength{\parskip}{0.5\baselineskip}}
  ]
  #1
  \end{tcolorbox}
}

\newcommand{\modificaciones}[1]{
  \begin{tcolorbox}[
    colback=white,
    colframe=red,
    title={Modificaciones a realizar},
    before upper={\setlength{\parindent}{0.5cm}\setlength{\parskip}{0.5\baselineskip}}
  ]
  #1
  \end{tcolorbox}
}

%% ------------------- Listas----------------------%%
    % --- Lista alfabética ---
    \newenvironment{la}
      {\begin{list}{\alph{alphaitem})}{%
          \usecounter{alphaitem}%
          \setlength{\leftmargin}{1cm}%
          \setlength{\parskip}{3pt}% 
          \setlength{\itemsep}{0pt}% ← espacio entre ítems
          \setlength{\parsep}{0pt}%  ← párrafos dentro de un ítem
      }}
      {\end{list}}
      
    % --- Lista numérica ---
    \newenvironment{lnum}
      {\begin{list}{\arabic{numitem}.}{%
          \usecounter{numitem}%
          \setlength{\leftmargin}{1cm}%
          \setlength{\parskip}{3pt}% 
          \setlength{\itemsep}{0pt}% ← espacio entre ítems
          \setlength{\parsep}{0pt}%  ← párrafos dentro de un ítem
      }}
      {\end{list}}
      
% ------------------- Imágenes----------------------%
    \graphicspath{{Img/}}% Rutas por defecto 
    % --- Imagen adaptativa ---
    % Registros de dimensión definidos una sola vez
    \newdimen\imgcap
    \newdimen\imgmin
    \newdimen\imgmax
    \newdimen\imgremain
    \newdimen\imgh
    \newdimen\imgneed
    
    \newcommand{\imagenfit}[3]{%
      \addvspace{10pt}% separación antes
      \par\begingroup
        % Parámetros de composición (asignaciones, no \newdimen)
        \imgcap=1\baselineskip            % alto estimado de título+fuente
        \imgmin=0.15\textheight
        \imgmax=0.15\textheight
        \imgremain=\pagegoal
        \ifdim\pagegoal=\maxdimen \imgremain=0pt\fi
        \advance\imgremain by -\pagetotal
        \advance\imgremain by -\imgcap
        % Decisión de altura destino
        \ifdim\imgremain<\imgmin
          \newpage
          \imgh=0.20\textheight
        \else
          \imgh=\imgremain
          \ifdim\imgh>\imgmax \imgh=\imgmax \fi
        \fi
        % Reservar espacio para evitar cortes del bloque completo
        \imgneed=\imgh \advance\imgneed by \imgcap
        \Needspace{\imgneed}
        % Cuerpo
        \noindent\begin{minipage}{\textwidth}
          \centering
          \captionof{figure}{\textemdash\ #2}\par\vspace{0.3em}% título arriba
          \includegraphics[width=\linewidth,height=\imgh,keepaspectratio]{\detokenize{#1}}\par
          \vspace{0.3em}\small\textit{Fuente: }#3% fuente abajo
        \end{minipage}%
      \endgroup
      \par\addvspace{10pt}%
    }

%------------ Bloque de RECUADRO ESPECIAL -------------%
    \newenvironment{specialblock}[1]{%
      \begin{quote} % crea sangría a izquierda y derecha
      \small % texto más pequeño
      \noindent\textbf{#1}\\[-0.3em] % título en negrita
      \rule{\linewidth}{0.4pt}\\[0.6em]}{
      \end{quote}}
      

%-------------------------- Author Font -----------------------%
    \newcommand{\authorfont}[1]{{\fontfamily{EBGaramond-TLF}\selectfont #1}}

%-------------------------- Anexos -----------------------%
    \makeatletter
    \newcounter{anexo}
    \renewcommand{\theanexo}{\Roman{anexo}}
    \newcommand{\anexo}[1]{%
      \clearpage
      \phantomsection
      \refstepcounter{anexo}%
      \noindent\textbf{Anexo \theanexo.- }#1\par\medskip
      \addcontentsline{stc}{section}{Anexo \theanexo.- #1}%
    }
    \makeatother

    %-------------------------- Lista de figuras (personal) -----------------------%
\makeatletter
\newcommand{\MyListOfFigures}{%
  \clearpage
  \phantomsection
  {\centering\bfseries\Large Índice de figuras\par}%
  \medskip
  % Entrada en nuestro índice de contenido (stc)
  \addcontentsline{stc}{section}{Índice de figuras}%
  % Recuperación del listado (sin cabecera propia)
  \begingroup
    \parindent=0pt\parskip=2pt
    \@starttoc{lof}%
  \endgroup
}
\makeatother

%-------------------------- Bibliografía -----------------------%
\makeatletter
% Contador para las referencias
\newcounter{myref}
% Cabecera de la bibliografía, entra en el índice 'stc'
\newcommand{\MyBibliography}{%
  \clearpage
  \phantomsection
  {\centering\bfseries\Large Bibliografía y referencias\par}%
  \medskip
  \addcontentsline{stc}{section}{Bibliografía y referencias}%
  \setcounter{myref}{0}%
}
\newcommand{\MyBibItem}[4]{%
  \refstepcounter{myref}%
  \noindent[\themyref]\ #1.\ \emph{#2}. #3. #4\par
}
\makeatother

%--------- Establecer cuatro niveles de índice --------%
    \setcounter{tocdepth}{4}   
    \setcounter{secnumdepth}{4}% numera hasta \paragraph

%%%%%%%%%%%%%%%%%%%%%%%%%%%%%%%%

\makeatletter

    %--------- Interlineado Global --------%
    \edef\GlobalBaseLineStretch{\baselinestretch}
    
    %--------- Espaciado de seccinones --------%
    \renewcommand\section{\@startsection{section}
      {1}{\z@}
      {2.5ex \@plus 1ex \@minus .2ex} % espacio antes
      {1.5ex \@plus .2ex}             % espacio después
      {\normalfont\Large\bfseries}    % formato del título
    }
    \renewcommand\subsection{\@startsection{subsection}{2}{\z@}
      {3ex \@plus 0.8ex \@minus .2ex}
      {1.5ex \@plus .2ex}
      {\normalfont\large\bfseries}
    }
    \renewcommand\subsubsection{\@startsection{subsubsection}{3}{\z@}
      {3ex \@plus 0.5ex \@minus .2ex}
      {1ex \@plus .2ex}
      {\normalfont\normalsize\bfseries}
    }
    \renewcommand\paragraph{\@startsection{paragraph}{4}{\z@}%
    {-2ex \@plus -0.5ex \@minus -.2ex}% espacio antes
    {0.8ex \@plus .2ex}% espacio después
    {\normalfont\normalsize\itshape}}% ← cursiva en lugar de negrita

    %--------- Ecuaciones --------%   
    % Variables globales para almacenar títulos y notas
    \gdef\leq@current@section{}%
    \gdef\leq@current@subsection{}%
    \gdef\leq@last@printed@section{}%
    \gdef\leq@last@printed@subsection{}%
    
    % Comando para añadir nota (invisible en el cuerpo)
    \newcommand{\eqnote}[1]{%
      \addcontentsline{leq}{leqnote}{#1}%
    }
    
    % Comando para ecuaciones
    \newcommand{\eqblock}[2]{%
      % Verificar si necesitamos imprimir el título de sección
      \ifx\leq@current@section\leq@last@printed@section\else
        \addcontentsline{leq}{leqsection}{\leq@current@section}%
        \global\let\leq@last@printed@section\leq@current@section
        \gdef\leq@last@printed@subsection{}% Reset subsección al cambiar sección
      \fi
      % Verificar si necesitamos imprimir el título de subsección
      \ifx\leq@current@subsection\leq@last@printed@subsection\else
        \ifx\leq@current@subsection\@empty\else
          \addcontentsline{leq}{leqsubsection}{\leq@current@subsection}%
          \global\let\leq@last@printed@subsection\leq@current@subsection
        \fi
      \fi
      % Ecuación
      \refstepcounter{eqnum}%
      \begin{equation}
        #1\tag{E.\theeqnum}
      \end{equation}%
      \addcontentsline{leq}{leqt}{[E.\theeqnum] -- #2}%
      \addcontentsline{leq}{leqm}{#1}%
    }
    
    %%% ----- Índice de ecuaciones ----------%%%
    % Tamaño para []
    \newcommand{\leqIndexFont}{\normalsize}
    \newcommand{\leqTitleFont}{\tiny}
    \newcommand{\leqNoteFont}{\tiny}
    % Cabecera de sección 
    \newcommand\l@leqsection[2]{%
      \addvspace{25pt}% Mayor separación antes de nueva sección
      \noindent\leqIndexFont
      \makebox[\linewidth]{\rule{.38\linewidth}{0.4pt}\ \ \textbf{  \MakeUppercase{#1}}\ \ \rule{.38\linewidth}{0.4pt}}%
      \par\vspace{-10pt}
    }
    % Cabecera de subsección 
    \newcommand\l@leqsubsection[2]{%
      \par\addvspace{15pt}%
      \noindent\leqIndexFont #1
      \par\vspace{-7pt}
    }
    % Título de ecuación 
    \newcommand\l@leqt[2]{%
    \par\addvspace{10pt}%
      \par\noindent\leqTitleFont #1\par
    }
    % Miniatura de ecuación
    \newcommand\l@leqm[2]{%
      \vspace{0.3\baselineskip}%
      \par\scriptsize\noindent\hspace{1.5em}\(#1\)\par%
    }
    % Nota de ecuación 
    \newcommand\l@leqnote[2]{%
      \vspace{0.2\baselineskip}%
      \par\noindent\leqNoteFont\hspace{1.5em}\textbf{Nota.--} #1\par\vspace{0.5\baselineskip}%
    }
    % Generación del índice
    \newcommand{\mylistofequations}{%
      \clearpage
      \phantomsection
      % Título visual de la lista (ajústalo a tu gusto):
      {\centering\bfseries\Large Índice de ecuaciones\par}
      % Entrada en nuestro índice 'stc':
      \addcontentsline{stc}{section}{Índice de ecuaciones}%
      % Llamada a tu lista real (usa el nombre exacto de tu macro):
      \@starttoc{leq}%
    }
    
    %--------- Captura de títulos de sección --------%
    \let\leq@original@section\section
    \let\leq@original@subsection\subsection
    % Flag para saber si estamos en una sección asterisco
    \newif\ifLeqStarSection
    \LeqStarSectionfalse
    
    % Redefinir \section CON CONTROL DE ASTERISCO
    \renewcommand{\section}{%
      \@ifstar{\leq@section@star}{\@dblarg\leq@section@nostar}%
    }
    \def\leq@section@star#1{%
      \global\LeqStarSectiontrue
      \gdef\leq@current@section{#1}%
      \gdef\leq@current@subsection{}%
      \leq@original@section{#1}%
      \global\LeqStarSectionfalse
    }
    \def\leq@section@nostar[#1]#2{%
      \global\LeqStarSectionfalse
      \gdef\leq@current@section{#2}%
      \gdef\leq@current@subsection{}%
      \leq@original@section[#1]{#2}%
    }
    % Redefinir \subsection
    \renewcommand{\subsection}{%
      \@ifstar{\leq@subsection@star}{\@dblarg\leq@subsection@nostar}%
    }
    \def\leq@subsection@star#1{%
      \global\LeqStarSectiontrue
      \gdef\leq@current@subsection{#1}%
      \leq@original@subsection{#1}%
      \global\LeqStarSectionfalse
    }
    \def\leq@subsection@nostar[#1]#2{%
      \global\LeqStarSectionfalse
      \gdef\leq@current@subsection{#2}%
      \leq@original@subsection[#1]{#2}%
    }

    % ----- Restauración de valores Globales de estilo ----------%
    % Flags
    \newif\ifInFootnote
    \InFootnotefalse
    \pretocmd{\@footnotetext}{\InFootnotetrue}{}{}
    \apptocmd{\@footnotetext}{\InFootnotefalse}{}{}
    % Profundidad de listas (soporta anidadas)
    \newcount\MyListDepth \MyListDepth=0
    \AtBeginEnvironment{la}{\advance\MyListDepth by 1}
    \AfterEndEnvironment{la}{\advance\MyListDepth by -1}
    \AtBeginEnvironment{lnum}{\advance\MyListDepth by 1}
    \AfterEndEnvironment{lnum}{\advance\MyListDepth by -1}
    \AtBeginEnvironment{itemize}{\advance\MyListDepth by 1}
    \AfterEndEnvironment{itemize}{\advance\MyListDepth by -1}
    \AtBeginEnvironment{enumerate}{\advance\MyListDepth by 1}
    \AfterEndEnvironment{enumerate}{\advance\MyListDepth by -1}
    % Restauración diferida: hazlo al INICIO del próximo párrafo real
    \newif\if@pendingRestore
    \@pendingRestorefalse
    \newcommand{\RequestRestore}{\global\@pendingRestoretrue}
    \everypar{%
      \if@pendingRestore
        \global\@pendingRestorefalse
        \ifInFootnote\else
          \ifnum\MyListDepth=0
            \setlength{\parskip}{\GlobalParskip}%
            \setstretch{\GlobalBaseLineStretch}%
          \fi
        \fi
      \fi
    }
    % Al final de bloques:
    \AfterEndEnvironment{equation}{\RequestRestore}
    \AfterEndEnvironment{equation}{\RequestRestore}
    \AfterEndEnvironment{la}{\RequestRestore}
    \AfterEndEnvironment{lnum}{\RequestRestore}
    % Al final de macros:
    \apptocmd{\eqblock}{\RequestRestore}{}{}
    \apptocmd{\imagenfit}{\RequestRestore}{}{}
    
\makeatother

% ===================== ToC propio con 4 niveles (único bloque) =====================
\makeatletter

% --- Escritores: añadimos una línea al .stc en cada cabecera numerada ---
% Guardamos las versiones actuales para llamar al comportamiento normal
\let\MyTOC@orig@section\section
\let\MyTOC@orig@subsection\subsection
\let\MyTOC@orig@subsubsection\subsubsection
\let\MyTOC@orig@paragraph\paragraph

% SECTION
\renewcommand{\section}{%
  \@ifstar{\MyTOC@section@star}{\@dblarg\MyTOC@section@nostar}%
}
\def\MyTOC@section@star#1{%
  \MyTOC@orig@section{#1}% sin entrada al .stc
}
\def\MyTOC@section@nostar[#1]#2{%
  \MyTOC@orig@section[{#1}]{#2}% comportamiento normal (escribe al .toc)
  % Escribimos también nuestra línea al .stc (texto con \numberline)
  \addcontentsline{stc}{section}{\protect\numberline{\thesection}#2}%
}

% SUBSECTION
\renewcommand{\subsection}{%
  \@ifstar{\MyTOC@subsection@star}{\@dblarg\MyTOC@subsection@nostar}%
}
\def\MyTOC@subsection@star#1{%
  \MyTOC@orig@subsection{#1}%
}
\def\MyTOC@subsection@nostar[#1]#2{%
  \MyTOC@orig@subsection[{#1}]{#2}%
  \addcontentsline{stc}{subsection}{\protect\numberline{\thesubsection}#2}%
}

% SUBSUBSECTION
\renewcommand{\subsubsection}{%
  \@ifstar{\MyTOC@subsubsection@star}{\@dblarg\MyTOC@subsubsection@nostar}%
}
\def\MyTOC@subsubsection@star#1{%
  \MyTOC@orig@subsubsection{#1}%
}
\def\MyTOC@subsubsection@nostar[#1]#2{%
  \MyTOC@orig@subsubsection[{#1}]{#2}%
  \addcontentsline{stc}{subsubsection}{\protect\numberline{\thesubsubsection}#2}%
}

% PARAGRAPH
\renewcommand{\paragraph}{%
  \@ifstar{\MyTOC@paragraph@star}{\@dblarg\MyTOC@paragraph@nostar}%
}
\def\MyTOC@paragraph@star#1{%
  \MyTOC@orig@paragraph{#1}%
}
\def\MyTOC@paragraph@nostar[#1]#2{%
  \MyTOC@orig@paragraph[{#1}]{#2}%
  % Si no numeras los \paragraph, cambia \theparagraph por texto plano sin \numberline
  \addcontentsline{stc}{paragraph}{\protect\numberline{\theparagraph}#2}%
}

% --- Impresor del índice propio (lee el .stc) ---
\newcommand{\MyTOC}{%
  \par\bigskip
  {\centering\bfseries\Large Índice de contenido\par}%
  \medskip
  \begingroup
    \parindent=0pt\parskip=2pt
    \@starttoc{stc}%
  \endgroup
}

\makeatother
% ===================== fin ToC propio =====================


%%%%%%%%%%%%%%


% --- Datos del tema ---
\title{Política Presupuestaria (VI): Saldo presupuestaria y deuda de las Administraciones Públicas en España\\
\large Análisis de su evolución reciente y comparación internacional. Política de financiación del Tesoro en la actualidad}
\author{Marcelo Ortega}
\date{Septiembre 2026}

\begin{document}
\maketitle
\newpage

% --- Página de resumen y modificaciones ---
\puntosclave{ %% Escribir puntos clave Apuntes CECO
     
    }
\vspace{1cm}

\modificaciones{
    
    }

\newpage
\MyTOC
\newpage

\mylistofequations
\clearpage

\MyListOfFigures
\clearpage


\pagenumbering{arabic}
\setcounter{page}{1}


\section{Introducción}



\subsection{Contextualización}


\subsubsection{Relevancia}




\subsubsection{Problemática}



\subsection{Estructura de la exposición}
Este tema es de gran importancia en el contexto actual. caracterizado en España por los elevados niveles de déficit y deuda pública. que hacen de la actuación del Tesoro una cuestión muy relevante. Es importante mantener el tema actualizado, para lo que se puede recurrir a dos documentos disponibles en la web del Tesoro: su estrategia anual de financiación y el chart pack, un compendio de gráficos que se actualiza mensualmente y que también puede ser útil para mantener actualizados los datos de todos los temas de economía española del cuarto ejercicio.
Por lo que se refiere a la estructura del tema, lo dividiremos en dos grandes secciones. Así, tras la introducción (a la que se pueden dedicar en torno a 2 minutos). el primer apartado analiza la evolución reciente del saldo presupuestario y de la deuda de las Administraciones Públicas (AAPP) en España (~16 minutos):
• Comenzaremos revisando qué se entiende por saldo presupuestario y por deuda, así como delimitando el sector institucional de las AAPP en España (~2 minutos).
• A continuación, repasaremos la evolución del saldo presupuestario y de la deuda de las AAPP en España siguiendo un enfoque temporal:
o Aunque el título del tema pide analizar la evolución reciente del saldo presupuestario y de la deuda, comenzaremos nuestro análisis en 1975, pues fue durante la Transición cuando el saldo presupuestario español empieza a mostrar déficits recurrentes. Así, en el primer subapartado analizaremos un periodo de más de tres décadas, siendo necesario un ejercicio de síntesis importante por parte del opositor (~4 minutos).
o Continuaremos con el estudio del periodo que va desde la crisis financiera internacional hasta la crisis de la covid-19 (~4 minutos). 
o Posteriormente, repasaremos la situación actual de las cuentas públicas españolas (~3 minutos), prestando especial atención al impacto de la pandemia sobre las mismas.
• Terminaremos con una comparación internacional de la evolución reciente del saldo presupuestario y de la deuda, prestando especial atención a las diferencias existentes dentro del grupo de las economías avanzadas al que pertenece España (~3 minutos).

El segundo apartado analiza la política de financiación del Tesoro en la actualidad ( ~ 1 O minutos), una cuestión clave dadas las necesidades de financiación del Estado, que se sitúan alrededor de los 250.000 millones de euros anuales, equivalentes a algo más del 20% del PIB. Comenzaremos repasando la evolución reciente de las necesidades de financiación del Tesoro (~3 minutos), seguiremos con el estudio de sus instrumentos (~3 minutos) y procedimientos de emisión (~3 minutos) y terminaremos revisando la última estrategia de financiación del Tesoro (~J minuto). El tema finaliza con una conclusión a la que, dada la naturaleza del tema, el opositor puede aprovechar para añadir alguna cuestión de actualidad (~2 minutos).

Finalmente, se hace notar que es de esperar que el opositor sea capaz de relacionar el contenido de este tema con el de otros del temario; en especial con el 3A39 y con el 3Al4 (en lo referente a subastas). También se recomienda mencionar en algún momento de la exposición que en este tema no se analiza en términos teóricos la sostenibilidad de la deuda pública, por ser objeto de estudio de otro tema

---------> Introducción

La evolución tanto del saldo presupuestario como de la forma de financiación del déficit público ha venido marcada a nivel internacional por el desarrollo de la teoría económica.

Así. hasta los años 30 del siglo XX las tesis dominantes defendían el mantenimiento del equilibrio presupuestario. en línea con los postulados de la teoría clásica de la Hacienda Pública. Autores como David Ricardo, de hecho, señalaban que la deuda es una de las más terribles plagas que nunca se han inventado para afligir a una nación

Además, el laissez.faire clásico implicaba limitar el papel del sector público a la provisión de los bienes públicos tradicionales Uusticia, defensa, infraestructuras ... ) y de las condiciones necesarias para el funcionamiento del libre mercado. lo que mantenía el gasto público en niveles reducidos. 

La Gran Depresión de los años 30 y el nuevo papel que adquirió la demanda en la teoría económica supusieron un cambio radical de paradigma. Siguiendo las tesis keynesianas, se empieza a defender el recurso a las políticas de demanda, especialmente a través de la política fiscal, para reactivar una economía en crisis. Un claro ejemplo de lo anterior fue el New D·eal de Roosevelt, que se pone en marcha en 1933.

También en este periodo, la teoría económica comienza a justificar la intervención del Estado en la economía por motivos de estabilidad, dada la tendencia a la inestabilidad de las economías de mercado defendida por el modelo de Harrod-Domar; o para evitar la posibilidad de llegar a un estancamiento secular como el que plantea Alvin Hansen ( 1938). P��r otro lado, el arraigo del concepto de equidad lleva a la creación del estado de bienestar, que supuso un aumento sin precedentes del rol del sector público en la economía. Como resultado de todo lo anterior, se rompe con el principio de equilibrio presupuestario y comienzan a pennitirse los saldos presupuestarios negativos discrecionales.

La reconstrucción tras la II Guerra Mundial y la bonanza económica que trajo consigo el sistema de Bretton Woods permitió que, pese al aumento del gasto público derivado de la creación del estado de bienestar, la mayoría de los países desarrollados consiguieran mantener unas cuentas públicas relativamente equilibradas durante los años 50 y 60.

Con las crisis del petróleo de los 70, no obstante, llega una coyuntura económica caracterizada por los bajos niveles de crecimiento y el elevado desempleo, que hace imposible recortar el gasto público y convierte el déficit público en un rasgo estructural de las economías avanzadas. Queda de manifiesto, además, la incapacidad de las políticas de demanda para domar el ciclo, como defendía la síntesis neoclásica, y se desarrolla la conocida como economía de la oferta.

La mayor inflación que trajeron consigo las crisis del petróleo llevó a articular políticas monetarias antiinflacionistas. Un claro ejemplo de ello fueron las políticas llevadas a cabo por la Reserva Federal de EE. UU. desde el nombramiento de Paul Volcker como presidente en 1979.

Lo anterior provocó que el déficit público dejase de ser monetizado y pasase a financiarse fundamentalmente vía deuda. El aumento de la deuda en los 80 y los 90 y la crisis de deuda de los 80 en muchos países en desarrollo, hacen evidente la posibilidad de insostenibilidad, lo que provoca la recuperación de las ideas de equilibrio presupuestario y la introducción de reglas fiscales para garantizar la credibilidad macroeconómica. En Europa, la firma del Tratado de Maastricht a principios de los 90 impone unos criterios de convergencia para poder acceder a la moneda única, lo que obliga a muchos países a acometer procesos de consolidación fiscal. 

España, con ciertas particularidades, no ha sido ajena a esta evolución, si bien este proceso ha sido más intenso y concentrado en el tiempo. Así, en España el déficit público es un fenómeno más reciente, pues fue durante la Transición cuando el saldo presupuestario español empieza a mostrar déficits recurrentes.

Hasta principios de los 80, dichos déficits se financiaban en buena medida mediante el recurso a la monetización, pero desde 1982 el Tesoro Público empieza a desarrollar una política de deuda pública que permite financiar el déficit público de una forma que hoy llamaríamos ortodoxa. El aumento de las emisiones provocó que la deuda pública llegara a suponer el 65% del PIB en 1996, una cifra superior a la prevista en el Tratado de Maastricht para entrar a formar parte de la Unión Económica y Monetaria (UEM). Por ello, ese mismo año comienza un proceso de consolidación presupuestaria que se extendió durante una década y que permitió, incluso, conseguir los primeros y únicos superávits de la democracia entre 2005 y 2007, año en el que la deuda pública se situó por debajo del 36% del PIB.

Sin embargo, el estallido de la crisis financiera internacional y la crisis de deuda soberana en la zona euro truncaron ese proceso, llegando a registrarse déficits anuales superiores al 10\% del PIB, que provocaron que la deuda pública alcanzase en 2014 el 105\% del PIB.\footnote{Como se verá, el aumento de la deuda pública española no se explica solo por la acumulación de déficits, sino que contribuyeron también la necesidad de aportar capital para la constitución del MEDE, o la creación del Fondo de Amortización del Déficit Eléctrico (FADE, ver anexo 1) y de la Sociedad de Gestión de Activos Procedentes de la Reestructuración Bancaria (Sareb, ver anexo 2), entre otras cosas.
}

En ese momento comienza un lento proceso de consolidación presupuestaria que se vio frenado en seco en 2020, cuando el estallido de la pandemia del covid- 19 provocó un enorme aumento del déficit y de la deuda.

A lo largo del tema vamos a repasar todos estos acontecimientos. Así, tras introducir algunos conceptos básicos, repasaremos la evolución del saldo presupuestario y de la deuda de las AAPP en España y realizaremos algunas comparaciones internacionales, prestando especial atención a las diferencias existentes dentro del grupo de las economías avanzadas al que pertenece España. En la segunda parte del tema, analizaremos la política de financiación del Tesoro en la actualidad, una cuestión clave dadas las elevadas necesidades de financiación del Estado, que se sitúan alrededor de los 250.000 millones de euros anuales, equivalentes a algo más del 20% del PIR



\clearpage
\section{Saldo presupuestario y la deuda de las Administraciones Públicas en España: evolución}
\puntosclave{
    En 2025, España registró un déficit público de unos 37.000 a 40.000 millones de euros, según se incluyan o no los gastos de la DANA. Ese mismo año, el Estado pagó por sí solo unos 34.900 millones de intereses. Es decir, casi todo el déficit de las AAPP se explica ya por la carga de la deuda heredada, y el saldo primario está cerca del equilibrio.    
}

En España, el déficit público es un fenómeno relativamente reciente. Hasta la Transición, el Estado mantuvo el equilibrio presupuestario e incluso pequeños superávits. Fue durante la Transición cuando el saldo empezó a mostrar déficits recurrentes, una situación que, con el breve paréntesis de 2005-2007, se ha mantenido hasta hoy. 

Podemos resumir el recorrido en tres puntos clave. El equilibrio previo a 1975 ocultaba el atraso del sistema tributario y solo era posible por el reducido tamaño del sector público. La creación del Estado de bienestar y del Estado autonómico trajo consigo un aumento estructural del gasto que no se acompañó de medidas equivalentes en los ingresos. Esos déficits han sido casi siempre estructuralmente negativos y, además, procíclicos, algo especialmente indeseable. La etapa de 1996 a 2007 fue muy favorable para el saneamiento y permitió los tres únicos superávits de la democracia. La recesión de 2008 provocó déficits superiores al 10% del PIB y llevó la deuda al 105% en 2014. Y la pandemia de 2020 frenó en seco la lenta consolidación posterior.

Este bloque reconstruye esa historia con una pregunta de fondo: ¿por qué España ha tenido déficit en casi todos los años de la democracia, incluidos muchos de expansión, y qué papel han desempeñado en su deuda los déficits, el crecimiento, los tipos de interés y las operaciones financieras? Para responderla, fragmentaremos la pregunta en cinco partes.

\textbf{¿Qué queremos decir?} En primer lugar, fijaremos el el marco de análisis. ¿Qué se mide, cómo se leen en España los indicadores del saldo, cómo se relacionan saldo y deuda y cómo se pasa del déficit de las AAPP a las necesidades del Tesoro? Después, recorreremos la evolución en cuatro etapas: del equilibrio aparente al déficit estructural (1975-1995); la consolidación, el euro y los superávits (1996-2007); la Gran Recesión y la crisis de la deuda soberana (2008-2019); y la pandemia, la inflación y la situación actual (2020-2026). Como veremos, la historia fiscal de la democracia española combina un déficit estructural persistente, porque los ingresos nunca se han ajustado del todo a un Estado de bienestar construido en muy poco tiempo, con grandes perturbaciones que han llevado la deuda a saltos (1992-1993, 2008-2013, 2020). Los periodos de reducción de la deuda se han debido más al crecimiento nominal y a los tipos bajos que a superávits primarios sostenidos. Esa combinación explica que España llegue a 2026 con una ratio de deuda cercana al 100\% del PIB, un saldo primario próximo al equilibrio y una carga de intereses que vuelve a crecer.

\subsection{Marco de análisis: qué se mide y cómo se lee en España}
El título del tema habla de saldo presupuestario y de deuda de las Administraciones Públicas como si fueran magnitudes evidentes. No lo son. Cada una tiene varias definiciones, y la elección entre ellas cambia el diagnóstico. 

No repiteremos aquí los desarrollos conceptuales, pues corresponden a otras partes del temario, pero si las presentaremos por su función operativa: proporcionarnos las claves para leer la evolución española y enlazarla con la política de financiación del Tesoro. Para ello, contestaremos a tres cuestion: ¿Qué se mide? ¿Cómo se lee el saldo en España? ¿Cómo se relacionan saldo y deuda? Con esto, podremos ver cómo se relaciona el déficit de las AAPP a las necesidades del Tesoro.

Como resumen, podemos decir que el saldo presupuestario y la deuda de las AAPP se miden en contabilidad nacional, en términos de devengo y con un perímetro móvil. La deuda se valora por su nominal y se consolida entre subsectores, lo que convierte al Estado en el prestamista de las comunidades y de la Seguridad Social y explica que su deuda represente casi el $90\%$ del total. Con esto, veremos como en España, el saldo primario está hoy cerca del equilibrio, pero la carga de intereses vuelve a crecer. Además, el saldo estructural revela un déficit persistente y enseña, con los falsos superávits de 2005-2007, la fragilidad de las estimaciones en tiempo real. Sin embargo, la deuda no varía en la cuantía del déficit: su evolución depende también del crecimiento nominal, de los tipos de interés y de numerosas operaciones financieras. Por ello, los periodos españoles de reducción de la deuda se han debido más al crecimiento nominal y a los tipos bajos que a superávits primarios. Por último, entre el déficit de las AAPP y las necesidades del Tesoro median los préstamos del Estado a otros subsectores, las diferencias entre caja y devengo y los vencimientos, lo que explica que el Tesoro emita cada año mucho más de lo que dice el déficit. 


\subsubsection{Qué se mide: el saldo y la deuda en su definición oficial}
El saldo presupuestario se define como la diferencia entre los ingresos y los gastos totales de las AAPP durante el ejercicio, es decir, su capacidad o necesidad de financiación. Según el SEC 2010, corresponde al saldo de la cuenta de capital del sector institucional AAPP y es el indicador esencial para las comparaciones internacionales y para valorar el cumplimiento de las reglas fiscales europeas.\footnote{En el SEC 2010, la capacidad ($+$) o necesidad ($-$) de financiación (B.9) es el saldo de la cuenta de capital del sector: el ahorro bruto más las transferencias de capital netas menos la formación bruta de capital y la adquisición neta de activos no financieros no producidos. Es el concepto de déficit público que utiliza el Protocolo de Déficit Excesivo (Reglamento (CE) 479/2009). ([\authorfont{Ver Tema 3.A.39}])
} 

Dos precisiones son relevantes para todo el tema. La primera: el saldo se mide en contabilidad nacional, con criterio de devengo, no en la contabilidad presupuestaria de caja con la que se ejecutan los presupuestos. De ahí que el déficit de las AAPP y las necesidades de financiación del Tesoro no coincidan.\footnote{La distinción entre la contabilidad presupuestaria (caja y derechos u obligaciones reconocidos) y la contabilidad nacional (devengo), y los ajustes que la IGAE aplica para pasar de una a otra, se desarrollan en ([\authorfont{Ver Tema 4.B.21}])
} La segunda: el saldo se refiere al conjunto de las AAPP, pero se publica también por subsectores, lo que permite ver cómo se reparte el déficit entre la Administración central, las comunidades, las entidades locales y la Seguridad Social.

Por el contrario, la deuda se define como el saldo vivo, en un momento dado, de los pasivos contraídos para financiar los gastos que no pueden cubrirse con ingresos corrientes y de capital. Existen distintos conceptos según los instrumentos incluidos y el método de valoración, y el más relevante para este tema es el del Protocolo de Déficit Excesivo (PDE). Este criterio valora la deuda por su valor nominal o facial, lo que evita las grandes oscilaciones que produciría la \textcolor{green}{valoración a precios de mercado}. La deuda PDE incluye el efectivo y los depósitos, los valores representativos de deuda (letras, bonos y obligaciones) y los préstamos, y excluye otros pasivos como los créditos comerciales, es decir, las facturas pendientes de pago, o los pasivos por pensiones. Por eso la literatura y los organismos internacionales complementan la deuda PDE con otras medidas: la deuda neta de activos financieros, el total de pasivos o los pasivos contingentes e implícitos, como los avales públicos o los compromisos de pensiones.\footnote{Los pasivos contingentes (avales, colaboración público-privada) y los implícitos (compromisos de pensiones, recogidos en la tabla complementaria 29 del SEC 2010) se tratan en ([\authorfont{Ver Tema 4.B.24}]) (sobre la deuda implícita del sistema de pensiones) y, para las empresas públicas y los pasivos fuera del perímetro, ([\authorfont{Ver Tema 4.B.22}]).
}

El criterio de valoración nominal tiene una consecuencia práctica que conviene adelantar, porque afecta a la lectura de la deuda. Cuando el Tesoro emite un bono por encima de la par, lo que ocurrió mucho entre 2015 y 2021 al reabrir referencias con cupones superiores a los tipos de mercado, ingresa más de lo que computa como deuda PDE, porque solo computa el nominal. La diferencia, la prima de emisión, financia el déficit de caja sin aumentar la deuda PDE, a cambio de pagar cupones más altos en el futuro. Desde 2022 ocurre lo contrario: con la subida de tipos, muchas reaperturas se emiten por debajo de la par, y el Tesoro necesita emitir más nominal para obtener el mismo efectivo. Por tanto, constituye un ejemplo claro de por qué la deuda no varía exactamente en la cuantía del déficit.

c) El perímetro: quién forma parte de las AAPP.

La Administración central comprende las unidades con competencias en todo el territorio nacional: el Estado (Casa del Rey, Gobierno, Cortes Generales, Tribunal Constitucional, etc.), los organismos de la administración central y las empresas clasificadas como administración central. Las comunidades autónomas incluyen sus órganos de gobierno y sus organismos administrativos, sus universidades y las empresas clasificadas como administración regional. Las corporaciones locales agrupan ayuntamientos, diputaciones, cabildos y consejos insulares, mancomunidades y agrupaciones, las ciudades autónomas de Ceuta y Melilla y sus organismos dependientes. Las administraciones de la Seguridad Social incluyen las unidades que proveen prestaciones sociales.\footnote{
Además, con carácter general, forman parte de las AAPP las unidades institucionales controladas por ellas que sean productores no de mercado, definidos por el SEC-2010. Los criterios de delimitación (control público, carácter de productor de mercado o no de mercado y prueba del $50\%$ de los costes cubiertos con ventas), así como las reclasificaciones de unidades como la Sareb, el ADIF o las empresas públicas, se desarrollan en el ([\authorfont{Ver Tema 4.B.23}]).
}

El perímetro no es una cuestión técnica menor. En España, varias reclasificaciones han alterado de forma relevante la deuda contabilizada. La más importante fue la de la SAREB. En marzo de 2021, siguiendo el criterio de Eurostat, el INE, la IGAE y el BdE la incluyeron en el sector de las AAPP, lo que añadió unos $35.000$ millones de euros a la deuda de 2020 y contribuyó a que la ratio superara el $120\%$ del PIB ese año.\footnote{La inclusión de la Sareb en las AAPP se basó en que el Estado controlaba la entidad a través del FROB y garantizaba su deuda sénior, y en que las pérdidas acumuladas, no previstas inicialmente, desvirtuaban su carácter de productor de mercado. Su vida limitada, con disolución prevista en 2027, había sido una de las condiciones para mantenerla fuera del perímetro. ([\authorfont{Ver Tema 4.B.23}])
} El episodio muestra que la frontera de las AAPP es móvil y que la deuda publicada depende de decisiones estadísticas, no solo de decisiones presupuestarias.

d) La consolidación entre subsectores.

Además, hay que tener en cuenta que la deuda de las AAPP se publica consolidada: se eliminan los pasivos de un subsector que son activos de otro. En España, esa eliminación es muy importante, porque una gran parte de la deuda de las comunidades autónomas y casi toda la de la Seguridad Social se deben a préstamos del Estado. 

Al cierre de 2025, la deuda consolidada de las AAPP ascendía a unos $1,698$ billones de euros ($100,7\%$ del PIB), mientras que la suma de los subsectores sin consolidar superaba esa cifra: Administración central, unos $1,56$ billones ($92,6\%$ del PIB); comunidades, $341.642$ millones ($20,2\%$); corporaciones locales, $20.729$ millones ($1,2\%$); y Seguridad Social, unos $136.200$ millones ($8,1\%$). 

En la práctica, como podemos ver, el Estado se ha convertido en el prestamista de los demás subsectores, y el Tesoro financia en el mercado no solo el déficit del Estado, sino también parte del de las comunidades y la Seguridad Social.


\subsubsection{Cómo se lee el saldo en España: total, primario y estructural}
Ahora que hemos repasado qué es el déficit y la deuda, conviene entender cómo se lee el saldo presupuestario, pues existen tres saldos que responden a tres preguntas distintas. Evidentemente, el \textbf{saldo total} responde a la pregunta de cuánto necesitan endeudarse las AAPP en un año. Pero para leer la política fiscal hacen falta otros dos indicadores. 

En primer lugar, el saldo primario, que excluye los intereses, responde a la pregunta de si los ingresos cubren el gasto corriente y de capital sin contar la carga de la deuda heredada. En segundo lugar, el saldo estructural, que corrige el efecto del ciclo y de las medidas excepcionales, responde a la pregunta de cuál sería el saldo con la economía en su nivel potencial.\footnote{Saldo primario $=$ saldo total $+$ intereses. Saldo ajustado de ciclo $=$ saldo total $−$ $\varepsilon x_t$, donde $\varepsilon$ es la semielasticidad presupuestaria al ciclo (en torno a $0,5-0,6$ para España en la metodología común de la Unión) y la brecha de producción ($x_t$) es la diferencia porcentual entre el PIB observado y el potencial. Saldo estructural $=$ saldo ajustado de ciclo $−$ medidas excepcionales y temporales. Para los conceptos básicos [\authorfont{Ver Tema 3.A.39}]. Para la estimación del producto potencial y de la brecha de producción, y sus problemas en tiempo real, véanse [\authorfont{Ver Tema 3.A.42, Tema 3.A.43}].
} 

¿Qué dicen estos dos indicadores en el caso de España?

\paragraph{Saldo primario: la carga de la deuda heredada.}
En 2025, el déficit de las AAPP fue del $2,2\%$ del PIB, unos $36.800$ millones, o del $2,4\%$, unos $40.300$ millones, si se incluyen los gastos extraordinarios de la DANA de Valencia de octubre de 2024. 

El Estado pagó ese año unos $34.900$ millones en intereses, y el conjunto de las AAPP, algo más de $40.000$ millones. La consecuencia es que el saldo primario de las AAPP estuvo muy cerca del equilibrio en 2025, lo que no ocurría desde antes de la Gran Recesión. No obstante, esa mejora se produce justo cuando el gasto en intereses vuelve a subir. 

Hasta julio de 2026, los intereses pagados por el Estado crecían en torno a un $14\%$ interanual, y la AIReF (2025) y Fedea (2025) proyectan un aumento sostenido de la carga de intereses en la próxima década, a medida que se refinancia a tipos más altos la deuda emitida en la etapa de tipos cero.\footnote{Según distintas proyecciones, la carga de intereses podría duplicarse en porcentaje del PIB hacia 2050, \textit{ceteris paribus}.
} Dicho de otro modo: el esfuerzo fiscal de los próximos años tendrá que hacerse para absorber el coste de la deuda heredada, no para corregir un exceso de gasto corriente.

\paragraph{Saldo estructural: el diagnóstico persistente y sus trampas}
No obstante, es la lectura del saldo estructural la que ofrece el diagnóstico más repetido. La AIReF, la Comisión, el Banco de España (BdE) y el FMI coinciden en que España tiene un déficit estructural persistente: incluso con la economía cerca de su nivel potencial, las AAPP registran un déficit relevante, que diversas estimaciones sitúan en los últimos años en torno a $2,5$ o $3$ puntos del PIB potencial 

Esta lectura presenta, también, un ejemplo histórico reciente. En 2005-2007, España registró superávits totales de hasta el $2\%$ del PIB, y las estimaciones de la época situaban también en superávit el saldo estructural. Con la información posterior, esas estimaciones se revisaron a fondo. 

Buena parte de los ingresos de aquellos años eran extraordinarios: procedían de la burbuja inmobiliaria (impuestos sobre transmisiones, IVA de la construcción, IRPF e Impuesto sobre Sociedades ligados al sector) y no desaparecían con el ciclo tal como lo medían las técnicas convencionales, sino con el pinchazo de la burbuja. Las estimaciones actuales sitúan el saldo estructural de 2007 en déficit. Es decir, España entró en la crisis con un déficit estructural oculto por los ingresos de la burbuja.\footnote{El fenómeno lo documentaron, entre otros, el BdE y la Comisión tras la crisis, con análisis de los ingresos extraordinarios (revenue windfalls) ligados al precio de los activos y a la composición de la demanda. Los métodos estándar de ajuste cíclico, basados en la brecha de producción, no los detectan.
} 

La lección, que retomaremos más adelante, es que el saldo estructural estimado en tiempo real puede ser muy engañoso, y que los superávits de 2005-2007 no fueron tanto un éxito de la política fiscal como un espejismo de ingresos.

\begin{specialblock}{\textbf{Medidas excepcionales}: el caso de la DANA}
    El saldo estructural excluye también las medidas excepcionales y temporales, lo que plantea en cada caso la cuestión de qué se considera excepcional. 
    
    El Gobierno presentó el déficit de 2025 sin los gastos de la DANA ($2,2\%$), mientras que Eurostat y la IGAE publicaron la cifra total ($2,4\%$). La diferencia es una convención: los gastos de reconstrucción tras una catástrofe se consideran temporales en el marco fiscal europeo. Pero muestra que incluso el saldo total admite lecturas distintas según el propósito. \textcolor{red}{(Pendiente de verificar el tratamiento exacto de la DANA en la evaluación de la Comisión)}    
\end{specialblock}

\subsubsection{Cómo se relacionan saldo y deuda: la contabilidad de la variación de la deuda en España}
Habida cuenta de todo lo expuesto, podemos ahora presentar la relación entre ambas. La variación de la ratio de deuda sobre el PIB puede descomponerse contablemente en tres componentes: (i) el saldo primario; ($sp_t$) (ii) el llamado efecto bola de nieve ($(i_t − g_t)/(1 + g_t)]·b_{t−1}$), que recoge la diferencia entre el tipo de interés medio de la deuda y el crecimiento nominal de la economía; y, (iii) los ajustes entre déficit y deuda (ADD), que agrupan todas las operaciones que aumentan o reducen la deuda sin pasar por el déficit.\footnote{Siendo $b$ la ratio de deuda sobre el PIB, $i$ el tipo de interés nominal implícito de la deuda, $g$ el crecimiento del PIB nominal, $sp$ el saldo primario en porcentaje del PIB y ADD los ajustes entre déficit y deuda: $\Delta b_t \approx [(i_t − g_t)/(1 + g_t)]·b_{t−1} − sp_t + \text{ADD}_t$. El primer término es el efecto bola de nieve. Si $i < g$, la ratio puede reducirse aun con déficits primarios moderados. Para más información [\authorfont{Ver Tema 3.A.39}].
} 

Esta descomposición es puramente contable y es la herramienta que nos permitirá entender la evolución de las finanzas de las AAPP en España. Para verlo claramente y fijar los términos, vamos a presentar tres episodios (sobre los que incidiremos más adelante) que lo ilustran perfectamente.

En 2008-2014, la deuda pasó del $36\%\to105\%$ del PIB. Contribuyeron los tres componentes: grandes déficits primarios por la caída de los ingresos y los estabilizadores; un efecto bola de nieve desfavorable, con tipos de interés que subieron con la crisis de la deuda soberana mientras el PIB nominal caía; y ajustes entre déficit y deuda muy importantes, por la recapitalización bancaria con el préstamo del MEDE, la creación del FADE, el mecanismo de pago a proveedores y los préstamos a los países rescatados de la zona del euro.

Más adelante, entre 2015-2019, la deuda bajó lentamente, del $105\%\to98\%$ del PIB, a pesar de que hubo déficits en todos los años. La razón fue el efecto bola de nieve favorable: el crecimiento nominal superó el tipo de interés medio de la deuda, que se reducía gracias a la política del BCE. Finalmente, entre 2021-2025, la deuda bajó del máximo cercano al $120\%$ del PIB (2020) al $100,7\%$ de 2025, de nuevo con déficits todos los años. ¿Por qué? De nuevo, la razón fue el efecto denominador de un crecimiento nominal muy fuerte, impulsado por la recuperación y por la inflación de 2022-2023, con un coste medio de la deuda todavía bajo gracias a la vida media larga.

La conclusión es simple: en la historia reciente de España, los periodos de reducción de la deuda se deben al crecimiento nominal y a los tipos bajos más que a superávits primarios. Por tanto, no es de extrañar que sea una situación considerada \textit{frágil} por la literatura, ya que depende de que la diferencia entre el tipo de interés y el crecimiento siga siendo favorable.\footnote{
No es de extrañar, por tanto, que la reducción de la deuda proyectada por la AIReF se apoye principalmente en el crecimiento del PIB nominal, con una contribución muy notable del deflactor, y que, cuando termine ese impulso, las proyecciones dibujen una dinámica desfavorable a políticas constantes.

La literatura sobre el diferencial entre el tipo de interés y el crecimiento (BLANCHARD, 2019) sostiene que, con $r < g$, la deuda es menos costosa en términos de bienestar, pero advierte de que el diferencial puede invertirse bruscamente, como ocurrió desde 2022.
} 

\begin{specialblock}{\textbf{Reusmen}: Principales ajustes entre déficit y deuda (ADD)}
    Para leer bien la historia conviene tener presente los principales ajustes entre déficit y deuda:
    \begin{lnum}
        \item La recapitalización del sistema bancario a través del FROB, financiada en parte con el préstamo del MEDE de 2012 (41.333 millones utilizados), una parte de la cual computó además como déficit por las pérdidas;
        \item Aportaciones al MEDE y los préstamos bilaterales a Grecia y a los fondos de rescate de la zona del euro;
        \item El Fondo de Amortización del Déficit Eléctrico (FADE), creado para titulizar el déficit de tarifa del sistema eléctrico, cuya deuda se integró en las AAPP;\footnote{El FADE, creado en 2010 y 2011, emitió deuda con garantía del Estado para financiar los derechos de cobro de las eléctricas por el déficit de tarifa acumulado. Tras las decisiones de Eurostat, su deuda se incluyó en la deuda de las AAPP. El déficit de tarifa se atajó con la reforma del sector eléctrico de 2013 (Ley 24/2013).
        }
        \item La Sareb, reclasificada en 2021;
        \item Mecanismos de pago a proveedores, que convirtieron deuda comercial, no computable como deuda PDE, en deuda financiera computable;
        \item Variaciones de los depósitos del Tesoro: el colchón de liquidez que mantiene el Tesoro, que aumenta la deuda sin aumentar el déficit;
        \item Primas y descuentos de emisión;
        \item Préstamos del Estado a las comunidades (FLA) y a la Seguridad Social, que no alteran la deuda consolidada de las AAPP, pero sí la del Estado y las necesidades del Tesoro.
    \end{lnum}

    \textcolor{red}{\textbf{NOTA.-} Convendría desarrollar un poco más que es lo que se experimentó en esos \textit{ajustes}. Ahora mismo solo se entiende que en cada uno de estos episodios se experimentó un ajuste brusco entre déficit y deuda pero no se dice por qué ni en qué dirección. Así como tampoco a raíz de qué componente.}
\end{specialblock}

\subsubsection{Del déficit de las AAPP a las necesidades de financiación del Tesoro}
Por último, y con todo sobre la mesa, debemos destacar que lo relevante para la política de financiación del Tesoro no será el saldo presupuestario de las AAPP, sino la \textbf{necesidad de endeudamiento del Estado}. 

El déficit de las AAPP es una magnitud de contabilidad nacional, en términos de devengo, para el conjunto de las AAPP. La necesidad de endeudamiento del Estado es una magnitud de caja y del Estado. Para pasar de una a otra hay que considerar cinco elementos:
\begin{la}
    \item \textbf{Emisión de otros subsectores}. La parte del déficit que corresponde a los otros subsectores, que el Estado no financia directamente salvo cuando les presta;
    \item \textbf{Criterio de caja}. Las diferencias entre caja y devengo en el propio Estado;\footnote{Sobre los criterios SEC de delimitación del sector AAPP (unidades de mercado y no de mercado, la prueba del 50%) y las reclasificaciones, remito al Tema 4B-23.
    }
    \item \textbf{Préstamos intersectoriales}. Los préstamos del Estado a las comunidades autónomas, a través del FLA y la Facilidad Financiera, y a la Seguridad Social, que el Tesoro debe financiar en el mercado aunque se eliminen al consolidar;
    \item \textcolor{red}{\textbf{Operaciones financieras}. Las operaciones financieras del Estado, como las aportaciones a fondos, los préstamos a otros países o la variación de su tesorería; (Qué quiere decir? Por qué hay que tenerlo en cuenta?)}
    \item \textbf{Vencimientos de deuda}. Los vencimientos de la deuda emitida en años anteriores, que no forman parte del déficit, pero sí de las necesidades brutas de financiación (emisión bruta, operaciones de refinanciación).
\end{la}

Con esto en mente podemos comprender por qué en 2025, el déficit de las AAPP fue de unos $37.000$ a $40.000$ millones de euros. Pero el Tesoro programó para 2025 y 2026 una emisión neta de unos $55.000$ millones y, para 2026, una emisión bruta de $285.693$ millones. 

La diferencia entre el déficit y la emisión neta se explica por los préstamos del Estado a los demás subsectores y por las operaciones financieras. Solo en 2025, la deuda de la Seguridad Social, casi toda con el Estado, aumentó en unos $10.000$ millones, y la de las comunidades, en unos $5.700$ millones. La diferencia entre la emisión neta y la bruta se explica por los vencimientos: el Tesoro debe refinanciar cada año en torno al $13\%$ del saldo vivo. 

He aquí, entonces, la clave para entender el Tema, la razón por la que el título une el saldo de las AAPP y la política de financiación del Tesoro: \textbf{el primero determina, a través de esta bisagra, la dimensión de la segunda}.


\subsection{Desarrollo histórico: ¿cómo ha evolucionado a lo largo de los años?}
Ahora que tenemos en mente lo necesario para comprender estas dos magintudes y la relación entre ellas, pasamos a examinar la evolución histórica de éstas.

Remontémonos a la génesis misma de nuestra Nación, tal y como la concebimos hoy en día: al periodo que se expande desde 1975 hasta 1995. 

\subsubsection{Del equilibrio aparente al déficit estructural (1975-1995)}
Dividiremos el periodo en cuatro partes. El primero trata la herencia: por qué el equilibrio presupuestario de 1975 era un síntoma de atraso y no de solidez. El segundo cubre los años 1975-1985: la construcción simultánea del Estado democrático, del Estado de bienestar y del Estado autonómico, y la omnipresencia del déficit. El tercero cubre la integración en la CEE y una consolidación que fue más cíclica que estructural. Acabaremos con la crisis, el SME y la bola de nieve de la deuda que se prolonga a nuestros días.

Durante estos veinte años, la hacienda pública española pasó de ser una de las más pequeñas de Europa occidental a tener un tamaño homologable a la media comunitaria. Lo hizo en muy poco tiempo y en las peores condiciones macroeconómicas: dos crisis del petróleo, una reconversión industrial, una crisis bancaria sistémica, una transición política y una descentralización territorial sin precedentes. Todo a la vez y en todas partes. Como podemos suponer, entonces, el déficit de este periodo no será un accidente coyuntural. Es el precio de un catch-up fiscal: la construcción del Estado de bienestar se hizo antes de que el sistema tributario pudiera financiarla, y los mecanismos de financiación del déficit (primero la monetización, después la deuda a corto plazo y, al final, la deuda de mercado) fueron cambiando por el camino. La deuda pasa de menos del $10\%$ a más del $60\%$ del PIB, y su composición explica muchas decisiones posteriores del Tesoro.

Identificaremos dos lógicas distintas de acumulación de deuda. Hasta mediados de los ochenta la deuda crece por déficits primarios, pero la represión financiera y la inflación alta mantienen un diferencial tipo-crecimiento favorable, y el déficit se financia primero con el BdE y después con instrumentos casi monetarios como los pagarés. En los noventa, con la liberalización financiera, la defensa del SME y la recesión, el diferencial se vuelve fuertemente positivo: la bola de nieve explica la mayor parte del salto de la deuda de $(\sim41\%) \to (\sim62-65\%)$.\footnote{Además, la autonomía del BdE (1994) pone fin a la monetización, y Maastricht introduce la restricción externa que disciplinará la siguiente década.
}

\paragraph{La herencia: un equilibrio sin Estado}
\textcolor{red}{Como en cualquier cambio, el marco de ideas es especialmente relavante para comprender las decisiones que se toman en cualquier ámbito de las ciencias sociales. Si repasamos brevemente la Historia del Pensamiento Económico, (hay que ampliar un poco para entender el contexto en el que nos situamos teniendo en cuenta que debe acabar hilando perfectamente con el párrafo siguiente ->)}

[...] España llega tarde a cada una de estas fases y las recorre comprimidas. Cuando los países de la OCDE construían sus Estados de bienestar con crecimiento alto y tipos reales bajos (1950-1973), España tenía un sector público mínimo. Cuando lo construyó (1977-1995), el contexto internacional ya era el contrario: crecimiento bajo, desempleo alto y, desde principios de los ochenta, tipos reales positivos y elevados. Esa desincronización es clave para entender por qué el déficit español aparece más tarde que en Europa pero crece más deprisa.

b) Las magnitudes de partida

[DOC] El documento describe las cuentas públicas de 1975 con cuatro rasgos:
- equilibrio presupuestario o incluso pequeños superávits;
- un gasto público en torno al 25% del PIB, unos veinte puntos por debajo de la media comunitaria;
- un sistema tributario de escasa capacidad recaudatoria, ineficiente y poco redistributivo;
- una deuda pública por debajo del 10% del PIB, frente a casi el 50% de media comunitaria.

Las cifras de gasto y deuda son coherentes con las series históricas de Comín y Díaz Fuentes (2005) para las Estadísticas históricas de España y con la reconstrucción de Lago Peñas (2024). Este último sitúa la ratio de deuda por debajo del 10% en 1976 y 1977, y otras series del BdE la dejan en torno al 7% en 1975-1976. Sobre la presión fiscal, la OCDE (Revenue Statistics) sitúa a España en torno al 18% del PIB a mediados de los setenta, frente a una media OCDE cercana al 29%. El diferencial era de diez puntos o más.

(OJO: la comparación casi el 50% de media en los países comunitarios del documento es discutible. La media de la CEE de los nueve en 1975 estaba más cerca del 35-40%, porque la alta deuda belga, italiana e irlandesa convivía con ratios bajas en Alemania y Francia (Francia rondaba el 20% a finales de los setenta según la base histórica del FMI). Si se cita, conviene decir muy superior y no dar una cifra.)

c) Por qué el equilibrio era un síntoma y no una virtud

No obstante, bajo la aparente solidez de las cuentas públicas de 1975 se escondía un sistema tributario poco eficiente y un gasto público insuficiente por tres razones, esencialmente.

La primera es de economía política. COMÍN (1996) y DÍAZ FUENTES (2023) describen la hacienda franquista como insuficiente e injusta. Su equilibrio no respondía a una preferencia social por la austeridad. Era el resultado de un pacto implícito entre el régimen y sus bases de apoyo: impuestos bajos y poco progresivos a cambio de un gasto social mínimo. La reforma tributaria propuesta en 1973 (Libro Verde, impulsado por FUENTES QUINTANA) fue archivada precisamente porque aumentar la progresividad y la transparencia era políticamente inasumible para el régimen. Por tanto, vemos como España llega a la democracia con treinta años de retraso tributario acumulado.

La segunda dirección es financiera. El equilibrio del presupuesto del Estado convivía con mecanismos de financiación del sector público que no pasaban por el presupuesto. ESCARIO, SABATÉ y GADEA (2011) reconstruyen la sombra monetaria del déficit a lo largo de toda la era de la peseta. Tras la pignoración automática de la deuda (hasta el Plan de Estabilización de 1959), el canal principal en los sesenta y setenta fue el redescuento especial: crédito privilegiado del BdE a sectores prioritarios. A él se sumaron los coeficientes de inversión obligatoria, que desde 1971 obligaban a la banca a mantener una parte de su activo en fondos públicos y en crédito privilegiado.\footnote{Caso de manual de la represión financiera de REINHART y SBRANCIA (2015): el sector público se financia a tipos reales bajos o negativos obligando a los intermediarios a absorber su deuda.
} La deuda explícita baja no significaba ausencia de financiación pública; significaba que ese coste lo soportaban de forma implícita los ahorradores, vía tipos de interés reprimidos e inflación.

La tercera dirección es perimetral. Una parte de la actividad pública (el INI, el crédito oficial y numerosos organismos autónomos) quedaba fuera del presupuesto del Estado y fuera de lo que entonces se consolidaba. Con los criterios actuales del SEC, parte de esa actividad se clasificaría dentro de las AAPP. No existe una reconstrucción homogénea que permita afirmar cuánto déficit oculto había. Pero el argumento cualitativo es sólido: el equilibrio de 1975 se mide con un perímetro más estrecho que el de hoy.

La lectura en términos de dinámica de la deuda es inmediata: $\Delta d_t ≈ (i − g)/(1 + g)·d_{t-1} − sp_t + \text{ADD}_t$. En 1960-1975 la economía española creció a tasas nominales muy altas (real en torno al $7\%$ más una inflación creciente), y el tipo de interés implícito de la deuda estaba reprimido. El diferencial entre tipo de interés y crecimiento nominal era fuertemente negativo, de modo que incluso pequeños déficits primarios habrían sido compatibles con una ratio de deuda decreciente. 

Por eso la deuda cae a mínimos históricos justo cuando empieza la Transición: el punto de partida no reflejaba una disciplina fiscal extraordinaria, sino un entorno de crecimiento nominal alto y represión financiera que desaparecería en la década siguiente.

\paragraph{1975-1985: tres Estados construidos a la vez y la omnipresencia del déficit}
La Constitución de 1978 consagra el Estado de bienestar y el Estado autonómico, lo que trae un aumento del gasto público de carácter estructural. A ese componente se suma otro coyuntural: la crisis económica internacional de los setenta obligó a una reconversión industrial que absorbió grandes cantidades de recursos públicos. El gasto supera el $40\%$ del PIB en 1985.

La distinción entre componente estructural y coyuntural es pertinente, pero conviene precisar cuatro motores del gasto con su base normativa. Los tres primeros son estructurales: (i) \textbf{pensiones}, su gasto crece por maduración del sistema, extensión de la cobertura y revalorizaciones por encima de la inflación en los primeros años de la democracia ([\authorfont{Ver Tema 4.B.24}]); (ii) \textbf{protección por desempleo}, que la Ley 31/1984 la reordena en un contexto en que el paro supera el $20\%$ a mediados de los ochenta; y, (iii) \textbf{sanidad y educación}, cuya universalización culmina con la Ley 14/1986, General de Sanidad, y la LOGSE de 1990. El cuarto motor es coyuntural pero persistente: la reconversión industrial, con la Ley 27/1984 de reconversión y reindustrialización y las pérdidas de las empresas del INI.

Pero hay un quinto elemento especialmente relevante para la deuda: la crisis bancaria de 1977-1985 y la expropiación de Rumasa de 1983. SUDRIÀ (2014) documenta que afectó a $58$ bancos, el $52\%$ del total, que representaban el $27\%$ del sector en recursos ajenos, recursos propios o empleo. Las aportaciones públicas rondaron $1,3$ billones de pesetas hasta 1985, frente a unos ingresos presupuestados del Estado de $4,6$ billones ese año. Una parte se canalizó a través del Fondo de Garantía de Depósitos (creado en noviembre de 1977) y del propio BdE, y otra a través del presupuesto. Atendiendo a nuestra descomposición, por tanto, una fracción de ese coste aparecería como transferencia de capital en el déficit y otra como ajuste déficit-deuda.\footnote{No existe una partición homogénea publicada, pero es el primer gran episodio de lo que conocemos como deuda que no es déficit. Por tanto, el paralelismo con 2012 (FROB, MEDE, Sareb) es claro y conviene explotarlo en el examen.
} 

A estos motores se suma la descentralización. La LOFCA (1980) y los traspasos de competencias iniciados a principios de los ochenta crean un nuevo nivel de gasto. En esta etapa el efecto no es tanto un aumento neto del gasto (las competencias se traspasan con su coste efectivo) como una pérdida de control centralizado sobre su dinámica futura, porque el gasto autonómico empieza a crecer con lógicas propias.\footnote{Sobre la evolución del sistema de financiación autonómica y sus incentivos fiscales (restricción presupuestaria blanda, corresponsabilidad), remito al Tema 4B-25 (ya redactado).
}

\textcolor{red}{b) La respuesta por el lado de los ingresos: los Pactos de la Moncloa y su lento rendimiento

[DOC] Según el documento, los Pactos de la Moncloa (1977) acordaron una profunda reforma del sistema fiscal:
- se revisaron el IRPF y el Impuesto sobre Sociedades, que adquieren su forma básica actual;
- se creó el Impuesto sobre el Patrimonio;
- se eliminaron figuras arcaicas: se derogaron los impuestos de producto y el impuesto general sobre la renta, sustituidos por retenciones en la fuente y un IRPF progresivo.

Añade que los ingresos tardaron mucho más en responder que el gasto, y de ahí la omnipresencia del déficit público.

Para que la exposición sea rigurosa conviene fijar la secuencia normativa:
- Ley 50/1977, de 14 de noviembre, de Medidas Urgentes de Reforma Fiscal. Crea el Impuesto Extraordinario sobre el Patrimonio de las Personas Físicas, que era extraordinario en el nombre pero estuvo en vigor hasta la Ley 19/1991. También tipifica el delito fiscal y levanta el secreto bancario a efectos tributarios.
- Ley 44/1978, del IRPF.
- Ley 61/1978, del Impuesto sobre Sociedades.

(OJO: el documento dice que se creó el Impuesto sobre el Patrimonio. Técnicamente se creó uno extraordinario; el impuesto ordinario es de 1991. Es una precisión menor pero fácil de señalar.)

La pregunta de fondo es por qué la reforma tardó tanto en rendir. La literatura ofrece tres explicaciones complementarias:
1. Administración tributaria. Faltaban medios, censos y cultura de cumplimiento; la AEAT no se crea hasta 1992. Un impuesto sintético y progresivo como el IRPF exige una administración que España no tenía.
2. Progresividad en frío con inflación alta. En un contexto de inflación del 15-25\%, la progresividad en frío aumentaba mecánicamente la recaudación del IRPF. Esto alivió en parte el problema, pero generó una fuerte presión política para deflactar la tarifa y erosionó la legitimidad del impuesto. Es el mismo debate que ha reaparecido en 2022-2024.
3. Composición del esfuerzo. Según los datos de la OCDE, la presión fiscal pasó de ~18\% a ~27\% del PIB entre 1975 y 1985, un aumento de unos nueve puntos que en términos comparados es enorme. Pero el gasto creció más de quince puntos en el mismo periodo. El déficit no fue consecuencia de que los ingresos no subieran. Fue consecuencia de que el gasto subió aún más deprisa.}


El resultado de todo lo anterior fueron déficits primarios persistentes con deuda todavía moderada. La base histórica de deuda del FMI sitúa la ratio en el $40,9\%$ en 1985 ($16,1\%$ en 1980, $29,5\%$ en 1983). ALBENTOSA y ROCAMORA (2016) dan un $44\%$ de deuda y un $5,8\%$ de déficit. 

No obstante, lo interesante no es el nivel sino la mecánica de la acumulación. En diez años la deuda sube más de $30$ puntos de PIB, de $7\%\to 41\%$. Con un crecimiento nominal medio del orden del $15\%$ anual e intereses implícitos todavía reprimidos, el diferencial entre tipo de interés y crecimiento fue negativo y restó deuda. Por lo tanto, la totalidad del aumento (y algo más) se explica por los \textbf{déficits primarios acumulados} y por los \textbf{ajustes déficit-deuda} vinculados a la crisis bancaria y a la asunción de pasivos de empresas públicas. 

Y, ¿cómo se financió? ESCARIO, SABATÉ y GADEA (2011) precisan que, entre 1978 y 1983 el sector público alcanzó su máxima contribución histórica a la creación de base monetaria, fundamentalmente mediante el crédito del BdE al Tesoro a través de su cuenta corriente. Los autores sitúan en 1984 el fin efectivo de la financiación por señoreaje. A partir de ese año la contribución del sector público a la base monetaria se vuelve negativa, y el déficit se financia con deuda colocada fuera del banco central. El instrumento de transición fueron los Pagarés del Tesoro, creados en 1981: títulos a corto plazo con un régimen fiscal opaco (no sujetos a retención), que los hacía atractivos para capitales que buscaban anonimato. Además eran computables en los coeficientes obligatorios de la banca. Según los mismos autores, entre 1984 y 1989 los pagarés y las letras explican cerca de una cuarta parte del crecimiento de los activos líquidos en manos del público (ALP), y en 1986 más de la mitad. Ahora bien, estos activos eran casi dinero: no ampliaban la base monetaria pero sí la liquidez de la economía. Por tanto, dejar de monetizar no fue un salto a la ortodoxia, sino una transición gradual por instrumentos intermedios.\footnote{La lectura teórica es la de la aritmética monetarista desagradable de SARGENT y WALLACE (1981). Si la autoridad fiscal es dominante y fija una senda de déficits que el público no absorbe como deuda de forma indefinida, la monetaria acabará teniendo que monetizarlos, y una política monetaria restrictiva hoy puede traducirse en más inflación mañana ([\authorfont{Ver Tema 3.A.39}]).

Mientras la autoridad fiscal fija la senda de déficits y la monetaria se acomoda, hay dominancia fiscal. Romperla exige que el banco central deje de financiar al Tesoro, y en España eso no se consolida legalmente hasta 1994. La reforma de 1983-1984, que acota la apelación del Tesoro al BdE, es el primer paso. En términos de credibilidad es incompleto, porque la separación dependía de acuerdos administrativos y no de una prohibición legal.
} 

\paragraph{1986-1990: la CEE, el IVA y una consolidación más cíclica que estructural}
El marco temporal que se expande de 1986 a 1990 fue una etapa de reducción del déficit gracias a cuatro factores: el crecimiento en torno al $5\%$ (que alivia los estabilizadores automáticos), la reducción de subvenciones exigida por la normativa comunitaria, la maduración de las reformas de los Pactos de la Moncloa y la introducción del IVA en 1986.\footnote{El IVA entra en vigor el 1 de enero de 1986 (Ley 30/1985), como condición de la adhesión (Sexta Directiva, 77/388/CEE). Sustituye al Impuesto General sobre el Tráfico de Empresas, un impuesto en cascada, y a otras figuras indirectas. Sus efectos recaudatorios fueron positivos pero menores de lo que suele suponerse, porque sustituía a impuestos existentes ([\authorfont{Ver Tema 4.B.20}]). Por el lado del IRPF, la STC 45/1989 declara inconstitucional la tributación conjunta obligatoria de los matrimonios. Esto obligó a una reforma que culmina en la Ley 18/1991. La sentencia tuvo un coste recaudatorio y abrió un periodo de inestabilidad normativa ([\authorfont{Ver Tema 4.B.16}]).
} 

Conforme a la base histórica del FMI, la deuda se estabiliza en un $38,5-42\%$ durante el periodo: $42,1\%$ en 1986, $38,5\%$ en 1988 y $41,3\%$ en 1990. Por lo que respecta al déficit, su mínimo del periodo se alcanza en 1989, en torno al $3\%$ (o algo por debajo según la serie). No obstante, en 1990 el saldo ya se deteriora hasta el $4\%$ aproximadamente. La cuestión que resulta interesante dentro del periodo es, ¿por qué? ¿Por qué ese mínimo es tan esporádico? Aquí está el punto de mayor calado analítico. 

Entre 1986 y 1990 la economía española creció por encima de su potencial: un $4,5-5,5\%$ anual real con inflación en aumento. En una fase así, el saldo total mejora casi mecánicamente por los estabilizadores automáticos aunque la política fiscal discrecional sea expansiva. La lectura dominante en la literatura es que la política fiscal española fue procíclica en la fase alcista: el saldo estructural no mejoró o empeoró mientras mejoraba el total.\footnote{Ver: OCDE, Economic Surveys de finales de los ochenta; el BdE en sus Informes Anuales de la época; y, con perspectiva, los análisis de ciclicidad de Galí y Perotti, 2003, para Europa.
} La prueba política fue la huelga general del 14 de diciembre de 1988. Tras ella el Gobierno aumentó el gasto social: mejora de las pensiones mínimas, extensión de la protección por desempleo (el documento lo recoge en nota) y, poco después, la Ley 26/1990 de pensiones no contributivas. Evidentemente, son decisiones legítimas en términos de equidad, pero se toman en el pico del ciclo y comprometen gasto estructural justo antes de la recesión.\footnote{ALESINA y DRAZEN (1991), con su modelo de guerra de desgaste, predicen que la estabilización se retrasa cuando hay conflicto distributivo sobre quién soporta el ajuste, lo que es plausible en el contexto de ruptura del diálogo social que culmina en el 14-D. Por su parte, ROUBINI y SACHS (1989) asocian los déficits a los gobiernos de coalición. Su predicción no encaja con España, que tuvo mayorías absolutas de un solo partido entre 1982 y 1993. La explicación española hay que buscarla más en la debilidad de las instituciones presupuestarias (VON HAGEN, 1992, sitúa a España en el grupo de procedimientos presupuestarios poco centralizados) y en la falta de reglas fiscales, que no existirán hasta Maastricht.
}

No obstante, el elemento estructural más duradero del periodo no fue fiscal sino institucional: la CEE, y desde 1989 el SME, introducen por primera vez un ancla externa. Con esto, España materializó lo que GIAVAZZI y PAGANO (1988) concibieron como \textit{atarse las manos}: un país con baja credibilidad antiinflacionista importa la del Bundesbank al fijar su tipo de cambio. 

No obstante, esto tuvo un problema: el ancla monetaria sin ancla fiscal genera tensiones. Con una política fiscal laxa (como la de la época), el BdE sostuvo tipos de interés muy altos para contener la inflación. Esos tipos atrajeron capitales y apreciaron la peseta, y tras la entrada en el SME en junio de 1989 la peseta se situó en la parte alta de la banda. Combinación (fiscal expansiva, monetaria restrictiva, tipo de cambio sobrevalorado) que se paga en 1992-1993.

Un patrón opuesto al que veremos a continuación.

\paragraph{1991-1995: la crisis, el SME y la bola de nieve}
La economía se desacelera desde 1990 y desemboca en la crisis de 1992-1993, en parte por las debilidades del SME. 

\textcolor{red}{\textbf{NOTA.- Reescribir sin referencias a errores, arreglar el documento original y prosa continua: no esquemas}. Varias economías sufrieron crisis cambiarias: el Reino Unido, Italia, Finlandia, Suecia y Noruega abandonaron el SME, y España devaluó frente al ECU en tres ocasiones, la última en marzo de 1993, antes de que en agosto de 1993 se ampliaran las bandas al $\pm1,5\%$.

Dos correcciones. [Seguro]
- Finlandia, Suecia y Noruega no pertenecían al Mecanismo de Cambios del SME. Mantenían vinculaciones unilaterales al ECU, que abandonaron en 1992. Los que salieron del mecanismo de cambios fueron el Reino Unido y la lira italiana, en septiembre de 1992.
- Las devaluaciones de la peseta fueron cuatro, no tres, y la tercera no fue en marzo: un 5\% en septiembre de 1992, un 6\% en noviembre de 1992, un 8\% en mayo de 1993 y un 7\% en marzo de 1995. La confusión probablemente viene de mezclar la de mayo de 1993 con la de marzo de 1995.

La ampliación de bandas al $\pm15\%$ en agosto de 1993 es correcta.

[DOC] El documento añade varios factores:
- las grandes inversiones de 1992: el AVE Madrid-Sevilla, los Juegos Olímpicos de Barcelona y la Expo de Sevilla;
- el aumento del gasto en desempleo, con un paro que alcanza el 24,5\% en 1994, en parte por la reforma laboral de 1984 (contratación temporal) y por las mejoras de prestaciones tras el 14-D;
- el aumento de la carga de intereses.

El resultado sería un gasto superior al 45\% del PIB en 1993, un déficit superior al 7\% y una deuda del 65% en 1995.

La descripción es correcta en su lógica, con matices de cifras:
- La tasa de paro de la EPA de la época superó el 24\% en 1994. Las series homogeneizadas posteriores la rebajan algo, así que conviene decir en torno al 24\%.
- La Ley 32/1984 introdujo el contrato temporal de fomento del empleo, y la literatura sobre dualidad (Bentolila y Dolado, 1994) le atribuye la extrema sensibilidad del empleo al ciclo.
- Sobre el déficit, las series difieren según la contabilidad: el máximo del periodo se sitúa entre el 6,5\% y el 7,5\% del PIB, en 1993 o en 1995 según la serie. Esteve y Prats (2024) identifican 1995 como el año del déficit récord hasta entonces (6,6\%).
- Sobre la deuda, la base del FMI da un 61,6\% en 1995 y un 65,4\% en 1996. El 65\% del documento corresponde en realidad a 1996, como el propio documento reconoce en su introducción. Hay que corregir la incoherencia interna.

b) La mecánica de la deuda: el diferencial toma el relevo

Aquí la descomposición de 1.1.3 es más reveladora que en ningún otro momento del periodo. Entre 1990 y 1995 la deuda sube unos 20 puntos (del 41,3\% al 61,6\% según el FMI), y más de 10 puntos solo en 1993 (del 44,1\% al 54,6\%). Ese salto de un año no se explica solo por el déficit. En orden de magnitud, contribuyen tres factores:
1. El déficit primario. Ronda el 2\% del PIB en 1993, con un saldo total de ~7\% y unos intereses de ~4,5-5\%.
2. El efecto bola de nieve. En 1993 el PIB real cae (≈ −1\%) y el crecimiento nominal apenas llega al 2-3\%, mientras el tipo implícito de la deuda supera el 10\%. Con una ratio en torno al 45\%, el diferencial aporta por sí solo en torno a 3-4 puntos.
3. Los ajustes déficit-deuda. Incluyen la revalorización en pesetas de la deuda en divisas tras las devaluaciones, la asunción de deudas de entes públicos y la acumulación de activos líquidos del Tesoro.

Estas cifras son una estimación aritmética mía a partir de las magnitudes agregadas, no una descomposición publicada. Sirven para la narrativa, pero no las citaría como dato exacto sin verificar con AMECO.

(OJO: Prats Albentosa y Rocamora (2016) dan un déficit primario del 0,5\% en 1993. Con un déficit total cercano al 7\% y unos intereses que, según los mismos autores, alcanzan el 5,3\% en 1996, esa cifra parece baja. Una diferencia de vintage contable es la explicación más plausible. Lo marco como pendiente.)}

La conclusión analítica es nítida. Entre 1975 y 1985 la deuda crece por déficits primarios con un diferencial favorable; entre 1991 y 1995 crece sobre todo por el diferencial, con déficits primarios moderados. Los intereses llegan a su máximo histórico en 1995-1996, en torno al $5,3\%$ del PIB, frente al $2,5\%$ de 2024-2025. El cambio de régimen tiene tres causas: la liberalización financiera, que acaba con la represión y eleva los tipos reales; la defensa de la peseta en el SME, que exige tipos muy altos; y la recesión, que hunde el crecimiento nominal. Este es el nexo con el tema del Tesoro. La estructura de la deuda de la época (predominio del corto plazo, vida media de pocos años) hacía que cada subida de tipos se trasladara rápidamente al coste medio. La vida media de la deuda del Estado rondaba entonces los $2-3$ años, frente a más de $8$ hoy. Alargar esa vida media será el objetivo central de la política de emisiones desde 1995-1996.

\textcolor{red}{\textbf{NOTA.- Lo mismo que antes, menudo desastre y continuo de cosas esquemáticas, con referencias a errores, etc. Depurar y reescribir} c) El marco institucional cambia: Maastricht y la autonomía del BdE

[DOC] El documento señala que el Tratado de Maastricht (1992) incluyó el requisito de situar el déficit por debajo del 3\% del PIB para entrar en la UEM, lo que obligó a iniciar la consolidación, favorecida desde 1995 por la mejora de la coyuntura.

Corrección. El año 1995 no es el inicio de la consolidación: es el año del déficit récord de la serie de Esteve y Prats. Los PGE para 1996 fueron rechazados en el Congreso en octubre de 1995, y los de 1995 hubieron de prorrogarse. El ajuste real empieza en 1996 con el nuevo Gobierno, sobre esos presupuestos prorrogados. Lo desarrollo en 1.3.

Por el lado institucional, dos hitos del periodo son decisivos para el tema:
- Programa de Convergencia de 1992. Es el primer ejercicio de planificación fiscal plurianual con objetivos numéricos vinculados a una restricción externa. Sus objetivos se incumplieron en 1992-1993 por la recesión y hubo que revisarlo en 1994. Su interés histórico está en ser el precedente de los Programas de Estabilidad\footnote{Evolución de la planificación fiscal plurianual europea (Programas de Convergencia y Estabilidad, semestre europeo, planes fiscales estructurales de medio plazo del marco de 2024), en el Tema 3B-45.
}
- Ley 13/1994, de Autonomía del BdE. Su artículo 13.2 prohíbe la autorización de descubiertos o la concesión de cualquier otro tipo de crédito por el BdE al Estado, en aplicación del artículo 104 del Tratado de Maastricht. El saldo deudor que el Tesoro mantenía con el BdE se reconvirtió en un préstamo a largo plazo, amortizable a lo largo de varias décadas. Es el cierre formal de la etapa de monetización: desde 1994, todo el déficit se financia en el mercado. Es la premisa de toda la tercera parte del tema.

Por último, la crisis de 1992-1993 deja un legado para la sostenibilidad del sistema de protección social. Obliga a recortar prestaciones por desempleo (RDL 1/1992 y Ley 22/1992) y lleva al Pacto de Toledo (1995), que separa las fuentes de financiación de la Seguridad Social. En esos mismos años el Estado empieza a conceder préstamos a la Seguridad Social para cubrir sus desequilibrios. Es el origen remoto de los préstamos que figuran en el inventario de ajustes déficit-deuda de 1.1.3 y que se desarrollan en el Tema 24.\footnote{Los préstamos del Estado a la Seguridad Social, su tratamiento contable (no computan como déficit del Estado sino como activo financiero, con impacto nulo en la deuda consolidada) y su condonación parcial, en el Tema 4B-24.
}

La deuda de 1995 es alta, cara y corta, y reducir su coste y alargar su vida media será la tarea del Tesoro en los años siguientes.}

\subsubsection{Consolidación de Maastricht y el espejismo del superávit (1996-2007)}
Entre 1996 y 2007 las AAPP españolas pasan de un déficit de más del $5\%$ del PIB a un superávit en torno al $2\%$, y la deuda pública cae a la mitad: de cerca del $65-70\%$ del PIB, según la serie, a en torno al $36\%$. En términos de saldo total es la mejora más larga y continuada de la historia fiscal española, y culmina con los tres únicos superávits de la democracia (2005, 2006 y 2007).

El periodo tiene, sin embargo, dos fases con lógicas distintas. La primera (1996-1998) es un ajuste deliberado con un objetivo externo: cumplir los criterios de Maastricht para entrar en la UEM desde el primer momento. La segunda (1999-2007) es una mejora del saldo impulsada sobre todo por dos fuerzas que no eran de política fiscal: la caída de la carga de intereses que trajo el euro, y una expansión de ingresos ligada al auge inmobiliario y crediticio, un auge transitorio. La mejora del saldo total, por tanto, ocultó un deterioro relativo del saldo estructural, que los instrumentos de medición y las reglas de la época no permitían ver. 

Por tanto, la conclusión a la que debemos llegar no es que la política fiscal de este periodo fuese irresponsable en nivel, porque la deuda baja fue un activo en 2008. Sino, más bien, que fue procíclica en composición y débil en estructura. Y que la lección institucional (reglas de gasto, desequilibrios macroeconómicos, instituciones fiscales independientes) se aprendió después de la crisis, no antes.

\paragraph{1996-1998: el ajuste de Maastricht}
Como se vio España llega a 1995 con el déficit registrado más elevado del periodo democrático hasta entonces. En esa linea, los PGE para 1996 fueron rechazados en el Congreso en otoño de 1995, lo que obligó a prorrogar los de 1995 y precipitó las elecciones de marzo de 1996. 

El nuevo Gobierno empezó el ajuste hacia Maastricht sobre un presupuesto prorrogado, mediante acuerdos de no disponibilidad de créditos, y lo intensificó con los PGE de 1997. El objetivo era la primavera de 1998, cuando el examen de convergencia se evaluaría según los datos de 1997. Teniendo en cuenta que el criterio fiscal exigía un déficit no superior al $3\%$ del PIB y una deuda por debajo del $60\%$ o disminuyendo suficientemente hacia ese valor, la labor era complicada.

Entre 1995 y 1997 el déficit logró bajar más de cuatro puntos. Por el lado del gasto, contribuyeron la contención de la masa salarial pública (congelación de 1997 y contención de la oferta de empleo), el recorte de la inversión pública, la partida más fácil de ajustar a corto plazo y la que más había crecido en 1992-1993, y el inicio de la caída de la carga de intereses, a medida que los mercados descontaban la entrada en el euro y se estrechaban los diferenciales con Alemania. Además, la Ley 24/1997, de consolidación y racionalización del sistema de Seguridad Social, primera traslación legislativa del Pacto de Toledo, reformó el Sistema para su contener el gasto, ampliando de $8$ a $15$ años el periodo de cálculo de la base reguladora, e iniciando la separación de fuentes de financiación.\footnote{Contenido de la Ley 24/1997, del Pacto de Toledo y del Fondo de Reserva: Tema 4B-24 (ya redactado).
}

Debemos tener en cuenta, también, que en 1996 se aprueba también el Programa de Modernización del Sector Público Empresarial del Estado: la privatización de hasta $36$ empresas en cuatro años, con unos ingresos de unos $30.000$ millones de euros, más del $6\%$ del PIB de 1996. Sabemos que en contabilidad nacional, la venta de acciones es una operación financiera: el Estado cambia un activo (participaciones) por otro (liquidez), que no reduce el déficit (no es ingreso no financiero). Sin embargo, las privatizaciones actuaron sobre la deuda a través del ajuste déficit-deuda: los ingresos destinados a amortizar deuda o a reducir la emisión neta.\footnote{Además, también reducen indirectamente la carga futura de intereses, aunque a cambio se pierden los dividendos.
} Por tanto, los ingresos de privatización fueron un factor relevante en la reducción de la ratio de deuda desde el máximo de 1996. 

En definitiva, según el Informe de Convergencia del Instituto Monetario Europeo (marzo de 1998), España registró en 1997 un déficit del $2,6\%$ del PIB y una deuda del $68,8\%$, habiendo alcanzado su máximo en 1996 ($70,1\%$). Por tanto, no estaba ligeramente por encima del $60\%$, pero sí la había reducido y podía considerarse en camino hacia el cumplimiento del criterio. De hecho, aunque el IME advirtió de que España necesitaba un sustancial progreso en la consolidación fiscal para situar la deuda por debajo del $60\%$, estimaba que, con presupuestos equilibrados desde 1999, ese nivel se alcanzaría hacia 2001.

Con todo lo expuesto, no debe extrañar que el caso español de 1996-1998 se haya citado a veces como ejemplo de consolidación fiscal compatible con un crecimiento elevado: el PIB crece en torno al $4\%$ en 1997-1998 mientras se reduce el déficit.\footnote{La hipótesis de la consolidación fiscal expansiva (GIAVAZZI y PAGANO, 1990; ALESINA y PEROTTI, 1995; ALESINA y ARDAGNA, 2010) sostiene que un ajuste creíble, sobre todo por el lado del gasto, puede estimular la demanda privada a través de las expectativas: menos impuestos futuros y menores primas de riesgo. El argumento teórico, a priori, es sólido: en un modelo con agentes ricardianos o con primas de riesgo endógenas, una consolidación creíble puede elevar la renta permanente esperada y reducir los tipos, y el efecto riqueza puede superar el efecto demanda directo, dando un multiplicador negativo. Además, hay cierta evidencia empírica que la respalda, aunque sujeta a cierta controversia metodológica (acción de política frente a variación del saldo ajustado de ciclo; GUAJARDO, LEIGH y PESCATORI, 2014). Para más información [\authorfont{Ver Tema 4.B.26}]
} No obstante, esta lectura cuantitativa debe ser también matizada por los efectos de segunda ronda que conllevaba entrar en la UEM. En la línea de PEROTTI (2011), para los casos europeos de los ochenta y noventa, la mayoría de estas experiencias de consolidación fiscal y crecimiento se explican sobre todo por canales: (i) el desplome de los tipos de interés por la convergencia hacia la UEM, que en España fue extraordinario (el diferencial a diez años frente a Alemania pasa de unos 500 puntos básicos a mediados de 1995 a prácticamente cero en 1998); (ii) la depreciación previa de la peseta en 1992-1995, que mejoró la competitividad; y, (iii) la moderación salarial.

La consolidación fue probablemente condición necesaria para que esa convergencia de tipos se materializase, porque sin ella España no habría entrado en el euro en 1999. Pero no es evidencia de que el ajuste fiscal fuese por sí mismo expansivo. 

\paragraph{1999-2007: el euro y el dividendo de los tipos}
Durante el periodo posterior a la entrada a la UEM hasta el estallido de la burbuja inmobiliaria, España experimenta un crecimiento económico sostenido en paralelo a una notable consolidación fiscal.

El déficit del periodo previo, que en 1999 se situaba algo por encima del $1\%$, se convirtió a un superávit durante el periodo, en 2007 se publicó en tiempo real en el $2,23\%$ del PIB ($23.368$ millones), con una deuda del 36,2\% (a cierre del Ministerio). Paralelamente, el saldo primario mejora del orden de $3,5$ puntos, de un valor cercano al equilibrio en 1996 a un superávit en torno al $3,5\%$ en 2007. Y buena parte de esa mejora primaria se concentra en 2005-2007 (precisamente cuando los ingresos ligados al ciclo inmobiliario alcanzan su máximo). 

La pregunta relevante, por tanto, es cuánta mejora del saldo se debió a la política fiscal y cuánta a factores externos. En primer lugar, los intereses de las AAPP pasan de un máximo de alrededor del $5,3\%$ del PIB en 1995-1996 a en torno al $1,6\%$ en 2007: un ahorro de más de $3,5$ puntos de PIB. Si la mejora total del saldo entre 1996 y 2007 es de unos $7$ puntos: casi la mitad de la mejora del saldo total no fue, por tanto, mejora del saldo primario, sino un dividendo del euro en forma de menor carga de intereses.\footnote{Esta descomposición tiene un paralelo comparado conocido. Los países de la periferia del euro recibieron el mismo dividendo de tipos, y la cuestión es qué hicieron con él. FERNÁNDEZ-VILLAVERDE, GARICANO y SANTOS (2013) sostienen que el acceso a financiación barata debilitó los incentivos a hacer reformas en la periferia. En su lectura, el crédito político del auge hizo menos urgente cualquier ajuste. España ahorró una parte del dividendo, a diferencia de Grecia o Portugal, pero no la suficiente para compensar el carácter transitorio de parte de sus ingresos.
}

Con la identidad $\Delta d \approx (i − g)/(1 + g)·d_{t-1} − sp_t + \text{ADD}_t$, la caída de la deuda de casi $30$ puntos entre 1996 y 2007 tiene dos motores. Unos dos tercios de la reducción se explican por los superávits primarios acumulados, y aproximadamente un tercio por el efecto bola de nieve favorable: con un crecimiento nominal medio en torno al $7\%$ y un tipo implícito que baja del $6-7\%\to4\%$, el diferencial es negativo durante casi todo el periodo. Los ajustes déficit-deuda operan en ambos sentidos: los ingresos de privatización reducen la deuda, mientras que la acumulación de activos y ciertas asunciones de deuda la aumentan. Por tanto, es la imagen especular de 1991-1995. Entonces el diferencial empujaba la deuda hacia arriba con déficits primarios moderados. Ahora la empuja hacia abajo aunque el esfuerzo primario sea limitado.

¿Quién generaba el superávit? La composición por subsectores de los superávits se suele olvidar, y es muy reveladora. En el cierre de 2007: la Administración Central registró un superávit del $1,29\%$ del PIB; la Seguridad Social, uno del $1,25\%$; las CCAA, un déficit del $0,17\%$; y las corporaciones locales, uno del $0,14\%$.

El superávit de la Seguridad Social era estructural en el sentido demográfico del momento: cohortes numerosas en edad de trabajar, una inmigración masiva que elevó la afiliación (más de $19$ millones de afiliados en 2007) y unas pensiones aún moderadas. Una parte importante de ese superávit se canalizaba al Fondo de Reserva, que alcanzó $45.716$ millones de euros a finales de 2007 y $57.223$ millones en 2008. El Fondo invertía mayoritariamente en deuda pública española. Como esa deuda está en manos de otra administración, desaparece en la deuda PDE consolidada. El Fondo reducía así la deuda consolidada y, cuando se dispuso de él en 2012-2017, esa deuda reapareció en manos del mercado. Por tanto, parte del superávit español de 2005-2007 era ahorro de la Seguridad Social para pensiones futuras, no un margen de maniobra de libre disposición. 

La entrada al euro, además, cambió la política de financiación. La vida media de la deuda del Estado pasó de unos $3$ años a mediados de los noventa a más de $6$ años en 2007. Desaparece la deuda en divisas, se redenomina toda la deuda en euros en 1999 y se consolidan los programas de creadores de mercado y de segregación de principal y cupones (\textit{strips}, 1997). El volumen de emisión se reduce y, en algunos años, la emisión neta de medio y largo plazo es negativa o casi nula. Como lo veremos más adelante, aquí basta retener que el Tesoro llega a 2008 con una deuda larga y barata, lo que amortiguó el impacto inicial de la crisis. Pero también llega con un mercado de deuda con poca liquidez y una base inversora cada vez más concentrada en no residentes (en torno al $50\%$ de la deuda negociable en 2007-2008), factor que pesará en 2010-2012.

\begin{specialblock}{\textbf{Material complementario:} ¿Cómo estaba \textbf{realmente} España durante este periodo?}
    Conviene entrar un poco, aunque no sea durante la exposición, a los determinantes del crecimiento económico español del periodo y sus consecuencias sobre los datos que acabamos de arrojar. Tratamos, por tanto, algunas cuestiones que son fundamentales para comprender como las finanzas españolas pudieron deteriorse tan rápidamente durante un lapso de tiempo tan corto.

    \textbf{Espejismo del superávit: ingresos transitorios y saldo estructural en tiempo real}. Lo primero que debería llamarnos la atención es: ¿cómo pasó España de un superávit de más del $2\%$ del PIB a un déficit del $11\%$, en solo dos años? Evidentemente, la caída de la actividad y los estabilizadores automáticos explican una parte. Pero también es evidente que no puede explicarlo todo.
    
    La magnitud del deterioro solo se entiende si se acepta que el saldo de 2007 incluía un componente transitorio muy superior al que estimaban los métodos convencionales. Es el llamado \textbf{falso superávit}.

    DE CASTRO, ESTRADA, HERNÁNDEZ DE COS y MARTÍ (2008), estimaron que el auge inmobiliario generó un aumento acumulado de ingresos tributarios de alrededor de \textcolor{re}{$2$ ¿puntos? ¿o $\%$?} de PIB entre 1998 y 2006. En ese mismo periodo, la presión fiscal total aumentó unos $4\%$: la mitad del aumento de la presión fiscal estaba ligado a la vivienda. Incorporando los efectos riqueza, su estimación del saldo estructural de 2007 bajaba del $1,78\%$ al $1,02\%$ del PIB. Ya advertían de que la desaceleración inmobiliaria había restado $0,6$ puntos de ingresos por IVA e ITP-AJD solo en 2007.\footnote{Las figuras más sensibles fueron: el ITP-AJD, un tributo cedido, lo que explica parte del deterioro autonómico posterior (Tema 25); el IVA de la vivienda nueva; el Impuesto sobre Sociedades de los sectores de la construcción y financiero; el IRPF de las ganancias patrimoniales; y, los impuestos locales vinculados al urbanismo: ICIO, plusvalía municipal y aprovechamientos urbanísticos.
    } MARTÍ y PÉREZ (2015) cifran en torno al $2,5\%$ del PIB el peso de los impuestos sobre la propiedad y las transmisiones en el pico del ciclo. Documentan además que el gasto primario, excluido el desempleo, creció en términos reales alrededor de un $4\%$ anual, por encima del crecimiento tendencial.
    
    Además, tenemos que tener en cuenta que este hecho permaneció oculto porque la metodología de la Unión para estimar el saldo estructural corrige el saldo por la brecha de producción, pero no por los precios de los activos ni por la composición del crecimiento. En España esto produjo dos errores que se sumaron:
    \begin{lnum}
        \item \textbf{Brecha de producción}. El primero fue la infravaloración de la brecha de producción. ALBEROLA, ESTRADA y SANTABÁRBARA (2014) muestran que, con la metodología de la función de producción de la Comisión, la brecha de 2007 estimada con datos hasta 2007 era ligeramente negativa ($-0,4\%$), mientras que con información hasta 2011 pasó a ser positiva ($+1,5\%$): una revisión de casi dos puntos. Su método alternativo de crecimiento sostenible, que incorpora los desequilibrios financieros y exteriores, estimaba una brecha del $5,4\%$ en 2007, y apenas se revisaba con información nueva. Con una semielasticidad presupuestaria de $0,5$ aproximadamente, esa diferencia de brecha supone entre $1$ y $3$ \textcolor{red}{¿puntos?} de saldo estructural. El superávit estructural que se veía en tiempo real se convierte, ex post, en equilibrio o déficit estructural según el método;
        \item \textbf{Falacia de composición}. El segundo es la omisión del componente de composición (DE CASTRO et al., 2008). Aunque la brecha de producción se hubiera estimado bien, el saldo ajustado de ciclo no captura los ingresos ligados a un crecimiento intensivo en construcción y en transacciones de activos: un crecimiento \textbf{nominal}. Estos generan mucha más recaudación por unidad de PIB que el consumo o la exportación. 
    \end{lnum}
    
    En definitiva, la conclusión que comparte hoy la literatura, del BdE a la AIReF y la Comisión, es que el saldo estructural español de 2007 estaba lejos del superávit que mostraba el saldo total. Las estimaciones ex post lo sitúan entre el equilibrio y un déficit de $1-2$ puntos, según el método.

    \textbf{El debate: ¿prudencia o prociclicidad?}. Dicho esto, lo que conviene determinar es si, dado que en realidad se experimentaba un déficit, el Gobierno actuaba con prudencia o con prociclicidad. Este debate conviene presentarlo desde dos posiciones, porque ambas tienen argumentos.
    
    La posición defensiva, que es la ortodoxia oficial de la época y la de una parte de la literatura posterior sobre la crisis del euro, sostiene cuatro cosas. En primer lugar, que España cumplió holgadamente el PEC, a diferencia de Alemania y Francia. En segundo lugar, que redujo su deuda a la mitad. En tercer lugar, que acumuló un Fondo de Reserva para la Seguridad Social. Y, en cuarto lugar, que la crisis española de 2008-2012 fue una crisis de balance privado (endeudamiento de hogares y empresas, burbuja inmobiliaria, cajas de ahorros) que acabó trasladándose al sector público, no una crisis de indisciplina fiscal. Esta linea es la lectura de DE GRAUWE (2011) y de buena parte de la literatura sobre la crisis del euro que distingue entre Grecia (crisis fiscal) e Irlanda y España (crisis bancaria y privada). Evidentemente, tiene un sólido apoyo empírico: con una deuda del $36\%$ en 2007, España tenía un margen que no tenían Italia o Grecia, y ese margen permitió absorber el choque inicial.

    No obstante, la posición crítica (BdE, Martí y Pérez, Hernández de Cos, y la OCDE y el FMI en sus revisiones posteriores) sostiene tres cosas. En primer lugar, que en una economía que crecía muy por encima de su potencial sostenible, con un déficit exterior cercano al $10\%$ del PIB y una inflación diferencial persistente, la política fiscal debió haber sido mucho más restrictiva. En segundo lugar, que el superávit era insuficiente para contrarrestar el recalentamiento de la economía. Y, en tercer lugar, que la política fiscal fue procíclica en su composición, pues los ingresos extraordinarios se usaron para financiar gasto primario permanente y rebajas fiscales permanentes (bonificaciones/gasto fiscal). Entre estas se cuentan las reformas del IRPF de 1998 (Ley 40/1998), 2002 (Ley 46/2002) y 2006 (Ley 35/2006), y la rebaja del tipo del Impuesto sobre Sociedades del $35\%$ al $30\%$ en 2007-2008. A ello le añaden la ausencia de política macroprudencial en un país sin política monetaria propia: la política fiscal era el único instrumento nacional de estabilización, y no se usó con esa función.

    Evidentemente, las dos posiciones son compatibles. El nivel de deuda de 2007 era efectivamente bajo y fue un activo en la crisis; y el saldo estructural y la estructura de ingresos eran mucho más débiles de lo que parecía. El problema, por tanto, no fue el punto de partida en deuda, sino el agujero estructural de ingresos que quedó cuando desaparecieron los transitorios: la presión fiscal cayó más de $6$ puntos de PIB.

    Además, existían una serie de pasivos (públicos) que se estaban acumulando fuera del saldo. Como ya hemos visto, en esos años se gestan varios pasivos que emergerán después como deuda o como déficit, y que figuran en el inventario de ajustes déficit-deuda: el déficit de tarifa eléctrica, que empieza a acumularse a principios de los 2000 y acabará titulizado a través del FADE; las facturas no contabilizadas de las CCAA y las entidades locales, sobre todo en sanidad, que aflorarán con el plan de pago a proveedores de 2012; la deuda de entidades del sector público empresarial clasificadas fuera de las AAPP, como ADIF, creada en 2005; y, el pasivo contingente del sistema bancario, concentrado en las cajas de ahorros, cuya materialización en 2012 (FROB, MEDE, Sareb) es el mayor ajuste déficit-deuda de la historia fiscal española.
\end{specialblock}

En definitiva, la lectura más extendida de este periodo (España llegó a 2008 con las cuentas saneadas) es, en el mejor de los casos, incompleta. Los superávits de 2005-2007 se sostenían sobre unos ingresos transitorios ligados al ciclo inmobiliario y sobre un ahorro de intereses que no dependía de España, mientras el gasto primario crecía por encima de la tendencia de la economía. Este periodo, 1996-2007, no fue una historia de disciplina fiscal porque no se hubiese evaporado en $18$ meses.

\subsubsection{Gran recesión, la crisis soberana y la consolidación incompleta (2008-2019)}
Como todos sabemos, el siguiente periodo es el de la Gran Recesión, la crisis de deuda soberana y el agónico letargo de la consolidación fiscal en 2019. Es la década más intensa de la historia fiscal española desde la Guerra Civil. 

La década tiene tres fases con lógicas distintas. La primer, 2008-2009, el desplome. Desaparecen los ingresos transitorios de la fase anterior, actúan los estabilizadores automáticos y se aplica un estímulo discrecional que acabó siendo en parte permanente. La segunda, 2010-2013, la crisis de deuda soberana. El ajuste fiscal se hace bajo presión de los mercados y de Bruselas, y el sector bancario se reestructura con apoyo europeo. La tercera, 2014-2019, la recuperación. El déficit total se reduce, sobre todo por el ciclo y por los tipos de interés, pero el déficit estructural queda prácticamente sin corregir.

Como ya sabemos, transforma la posición fiscal española: del $36\%$ al $98\%$ del PIB de deuda, con un déficit estructural cercano al $3\% $al final.

\textcolor{red}{Esta década tiene dos lecturas oficiales, y ninguna de las dos resiste bien el análisis. El documento base sugiere que las reformas estructurales y el ajuste de 2010-2013 explican la mejora del saldo desde 2015. Pero la mayor parte de la reducción del déficit entre 2014 y 2019 fue cíclica y financiera: crecimiento, empleo y caída de los intereses gracias al BCE.

La verdad incómoda es esta: España cerró 2019, en la fase alta del ciclo y con la brecha de producción ya positiva, con un déficit del 3,1\% del PIB. Eso es prácticamente el mismo déficit estructural con el que salió de la recesión. Si en el examen presentas 2008-2019 como crisis, ajuste y saneamiento, el tribunal puede preguntarte por qué el procedimiento de déficit excesivo se cerró en 2019 y el déficit volvió al 3% ese mismo año.

Hay otro problema, de datos. Las cifras de esta década se han revisado varias veces, sobre todo por la reclasificación retroactiva de la Sareb en 2021. El déficit de 2012 ya no es el 10,6\% que se publicó en su día, sino el 11,5\%. En cada caso indico qué serie uso.

Con la serie vigente, el saldo pasa de un superávit del 1,9\% del PIB en 2007 a un déficit del 11,2\% en 2009, y la deuda se triplica: del 36\% en 2007 a 105\% en 2014. Por primera vez desde la entrada en el euro, España pierde en la práctica el acceso normal a los mercados. En el verano de 2012 la prima de riesgo supera los 600 puntos básicos y el Gobierno solicita asistencia financiera europea para el sistema bancario.
}


\paragraph{2008-2009: el desplome}
El desplome de 2008-2009 fue, a priori, una crisis de ingresos: desaparecieron los ingresos transitorios del auge inmobiliario, y el estímulo discrecional, en parte con medidas permanentes, se aplicó sobre un diagnóstico erróneo de problema cíclico.

Durante la época, el saldo pasa de un superávit del $1,9\%$ en 2007 a un déficit del $4,1\%$ en 2008 y de más del $11\%$ en 2009, que no bajará del $9,5\%$ hasta 2012. Por su parte, la deuda sube casi $70$ puntos, del $36\%$ en 2007 al $105\%$ en 2014.

En términos generales, este deterioro se puede atribuir a cuatro factores: los estabilizadores automáticos (con un paro que llegó a superar el $25\%$ y una caída de la recaudación de más de $6$ puntos de PIB, sin paralelo en el entorno), las políticas expansivas (Plan E), las ayudas al sector bancario y la crisis de deuda soberana.

Con la serie vigente en SEC-2010, el déficit de 2008 es del 4,6\%, no del 4,1\%. El diagnóstico de los cuatro factores es correcto, pero mezcla fases distintas. Las ayudas bancarias y la crisis soberana pertenecen a 2010-2013; en 2008-2009 dominan los dos primeros factores. La tasa de paro de la EPA alcanzó su máximo en el primer trimestre de 2013, con un 26,9%.

La variación del saldo entre 2007 y 2009 es de unos $13$ puntos de PIB. Más de la mitad procede de los ingresos. La presión fiscal, incluidas las cotizaciones, cae más de $6$ puntos en dos años: del $41\%$ al $35\%$ del PIB. No es un simple efecto de los estabilizadores automáticos, porque el PIB nominal cayó mucho menos: es la reversión de los ingresos transitorios y responde al esquema de DE CASTRO et al. (2008): el ITP-AJD y el IVA de la vivienda se desploman con las transacciones; el Impuesto sobre Sociedades pierde las bases del sector inmobiliario y financiero, y las pérdidas de ejercicios anteriores alimentan créditos fiscales; y el IRPF pierde las ganancias patrimoniales.

Martí y Pérez (2015) documentan que el saldo empeoró $6,4$ puntos solo en 2008. Por el lado del gasto actúan dos fuerzas. La primera es el aumento del desempleo: las prestaciones pasan de $1,5\%\to3\%$ del PIB. La segunda es la inercia del gasto primario comprometido en la fase alcista, que siguió creciendo en términos nominales en 2008-2009 mientras el PIB se contraía. El efecto denominador amplifica todas las ratios.

La respuesta fiscal española se inscribió en el Plan Europeo de Recuperación Económica (Comisión, noviembre de 2008), que recomendaba un estímulo coordinado de en torno al $1,5\%$ del PIB de la Unión. Las principales medidas fueron: la deducción de $400$ euros en el IRPF (RDL 2/2008); la bonificación del $100\%$ del Impuesto sobre el Patrimonio, que supuso su supresión efectiva (Ley 4/2008); el Fondo Estatal de Inversión Local, dotado con $8.000$ millones de euros (RDL 9/2008) y, su continuación, el Fondo Estatal para el Empleo y la Sostenibilidad Local, con $5.000$ millones (RDL 13/2009). A ello se sumaron ayudas sectoriales, como el Plan 2000E de automoción, y medidas de liquidez. Todo ello se conoce como Plan E.

El debate sobre este estímulo tiene dos niveles. El primero es su eficacia. Los fondos locales de inversión se criticaron por su baja calidad y por un multiplicador probablemente reducido: eran proyectos pequeños, intensivos en obra menor y sin continuidad. Su defensa se basa en que actuaron rápido en el sector más castigado por el paro, la construcción, cuando el multiplicador en recesión y con política monetaria restringida debía ser alto ([\authorfont{Ver Tema 3.A.40}])

El segundo nivel, más importante para este tema, es el error de diagnóstico. El estímulo se diseñó como si el deterioro fuese cíclico. Pero, como vimos, buena parte de él era estructural: los ingresos perdidos no iban a volver. Varias medidas del estímulo, como la deducción de 400 euros o la supresión del Impuesto sobre el Patrimonio, eran rebajas fiscales permanentes. En la práctica, España utilizó en 2008-2009 un margen fiscal que creía tener por el bajo nivel de deuda, cuando en realidad tenía un déficit estructural que se revelaría como tal en cuanto se estimó bien la brecha de producción.\footnote{Esta es la lectura del BdE y de la OCDE ex post. Una lectura alternativa sostiene que el estímulo fue razonable en el contexto europeo de 2008-2009, porque el error de diagnóstico fue generalizado y era difícil detectarlo en tiempo real.
}

\paragraph{2010-2013: la crisis soberana, el ajuste y el rescate bancario}



La crisis soberana de 2010-2013 impuso un ajuste concentrado y probablemente excesivo en su calendario, en un contexto de multiplicadores altos y sin prestamista de última instancia. Se combinó con el círculo vicioso entre bancos y soberano, que acabó con la asistencia del MEDE (41.333 millones) y un coste fiscal de la reestructuración bancaria que el Tribunal de Cuentas estima en unos 72.000 millones. En esos años los tres motores de la deuda (déficit primario, intereses y ajustes déficit-deuda) sumaron en la misma dirección.

a) El punto de inflexión de mayo de 2010

[DOC] El documento sitúa la apertura del procedimiento de déficit excesivo (PDE) en 2009 y un cambio radical de estrategia presupuestaria en 2010, con el objetivo de situar el déficit por debajo del 3% lo antes posible, algo que no se conseguiría hasta 2018. En nota añade que el procedimiento se cerró en junio de 2019.

La cronología exacta es la siguiente:
- el Consejo abrió el PDE a España en abril de 2009, con plazo de corrección en 2012;
- el plazo se amplió a 2013 (diciembre de 2009), a 2014 (julio de 2012) y a 2016 (junio de 2013);
- en julio de 2016 el Consejo declaró que España no había adoptado medidas eficaces, y en agosto de ese año decidió anular la multa;
- el PDE se cerró en junio de 2019, con los datos de 2018.

Las ampliaciones sucesivas de plazo son en sí mismas un dato relevante: el ajuste fue más lento de lo previsto porque la recesión se prolongó y los multiplicadores resultaron mayores de lo supuesto (ver c).

El giro se produce en mayo de 2010. Tras el primer programa griego y la creación del FEEF (9 de mayo de 2010), el RDL 8/2010 aprobó:
- una reducción media del 5% de las retribuciones del sector público;
- la congelación de las pensiones contributivas en 2011, salvo las mínimas;
- la supresión del cheque bebé;
- un recorte de la inversión pública.

Es la primera vez que el ajuste se impone desde fuera: la presión de los socios europeos y de los mercados obliga a abandonar el estímulo.

b) El ajuste: composición y cronología

[DOC] El documento enumera las medidas de gasto (reducción y congelación de salarios públicos, prestaciones por desempleo, inversión pública, educación, sanidad y ayuda oficial al desarrollo) y de ingreso:
- ingresos extraordinarios, con la amnistía fiscal de 2012 y la privatización parcial de AENA en 2015;
- subidas de impuestos: recargo de solidaridad en el IRPF en 2012-2014, subida del gravamen del ahorro, restablecimiento del Impuesto sobre el Patrimonio, límites a las deducciones del Impuesto sobre Sociedades y subidas de los tipos del IVA y de los impuestos especiales.

En nota, el documento cita la caída del límite de gasto no financiero del Estado de 182.439 millones de euros (2010) a 117.353 millones (2012).

Tres correcciones.
1. La venta de AENA no reduce el déficit. Como las privatizaciones de 1996-2000 (1.3.1), la salida a bolsa del 49% de AENA en febrero de 2015 es una operación financiera. No reduce el déficit, sino la deuda, a través del ajuste déficit-deuda. No debe presentarse como ingreso extraordinario en el sentido del saldo.
2. La amnistía fiscal fue anulada. La regularización del RDL 12/2012 fue declarada inconstitucional por la STC 73/2017, por haberse aprobado mediante decreto-ley en materia que afecta al deber de contribuir (art. 31 CE). La sentencia no afectó a las situaciones ya consolidadas. Recaudó en torno a 1.200 millones de euros, aproximadamente la mitad de lo previsto. El documento lo omite y es un dato fácil de preguntar.
3. La caída del techo de gasto no mide un recorte. La reducción del límite de gasto no financiero del Estado entre 2010 y 2012 no es comparable en términos homogéneos. Una parte sustancial se debe a la entrada en vigor del modelo de financiación autonómica de la Ley 22/2009, que amplió la cesión de tributos a las CCAA (IRPF al 50%, IVA al 50%, impuestos especiales al 58%) y redujo, en contrapartida, las transferencias del Estado que computan en su techo.\footnote{Modelo de financiación autonómica de 2009 (Ley 22/2009), cesión de tributos y fondos: Tema 4B-25 (ya redactado).
} Presentar el paso de 182.000 a 117.000 millones como recorte sobreestima el ajuste del Estado.

La secuencia normativa del ajuste que conviene retener es la siguiente:
- el IVA sube del 16% al 18% en julio de 2010 y al 21% en septiembre de 2012; el tipo reducido pasa del 7% al 8% y después al 10%;
- el RDL 20/2011 (diciembre de 2011) introduce el recargo temporal del IRPF y del IBI;
- el RDL 13/2011 restablece temporalmente el Impuesto sobre el Patrimonio;
- el RDL 20/2012 (julio de 2012) suprime la paga extraordinaria de diciembre de los empleados públicos, reduce la prestación por desempleo a partir del día 181 (del 60% al 50% de la base) y aprueba la subida del IVA;
- en pensiones, la Ley 27/2011 y la Ley 23/2013, con el índice de revalorización y el factor de sostenibilidad, se analizan en el Tema 24.\footnote{Reformas de pensiones de 2011 y 2013, IRP y factor de sostenibilidad: Tema 4B-24 (ya redactado).
}

c) El gran debate: ¿austeridad autodestructiva o ajuste inevitable?

Es el debate central de la década, y conviene presentarlo con sus dos posiciones.

La posición crítica se apoya en dos hallazgos:
- Blanchard y Leigh (2013), entonces en el FMI, mostraron que las previsiones de crecimiento de 2010-2011 infravaloraron sistemáticamente el impacto de la consolidación: los multiplicadores efectivos estaban por encima de 1, frente al ~0,5 supuesto.
- De Grauwe y Ji (2013) mostraron que las primas de riesgo de la periferia en 2010-2012 no se explicaban por los fundamentales fiscales, sino por el pánico del mercado en una unión monetaria sin prestamista de última instancia. Su comparación entre España y el Reino Unido es ilustrativa: con peores ratios de deuda y déficit, el Reino Unido pagaba mucho menos porque tenía su propio banco central.

De ahí la tesis de la austeridad autodestructiva: en recesión, con multiplicadores altos y sin política monetaria propia, la consolidación reduce el PIB más que el déficit, y la ratio de deuda puede aumentar. En España, la ratio de deuda siguió subiendo con fuerza en 2012-2014 pese al ajuste, lo que es coherente con esta tesis.

La posición defensiva, que es la de la Comisión, el BdE en su momento y una parte de la literatura, sostiene que la alternativa no era un ajuste más lento sino la pérdida completa del acceso al mercado. Con necesidades de financiación del Estado en torno al 20% del PIB anual y una prima de riesgo por encima de 600 puntos básicos, un ajuste más gradual no era financiable. La credibilidad del ajuste fue, además, condición política para que el BCE anunciase el OMT. Alesina, Favero y Giavazzi (2019) defienden que los ajustes por el lado del gasto son mucho menos recesivos que los basados en impuestos. En España predominaron las subidas de impuestos, lo que en su marco explicaría parte del coste.

La síntesis más aceptada hoy es que el ajuste era inevitable en su dirección pero probablemente excesivo en su concentración temporal (2010-2013), por la combinación de multiplicadores altos, recesión y ausencia de un prestamista de última instancia hasta julio de 2012. El giro del BCE, más que el ajuste, fue el factor decisivo para romper el círculo vicioso.\footnote{La hipótesis de la consolidación expansiva (canal de expectativas y de primas de riesgo) frente a la austeridad autodestructiva (multiplicador alto con tipos en su límite inferior e histéresis). La evidencia de Blanchard y Leigh (2013) y de Guajardo, Leigh y Pescatori (2014), en el Tema 3A-40. El OMT y la política monetaria no convencional, en los Temas 3A-36/37.
}

d) El rescate bancario y el círculo vicioso entre bancos y soberano

[DOC] En junio de 2012 España solicitó asistencia del MEDE para recapitalizar su sistema bancario. El Memorando de Entendimiento de julio de 2012 abrió una línea de hasta 100.000 millones de euros, de los que España utilizó 41.333. La condicionalidad era de tres tipos: sobre las entidades, horizontal sobre el sistema financiero y recordatorios de compromisos fiscales y estructurales. España devolvió anticipadamente 17.612 millones en 2014-2018, y los 23.721 restantes vencen hasta 2027. La prima de riesgo alcanzó los 638 puntos básicos en julio de 2012, con el bono a diez años por encima del 7,5%, y Moody's rebajó la calificación en nueve escalones en menos de dos años (de Aaa a Baa3 en junio de 2012).

Los datos del documento son correctos. Lo que falta es la explicación de por qué 2012 es el momento crítico: el círculo vicioso entre bancos y soberano. Las dudas sobre la solvencia de los bancos, en particular de las cajas expuestas al sector inmobiliario, elevaban el riesgo del soberano, que era su garante implícito. A su vez, las dudas sobre el soberano deterioraban el valor de la deuda pública que los bancos tenían en balance. Los préstamos a largo plazo del BCE de diciembre de 2011 y febrero de 2012 (LTRO a tres años) reforzaron el vínculo: los bancos españoles usaron esa liquidez para comprar deuda del Estado. La proporción de deuda en manos de no residentes cayó del ~50% a en torno al 33% a mediados de 2012, según el BdE. La deuda se renacionalizó en los balances bancarios. Es el diagnóstico que llevó a la Unión Bancaria (2012-2014).\footnote{FEEF, MEDE, Unión Bancaria, MUS y MUR: Tema 3B-45.
}

Sobre el coste fiscal, el Tribunal de Cuentas actualizó en marzo de 2025 el coste de la reestructuración bancaria a 71.833 millones de euros (datos a 1 de enero de 2022). De ellos, 50.622 millones corresponden al FROB y 21.273 millones a las propias entidades de crédito a través del Fondo de Garantía de Depósitos. El coste no es definitivo. En contabilidad nacional, las aportaciones de capital que no se esperaba recuperar se registraron como transferencias de capital, con impacto en el déficit (en torno a 3,7 puntos de PIB en 2012). El resto se contabilizó como activos financieros, con impacto solo en la deuda.

e) Los mecanismos de liquidez para las administraciones territoriales

En 2012 se crean el Fondo para la Financiación de los Pagos a Proveedores (RDL 4/2012 y RDL 7/2012) y el Fondo de Liquidez Autonómico (RDL 21/2012), después integrados en el Fondo de Financiación a CCAA (RDL 17/2014). Ya se analizaron en el Tema 25 y en el inventario de 1.1.3.\footnote{FFPP, FLA y FFCCAA, volumen de recursos (126.895 millones en 2012-2024) y condonación de 2025: Tema 4B-25.
} Aquí importa su efecto sobre la deuda consolidada. El plan de pago a proveedores convirtió deuda comercial, que no computa como deuda PDE, en deuda financiera del Estado, que sí computa. Aumentó la deuda de las AAPP sin aumentar el déficit del año, salvo en la parte de las facturas no registradas en su momento. Ese aflorar de facturas en el cajón confirma lo apuntado en 1.3.3e: parte del déficit de los años anteriores no se había registrado.

f) La reforma del marco institucional

[DOC] El documento recoge la reforma del art. 135 CE (septiembre de 2011), la LOEPSF (2012), con sus tres reglas (equilibrio estructural, límite de deuda del 60% y regla de gasto), y la creación de la AIReF.

Conforme a la remisión acordada, aquí solo fijo lo imprescindible. La reforma constitucional de 27 de septiembre de 2011 introduce el principio de estabilidad presupuestaria y la prioridad absoluta del pago de intereses y capital de la deuda. La LO 2/2012 desarrolla esa reforma y traslada el six-pack y el TECG. La LO 6/2013 crea la AIReF, que empieza a funcionar en 2014. Una precisión sobre el documento: la regla de gasto de la LOEPSF no excluye formalmente las pensiones. Lo que ocurre es que no se aplicaba a la Seguridad Social, sino al Estado, las CCAA y las entidades locales, y por eso el gasto en pensiones quedaba fuera en la práctica.\footnote{Reglas de la LOEPSF (arts. 11 a 13), su relación con el TECG y el six-pack, las cláusulas de escape y su suspensión en 2020-2023: Tema 3B-45; articulación con la LGP y el ciclo presupuestario: Tema 4B-21.
}

g) La mecánica de la deuda, 2008-2014

El BdE (Boletín Económico, julio de 2015) descompone el aumento acumulado de la deuda entre 2008 y 2014 así:
- déficit primario, unos 37 puntos de PIB;
- carga de intereses, unos 13 puntos, parcialmente compensados por un crecimiento nominal prácticamente nulo en el periodo;
- ajustes déficit-deuda, unos 6 puntos netos. Incluyen la parte financiera de las ayudas bancarias, la contribución española a los programas europeos (préstamos bilaterales a Grecia y participación en el FEEF), el capital desembolsado en el MEDE y la acumulación de depósitos del Tesoro, con reversiones parciales.

A diferencia de 1991-1995 (dominado por el diferencial) y de 1975-1985 (dominado por el déficit primario con diferencial favorable), aquí los tres componentes suman en la misma dirección. Es el único episodio de la serie en que eso ocurre.

(OJO sobre las cifras de deuda: la ratio de 2014 se publicó en 2015 como 97,7%, pasó a 100,7% en revisiones posteriores y hoy es de ~105%. Las diferencias se deben sobre todo a la reclasificación de la Sareb en las AAPP con efectos desde 2012, decidida por Eurostat en 2021, y a revisiones del PIB. Es el mismo problema de perímetro de 1.1.1, y conviene citar la serie vigente, con la advertencia si procede).\footnote{Reclasificación de la Sareb por Eurostat (2021) y criterios SEC de control y de productor de mercado: Tema 4B-23.
}

\paragraph{2014-2019: la recuperación y la consolidación incompleta}

La recuperación de 2014-2019 redujo el déficit total gracias al ciclo y al BCE, no gracias a un esfuerzo estructural. La fragmentación política y las decisiones de gasto e ingreso de la fase alcista volvieron a hacer procíclica la política fiscal en la expansión, igual que en 1986-1990 y en 2005-2007. España llega a 2020 con una deuda del 98% del PIB, un déficit estructural del orden del 3% y una dependencia creciente del Eurosistema como tenedor de su deuda. Ese es el margen, muy estrecho, con el que afrontará la pandemia.

a) Lo que dice el documento

[DOC] Según el documento, las reformas estructurales permitieron que España volviera a crecer desde 2013 y se convirtiera en la gran economía occidental con mayor crecimiento. Cita la reforma de pensiones de 2011, la laboral de 2012, la Ley de Garantía de la Unidad de Mercado (2013) y la del sector eléctrico (2013). Desde 2015 el déficit se reduce hasta situarse por debajo del 3% en 2018, tras más de diez años por encima, y la deuda baja del 105% en 2014 al 98% en 2019.

Los datos son correctos con la serie actual: un déficit del 2,6% en 2018 y del 3,1% en 2019, y una deuda del 105,1% en 2014 y del 98,2% en 2019. La atribución causal a las reformas, en cambio, necesita matices.

b) ¿Qué redujo el déficit? Ciclo, intereses y (poco) esfuerzo estructural

Entre 2014 y 2019 el déficit total baja unos 3 puntos de PIB, del ~6% al ~3%. Hay tres componentes:
1. El ciclo. La economía crece por encima del 2,5% anual en 2015-2018 y se crean más de 2,5 millones de empleos. La brecha de producción pasa de muy negativa a positiva hacia 2018-2019, y la mejora cíclica del saldo explica la mayor parte de la reducción.
2. Los intereses. Tras el máximo de 2013-2014, en torno al 3,5% del PIB, la carga de intereses baja al 2,3% en 2019, según el BdE. El motor es la política monetaria: el OMT (2012) eliminó el riesgo de redenominación, y el programa de compra de deuda pública del BCE (PSPP), desde marzo de 2015, llevó los tipos a mínimos. A finales de 2019 el Eurosistema tenía el 18,7% de la deuda pública española.\footnote{El programa de compras del BCE (APP/PSPP y PEPP), la regla de reparto por la clave de capital y sus efectos sobre las primas: Temas 3A-36/37; su efecto sobre el Tesoro, en 3.5 de este tema.
}
3. El esfuerzo estructural primario. Fue prácticamente nulo o negativo a partir de 2015. Hubo rebajas del IRPF en 2015-2016 y 2018, revalorizaciones de las pensiones por encima del índice de revalorización en 2018-2019 y subidas salariales públicas. Las estimaciones de la Comisión y de la AIReF sitúan el déficit estructural en torno al 3% del PIB en 2019, un nivel similar al de 2014-2015.

La conclusión es incómoda: el ajuste estructural de la década se concentró en 2010-2013, en plena recesión y cuando más costaba. En la recuperación no se aprovechó para completarlo. La política fiscal volvió a ser procíclica en la fase expansiva, igual que en 1986-1990 y en 2005-2007. Es el diagnóstico compartido por el FMI (consultas del artículo IV de 2018 y 2019), la Comisión (Semestre Europeo) y la AIReF.

c) Por qué no se completó el ajuste: la economía política

La explicación más directa es política. Entre 2015 y 2019 España atraviesa el periodo de mayor fragmentación parlamentaria de su democracia:
- dos elecciones generales en seis meses (2015-2016) y otras dos en 2019;
- una moción de censura (2018);
- presupuestos prorrogados: los de 2018 rigieron en 2019 y 2020.

Aquí la hipótesis de Roubini y Sachs (1989), que no encajaba con las mayorías absolutas de los ochenta (1.2.3), sí encuentra respaldo: la fragmentación dificulta el ajuste y favorece las concesiones de gasto a los socios de investidura. El caso más claro es la decisión de revalorizar las pensiones con el IPC en 2018-2019, abandonando el índice de revalorización de 2013, que el Tema 24 analiza como el principio del fin del factor de sostenibilidad.

d) La mecánica de la deuda, 2014-2019: el crecimiento nominal como protagonista

Según el BdE, la reducción de la deuda en 2019 (2,1 puntos) se descompone en tres partes:
- el déficit total sumó 2,8 puntos (0,5 de déficit primario y 2,3 de intereses);
- el crecimiento nominal restó 3,4;
- los ajustes déficit-deuda restaron 1,6.

Es el patrón de todo 2015-2019. La deuda baja muy despacio, unos 7 puntos en cinco años, y lo hace sobre todo porque el crecimiento nominal compensa un déficit total que todavía aumenta la deuda. Con un saldo primario casi equilibrado y un diferencial negativo entre tipo y crecimiento, la deuda baja; pero a este ritmo habrían hecho falta décadas para volver al 60%. Es el punto de partida de la pandemia.

e) El Tesoro llega a 2020 con una deuda más larga y un inversor distinto

A finales de 2019 la vida media de la deuda era de unos 7,5 años, el 93,5% era a largo plazo y los no residentes tenían el 49% (frente al 33% de mediados de 2012). El dato relevante para la tercera parte del tema es que un inversor nuevo, el Eurosistema, se ha convertido en el tenedor individual más importante. Esa dependencia tiene consecuencias. Cuando el BCE deje de reinvertir los vencimientos de sus programas (2023-2024 para el APP y 2025 para el PEPP), la demanda que el Tesoro tendrá que sustituir en el mercado será muy grande. Lo desarrollo en 3.5.


\subsubsection{Pandemia, inflación y nuevo marco fiscal (2020-2026)}
El último periodo que nos queda por recorres son seis años en los que la hacienda española sufre tres shocks de naturaleza distinta. El primero, en 2020, es un choque de oferta y de demanda sin precedentes en tiempos de paz, que hunde el PIB, dispara el déficit por encima del $10\%$ y lleva la deuda a su máximo histórico. El segundo, en 2022-2023, es un choque inflacionario de origen energético. Cambia por completo la aritmética de la deuda: primero la alivia a través del denominador y después la encarece a través de los tipos de interés. El tercero es un conjunto de shocks que se superponen en 2024-2026: la DANA de octubre de 2024, las borrascas de principios de 2026 y la guerra en Oriente Medio. 

El relato habitual del periodo es que, tras la pandemia, España ha reducido su deuda en casi 20 puntos. Como veremos, es cierto en la cifra y engañoso en la causa. Según la descomposición del BdE, entre 2021 y 2024 el déficit (primario más intereses) añadió unos $18$ puntos de PIB a la ratio de deuda y el crecimiento del PIB nominal restó unos $37$. La deuda bajó por el denominador, con inflación y crecimiento, no porque las administraciones generasen excedentes. No obstante, la hacienda española sale de este periodo con un saldo primario cercano al equilibrio y una deuda ligeramente por encima del $100\%$ del PIB: una posición claramente mejor que en 2020, pero gracias a un entorno de crecimiento nominal que no se va a repetir, y no ha corregido el componente estructural del déficit que arrastra desde 2008. 

En definitiva, que ambas lecturas son compatibles, tanto la positiva como la negativa. La posición de partida es mejor que en 2020 en el saldo primario y en la estructura de la deuda. Pero la sostenibilidad de la trayectoria depende de nuevo de que el diferencial entre el tipo de interés y el crecimiento siga siendo favorable, algo que la experiencia española de 1991-1995 y 2010-2013 enseña que no puede darse por descontado.

\textcolor{red}{Entre 2020 y 2026 la hacienda española sufre tres shocks de naturaleza distinta.

La inflación de 2022-2023 invirtió la aritmética de la deuda. Entre 2021 y 2024 la ratio bajó unos 17,5 puntos, pero lo hizo porque el crecimiento nominal restó unos 37 puntos mientras el déficit seguía añadiendo unos 18. La recaudación alcanzó máximos históricos, en parte por la progresividad en frío, y el Plan de Recuperación financió inversión sin aumentar la deuda española en la parte de subvenciones.

En 2025-2026, España alcanza un saldo primario cercano al equilibrio y una deuda en torno al 100% del PIB, en el primer año de un nuevo marco fiscal europeo basado en una senda de gasto neto. Pero lo hace con presupuestos prorrogados desde 2023, con el retorno de los shocks (la DANA, las borrascas y la guerra en Oriente Medio) y con un BCE que sube tipos por primera vez en tres años. El gasto neto crece muy por encima de lo comprometido. La lección del bloque se repite: los periodos favorables han reducido la deuda a través del denominador, pero no han corregido el déficit estructural, y la sostenibilidad depende de nuevo de condiciones que España no controla.}


Además, 2026 es el año en que las condiciones que hicieron posible esa reducción empiezan a revertirse. El conflicto de Oriente Medio iniciado el 28 de febrero de 2026 ha devuelto la inflación al 4,3% en agosto y ha llevado al BCE a subir los tipos en junio y septiembre, hasta el 2,50% en la facilidad de depósito. La AIReF, por su parte, prevé un crecimiento del gasto neto del 6,4% en 2026, frente al 3,5% comprometido con Bruselas. El apartado situación actual del documento base describe el mundo de 2022, y en un examen de 2026 sería lo primero que un tribunal detectaría. Muchas de estas cifras de 2026 son muy recientes y pueden revisarse antes de tu examen, así que las he marcado para verificarlas de nuevo.




\paragraph{2020-2021: la pandemia y la respuesta fiscal}
El primer periodo que veremos será la pandemia. En 2020 el déficit volvió a situarse por encima del $10\%$ del PIB y la deuda superó el $120\%$: un máximo histórico.

La revisión de la contabilidad nacional publicada por el INE en septiembre de 2024 elevó el nivel del PIB nominal, y con ella la caída real de 2020 queda en torno al $10,9\%$. El BdE sitúa hoy la ratio de deuda de 2020 en el $119,3\%$, no por encima del $120\%$. El déficit de 2020 se publicó en marzo de 2021 en el $11\%$ del PIB, incluida la reclasificación de la Sareb, y con la serie vigente queda cerca del $10\%$.\footnote{La diferencia entre las cifras de hoy y las que se publicaron en su momento es otro ejemplo de los desajustes: las cifras fiscales dependen del perímetro, en este caso la Sareb, y del denominador, en este caso la revisión del PIB.
} La lectura con la descomposición del BdE es especialmente clara para 2020. La deuda aumentó un $21,6\%$. El déficit primario aportó $7,7\%$, los intereses $2,2\%$, los ajustes déficit-deuda $0,9\%$ y la caída del PIB nominal $10,8\%$. Por primera vez en la serie histórica, el efecto denominador explica prácticamente la mitad del aumento de la ratio en un solo año: la imagen inversa de lo que ocurriría en 2021-2024.

¿A qué se deben estas variaciones? Evidentemente a la respuesta a la Pandemia, que fue clara: mucha liquidez y un gasto directo relativamente contenido. Según la actualización del Programa de Estabilidad 2021-2024, se movilizaron unos $231.000$ millones de euros en 2020-2021, el $20,6\%$ del PIB de 2020. Algo más de $73.000$ millones fueron ayudas directas y unos $158.000$ millones, financiación movilizada mediante avales, garantías y moratorias.

La flexibilización de los expedientes de regulación temporal de empleo exoneró de cotizaciones a las empresas a cambio de mantener el empleo y se prolongó, con modificaciones, hasta marzo de 2022. Se crearon dos fondos de recapitalización: uno gestionado por la SEPI para empresas estratégicas, con $10.000$ millones de los que se utilizaron $3.255,8$, y otro gestionado por COFIDES para medianas empresas, con $1.000$ millones de los que se utilizaron $779$. Hubo además un fondo de ayudas directas de $7.000$ millones, gestionado por las CCAA, del que se utilizaron algo más de $6.000$, y se aprobó el Ingreso Mínimo Vital. Entre las medidas sin impacto inmediato en el gasto destacan las líneas de avales, por unos $143.000$ millones y en su mayoría gestionadas por el ICO, las moratorias de créditos y los aplazamientos tributarios. Según el ICO, a junio de 2022 se habían desplegado avales por $107.187$ millones, que movilizaron $140.737$ millones de financiación en cerca de $1,2$ millones de operaciones.

En definitiva, la respuesta fiscal española fue desproporcionadamente de liquidez y comparativamente contenida en gasto directo. El Fiscal Monitor del FMI de 2021 situaba el apoyo presupuestario directo (por encima de la línea) de España claramente por debajo del de Alemania, el Reino Unido o Estados Unidos, y el apoyo de liquidez (avales, préstamos, capital) en niveles comparables o superiores a los de sus pares. La razón es evidente, España entró en la pandemia con una deuda cercana al $98\%$ y un déficit estructural del $3\%$, y su espacio fiscal era menor. Por eso optó por instrumentos que no computan como déficit en el momento, sino como pasivos contingentes.\footnote{
Como sabemos ([\authorfont{Ver Tema 4.B.23}]), los avales no generan gasto salvo que se ejecuten, aunque se dotó una provisión por pérdida esperada con impacto en el déficit de 2020. La experiencia posterior, con las ejecuciones de avales y el marco de reestructuración de créditos avalados del RDL 5/2021, muestra que la pérdida efectiva fue moderada en relación con el volumen avalado. Esto avala ex post el diseño, aunque la cifra definitiva está pendiente de verificar.
}

No obstante, la diferencia decisiva frente a 2010-2012 fue la respuesta europea. En marzo de 2020 se activó la cláusula general de salvaguardia del Pacto de Estabilidad, que suspendía en la práctica los requisitos de ajuste y se mantuvo hasta finales de 2023. En España, el Congreso apreció la existencia de una situación de emergencia extraordinaria, lo que suspendía a su vez las reglas de la LOEPSF. El instrumento SURE financió los ERTE y mecanismos similares. España fue uno de sus principales beneficiarios, con unos 21.300 millones en préstamos. En julio de 2020 se acordó el instrumento Next Generation EU.

Po otro lado, y en términos financieros, la pieza más importante fue la del BCE. Cuando la presidenta del BCE declaró el 12 de marzo que su función no era cerrar diferenciales, la rentabilidad del bono español a diez años subió casi 100 puntos básicos en una semana. El anuncio del programa de compras de emergencia frente a la pandemia (PEPP) el 18 de marzo y la flexibilización del colateral en abril devolvieron la rentabilidad por debajo de su nivel previo a la pandemia en julio de 2020. Lo que refleja la absoluta dependencia de los tipos de la deuda respecto de las decisiones de política monetaria. Lección que la crisis de 2010-2012 había enseñado a un coste muy alto y que en 2020 se aplicó desde el primer momento.

Dicho en los términos de DE GRAUWE, en 2020 los países de la zona del euro sí tuvieron un prestamista de última instancia de facto. Sin él, la respuesta fiscal española habría sido mucho más limitada. Es también una advertencia para el presente: la dependencia se ha vuelto a manifestar, en sentido contrario, desde que el BCE dejó de comprar deuda.

d) El debate de 2020-2021: ¿había que haber hecho más?

Hubo un debate genuino sobre si la respuesta española fue suficiente. Una posición, sostenida por parte de la literatura y por organismos como el FMI en 2020, sostenía que en una crisis de esta naturaleza el coste de hacer demasiado poco (quiebras de empresas viables, histéresis en el empleo) era mayor que el de hacer demasiado. En consecuencia, España debió aprobar más ayudas directas a empresas y antes, sobre todo en los sectores del turismo y la hostelería. El fondo de ayudas directas de $7.000$ millones no llegó hasta marzo de 2021 (RDL 5/2021), un año después del inicio de la crisis.

La posición contraria recordaba que la recuperación del empleo fue relativamente rápida gracias a los ERTE, que el tejido productivo no sufrió la ola de quiebras temida y que, dada la posición de partida, un gasto directo mayor habría elevado todavía más la deuda. El balance que hoy parece más razonable es que el diseño de los ERTE fue un acierto, porque evitó la destrucción masiva de empleo de 2008-2013. El apoyo directo a empresas, en cambio, llegó tarde y se ejecutó por debajo de lo dotado, como muestra la infrautilización de los fondos de la SEPI y de COFIDES que recoge el propio documento.

\paragraph{2022-2024: inflación, Plan de Recuperación y reducción de la deuda por el denominador}
El segundo periodo viene caracterizado por la inflación. Sin embargo, antes de antrar a valorar su impacto, veamos las principales medidas de Política Fiscal y Monetaria que se adoptaron.

Las medidas que se adoptaron como respuesta a la crisis energética de 2022-2023 fue diversa y el coste conjunto en 2022-2023 se sitúa en torno a los $45.000$ millones de euros según distintas estimaciones, lo que queda pendiente de verificar con la AIReF.\footnote{Incluyó la bonificación de 20 céntimos por litro de combustible en 2022, las rebajas del IVA y del impuesto especial de la electricidad y del gas, las rebajas del IVA de los alimentos básicos, las ayudas al transporte público y el llamado mecanismo ibérico de limitación del precio del gas para la generación eléctrica, este último sin coste presupuestario directo.

La literatura, con la OCDE, el FMI y el BdE a la cabeza, criticó que la mayoría de estas medidas no estaban focalizadas: beneficiaban a todos los consumidores, con independencia de su renta, y mantenían la señal de precios de la energía en lugar de corregirla. Una transferencia de renta a los hogares vulnerables habría tenido menor coste y habría preservado el incentivo al ahorro energético. La defensa de las medidas generales se basaba en la rapidez de aplicación y en su efecto directo sobre el IPC, que a su vez moderaba la indexación de pensiones y salarios.
} Además, por el lado de los ingresos, la Ley 38/2022 creó gravámenes temporales sobre las empresas energéticas y las entidades de crédito, así como el Impuesto Temporal de Solidaridad de las Grandes Fortunas. El gravamen bancario se transformó a finales de 2024 en un impuesto sobre el margen de intereses y comisiones de las entidades financieras (Ley 7/2024), mientras que el energético decayó. 

Esa recaudación merece ser comentada. Los ingresos tributarios del Estado superaron los $325.000$ millones de euros en 2025, con un crecimiento del $10,4\%$. La presión fiscal alcanzó máximos históricos, en torno al $37-38\%$ del PIB en 2024-2025. Es la contracara del efecto denominador: la inflación y la progresividad en frío del IRPF, que no se ha deflactado con carácter general en el ámbito estatal, elevaron la recaudación por encima del crecimiento de las bases. Una parte de este aumento es estructural (el empleo récord, más de $21$ millones de afiliados, y los cambios normativos), pero otra parte es transitoria o depende del nivel de inflación. Estimar cuál es cuál es de nuevo el problema central para leer el saldo estructural.

Por otro lado, el Mecanismo de Recuperación y Resiliencia, núcleo de Next Generation EU, asignó a España el mayor volumen de subvenciones tras Italia. Tras la adenda de 2023, el Plan de Recuperación español alcanzaba unos $163.000$ millones de euros, aproximadamente la mitad en transferencias no reembolsables y la otra mitad en préstamos. El Gobierno solicitó finalmente solo una fracción de los préstamos disponibles, en torno a $21.500$ millones, al considerar que el coste de financiarse directamente en el mercado era similar o inferior. A 31 de agosto de 2026, España había recibido unos $78.000$ millones en seis desembolsos. El sexto, de $6.234$ millones, llegó el 11 de agosto de 2026.\footnote{Quedaban pendientes unos 25.862 millones del último tramo, condicionados al cumplimiento de 148 hitos y objetivos, con plazo para solicitar el desembolso hasta el 30 de septiembre de 2026.
} Por lo que a nuestro tema respecta, las subvenciones no generan deuda española. La deuda la emite la Comisión en nombre de la Unión y se amortizará con cargo al presupuesto europeo a partir de 2028. En contabilidad nacional, las subvenciones se registran como ingreso cuando se realiza el gasto que financian, de modo que su efecto sobre el déficit español es en principio neutro: financian gasto adicional sin aumentar el déficit. Los préstamos, en cambio, computan como deuda de las AAPP españolas igual que cualquier otra financiación del Tesoro.

Esto tiene dos implicaciones. La primera es que el Plan permitió mantener niveles de inversión pública que de otro modo habrían exigido más déficit. La segunda es que el fin del Plan en 2026 plantea el llamado precipicio fiscal: si la inversión financiada con fondos europeos debe mantenerse, tendrá que pasar a financiarse con recursos propios, con impacto en el déficit a partir de 2027.

Por otro lado, entre julio de 2022 y septiembre de 2023 el BCE elevó la facilidad de depósito del $−0,5\%$ al $4\%$, el ciclo de subidas más rápido de su historia. Dejó además de reinvertir los vencimientos del programa de compra de activos (APP) desde julio de 2023 y del PEPP a finales de 2024. Para el Tesoro, el efecto fue doble: el coste de las nuevas emisiones subió de forma abrupta, y el Eurosistema, que había sido el principal comprador neto de deuda española en 2015-2022, pasó a reducir sus tenencias. El Tesoro tuvo que colocar en el mercado no solo sus necesidades netas, sino también los vencimientos que antes absorbía el banco central: la ratio de no residentes subió al $44,9\%$ en 2024. La prima de riesgo española se mantuvo en niveles moderados, e incluso por debajo de la francesa en algunos momentos de 2024-2025, lo que sugiere que los mercados valoraban positivamente la evolución del saldo primario y del crecimiento.\footnote{El rápido aumento de los tipos desde 2022 puso en entredicho la hipótesis, cada vez más aceptada antes de la pandemia, de que cambios estructurales en las economías avanzadas (demografía, exceso de ahorro, caída de la productividad) mantendrían los tipos reales bajos o negativos de forma persistente, lo que haría sostenibles niveles de deuda más altos. Cita a GALESI, NUÑO y THOMAS (2017).

La formulación más influyente de esa tesis es la de BLANCHARD (2019): con $r < g$, la deuda pública puede tener un coste fiscal nulo y un coste en bienestar menor del supuesto. El propio BLANCHARD advertía de que la condición podía invertirse, y la subida de tipos de 2022-2023 la puso a prueba. La evidencia del periodo es matizada. Los tipos nominales subieron, pero la inflación subió más, de modo que el diferencial entre tipo de interés y crecimiento siguió siendo favorable en España durante 2022-2024. Lo que sí se ha demostrado es que la condición puede cambiar rápidamente y que depende de la política monetaria. El debate sigue abierto sobre si el tipo de interés natural ha vuelto a los niveles previos a 2008 o si se mantiene bajo, con una inflación transitoria que ha alterado temporalmente los tipos observados.
}

Con todo lo expuesto, sorprenderá la evolución de la deuda. La ratio de deuda pasa del $119,3\%$ en 2020 al $115,7\%$ en 2021, al $109,5\%$ en 2022, al $105,1\%$ en 2023 y al $101,8\%$ en 2024. En esos cuatro años se redujo en unos $17,5$ puntos. La suma de los déficits primarios aportó unos $8,6$ puntos al alza ($4,5$ en 2021, $2,3$ en 2022, $1,1$ en 2023 y $0,7$ en 2024) y los intereses, unos $9,2$ (entre 2,1 y 2,4 puntos cada año). Los ajustes déficit-deuda añadieron en conjunto algo menos de $2$ puntos. ¿Cómo es posible, entonces, que experimentáramos semejante descenso?

Como es de esperar, el crecimiento del PIB nominal restó unos $37$ puntos: $10,3$ en 2021, $11,6$ en 2022, $9,1$ en 2023 y $6,2$ en 2024. La conclusión es inequívoca. Sin el crecimiento nominal, la deuda habría aumentado unos $20$ puntos en ese periodo, porque las administraciones siguieron teniendo déficit todos los años. La mejora del saldo primario es real, de un déficit del $7,7\%$ a uno del $0,7\%$ en cuatro años. Pero un déficit primario casi nulo solo estabiliza la deuda si el crecimiento nominal supera al tipo de interés, que es precisamente la condición que empieza a debilitarse en 2026.

Ahora bien, la inflación afecta a las finanzas públicas de forma asimétrica en el tiempo. A corto plazo, la inflación favorece las cuentas públicas por tres vías. La primera es que los ingresos se ajustan automáticamente a los precios: el IVA recauda sobre bases nominales y el IRPF, que no se deflacta, sufre la llamada progresividad en frío, porque los contribuyentes saltan de tramo sin haber ganado poder adquisitivo. La segunda es que buena parte del gasto se ajusta con retraso. La tercera es que el tipo de interés implícito de la deuda apenas se mueve, porque la vida media es larga (unos 8 años). El resultado es que la ratio de deuda baja con fuerza por el denominador mientras la carga de intereses sigue siendo baja.

No obstante, a medio plazo el efecto se invierte cuando el gasto indexado alcanza a la inflación: las pensiones se revalorizaron un $8,5\%$ en 2023 por aplicación de la Ley 21/2021, y los salarios públicos y la contratación pública se actualizan. Al mismo tiempo, la deuda se va refinanciando a los nuevos tipos. Por tanto, la reducción de la deuda se sustentaba principalmente en el crecimiento del PIB nominal, con una contribución muy notable del deflactor. Una vez agotado ese impulso, las proyecciones de la AIReF dibujaban una dinámica desfavorable a políticas constantes, agravada por la subida de tipos. 

\paragraph{2025-2026: la situación actual}
Como anticipimamos, España cerró 2025 con un déficit del $2,2\%$ del PIB sin los gastos extraordinarios de la DANA de Valencia, o del $2,4\%$ incluyéndolos, y con una deuda del $100,7\%$, la ratio más baja desde el inicio de la pandemia. 

El Estado pagó unos $34.900$ millones en intereses y el conjunto de las AAPP algo más de $40.000$, en torno al $2,4\%$ del PIB. El saldo primario estuvo, por tanto, muy cerca del equilibrio, la primera vez desde 2007 que se da esta situación.

El año 2025 es también el primero de aplicación del nuevo marco de gobernanza económica de la Unión (Reglamento (Unión) 2024/1263), que sustituye los objetivos de saldo estructural por una senda plurianual de gasto primario neto negociada con cada Estado miembro.\foonote{Contenido del nuevo marco de gobernanza económica (Reglamento (Unión) 2024/1263 y Reglamento (Unión) 2024/1264 de la vertiente correctiva, con el análisis de sostenibilidad de la deuda, la trayectoria técnica, la cuenta de control y las cláusulas de salvaguardia general y nacional): Tema 3B-45.
}El Plan Fiscal y Estructural de Medio Plazo de España, aprobado por el Consejo el 21 de enero de 2025, fija un periodo de ajuste de siete años (2025-2031), ampliado en contrapartida por compromisos de reformas e inversiones. Establece un crecimiento máximo del gasto neto del 3,7% en 2025, del 3,5% en 2026, del 3,2% en 2027 y del 3% en 2028. Según el Plan, la deuda bajaría del 102,5% del PIB en 2024 al 90,6% en 2031 y al 76,8% en 2041.

La traslación al ordenamiento interno está pendiente. La LOEPSF no se ha adaptado al nuevo marco, y la regla de gasto nacional, el objetivo de estabilidad estructural y el límite de deuda de la ley orgánica siguen formalmente vigentes pero conviven con una regla europea de naturaleza distinta. A ello se añade una anomalía institucional sin precedentes en la democracia: España funciona desde 2024 con los Presupuestos Generales del Estado de 2023 prorrogados, porque ni en 2024, ni en 2025, ni en 2026 se aprobaron unos nuevos. Las sendas de estabilidad presentadas en 2024 y 2025 no obtuvieron el respaldo parlamentario necesario. La prórroga presupuestaria limita el margen de actuación discrecional, lo que paradójicamente ha contribuido a contener el gasto en algunos capítulos. Pero impide a la vez aplicar una estrategia fiscal de medio plazo y obliga a canalizar las nuevas medidas por decreto-ley y créditos extraordinarios.

El retorno de los shocks fue 2026. El 28 de febrero de 2026 comenzó la operación militar conjunta de Estados Unidos e Israel contra el régimen iraní, seguida de la respuesta de Irán y de un bloqueo del estrecho de Ormuz que disparó los precios del petróleo y del gas. El Gobierno aprobó el 20 de marzo de 2026 un Plan de Respuesta que movilizaba más de $5.000$ millones de euros, y la inflación española alcanzó el $4,3\%$ en agosto de 2026, la tasa más alta en más de tres años.\footnote{Incluía rebajas escalonadas del impuesto sobre hidrocarburos, descuentos en el bono social eléctrico, ayudas a los transportistas y al sector primario y la supresión progresiva del impuesto sobre la producción de energía eléctrica. Ayer mismo, el 29 de septiembre de 2026, el Real Decreto-ley 25/2026 prorrogó buena parte de estas medidas hasta el 31 de diciembre, incluida una rebaja de 20 céntimos por litro en el impuesto sobre hidrocarburos en octubre, con reducciones condicionadas en noviembre y diciembre.
} 

No obstante, la consecuencia monetaria es la más relevante. El BCE, que había llevado la facilidad de depósito al $2\%$ en junio de 2025, la subió al $2,25\%$ el 11 de junio de 2026 (la primera subida en tres años) y al $2,50\%$ el 10 de septiembre de 2026. Justificó la decisión porque el conflicto seguía generando presiones inflacionistas y la inflación se mantendría claramente por encima del objetivo durante un periodo prolongado. Sus proyecciones sitúan la inflación en el $3\%$ en 2026 y en el 2,5\% en 2027. Además, las carteras del APP y del PEPP siguen reduciéndose al no reinvertirse los vencimientos. Para España esto significa que \textbf{el coste marginal de la deuda vuelve a subir justo cuando se refinancia una parte creciente de la deuda emitida en la etapa de tipos cero}.

En su informe de julio de 2026, la AIReF mantenía su previsión de déficit para 2026 en el 2,6\% del PIB. Precisaba que, sin las medidas de emergencia (DANA, borrascas y conflicto de Oriente Medio), el déficit se situaría en el 1,8\%. Rebajaba su previsión de deuda al 99,3\% y elevaba la de crecimiento al 2,5\%. Pero advertía de que el crecimiento del gasto neto previsto para 2026 subía del 5,8\% al 6,4\%, muy por encima del 3,5\% comprometido en el Plan Fiscal, y reclamaba un ajuste del 0,6\% del PIB para no incumplir las reglas nacionales. Una parte de esa desviación puede quedar justificada por el carácter excepcional de las medidas de emergencia. La valoración final dependerá de cómo las trate la Comisión en el Semestre Europeo y de la cuenta de control del Plan, lo que queda pendiente de verificar.

A las presiones inmediatas se suman otras de medio plazo que se tratan en otros temas pero que condicionan la lectura de la situación actual. Una es el gasto en pensiones, cuyo crecimiento por el envejecimiento se analizó en el Tema 24 con la cláusula de salvaguarda y el informe de la AIReF de mayo de 2026. Otra es el compromiso de gasto en defensa, que España elevó al 2\% del PIB en 2025 mediante un plan de más de 10.000 millones de euros sin recurrir a la cláusula nacional de salvaguardia prevista en el marco europeo. La tercera es la condonación de deuda autonómica: el proyecto de ley que prevé que el Estado asuma 83.252 millones de deuda de las CCAA, enviado al Congreso en diciembre de 2025, superó el debate de totalidad en julio de 2026 y reanudó su tramitación el 29 de septiembre.\footnote{El contenido de la condonación, su distribución por comunidades y el análisis de la AIReF de junio de 2025: Tema 4B-25 (ya redactado).
} Su efecto sobre la deuda consolidada de las AAPP es nulo, porque se trata de una transferencia entre subsectores. Pero sí altera la distribución de la deuda entre el Estado y las CCAA y los incentivos de estas.





\clearpage
\section{España en perspectiva comparada}
\puntosclave{
    
}

El primer bloque ha reconstruido la trayectoria del saldo y de la deuda de las AAPP españolas desde 1975. Cerrado el recorrido, conviene presentar las dos lecturas de la situación actual: ¿son hoy sanas las cuentas públicas españolas? La lectura favorable subraya tres cosas: la reducción de la deuda en unos $19$ puntos de PIB desde 2020; un saldo primario cercano al equilibrio, que no se veía desde 2007; y un crecimiento económico por encima de la media de la zona del euro, que ha permitido a España financiarse con una prima de riesgo moderada, en ocasiones inferior a la de Francia. A ello añade una estructura de deuda larga (unos 8 años de vida media), que amortigua el impacto de la subida de tipos, y una base inversora diversificada. La lectura crítica subraya que el nivel de deuda sigue siendo uno de los más altos de la Unión. Añade que su reducción se ha debido casi por completo al crecimiento nominal, que el déficit estructural se mantiene en torno al 2,5-3\% del PIB según la Comisión y la AIReF (lo que queda pendiente de verificar para 2026), que el gasto neto crece muy por encima de lo comprometido y que España lleva tres años sin presupuestos aprobados. Sostiene que, con el BCE subiendo tipos, una nueva perturbación energética y presiones estructurales de gasto (pensiones, defensa, fin del Plan de Recuperación), la ventana favorable de 2021-2025 se ha cerrado sin haber completado la consolidación estructural. Es la misma crítica que se hizo a 1986-1990, a 2005-2007 y a 2015-2019: no se aprovechó la fase expansiva para crear margen.

Decantarse por una u otra valoración, es una cuestión que se reduce al sistema de creencias de cada individuo. Como sabemos (o deberíamos), las ciencias sociales parten de supuestos (implícitos u explícitos) sin los que no se podría construir un sistema de relaciones tan complejo como el que tenemos. A diferencia de las ciencias puras: dan argumentos, pero no ofrecen verdades.

Por ello puede ser interesante alejarse de la disyuntiva y optar por una vía más empírica, comparada. Al fin y al cabo, el título de deuda de una Gobierno es, por definición, el activo más seguro que se ofrece en el mercado nacional; por tanto, parece razonable que la situación de éste se extraiga por comparación por pares. 

Es frecuente que este ejercicio comparativo se realice sobre todo por niveles de déficit y deuda, acabando por definir un ránking que dice poco sobre la situación real del activo de España. Precisamente, la lección más importante que ofrece la comparación histórica es precisamente que el nivel de deuda pública es un mal predictor de la vulnerabilidad fiscal. En 2007, los dos países de la zona del euro con menos deuda eran Irlanda (24\%) y España (36\%), y fueron los que sufrieron los mayores deterioros en 2008-2012. 

Por ello, optaremos por restringir nuestro análisis y preguntarnos si la trayectoria presentada es peculiar de España o común a otras economías avanzadas, y en qué se diferencia de las de los países con los que se compara habitualmente: los grandes países de la zona del euro, las economías de la periferia europea y el conjunto de la OCDE.\footnote{
Hay que advertir de un problema de método que condiciona la lectura de todos los datos del bloque: la comparabilidad de las estadísticas fiscales entre países. Los países de la Unión publican el déficit y la deuda según el SEC 2010 y el Protocolo de Déficit Excesivo, con la deuda valorada por su nominal y consolidada dentro del sector AAPP. Las bases del FMI y de la OCDE, que son las que permiten comparar con Estados Unidos, Japón o Canadá, siguen el Manual de Estadísticas de Finanzas Públicas (GFSM 2014). Este manual incluye en la deuda bruta instrumentos que la deuda PDE excluye, como los créditos comerciales, y en algunos países los pasivos de los planes de pensiones de los empleados públicos. Por eso la deuda de Estados Unidos que publica el FMI es sensiblemente superior a la que resultaría con el criterio europeo.

A ello se suma la distinción entre deuda bruta y deuda neta de activos financieros. Para países con grandes fondos públicos, como Noruega, Japón o los nórdicos, la deuda neta puede ser muy inferior a la bruta, e incluso negativa. En el caso de Japón, la deuda bruta supera el $230\%$ del PIB, pero la neta ronda la mitad, porque el sector público japonés tiene muchos activos financieros, entre ellos los del fondo de pensiones público y las tenencias cruzadas con el banco central. En España la diferencia es moderada, porque las AAPP tienen relativamente pocos activos financieros líquidos. Estas cautelas obligan a comparar siempre series homogéneas de la misma fuente y a no mezclar datos de Eurostat con datos del FMI en una misma frase.
}

Como veremos, España llega tarde al ciclo de expansión del sector público de las economías avanzadas y lo recorre comprimido en el tiempo. Por eso su deuda es inferior a la media hasta 2010, pero su déficit estructural es persistentemente mayor. Las comparaciones muestran también que la vulnerabilidad fiscal de un país depende menos de su nivel de deuda que de la estructura de sus ingresos, de sus desequilibrios privados y exteriores, de sus instituciones presupuestarias y de su pertenencia a una unión monetaria sin prestamista de última instancia. En todas esas dimensiones España ha sido, hasta fechas recientes, más vulnerable de lo que sugería su nivel de deuda.

\subsection{Trayectoria histórica comparada (1970-2007)}
Empecemos por comparar la trayectoria española con la de las economías avanzadas en las cuatro décadas que van de las crisis del petróleo a la víspera de la Gran Recesión. 

La comparación persigue tres objetivos. El primero es situar el punto de partida español, cuya anomalía (equilibrio con un sector público pequeño) solo se aprecia bien frente a los demás. El segundo es identificar los mecanismos que explican por qué, entre 1975 y 1995, algunos países duplicaron o triplicaron su deuda y otros la contuvieron, lo que permite leer el caso español a la luz de la literatura comparada de economía política. El tercero es mostrar cómo la convergencia de Maastricht y los primeros años del euro produjeron la paradoja con la que se abre 2008: los países con menos deuda no eran los menos vulnerables.

Como decíamos, la comparación muestra a España como un caso de llegada tardía y trayectoria comprimida. En 1975 tenía una de las deudas más bajas de las economías avanzadas, pero con un sector público veinte puntos de PIB menor que la media comunitaria, lo que hacía de su equilibrio un síntoma de atraso y no de virtud. En los ochenta, mientras la subida de los tipos reales provocaba la gran divergencia de las deudas (Bélgica, Italia, Irlanda, Grecia), España multiplicaba su ratio por dos veces y media, pero desde niveles tan bajos que en 1990 seguía por debajo de Alemania o Francia. Su déficit, sin embargo, era ya superior a la media de la OCDE. La crisis de principios de los noventa tuvo epicentros más graves que España (los nórdicos, Italia, Bélgica, Canadá y, a partir de entonces, Japón). Las consolidaciones más duraderas que siguieron, como las de Suecia, Canadá y Bélgica, combinaron ajuste y reforma institucional, sobre todo reglas de gasto. Maastricht ancló la convergencia de los déficits, con una contabilidad creativa cuya huella estadística documentó la literatura. Los primeros años del euro produjeron la gran paradoja con la que se abre el 2008: Irlanda y España, los países con menos deuda y con superávit, eran los más vulnerables, porque su posición fiscal dependía de ingresos transitorios y de un endeudamiento privado y exterior insostenible.

\subsubsection{Los años setenta: la ruptura del equilibrio en las economías avanzadas y la anomalía española}
Las economías avanzadas llegan a los años setenta con ratios de deuda pública históricamente bajas. La deuda heredada de la Segunda Guerra Mundial, que en el Reino Unido superaba el 200\% del PIB en 1946 y en Estados Unidos rozaba el 120\%, se redujo durante los treinta gloriosos (1945-1973) sin grandes superávits primarios. 

La literatura (ABBAS et al., 2010; REINHART y SBRANCIA, 2015) atribuye esa reducción sobre todo a un diferencial muy negativo entre el tipo de interés y el crecimiento. El crecimiento real fue alto, la inflación moderada pero persistente, y en muchos países hubo represión financiera, con controles de capital, coeficientes obligatorios y techos a los tipos de interés, que mantenía los tipos reales bajos o negativos.\footnote{La misma mecánica que vimos para la España franquista. La diferencia es que en las demás economías avanzadas esa mecánica convivía con una rápida expansión del Estado de bienestar, y en España no.
} 

TANZI y SCHUKNECHT (2000) documentan que el gasto público de los países industrializados pasó de un promedio de en torno al 28\% del PIB en 1960 a más del 40\% en 1980. Fue la etapa de mayor crecimiento del sector público en tiempos de paz. 

\paragraph{Crisis del petróleo y la aparición del déficit estructural}
Las crisis del petróleo convierten el déficit en un rasgo estructural de las economías avanzadas. En los setenta el déficit medio de la OCDE siguió siendo bajo, entre el 1\% y el 2\% del PIB, pero ya por encima del casi equilibrio español.

No obstante, lo relevante de los setenta no es tanto el nivel medio del déficit como el cambio de régimen. Los shocks de oferta de 1973 y 1979 redujeron el crecimiento potencial, mientras que los compromisos de gasto social adquiridos en la fase de crecimiento alto se mantuvieron. Los gobiernos intentaron además amortiguar el choque con políticas de demanda y el resultado fue la aparición de déficits que ya no se corregían en la fase alcista.\footnote{Por eso la literatura habla del fin del presupuesto equilibrado a lo largo del ciclo como norma implícita de la posguerra (BUCHANAN y WAGNER, 1977).
}

La inflación de los setenta, sin embargo, siguió conteniendo la deuda. En 1975, la deuda era del 23,7\% del PIB en Alemania, del 20,9\% en Japón, del 44,4\% en Estados Unidos, del 49,6\% en el Reino Unido y del 49,9\% en Canadá. Los casos más altos de Europa continental eran Bélgica (59,5\%) e Italia (57,9\%), y en el extremo opuesto estaban Dinamarca (11,2\%) y Finlandia (6,8\%). Con los datos, la media de la CEE de los nueve, ponderada por el PIB, se situaba en torno al 35-40\%, porque el mayor peso de Alemania y Francia (esta con una deuda en torno al 20\% a finales de los setenta) compensaba los niveles altos de Bélgica, Italia, Irlanda y el Reino Unido. 

España, con una deuda inferior al 10\% del PIB, se situaba en 1975 en el grupo de menor endeudamiento, junto a Finlandia y Dinamarca. Pero su caso era muy distinto del de estos dos países. Finlandia y Dinamarca tenían deuda baja con un sector público ya grande y un sistema tributario moderno. España la tenía con un sector público pequeño y un sistema tributario atrasado. De nuevo, el equilibrio español no era una virtud fiscal sino un síntoma de atraso.


\subsubsection{Los años ochenta: la gran divergencia de las deudas}
Los años ochenta son el periodo en que la deuda de las economías avanzadas diverge como nunca lo había hecho en tiempos de paz. El detonante común es el cambio de política monetaria: el giro antiinflacionista de la Reserva Federal desde 1979, seguido por los demás bancos centrales, elevó los tipos de interés reales a niveles positivos y altos. 

Al mismo tiempo, la liberalización financiera fue desmontando la represión financiera. El resultado fue que el diferencial entre tipo de interés y crecimiento, que había sido negativo durante tres décadas, se volvió positivo en la mayoría de las economías avanzadas: $\Delta d \approx [(i − g)/(1 + g)]·d_{t-1} − sp_t + \text{ADD}_t$, el paso de $i < g$ a $i > g$, convierte la deuda heredada en un motor autónomo de endeudamiento, tanto mayor cuanto mayor sea $d$. Por tanto, con la liberalización apareció el efecto bola de nieve: un país con déficits primarios moderados y una deuda ya relevante podía ver crecer su ratio de forma casi automática.

El déficit medio de la OCDE de los ochenta se situaba en torno al 3\% del PIB y el de España algo por encima del 4\%. A diferencia de los setenta, por tanto, España ya no estaba por debajo de la media sino por encima, y esa diferencia persistirá hasta 2005.

\paragraph{Gran Divergencia}
Entre 1975 y 1990 la deuda belga pasó del 59,5\% al 130,3\% del PIB, y la italiana del 57,9\% al 95,3\%. Irlanda llegó al 108,3\% en 1987, desde el 64,6\% de 1980. Grecia pasó del 22,8\% en 1980 al 74,2\% en 1990, y los Países Bajos del 36\% en 1975 al 75\% en 1990. Canadá pasó del 44,6\% en 1980 al 73,7\% en 1990. Estados Unidos también subió, del 41,3\% en 1980 al 62,2\% en 1990, con la combinación de rebajas de impuestos y aumento del gasto en defensa de la primera administración Reagan. 

En el extremo opuesto, el Reino Unido redujo su deuda del 43,7\% en 1980 al 28,3\% en 1990, con superávits a finales de la década e ingresos de privatizaciones. Alemania la contuvo en torno al 40\%. Japón la subió del 47,8\% en 1980 al 68,3\% en 1985 y la estabilizó después con la burbuja de finales de la década.

En este contexto, España pasa del 16,1\% en 1980 al 40,9\% en 1985 y al 41,3\% en 1990. Es uno de los mayores aumentos relativos del periodo (la deuda se multiplica por más de dos veces y media en cinco años), pero desde un nivel tan bajo que en 1990 la ratio española seguía por debajo de la de Alemania, Francia o Estados Unidos.

La pregunta que le surgiría a cualquiera es por qué. ¿Por qué divergieron? La explicación económica más simple (distintos shocks, distinto crecimiento) no basta, porque los países que más aumentaron su deuda no fueron necesariamente los que sufrieron los peores shocks. Por eso, al buscar en otros ámbitos, la divergencia de los ochenta dio lugar a una de las literaturas más fecundas de la economía política de la hacienda pública.

ROUBINI y SACHS (1989) mostraron que los déficits eran mayores y más persistentes en los países con gobiernos de coalición amplios, frágiles y de corta duración, como Bélgica, Italia, los Países Bajos, Irlanda o Dinamarca en ciertos periodos. En estos gobiernos cada socio tiene capacidad de veto sobre los recortes que afectan a su electorado y el ajuste se retrasa. GRILLI, MASCIANDARO y TABELLINI (1991) ampliaron el análisis a la estructura institucional. Encontraron que los países con sistemas electorales proporcionales y gobiernos inestables acumularon más deuda que los de sistemas mayoritarios, y que la independencia del banco central estaba asociada a menor inflación sin coste en crecimiento. 

ALESINA y DRAZEN (1991) formalizaron el mecanismo con su modelo de guerra de desgaste: cuando los grupos sociales discrepan sobre quién debe soportar el coste del ajuste, cada uno espera que el otro ceda, y la estabilización se retrasa hasta que el coste de no hacer nada se vuelve insoportable. PERSSON y SVENSSON (1989) y ALESINA y TABELLINI (1990) añadieron el argumento de la deuda estratégica: un gobierno que prevé perder el poder puede endeudarse para limitar las opciones de gasto de su sucesor.

Por último, la línea iniciada por VON HAGEN (1992) mostró la importancia de las instituciones presupuestarias. Los países con procedimientos centralizados, un ministro de Hacienda con poder de agenda frente a los ministerios de gasto y restricciones en la fase parlamentaria tenían déficits significativamente menores que los de procedimientos fragmentados.\footnote{El caso español encaja parcialmente en este cuadro. Evidentemente, no encaja en la hipótesis de las coaliciones, porque entre 1982 y 1993 España tuvo gobiernos de mayoría absoluta de un solo partido. Sin embargo, encaja en la de las instituciones presupuestarias débiles, porque VON HAGEN situaba a España en el grupo de procedimientos poco centralizados, y en la de la guerra de desgaste, a la luz de la ruptura del diálogo social que culminó en la huelga general de 1988. Encaja también, y sobre todo, en una explicación que la literatura comparada menciona menos: el catch-up del Estado de bienestar. España estaba construyendo en los ochenta lo que los demás países habían construido en los sesenta y setenta. Su déficit no era de un Estado de bienestar maduro que no quería ajustarse, sino de uno en construcción que no había terminado de ajustar sus ingresos.
}

\paragraph{Consolidaciones: ¿las consolidaciones expansivas?}
El contrapunto de la divergencia son los países que consiguieron corregir la trayectoria en los ochenta, porque de su experiencia nace una tesis que se invocará en todas las consolidaciones posteriores, incluida la española de 2010-2013. 

Dinamarca (1983-1986) e Irlanda (1987-1989) redujeron fuertemente sus déficits y, al mismo tiempo, crecieron con fuerza. GIAVAZZI y PAGANO (1990) interpretaron estos casos como consolidaciones fiscales expansivas. Un ajuste creíble y centrado en el gasto habría mejorado las expectativas de los hogares sobre sus impuestos futuros y reducido las primas de riesgo, estimulando la demanda privada lo bastante para compensar el efecto contractivo directo.

La relectura posterior, cuyo exponente más citado es Perotti (2011), ha mostrado que en esos casos pesaron mucho otros factores. En Irlanda, la devaluación de 1986 y el pacto social de moderación salarial de 1987. En Dinamarca, la fijación del tipo de cambio al marco alemán y la caída de los tipos de interés. En ambos, también, la recuperación de la demanda de sus respectivos socios comerciales. 

Con todo, lo que debemos retener es que el canal de las expectativas existe, pero rara vez es suficiente por sí solo, y requiere circunstancias (política monetaria acomodaticia, tipo de cambio o salarios flexibles) que no se dan en una unión monetaria en recesión. 

\subsubsection{1990-1997: la crisis de principios de los noventa y la convergencia de Maastricht}
La crisis de 1991-1993 desencadenó un periodo de fuerte aumento generalizado del déficit pero con epicentros muy distintos.

El caso más grave de deterioro fiscal fue el de los países nórdicos, que sufrieron crisis bancarias sistémicas tras la liberalización financiera de los ochenta. La deuda de Suecia pasó del 39\% del PIB en 1990 al 66,1\% en 1993, y la de Finlandia del 13,9\% al 54,2\% en solo tres años. El déficit público sueco superó el 10\% del PIB en 1993. En Italia, la crisis del SME de septiembre de 1992 coincidió con un sistema político en descomposición, y la deuda superó el 115\% en 1993. Bélgica alcanzó su máximo en 1993 (138,9\%) y Grecia superó el 100\%. Fuera de Europa, Canadá llegó al 100\% en 1995, y Japón inició, tras el estallido de su burbuja en 1990-1991, un aumento de la deuda que no se ha detenido desde entonces: del 63,2\% en 1990 al 92,5\% en 1995, al 135,6\% en 2000 y al 173\% en 2007. El déficit medio de la OCDE fue del 4,9\%.

En ese cuadro, España (del 41,3\% en 1990 al 61,6\% en 1995) muestra un deterioro intenso pero no extremo. La comparación con Suecia y Finlandia es la más instructiva: en los tres casos una crisis de balance privado, bancaria en los nórdicos y de competitividad y tipos de cambio en España, se tradujo en una crisis fiscal. Un anticipo de lo que ocurriría en 2008-2012.

\paragraph{Consolidaciones de los noventa: el papel de las instituciones}
La segunda mitad de los noventa es la etapa de las grandes consolidaciones, cuya comparación ilustra cuán importante es el diseño institucional.

Suecia reformó a fondo su marco presupuestario tras la crisis: techos de gasto nominales plurianuales aprobados por el Parlamento, un procedimiento presupuestario de arriba abajo y un objetivo de superávit a lo largo del ciclo (1997). Redujo su deuda del 68,7\% en 1995 al 39,2\% en 2007. 

Estados Unidos, con la Budget Enforcement Act de 1990 y sus reglas de compensación (PAYGO), el crecimiento de finales de los noventa y los ingresos de la burbuja bursátil, registró superávits federales entre 1998 y 2001. Canadá aplicó en 1995 su \textit{Program Review}, una revisión sistemática de todos los programas de gasto federal, con recortes muy significativos y una regla de previsiones prudentes. Su deuda bajó del 100\% al 67,2\% en 2007. 

Bélgica, el caso más extremo de Europa, mantuvo durante más de una década superávits primarios de en torno al 5-6\% del PIB, lo que le permitió bajar del 138,9\% en 1993 al 87,3\% en 2007. 

Como vemos, las consolidaciones más duraderas combinaron un ajuste sustancial con reglas de gasto y procedimientos presupuestarios reforzados, no solo con objetivos de saldo. 

\paragraph{Maastricht como ancla de la convergencia}
Además, dentro de este periodo destacan en la literatura comparada los efectos que tuvo Maastricht entre sus integrantes. Como sabemos, Maastricht fijó unos criterios de déficit (3\% del PIB) y de deuda (60\%, o disminuyendo suficientemente) uniformes para acceder a la moneda única.\footnote{Los criterios de Maastricht fue objeto de críticas desde su origen. Buiter, Corsetti y Roubini (1993) cuestionaron tres aspectos. En primer lugar, la arbitrariedad de las cifras de referencia (el 60\% era aproximadamente la media de la época, y el 3\% es el déficit que estabiliza la deuda en el 60\% con un crecimiento nominal del 5\%). En segundo lugar, su carácter procíclico, porque exigía ajustes en recesión. Y, por último, la ausencia de fundamento en una teoría de la sostenibilidad. 

La ratio de deuda se estabiliza cuando el déficit total es $d_t·g_t/(1 + g_t)$. Con $d = 60\%$ y $g = 5\%$ nominal ($3\%$ real y $2\%$ de inflación), el déficit estabilizador es $\sim 2,9\%$, de donde procede la coherencia aritmética entre el 3\% y el 60\%.
} Los once países que formaron la primera oleada de la UEM pasaron de un déficit medio del 5,7\% en 1993 a uno inferior al 2\% a finales de los noventa.

En la evaluación de 1998, con datos de 1997, todos los candidatos salvo Grecia cumplieron el criterio de déficit. En cuanto a la deuda, se aplicó una interpretación flexible: Bélgica e Italia, con ratios superiores al 115\%, fueron admitidas por la tendencia decreciente. España, con un 68,8\%, se situaba en una posición intermedia, lejos de los casos más problemáticos. Grecia se incorporó en 2001.

Ahora bien, la otra cara de la convergencia de Maastricht fue el recurso y proliferación de operaciones que mejoraban el saldo del año de referencia sin mejorar la posición fiscal subyacente (con efectos perversos sobre el periodo subsiguiente). Los ejemplos más conocidos son los de Italia y Francia.

Italia aprobó en 1997 una eurotasa, un impuesto extraordinario que luego se devolvió parcialmente. Francia contabilizó en 1997 como ingreso el pago de France Télécom al Estado a cambio de que este asumiera sus compromisos de pensiones, de alrededor de medio punto de PIB, lo que Eurostat aceptó. KOEN y VAN DEN NOORD (2005), para la OCDE, documentaron el uso sistemático de medidas puntuales en el periodo.

VON HAGEN y WOLFF (2006) aportaron la evidencia más sólida. \textcolor{red}{Los ajustes déficit-deuda, es decir, la diferencia entre la variación de la deuda y el déficit, eran significativamente mayores en los años en que el déficit de un país se acercaba al límite del 3\%}. Es la huella estadística de la contabilidad creativa: se reducía el déficit declarado desplazando operaciones al ajuste déficit-deuda. MILESI-FERRETTI (2003) había modelizado ese incentivo: una regla basada en un indicador manipulable induce la manipulación del indicador, no el ajuste.

El caso extremo, conocido después, fue el de Grecia. La revisión de sus estadísticas fiscales, iniciada en 2004 tras el cambio de gobierno, mostró que sus déficits habían superado el 3\% en todos los años desde 1997, incluido el de referencia para su entrada en el euro.\footnote{En el caso de España no hay evidencia de manipulación comparable en la convergencia. Pero sí se acumularon pasivos fuera del saldo, lo que pertenece a la misma familia de problemas: la distancia entre el indicador vigilado y la posición fiscal real.
}

\subsubsection{1999-2007: los primeros años del euro y la paradoja de los países con menos deuda}
La década previa a la crisis fue muy dispar. 

Empezando por el ámbito extra europeo, Estados Unidos cerró los superávits de 1998-2001 con las rebajas fiscales de 2001-2003, el gasto militar posterior al 11 de septiembre y la recesión de 2001. Su deuda subió del 53,2\% en 2000 al 64,7\% en 2007. Japón, por otro lado, representa el caso extremo de deuda creciente sin crisis de financiación: su ratio superó el 170\% en 2007. La explicación habitual (una deuda en moneda propia, en manos de residentes, con un banco central que puede comprarla y un exceso de ahorro privado que mantiene los tipos en mínimos).\footnote{Confirma, por contrafactual, la tesis de DE GRAUWE: lo que hace vulnerable a un país no es solo cuánto debe, sino en qué moneda, a quién y con qué respaldo del banco central.
} 

En el ámbito europeo, con la llegada de la moneda única, la evolución del déficit fue muy dispar. España e Irlanda continuaron la consolidación, mientras que Francia, Alemania, Italia y Portugal tuvieron grandes dificultades para cumplir el límite del 3\% del PEC. Fuera del euro, Estados Unidos registró superávits en 1999-2000, y Japón transformó sus superávits de principios de los noventa en déficits persistentes.

Las cifras de deuda confirman y completan esa descripción. Entre 2000 y 2007, la deuda de Irlanda bajó del 36,4\% al 24\% y la de España del 57,8\% al 35,7\%, y fueron los únicos grandes descensos de la zona del euro en ese periodo. Bélgica siguió reduciendo su deuda (del 109,6\% al 87,3\%), pero desde niveles muy altos. En cambio, Alemania subió del 59,2\% al 63,7\% y Francia del 59,7\% al 65,5\%, ambas por encima del 60\%. Italia se estancó en torno al 104\%, Portugal pasó del 54,2\% al 72,7\% y Grecia se mantuvo por encima del 100\%. Alemania y Francia incumplieron el límite de déficit a partir de 2002-2003, y la decisión del Ecofin de noviembre de 2003 de no aplicarles el procedimiento llevó a la sentencia del TJUE de 2004 y a la reforma del PEC de 2005.

Alemania, por tanto, es el contraejemplo del relato habitual. En los primeros años del euro no fue un modelo de disciplina fiscal: soportaba todavía el coste de la reunificación, crecía muy poco y era considerada el enfermo de Europa. Su consolidación fiscal y la mejora de su competitividad, con las reformas del mercado laboral de 2003-2005 y la moderación salarial, son posteriores. En 2007 tenía una deuda superior en casi treinta puntos a la española.

\paragraph{Paradoja: }
Con todo lo expuesto es prácticamente imposible dar respuesta a una pregunta. ¿Por qué los países con menos deuda fueron los más vulnerables? Aquí está la lección central de la comparación histórica. 

En 2007, Irlanda y España eran, con diferencia, los grandes países de la zona del euro con menos deuda pública y los únicos con superávit presupuestario significativo. Sin embargo, entre 2008 y 2012 fueron dos de los países con mayor deterioro fiscal de toda la OCDE, y ambos necesitaron asistencia financiera europea: Irlanda, un programa completo en 2010, y España, la asistencia para su sector bancario en 2012.

La literatura ha explicado la paradoja con tres argumentos complementarios. El primero, como habíamos visto para el caso español, es que sus saldos de 2005-2007 contenían un gran componente de ingresos transitorios ligados a burbujas inmobiliarias. El saldo estructural real era mucho peor que el publicado, y ese problema de medición afectaba especialmente a los países con auges de activos: \textbf{grandes economías financieras}.

El segundo argumento es que la vulnerabilidad procedía del sector privado y del sector exterior. LANE (2012) y BALDWIN y GIAVAZZI (2015) han mostrado que la crisis de la zona del euro fue en su origen una crisis de balanza de pagos. Los países de la periferia acumularon grandes déficits por cuenta corriente, financiados con entradas de capital privado del centro de la zona, que se detuvieron bruscamente en 2008-2010. En 2007 el déficit por cuenta corriente español rondaba el 9-10\% del PIB, y la deuda privada de hogares y empresas superaba el 200\%. Cuando las entradas se detuvieron, la deuda privada, sobre todo la bancaria, se convirtió en deuda pública por los rescates bancarios y la caída de ingresos.

El tercer argumento es la pertenencia a una unión monetaria sin prestamista de última instancia, que DE GRAUWE (2011) formuló como la fragilidad de la zona del euro: un país que se endeuda en una moneda que no controla puede sufrir una crisis de liquidez autocumplida aunque sus fundamentales sean sólidos.

La consecuencia para el marco de análisis es de primer orden. La comparación de los niveles de deuda de 2007 habría situado a España entre los países más sólidos de la OCDE. Una comparación que hubiera incluido el saldo estructural corregido por los ingresos transitorios, la posición de inversión internacional neta y la deuda privada la habría situado entre los más vulnerables. Evidentemente, esta lección está detrás del Procedimiento de Desequilibrios Macroeconómicos de 2011 y de la inclusión del análisis de sostenibilidad de la deuda en el nuevo marco fiscal de 2024, que se examinarán más adelante.\footnote{Procedimiento de Desequilibrios Macroeconómicos (Reglamento (Unión) 1176/2011) y su cuadro de indicadores, y el análisis de sostenibilidad de la deuda del Reglamento (Unión) 2024/1263: Tema 3B-45.
}


\subsection{Las crisis de 2008-2023 en perspectiva comparada}
Entre 2008 y 2023 las economías avanzadas atraviesan tres crisis que ponen a prueba sus finanzas públicas de formas distintas. La primera, la Gran Recesión de 2008-2009, es un choque financiero global que deteriora los saldos en todas partes, pero con intensidades muy diferentes. La segunda, la crisis de la deuda soberana de 2010-2013, es un fenómeno esencialmente europeo, y más precisamente de la zona del euro, que separa a los países de la moneda única en dos grupos: los que pierden el acceso a los mercados o se acercan a perderlo, y los que actúan como refugio. La tercera, la pandemia de 2020 seguida del choque inflacionario de 2022-2023, vuelve a ser global y sincronizada, pero llega con una respuesta institucional europea radicalmente distinta.

La comparación de las crisis de 2008-2023 muestra tres rasgos del caso español. El primero es la \textbf{elevada sensibilidad de su saldo al ciclo}. En 2008-2009, con una caída del PIB menor que la alemana o la italiana, España registró un deterioro del saldo del doble de la media de la zona del euro, comparable solo al de Irlanda y Grecia. La causa fue la desaparición de los ingresos del auge inmobiliario y la intensidad de los estabilizadores del desempleo, propia de su mercado de trabajo dual.

El segundo rasgo es su \textbf{posición intermedia en la crisis soberana}. España no recibió un programa completo como Grecia, Irlanda y Portugal, pero necesitó asistencia europea para su banca, y sufrió primas de riesgo que la comparación con el Reino Unido, Estados Unidos o Japón muestra que respondían más al diseño de la zona del euro que a sus fundamentales. El coste de su crisis bancaria fue mucho menor que el irlandés en porcentaje del PIB, pero el agujero de ingresos produjo un aumento de la deuda comparable.

El tercer rasgo, y el más revelador, es su \textbf{incapacidad para generar superávits primarios en las recuperaciones}. En 2014-2019 y, sobre todo, en 2021-2023, países que habían pasado por programas de ajuste, como Portugal, Grecia, Chipre e Irlanda, redujeron su deuda por debajo del nivel previo a la pandemia gracias a superávits primarios. España, con el mismo viento de cola del crecimiento nominal, salió de 2023 con más deuda que en 2019. 

El resultado de los quince años es un cambio de posición completo: de ser uno de los países con menos deuda de la zona del euro en 2007 (36\% frente a una media del 66\%) a estar entre los cinco con más. 

\subsubsection{La Gran Recesión (2008-2009): un deterioro general con intensidades muy distintas}
Tras la crisis financiera internacional prácticamente todos los países de la OCDE incurrieron en déficit o lo agravaron, y el déficit medio pasó de poco más del 1\% del PIB en 2007 a más del 8\% en 2009, con grandes diferencias entre países.

Las cifras de Eurostat de la época permiten concretar esas diferencias. La zona del euro pasó de un déficit del 0,7\% en 2007 al 6,4\% en 2009, un deterioro de unos seis puntos. España pasó de un superávit del 1,9\% a un déficit del 11,2\%, un deterioro de unos $13$ puntos, el doble que la media. Solo Irlanda (del 0,1\% de superávit al 14,2\% de déficit) y Grecia (del 6,5\% al 15,8\% de déficit) registraron deterioros comparables. Fuera de la zona del euro, el Reino Unido pasó del 2,7\% al 11,5\% de déficit. En el otro extremo, Alemania pasó de un ligero superávit a un déficit del 3,2\% e Italia del 1,6\% al 5,4\%.

\paragraph{Por qué España se deterioró el doble que la media: los ingresos}
La comparación permite descartar una explicación simple. La caída del PIB español en 2009 (en torno al −3,8\%) fue menor que la de Alemania (−5,7\%), Italia (−5,5\%) o el Reino Unido (-4,5\%). Si el deterioro del saldo dependiera solo del tamaño de la recesión, el de España debería haber sido menor que el alemán, y fue cuatro veces mayor.

La explicación está en la composición del deterioro. La recaudación española cayó más de seis puntos de PIB, una caída que no se dio en ningún otro país de nuestro entorno. Es el efecto, ya analizado, de la desaparición de los ingresos transitorios del auge inmobiliario, que tenía su paralelo en Irlanda y, en menor medida, en el Reino Unido y en Estados Unidos. Los tres países de la OCDE con mayor deterioro de los ingresos en 2008-2009 fueron precisamente los que habían vivido los mayores auges inmobiliarios.

Un segundo factor comparado es la mayor intensidad de los estabilizadores automáticos del gasto en España. DOLLS, FUEST y PEICHL (2012) estimaron que, en Europa, los estabilizadores automáticos absorbían una proporción del choque de renta de los hogares bastante mayor que en Estados Unidos. Dentro de Europa, España destacaba por la reacción del gasto en desempleo, porque su mercado laboral dual, con una elevada proporción de temporales, destruye empleo con mucha más intensidad que el de otros países ante una misma caída del PIB. Entre 2007 y 2013 la tasa de paro española pasó del 8\% al 26,9\%, mientras que la alemana simplemente bajó. Esa elasticidad del empleo al ciclo, que es un rasgo estructural del mercado de trabajo español, se traduce en una elasticidad del saldo al ciclo mayor que la de sus pares.\footnote{Son cuestiones importantes. La elasticidad presupuestaria al ciclo en la metodología común de la Unión y su mayor valor para España por la reacción del desempleo. Sobre la dualidad del mercado de trabajo español y su efecto en la destrucción de empleo, remito a los temas de mercado de trabajo del temario.
}

La respuesta discrecional fue coordinada a escala del G20 (cumbres de Washington y Londres, 2008-2009) y, en la Unión, a través del Plan Europeo de Recuperación Económica, que recomendaba un estímulo de en torno al 1,5\% del PIB. Los mayores estímulos discrecionales correspondieron a Estados Unidos (la American Recovery and Reinvestment Act de 2009, en torno al 5,5\% del PIB repartido en varios años) y a China, con un paquete de cuatro billones de yuanes (en torno al 12\% de su PIB) concentrado en inversión en infraestructuras. En Europa, Alemania aprobó dos paquetes coyunturales de tamaño relevante con la economía más saneada. El de España, con el Plan E, las rebajas del IRPF y la supresión efectiva del Impuesto sobre el Patrimonio, fue de tamaño similar en porcentaje del PIB. La diferencia es que España lo aplicó sobre un saldo estructural que, como vimos, era mucho peor de lo que parecía.

\subsubsection{La crisis de la deuda soberana (2010-2013): la fractura de la zona del euro}
Otro rago particular que muestrala comparación es que la crisis de 2010-2013 no fue una crisis de las economías con más deuda o más déficit. Fue una crisis de los países de la zona del euro con desequilibrios exteriores y problemas bancarios, como ya sabemos. 

El caso más claro es el del Reino Unido. En 2010 tenía un déficit del 10,3\% del PIB, superior al español (9,3\%), y Estados Unidos y Japón tenían déficits y deudas mayores. Sin embargo, sus rentabilidades a diez años bajaron durante la crisis mientras las de la periferia del euro se disparaban. Japón, con una deuda que superaba el 220\% del PIB en 2012, se financiaba a tipos inferiores al 1\%.\footnote{
De nuevo, la base empírica de la tesis de DE GRAUWE (2011). Los países que se endeudan en su propia moneda y cuentan con un banco central que puede actuar como prestamista de última instancia no sufren crisis de liquidez autocumplidas. Los que se endeudan en una moneda que no controlan, como los miembros de la zona del euro, sí pueden sufrirlas. DE GRAUWE y JI (2013) mostraron que el aumento de las primas de la periferia en 2010-2012 no se explicaba por sus fundamentales fiscales. La fuerte caída que siguió al anuncio del OMT en 2012, sin cambios en esos fundamentales, confirma que buena parte de la prima era un componente de pánico y de riesgo de redenominación.
}

Ante el sucesivo deterioro de las finanzas públicas, la zona del euro respondió con una escala de programas de intensidad creciente. Grecia recibió un primer programa en mayo de 2010, en forma de préstamos bilaterales de los países de la zona del euro y del FMI por 110.000 millones de euros. El segundo (2012) incluyó la mayor reestructuración de deuda soberana de la historia: el canje con participación del sector privado (PSI) de marzo de 2012 impuso a los acreedores privados una quita nominal del 53,5\%. Hubo un tercero en 2015, ya con el MEDE. Irlanda recibió un programa de 85.000 millones en noviembre de 2010, y Portugal uno de 78.000 millones en mayo de 2011. Chipre recibió en 2013 un programa de 10.000 millones condicionado a un bail-in de los depositantes de sus dos principales bancos, un precedente de la resolución bancaria europea. España solicitó en junio de 2012 una línea para su sector bancario de hasta 100.000 millones, de la que utilizaron 41.333. Italia no pidió asistencia, pero sufrió tensiones comparables a las españolas en 2011-2012, y el BCE compró deuda de ambos países con el programa SMP a partir de agosto de 2011.

En esta escala, España ocupa una posición intermedia que se suele malinterpretar. No fue rescatada en el sentido de los programas completos de Grecia, Irlanda y Portugal, que incluían financiación del Estado y condicionalidad macroeconómica, fiscal y estructural sujeta a la supervisión de la troika. Pero tampoco quedó al margen. La asistencia fue sectorial, con condicionalidad financiera, y en términos de deuda pública tuvo el mismo efecto que un préstamo al Estado, porque el FROB era la entidad receptora y el Estado el garante.

La comparación más instructiva en este contexto es con Irlanda, porque ambos países partían en 2007 de una deuda baja y un auge inmobiliario financiado por el crédito bancario. La decisión irlandesa de garantizar en septiembre de 2008 prácticamente todos los pasivos de sus bancos llevó a un déficit del 31,3\% del PIB en 2010, del que unos veinte puntos correspondían a la recapitalización bancaria. La deuda irlandesa pasó del 24\% en 2007 al 118,7\% en 2012. Según la base de crisis bancarias de LAEVEN y VALENCIA (2018), el coste fiscal bruto directo de la crisis bancaria irlandesa se situó en torno a un tercio del PIB. El español, en torno al 4-6\% según las estimaciones (5,5\% del PIB en el cómputo de fondos utilizados). Aun así, el coste de la reestructuración bancaria española, que el Tribunal de Cuentas cifra en 71.833 millones, es uno de los mayores de Europa en términos absolutos, solo por detrás de Alemania y del Reino Unido.

La diferencia de escala se explica por el tamaño relativo del sistema bancario (en Irlanda superaba varias veces su PIB), por la garantía general de 2008 y por el hecho de que en España el problema se concentró en las cajas de ahorros, mientras los grandes bancos resistían. Pero la consecuencia fiscal agregada fue parecida por otra vía. En España pesó mucho menos el rescate bancario y mucho más la caída estructural de los ingresos. Ambos países llegaron a ratios de deuda cercanas al 100\%-120\%: Irlanda por el rescate bancario y España por el agujero de los ingresos.

\paragraph{Consolidación: Ajuste comparado y el debate sobre sus costes}

Los países de la periferia aplicaron entre 2010 y 2013 consolidaciones fiscales de gran tamaño, en buena parte simultáneas y en plena recesión. La consolidación discrecional acumulada, medida por la mejora del saldo estructural primario, fue del orden de 15-20 puntos de PIB en Grecia, de 7-9 en Portugal y en Irlanda y de 6-8 en España en ese periodo, según las estimaciones de la Comisión y del FMI. Es una magnitud sin precedentes en economías avanzadas en tiempos de paz.

La literatura comparada ha evaluado el coste de este ajuste con gran intensidad. BLANCHARD y LEIGH (2013), como vimos, mostraron que los multiplicadores eran mayores de lo supuesto. HOUSE, PROEBSTING y TESAR (2020) estiman que la austeridad explica una parte muy importante de la caída del PIB de la periferia europea entre 2010 y 2014, y que en algunos casos la ratio de deuda habría aumentado menos sin ella. FATÁS y SUMMERS (2018) añaden el argumento de la histéresis: si la recesión reduce de forma permanente el producto potencial, por la pérdida de capital humano y la caída de la inversión, la consolidación en recesión puede empeorar la sostenibilidad a largo plazo.

La defensa del ajuste, se basa en que sin él no habría habido financiación, ni europea ni de mercado. La comparación internacional ofrece un matiz relevante. El Reino Unido, fuera del euro, también aplicó un programa de consolidación desde 2010, pero con su propio banco central comprando deuda y con una libra que se había depreciado en torno a un 25\% desde 2007. El coste en crecimiento fue menor. 

La diferencia no está tanto en el ajuste fiscal como en la ausencia de amortiguadores monetarios y cambiarios en la zona del euro hasta 2012.

\subsubsection{La recuperación (2014-2019): dos modelos de salida}
Entre 2014 y 2019, las trayectorias de deuda de la zona del euro divergen de nuevo. Alemania, con superávits presupuestarios desde 2014 y su regla constitucional de freno de la deuda (Schuldenbremse, 2009), redujo su deuda del 79,8\% del PIB en 2012 al 58,7\% en 2019, por debajo del umbral de Maastricht. Irlanda la redujo del 118,7\% en 2012 al 55,9\% en 2019. Portugal, del 132,5\% en 2014 al 116,1\%, y los Países Bajos, del 67,2\% al 47,6\%. En el grupo contrario, la de Francia subió del 96,2\% al 98,2\% y la de Italia se estancó en torno al 134\%. España pasó del 104,4\% al 97,6\%, un descenso de unos siete puntos en cinco años de crecimiento superior a la media.\footnote{
\textbf{NOTA.-} La reducción irlandesa está exagerada por un problema de denominador. En 2015, la relocalización de activos intangibles de multinacionales elevó el PIB irlandés un 25\% en un solo año, lo que redujo mecánicamente su ratio de deuda. Las autoridades irlandesas utilizan desde entonces un indicador alternativo, la renta nacional bruta modificada (GNI), sobre el que su deuda es muy superior a la que resulta sobre el PIB. 
}

La comparación más reveladora de este periodo, y la que anticipa el contraste de 2019-2023, es la de España con Portugal. Ambos países crecieron en 2015-2019 a tasas parecidas, y ambos se beneficiaron del programa de compras del BCE. Pero Portugal, que salió de su programa en 2014 con compromisos de consolidación supervisados, mantuvo superávits primarios de en torno al 3\% del PIB a partir de 2016, y cerró 2019 con un superávit total, el primero de su democracia. España mantuvo un déficit primario hasta 2018 y un saldo total cercano al −3\% al final del periodo.

La explicación que ofrece la economía política comparada combina dos factores. Por un lado, los programas de ajuste dejaron en los países rescatados un anclaje institucional más fuerte: consejos fiscales reforzados, supervisión posprograma y una memoria social del coste de perder el acceso a los mercados. Por otro, España vivió en 2015-2019 un periodo de fragmentación política sin precedentes, con presupuestos prorrogados, que la literatura (ROUBINI y SACHS, 1989) asocia a una menor capacidad de ajuste. 

En 2019 España tenía el cuarto déficit público más elevado de las economías avanzadas, solo superado por Estados Unidos, Australia e Israel, e igualado por Francia. La conclusión es robusta: en plena fase alta del ciclo y con tipos de interés en mínimos históricos, España estaba entre las pocas economías avanzadas con un déficit en torno al 3\% del PIB. 

Es la misma posición que en 2007 con respecto a su saldo estructural oculto. La diferencia es que ahora el déficit era visible, porque ya no había ingresos transitorios que lo enmascararan.

\subsubsection{La pandemia y la inflación (2020-2023): una crisis sincronizada con respuestas desiguales}
Según los datos del FMI, entre 2019 y 2023 la ratio de deuda cambió de forma muy desigual. Portugal la redujo del 116,1\% al 97,7\% (−18,4 puntos), Grecia del 183,7\% al 165,2\% (−18,5), Irlanda del 55,9\% al 43,3\% (−12,6) y Chipre del 92,3\% al 73,7\% (−18,6). En el otro extremo, España la aumentó del 97,6\% al 105\% (+7,4), Francia del 98,2\% al 109,8\% (+11,6), Estados Unidos del 108,2\% al 119\% (+10,8) y el Reino Unido del 85,7\% al 100,4\% (+14,7). Alemania (+4,2) e Italia (+0,7) se situaron en posiciones intermedias.

Si se comparan 2019 y 2022, España fue la quinta economía avanzada con mayor aumento de su ratio (unos 15 puntos), solo superada por San Marino, Japón, Nueva Zelanda y Malta. Su ratio, del 113,1\% en 2022, era la séptima más alta de las economías avanzadas y la cuarta de la zona del euro, unos 20 puntos por encima de la media.

El contraste no se explica por las condiciones macroeconómicas, que fueron parecidas. Todos los países de la periferia se beneficiaron del crecimiento nominal de 2021-2023, de la inflación y de los fondos europeos, y todos soportaron la subida de tipos de 2022-2023. Se explica por el saldo primario. Portugal, Grecia y Chipre volvieron a superávits primarios en 2022-2023: Grecia del orden del 2\% del PIB, Portugal por encima del 2\% y Chipre cerca del 3\%. España mantuvo déficits primarios en todos los años hasta 2024, y su saldo primario no se acercó al equilibrio hasta 2025. Con el mismo viento de cola del denominador, los países que generaban superávits primarios redujeron su deuda por debajo del nivel de 2019, y España no llegó a recuperarlo. En 2025 su ratio (100,7\%) sigue por encima de la de 2019.

La comparación con Francia merece un comentario, porque es el país con el que España ha convergido. Francia tenía en 2007 una deuda casi treinta puntos superior a la española y en 2023 era solo cinco puntos superior. Ambos países comparten un déficit estructural persistente, un gasto público elevado en relación con sus ingresos y dificultades políticas para aprobar presupuestos. Francia entró en procedimiento de déficit excesivo en 2024, y España no, porque su déficit de 2023 era inferior y se preveía por debajo del 3\% en 2024. En 2024-2025, la prima de riesgo española se situó en algunos momentos por debajo de la francesa, lo que indica que los mercados valoraban mejor la dinámica española que el nivel. 

Analicemos más en concreto los shocks del periodo para ver que nos dicen los datos.


\paragraph{Pandemia}
Evidentemente, el cataclismo del periodo fue la pandemis del COVID-$19$, donde la capacidad de respuesta fiscal en 2020 dependió del acceso a la financiación, y ese acceso dependió a su vez del respaldo del banco central. 

En las economías avanzadas, el déficit medio pasó de menos del 3\% del PIB en 2019 a casi el 11\% en 2020, y dos años después había vuelto a niveles próximos a los de 2019. En las economías emergentes (excluida China) pasó de casi el 4\% a casi el 8\%, y había vuelto a niveles cercanos al 4\% en 2022. China, que ya tenía un déficit cercano al 6\% en 2019, llegó a casi el 10\% en 2020 y a un 9\% en 2022, por la prolongación de su política de covid cero. Las economías de bajo ingreso pasaron de casi el 4\% a casi el 5\%. La deuda pública mundial aumentó unos $15$ puntos de PIB en 2020, el mayor aumento anual desde la Segunda Guerra Mundial, hasta alcanzar el 99\% del PIB. Más del 90\% del aumento se concentró en las economías avanzadas y en China, que pudieron apoyarse en sus bancos centrales y en mercados financieros más desarrollados. En las economías avanzadas la deuda subió $19$ puntos, hasta el 124\% del PIB, con Estados Unidos cerca del 135\%. En las emergentes y de bajo ingreso subió unos 7 puntos, sobre todo por la caída del PIB.

España fue la séptima economía avanzada con mayor aumento de su deuda en 2020, con más de $21$ puntos de PIB. El mayor aumento fue el de Canadá, con más de $30$ puntos, y Taiwán incluso la redujo ligeramente. La deuda española pasó del 97,6\% en 2019 al 119,2\% en 2020. 

La comparación con la zona del euro muestra que el aumento español se explica en gran medida por el denominador. La caída del PIB español en 2020 (en torno al −11\%) fue la mayor de la zona del euro, casi el doble de la media (en torno al −6\%), por el peso del turismo, de la hostelería y de los servicios de contacto personal en su estructura productiva. Como vimos, el efecto denominador explicó en torno a la mitad del aumento de la ratio española en 2020. La respuesta discrecional española fue, en comparación, contenida en gasto directo y amplia en liquidez. La base de datos de medidas fiscales del FMI sitúa el apoyo directo (por encima de la línea) de España por debajo del de Alemania, el Reino Unido, Estados Unidos o Japón, y su apoyo de liquidez (avales, préstamos y capital) en niveles comparables a los de Francia, aunque por debajo de Italia y Alemania. 

\paragraph{Crisis energética de 2022: el mismo patrón}
La respuesta a la crisis energética de 2022 muestra de nuevo una diferencia de espacio fiscal. Según la base de datos de Bruegel, los países europeos destinaron casi 800.000 millones de euros a proteger a hogares y empresas frente a la subida de los precios de la energía entre septiembre de 2021 y enero de 2023. 

Alemania encabezaba la lista con unos 270.000 millones, seguida del Reino Unido, Italia y Francia, con menos de 150.000 millones cada uno. España movilizó un volumen menor en porcentaje del PIB, en parte gracias al mecanismo ibérico, que redujo el precio de la electricidad sin coste presupuestario directo. La crítica de Bruegel, extensible a España, es que la mayor parte del apoyo europeo no estaba focalizado y actuaba como una subvención a los combustibles fósiles.

\subsection{Posición actual de España en perspectiva comparada (2024-2026)}
La verdad, con todo lo expuesto, es que la foto comparada de 2025-2026 favorece a España más de lo que lo explicado haría esperar, y eso encierra una trampa. 

Con datos de Eurostat de abril de 2026, España tuvo en 2025 un déficit (2,4\%) inferior a la media de la zona del euro (2,9\%) y muy inferior al de Francia (5,1\%) o Bélgica (5,2\%). Su prima de riesgo cotiza por debajo de $60$ puntos básicos, el diferencial de Francia frente a España está en máximos históricos y S\&P le da la misma calificación que a Francia. Pero España sigue teniendo la quinta deuda más alta de la Unión, con un déficit estructural que no ha corregido y unas necesidades brutas de financiación entre las mayores de Europa. 

Hasta ahora, hemos reconstruido la trayectoria comparada histórica. Sin embargo, para ver el por qué de la afirmación anterior debemos analizar su posición actual con tres grupos de indicadores. Los de nivel, que comparan el déficit, la deuda, el saldo primario y el tamaño del sector público. Los de dinámica y sostenibilidad, proyecciones de deuda, necesidades de financiación, envejecimiento y los indicadores de vulnerabilidad macrofinanciera. Y, por último, los de percepción recogen las primas de riesgo y las calificaciones crediticias.

Como veremos, España ocupa hoy una posición comparada mejor en dinámica y peor en nivel. 


En 2025-2026, España ocupa una posición comparada ambivalente. Por nivel de déficit, está por debajo de la media de la zona del euro (2,4% frente a 2,9%) y muy por debajo de Francia, Bélgica o Austria. Por nivel de deuda, tiene la quinta ratio más alta de la Unión (100,7%), unos 13 puntos por encima de la media, detrás de Grecia, Italia, Francia y Bélgica, y ha sido rebasada a la baja por Portugal. Su saldo primario está cerca del equilibrio, pero lejos de los superávits de la antigua periferia. Sus intereses pesan más que la media en el gasto público, y sus ingresos siguen varios puntos por debajo de la media europea.

Los indicadores prospectivos muestran un riesgo alto a medio plazo según la Comisión y necesidades brutas de financiación entre las más altas de la Unión. En cambio, los indicadores de vulnerabilidad macrofinanciera que fallaron en 2007 (cuenta corriente, posición exterior, deuda privada, solvencia bancaria) muestran hoy una posición mucho más sólida. La configuración de 2007 se ha invertido: deuda pública alta y deuda privada moderada.

Los mercados y las agencias premian la dinámica: la prima de riesgo está por debajo de 60 puntos básicos, por debajo de la francesa, y la calificación está en A+ con S&P y Scope. Pero las primas son hoy una señal menos fiable de disciplina por el respaldo del BCE. La comparación, en suma, confirma la tesis del Bloque I con otro ángulo: España ha mejorado su posición relativa gracias al crecimiento y al deterioro de otros, pero no ha resuelto el problema estructural que la comparación con Portugal pone de relieve.

Cierre del bloque. El Bloque II ha mostrado que la trayectoria fiscal española es una versión tardía y comprimida del ciclo de las economías avanzadas. España llegó tarde a la expansión del sector público y la recorrió con un déficit persistentemente superior a la media. Su deuda fue inferior a la de sus pares hasta 2010, y su vulnerabilidad, mayor de lo que esa deuda sugería. Las crisis de 2008-2023 la han llevado de estar entre los países con menos deuda de la zona del euro a estar entre los cinco con más. El rasgo que la distingue de los países que han reducido su deuda es la ausencia de superávits primarios en las recuperaciones. Con este diagnóstico, el Bloque III examina cómo el Tesoro Público financia una deuda de esta dimensión: sus objetivos, sus necesidades, sus instrumentos, sus procedimientos y su estrategia en un contexto en que el BCE ha dejado de comprar y vuelve a subir los tipos.



\subsubsection{Nivel: déficit, deuda y tamaño del sector público en 2025}
Según la notificación de Eurostat de abril de 2026, el déficit de la zona del euro se situó en el 2,9\% del PIB en 2025, frente al 3\% de 2024, y el de la Unión se mantuvo en el 3,1\%. En ese cuadro, España (2,4\%, incluidos los gastos de la DANA) está claramente por debajo de la media y muy por debajo de los grandes países con déficit elevado: Francia (5,1\%), Bélgica (5,2\%), Austria (4,2\%), Finlandia (3,4\%) e Italia (3,1\%). También está por debajo de Alemania, cuyo déficit (2,7\%) empieza a crecer con la reforma de su freno de la deuda para financiar defensa e infraestructuras. Los mayores déficits de la Unión fueron los de Rumanía (7,9\%) y Polonia (7,3\%). 

Solo cinco Estados miembros registraron superávit: Chipre (3,4\%), Dinamarca (2,9\%), Irlanda (1,8\%), Grecia (1,7\%) y Portugal (0,7\%).

La comparación muestra dos grupos en la zona del euro. En uno están los países de la antigua periferia (Grecia, Portugal, Chipre e Irlanda), con superávits y deudas en rápida reducción. En el otro, varios países del centro y del norte (Francia, Bélgica, Austria, Finlandia) con déficits superiores al 3\% y procedimientos de déficit excesivo abiertos o en perspectiva. España se sitúa entre ambos: sin superávit, pero con un déficit por debajo del 3\% y decreciente. Una inversión casi completa del mapa de 2010-2012, que conviene señalar.

Por otro lado, al cierre de 2025, la deuda de la zona del euro era del 87,8\% del PIB y la de la Unión del 81,7\%. España (100,7\%) tiene la quinta ratio más alta de la Unión, detrás de Grecia (146,1\%), Italia (137,1\%), Francia (115,6\%) y Bélgica (107,9\%). Supera en unos 13 puntos la media de la zona del euro y en más de 10 a Portugal (89,7\%), que en 2014 tenía casi treinta puntos más que España. Queda muy por encima de Alemania (63,5\%), de los Países Bajos (44,4\%) y de Irlanda (32,9\%). Las ratios más bajas de la Unión son las de Estonia (24,1\%), Luxemburgo (26,5\%) y Dinamarca (27,9\%).

Frente a 2022, la distancia con la media de la zona del euro se ha reducido de unos 20 a unos 13 puntos, y España ha pasado del cuarto al quinto lugar de la zona del euro, al superarla Bélgica. Pero la reducción de esa distancia se debe sobre todo a que la ratio española ha caído por el denominador, no a un cambio en su posición estructural. Además, varios países que estaban por debajo de España, como Portugal o Chipre, la han rebasado a la baja en estos años.

Por último, estaría la comparación del saldo primario, más informativa que la del saldo total. En 2025 España tiene un saldo primario cercano al equilibrio. Es mejor que el de Francia o Bélgica, que tienen déficits primarios de varios puntos, pero peor que el de Italia, que viene registrando superávits primarios, y claramente peor que los de Grecia, Portugal o Chipre, con superávits primarios del 2-4\%. 

Eurostat añade un indicador de coste de la deuda. En 2025 los intereses pagados por las AAPP representaron el 5,3\% del gasto público en España, frente al 3,8\% de media en la zona del euro. Solo Hungría (8\%), Italia (7,6\%), Rumanía (6,6\%) y Grecia (6,5\%) tuvieron una proporción mayor. En porcentaje del PIB, los intereses españoles (en torno al 2,4\%) superan la media de la zona del euro. Son una carga que otros países con menos deuda no soportan, y que en 2026 vuelve a crecer.

Por último, el tamaño del sector público de la zona del euro en 2025 fue del 49,8\% del PIB, por gastos, y del 46,9\%, por ingresos. España se sitúa varios puntos por debajo de la media en ambas magnitudes. Su gasto está en torno al 45-46\% del PIB y sus ingresos en torno al 42-43\%, siendo la brecha de ingresos frente a la media es mayor: una constante desde la Transición. España ha homologado el tamaño de su Estado de bienestar con Europa más que su capacidad recaudatoria, y su déficit estructural refleja en buena parte esa brecha de ingresos.\footnote{La brecha fiscal de España frente a la Unión (tipos efectivos, bases imponibles, beneficios fiscales y economía sumergida) y el debate sobre la suficiencia del sistema tributario corresponden a los temas de sistema fiscal del temario. Ver también el Libro Blanco sobre la Reforma Tributaria (2022) del Comité de personas expertas.

La literatura española discute si la brecha procede de tipos más bajos o de bases más estrechas (economía sumergida, beneficios fiscales, fraude); la mayoría de los análisis de la AIReF, del BdE y de la Comisión apuntan más a lo segundo.
}

\subsubsection{Sostenibilidad y vulnerabilidad: lo que miden los indicadores prospectivos}

a) El análisis de sostenibilidad de la Comisión

El nuevo marco fiscal europeo ha convertido el análisis de sostenibilidad de la deuda (DSA) en el eje de la supervisión. Las sendas de gasto neto se derivan de una trayectoria que debe situar la deuda en una senda descendente y plausible a medio plazo.\footnote{El análisis de sostenibilidad de la deuda en el Reglamento (Unión) 2024/1263 (trayectoria técnica, salvaguardias de resiliencia del déficit y de sostenibilidad de la deuda) y los indicadores S0, S1 y S2 de la Comisión: Tema 3B-45 para el marco y Tema 3A-39 para los indicadores de sostenibilidad.
} La Comisión publica anualmente su Debt Sustainability Monitor, que clasifica a los Estados miembros por su riesgo a corto, medio y largo plazo.

En la edición de 2024, la Comisión situó a España entre los nueve Estados miembros con riesgo alto de sostenibilidad a medio plazo. En un escenario sin cambios de política, proyectaba una deuda de en torno al 109% del PIB en 2030 y al 118% en 2034, con un ajuste anual necesario de entre el 0,4% y el 0,6% del PIB. Identificaba además unas necesidades brutas de financiación cercanas al 20% del PIB, entre las más altas de la Unión junto con Italia, Francia y Bélgica. La edición de 2025, publicada el 12 de febrero de 2026, eleva a doce el número de países con riesgo alto a medio plazo. España sigue, según todos los indicios, en ese grupo, aunque no he podido verificar su clasificación en la ficha país. En la dimensión a corto plazo, la Comisión no identifica a España entre los países con necesidades de financiación preocupantes, que en esta edición son Bélgica, Francia, Italia y Finlandia.

La diferencia entre la senda del Plan Fiscal y Estructural español, con la deuda bajando al 90,6% en 2031, y la proyección de la Comisión sin cambios de política, con la deuda subiendo, mide precisamente el esfuerzo que España se ha comprometido a hacer. Ese esfuerzo exige respetar la senda de gasto neto que la AIReF ve en riesgo en 2026 (1.5.3d). La Comisión clasifica el riesgo a largo plazo, vinculado al envejecimiento, como medio o alto según el escenario de pensiones. La interacción entre la reforma de 2021-2023 y la cláusula de salvaguarda se analizó en el Tema 24.\footnote{Proyecciones del gasto en pensiones (Ageing Report 2024), la reforma de 2021-2023, la cláusula de salvaguarda y el informe de la AIReF de mayo de 2026: Tema 4B-24 (ya redactado).
}

b) La lección de 2008: los indicadores de vulnerabilidad más allá de la deuda pública

La lección central de 2.1.4 era que el nivel de deuda pública, por sí solo, no mide la vulnerabilidad fiscal. En 2007, España e Irlanda tenían deuda baja pero enormes desequilibrios privados y exteriores. Conviene aplicar esa lección a la foto de 2025, porque su resultado es simétrico. Los indicadores que en 2007 señalaban una vulnerabilidad extrema muestran hoy una posición mucho más sólida. España tiene superávit por cuenta corriente desde 2013, que en 2024-2025 rondaba el 3% del PIB. Su posición de inversión internacional neta ha mejorado de en torno al −98% del PIB en 2014 a en torno al −40/−45% en 2025. La deuda de hogares y empresas se ha reducido desde más del 200% del PIB en 2010 a en torno al 125%, por debajo de la media de la zona del euro. El sistema bancario tiene más capital y menos exposición inmobiliaria, bajo la supervisión única del BCE.

En otras palabras, la configuración de 2007 se ha invertido. España tenía entonces una deuda pública baja con deuda privada y exterior muy alta, y tiene ahora una deuda pública alta con deuda privada y exterior moderada. La segunda configuración es menos vulnerable a una crisis de balanza de pagos como la de 2010-2012, porque ya no hay entradas de capital privado cuya interrupción pueda provocarla. Pero es más vulnerable a una subida de tipos y a un choque que exija gasto público, porque el espacio fiscal para absorberlo es menor. Por eso la mejor valoración de los mercados es compatible con la advertencia de la Comisión: miden riesgos distintos.\footnote{Los indicadores del Procedimiento de Desequilibrios Macroeconómicos (cuenta corriente, posición de inversión internacional neta, deuda privada, crédito, precios de la vivienda) y la salida de España de la categoría de desequilibrios: Tema 3B-45.
}

c) La estructura de la deuda: un amortiguador comparativo

Según el BdE, la vida media de la deuda española era de unos 8 años en 2024, frente a 8,9 años en la zona del euro. El 13,9% de la deuda vencía en menos de un año, y los no residentes tenían el 44,9%. Frente a Italia o Francia, España no tiene una estructura de vencimientos especialmente desfavorable. La vida media española es algo menor que la media de la zona del euro, lo que hace que las subidas de tipos se trasladen algo más deprisa al coste medio. Lo desarrollaré en 3.2 y 3.5, porque es el terreno de la política del Tesoro.

\subsubsection{Percepción de los mercados y de las agencias de calificación}

a) Las primas de riesgo: la inversión del mapa europeo

En septiembre de 2026, la prima de riesgo española frente al bono alemán a diez años ha bajado de 60 puntos básicos, un nivel que no se veía desde antes de la crisis financiera. Y el diferencial entre la deuda francesa y la española se sitúa en máximos históricos: Francia paga más que España a diez años, con su bono cerca del 4% en el verano de 2026. Italia, que vivió en septiembre la mayor caída de primas de la zona del euro, cotiza también con primas históricamente bajas. Es un cambio radical respecto a julio de 2012, cuando la prima española alcanzó los 638 puntos básicos (1.4.2).

La literatura y los analistas ofrecen tres interpretaciones complementarias. La primera es que los mercados valoran la dinámica más que el nivel: el crecimiento español superior a la media, un déficit decreciente y la estabilidad del saldo primario, frente a la parálisis presupuestaria y política francesa. La segunda es el respaldo institucional del BCE. Desde julio de 2022 existe el Instrumento para la Protección de la Transmisión (TPI), que permite al BCE comprar deuda de un país que sufra un aumento de la prima no justificado por sus fundamentales, siempre que cumpla el marco fiscal europeo. Ese respaldo reduce el componente de pánico que De Grauwe identificó en 2010-2012 (2.2.2). La tercera es la escasez relativa de activos seguros en euros y la diversificación de los inversores, que favorecen a los emisores de la zona del euro con fundamentales en mejora.

La crítica a esta lectura sostiene que los mercados pueden estar siendo complacientes. Si el TPI actúa como un seguro implícito, las primas han dejado de ser una señal fiable de disciplina fiscal: es el problema de riesgo moral que se discutió al diseñar el instrumento. Además, las primas de riesgo cayeron también en 2005-2007, cuando la deuda española bajaba pero su vulnerabilidad aumentaba (2.1.4).\footnote{El Instrumento para la Protección de la Transmisión (julio de 2022), sus criterios de elegibilidad y el debate sobre el riesgo moral: Temas 3A-36/37. Su efecto sobre la política de emisión del Tesoro, en 3.5 de este tema.
}

b) Las calificaciones crediticias

Tras las mejoras de 2025, las calificaciones de España son las siguientes. S\&P le asigna A+ con perspectiva estable. Fitch le asigna A y Moody's A3, ambas estables. Morningstar DBRS le asigna A (high) y Scope la elevó a A+ con perspectiva estable el 25 de septiembre de 2026. En el caso de S\&P, la calificación de España coincide hoy con la de Francia, que perdió notas con Fitch y S&P en septiembre y octubre de 2025, en ambos casos hasta A+. Italia sigue en la zona BBB+/Baa, e Italia y España son los dos grandes países cuyas notas han mejorado de forma más continuada desde 2023. Las agencias atribuyen las mejoras de España a su crecimiento, a la corrección de los desequilibrios exteriores y a la reducción del déficit. Condicionan nuevas subidas, que llevarían a la banda AA, a una reducción clara de la ratio de deuda, que sigue siendo su principal debilidad.

(OJO: la historia de las calificaciones españolas sirve para ilustrar el ciclo completo. España tuvo la máxima calificación, AAA/Aaa, hasta 2009-2010. Moody's la rebajó nueve escalones en menos de dos años, hasta Baa3 en junio de 2012, a un escalón del grado especulativo. Desde entonces ha recuperado posiciones de forma gradual, pero sigue lejos de la AAA que tenía con una deuda del $36\%$. La diferencia entre la nota de 2007 y la actual resume la lección del Bloque I: las agencias, como los mercados, no detectaron la vulnerabilidad de 2007 y la penalizaron después con dureza.)

\subsubsection{Una valoración comparada}

Cerrado el recorrido, conviene presentar al tribunal una valoración que integre las dimensiones anteriores.

La lectura favorable sostiene que, en términos relativos, España es hoy uno de los grandes países de la zona del euro con mejor trayectoria fiscal. Su déficit está por debajo de la media y bajando, su saldo primario está cerca del equilibrio y su crecimiento es de los más altos. Sus desequilibrios privados y exteriores se han corregido, su prima de riesgo está en mínimos de casi dos décadas y su calificación mejora. Frente a Francia, Bélgica o Austria, su posición ha mejorado de forma nítida, y la comparación con 2012 muestra un cambio de régimen.

La lectura crítica sostiene que esa mejora es relativa y en parte coyuntural. España tiene la quinta deuda más alta de la Unión y una de las mayores cargas de intereses en proporción al gasto. La Comisión la clasifica con riesgo alto a medio plazo, y sus necesidades brutas de financiación están entre las más elevadas de Europa. Mantiene un déficit estructural, una brecha de ingresos frente a Europa y un gasto neto que crece por encima de lo comprometido. Y, a diferencia de Portugal, Grecia o Chipre, no ha generado superávits primarios en la fase favorable del ciclo. Frente a la antigua periferia, su posición ha empeorado de forma igualmente nítida.

Mi síntesis, para presentarla como valoración razonada, es que ambas lecturas describen realidades distintas: la posición relativa de España ha mejorado porque otros países han empeorado (Francia, Bélgica) y porque su crecimiento ha sido alto, y su posición absoluta sigue siendo vulnerable porque no ha corregido su déficit estructural. La comparación más exigente, y la más útil para la política fiscal, no es con Francia sino con Portugal. Con una historia fiscal parecida y un choque mucho mayor en 2011-2014, Portugal ha reducido su deuda por debajo de la española generando superávits primarios. Esa es la vara de medir que el nuevo marco europeo, basado en la dinámica de la deuda, va a aplicar.



\clearpage
\section{Política de financiación del Tesoro Público}
El riesgo de este bloque es que el documento base trata la política de financiación del Tesoro como un catálogo de instrumentos y procedimientos. No dice qué problema resuelve el Tesoro ni con qué criterios decide. Además, contiene un error institucional de partida. Sitúa la gestión de la deuda en una Dirección General del Tesoro y Política Financiera del Ministerio de Asuntos Económicos y Transformación Digital. Ese ministerio ya no existe con ese nombre, y la competencia corresponde a la Secretaría General del Tesoro y Financiación Internacional, integrada en el Ministerio de Economía, Comercio y Empresa, de la que depende esa dirección general.

La pregunta que un tribunal puede hacer es sencilla: ¿qué objetivo persigue el Tesoro y cómo se mide si lo cumple?. Si la respuesta se limita a financiarse al menor coste, está incompleta, porque omite el riesgo, que es la mitad del problema. Por eso 3.1 se centra en el marco jurídico, la organización y, sobre todo, en los objetivos y su fundamento teórico. El catálogo de instrumentos queda para 3.3 y 3.4.

Dos datos quedan pendientes de cierre por problemas de acceso a las fuentes. La web del Tesoro no ha sido accesible desde aquí, y los datos de vida media y coste medio de 2025 proceden de fuentes secundarias. Y en 1.1.4 di una emisión bruta para 2026 de 285.693 millones, mientras que la nota oficial del Gobierno sobre la Estrategia 2026 da 285.677 millones. La cifra correcta es la oficial, y conviene corregirla allí.


## Introducción al bloque

Los dos primeros bloques han mostrado que el déficit es un rasgo estructural de las cuentas públicas españolas y que la deuda de las AAPP ronda hoy el 100% del PIB. Este tercer bloque estudia cómo se financia esa deuda.

[DOC] El documento justifica el estudio de la forma de financiación por dos razones. La primera son los efectos económicos que genera. La segunda es que una política de financiación óptima ayuda a reducir las necesidades de financiación. Añade que esta segunda razón es especialmente importante en la actualidad, dadas las grandes necesidades de financiación de las AAPP españolas, y recuerda que la deuda del Estado supone casi el 90% de la deuda total de las AAPP.

Ambas razones son correctas y pueden precisarse con lo que hemos visto. La primera razón, los efectos económicos, remite al vínculo entre la política de deuda, la política monetaria y la estabilidad financiera. La deuda pública es el activo sin riesgo de referencia para la economía nacional. Su curva de rentabilidades sirve de base para fijar el precio de la deuda privada. Y el volumen y los plazos de su emisión afectan a las condiciones financieras de toda la economía. La segunda razón, la eficiencia, remite a la carga de intereses. En 2025 las AAPP pagaron algo más de 40.000 millones de euros en intereses (1.1.2), una cifra comparable a la recaudación del Impuesto sobre Sociedades. Cada punto básico de ahorro en el coste medio de una cartera de más de un billón y medio de euros equivale a unos 150 millones anuales, una vez que toda la cartera se ha renovado a ese coste.

El bloque parte de la bisagra que se estableció en 1.1.4. El Tesoro no financia el déficit de las AAPP en contabilidad nacional, sino las necesidades de endeudamiento del Estado en términos de caja. Estas incluyen su propio déficit, los préstamos a las comunidades autónomas y a la Seguridad Social, las operaciones financieras y, sobre todo, los vencimientos de la deuda emitida en años anteriores.

Este bloque se organiza en cinco subapartados. El subapartado 3.1 presenta el marco institucional y los objetivos de la política de financiación: quién gestiona la deuda, con qué normas y con qué criterios. El 3.2 analiza la evolución de las necesidades de financiación y su descomposición. El 3.3 estudia los instrumentos: letras, bonos y obligaciones, bonos indexados, bonos verdes y otras modalidades. El 3.4 examina los procedimientos de emisión y el funcionamiento del mercado: subastas, sindicaciones, creadores de mercado y mercado secundario. El 3.5 analiza la interacción con el BCE, desde las compras masivas de 2015-2022 hasta la retirada actual, y la estrategia de 2026.

La tesis del bloque puede anticiparse. Desde la crisis de 2010-2012, el Tesoro español ha seguido una estrategia coherente y reconocida internacionalmente. Ha alargado la vida media de la deuda para reducir el riesgo de refinanciación, ha estandarizado y hecho predecibles sus procedimientos, y ha diversificado su base inversora con instrumentos nuevos. Esa estrategia se desarrolló en un entorno excepcional, en el que el BCE absorbía la mayor parte de la emisión neta y los tipos estaban en mínimos. El reto actual es mantenerla sin ese respaldo, con tipos que vuelven a subir en 2026 y con necesidades brutas cercanas al 20% del PIB, de modo que el mercado privado absorba lo que antes absorbía el banco central sin encarecer en exceso la financiación.

\subsection{Marco institucional y los objetivos de la política de financiación}

Introducción

Antes de estudiar las necesidades, los instrumentos y los procedimientos, hay que responder a tres preguntas previas que el documento apenas aborda. La primera es cuál es el marco jurídico que autoriza y limita el endeudamiento del Estado. La segunda es qué órgano gestiona la deuda y con qué grado de autonomía, una cuestión con un debate propio en la literatura sobre gestión de la deuda. La tercera es qué objetivos persigue esa gestión y cómo se concilian, que es el núcleo conceptual del bloque. A ellas se añade una cuarta, más aplicada: con qué indicadores se mide el cumplimiento de esos objetivos en España.

El subapartado se organiza en cuatro epígrafes: el marco jurídico y su evolución histórica (3.1.1), la organización y el debate sobre las agencias de deuda (3.1.2), los objetivos y su fundamento teórico (3.1.3) y los indicadores de coste y riesgo de la cartera española (3.1.4).

La política de financiación del Tesoro se apoya en un marco jurídico de cuatro niveles. El constitucional y europeo establece la competencia del Estado, la prioridad absoluta del pago de la deuda (art. 135 CE) y la prohibición de la financiación monetaria y del acceso privilegiado (arts. 123 y 124 TFUE). La Ley General Presupuestaria atribuye la creación de deuda (arts. 94 y 98). La Ley de PGE fija el límite anual, que desde 2024 es el de 2023 prorrogado, 96.022 millones. Y la orden anual de creación de deuda, para 2026 la Orden ECM/2/2026, desarrolla instrumentos y procedimientos.

Ese marco es el resultado de una transición histórica. España pasó de la monetización de los setenta y principios de los ochenta a una financiación de mercado, primero apoyada en demanda obligatoria y después, desde los noventa, plenamente de mercado.

La gestión corresponde a la Secretaría General del Tesoro y Financiación Internacional, en el Ministerio de Economía, Comercio y Empresa, no al ministerio que cita el documento. Es un modelo de unidad ministerial con alta autonomía técnica, que la literatura comparada no considera inferior al de agencia si hay mandato claro, transparencia y capacidad técnica. La separación respecto de Hacienda y la relación con el BCE plantean problemas de coordinación que la expansión cuantitativa ha hecho más visibles.

El objetivo canónico, el menor coste a medio y largo plazo con un riesgo prudente, encierra una disyuntiva entre coste y riesgo. Su fundamento teórico está en la nivelación impositiva y la deuda como cobertura fiscal (Barro, Lucas y Stokey, Angeletos) y en las crisis de refinanciación autocumplidas (Cole y Kehoe). Esta última línea justifica el alargamiento de la vida media que el Tesoro ha perseguido desde 2012. Los objetivos complementarios (liquidez, previsibilidad, diversificación y sostenibilidad) entran en tensión entre sí, y su resolución define la estrategia anual.

Los indicadores muestran una cartera larga (unos 8 años de vida media), con poca deuda a corto plazo y un coste medio aún bajo pero en aumento. El coste de emisión (3,16% en 2024) supera ampliamente al coste medio de la cartera, lo que anticipa el aumento de la carga de intereses que analizaré en 3.2 y 3.5.

\subsubsection{Marco jurídico: de la monetización a la financiación de mercado}

a) La evolución histórica: la construcción de una política de deuda

[DOC] El documento describe con detalle el origen de la política de deuda moderna en España. Hasta 1982, la financiación del déficit recayó en buena medida en el BdE: ese año, más del 80% del déficit del Estado se financió con su crédito. El fuerte crecimiento del déficit complicaba una política monetaria orientada a reducir una inflación que superaba el 14% en 1981. Por eso, desde 1982 se reorienta la financiación del Estado hacia mecanismos de mercado, lo que obliga al Tesoro a desarrollar una política de deuda: un conjunto de instrumentos que le permitieran financiarse sin monetización y al menor coste posible.

Según el documento, esa transición exigió además crear demanda donde no la había. En 1984 se estableció un coeficiente de inversión obligatoria en pagarés del Tesoro, del 12% de los pasivos computables de bancos y cajas, que se redujo desde 1989 y se eliminó en 1992. Durante los primeros años el sistema bancario sustituyó al BdE como principal financiador del déficit. En 1990 se restringió legalmente el acceso del Tesoro al crédito del BdE, y en 1994 se prohibió.

Esta secuencia coincide con la que reconstruí en 1.2.2 y 1.2.4 a partir de Escario, Sabaté y Gadea (2011). La monetización alcanzó su máximo en 1978-1983 y terminó en torno a 1984. La transición pasó por instrumentos casi monetarios, como los pagarés, y por la represión financiera, con coeficientes obligatorios de inversión en deuda pública. La prohibición legal llegó con la Ley 13/1994, de Autonomía del BdE, en aplicación del Tratado de Maastricht. Sin embargo, el mecanismo exacto del requisito de 1984 debe verificarse. Parte de la literatura lo describe como un tramo del coeficiente de caja que debía mantenerse en pagarés y no como un coeficiente de inversión independiente. Lo dejo pendiente. Lo esencial para el tema es que la financiación de mercado nació en España apoyada en una demanda obligatoria. Solo en los noventa, con la liberalización financiera, los incentivos a los no residentes y la creación de los creadores de mercado, pasó a ser una financiación verdaderamente de mercado.

En paralelo se construyó la infraestructura del mercado. En 1987 se crearon el Mercado de Deuda Pública en Anotaciones y la Central de Anotaciones del BdE, que eliminaron los títulos físicos. En 1988 se creó la figura de los creadores de mercado, y en 1995 empezó a anunciarse el rango objetivo de colocación de las subastas, según recoge el documento. En 2003, con la Ley 44/2002, el registro de la deuda anotada pasó a Iberclear. El detalle de esta infraestructura se estudia en 3.4.

b) Las normas que autorizan y limitan el endeudamiento

El marco jurídico vigente se articula en cuatro niveles, que conviene conocer porque el tribunal puede preguntar quién autoriza la deuda y dónde está su límite.

El primer nivel es el constitucional y europeo. El artículo 149.1.14.ª de la Constitución atribuye al Estado la competencia exclusiva sobre la Hacienda general y la deuda del Estado. El artículo 135, tras la reforma de 2011, exige que la emisión de deuda esté autorizada por ley y establece que los créditos para pagar los intereses y el capital de la deuda se entenderán siempre incluidos en el estado de gastos de sus presupuestos y su pago gozará de prioridad absoluta. En el ámbito europeo, el artículo 123 del TFUE prohíbe al BCE y a los bancos centrales nacionales conceder crédito a las administraciones públicas o adquirir directamente su deuda en el mercado primario. El artículo 124 prohíbe el acceso privilegiado de las administraciones a las entidades financieras, lo que impide reproducir coeficientes obligatorios como los de los ochenta. Y el artículo 125 recoge la cláusula de no corresponsabilidad financiera (no bail-out).\footnote{Prohibición de financiación monetaria (art. 123 TFUE), su interpretación por el TJUE en los asuntos Gauweiler (C-62/14, 2015, sobre el OMT) y Weiss (C-493/17, 2018, sobre el PSPP), y cláusula de no corresponsabilidad (art. 125): Temas 3A-36/37 y 3B-45.
}

El segundo nivel es la Ley General Presupuestaria. La Ley 47/2003 regula la deuda del Estado en su título IV. Sus artículos 94 y 98 atribuyen al Gobierno, y por delegación al ministro competente en economía, la creación de deuda dentro del límite fijado en la Ley de Presupuestos, y definen las modalidades de emisión. La misma ley regula la gestión de la tesorería del Estado.\footnote{La Ley General Presupuestaria, el ciclo presupuestario y la prórroga de los presupuestos (art. 134.4 CE y art. 38 LGP): Tema 4B-21.
}

El tercer nivel es la Ley de Presupuestos Generales del Estado, que fija cada año el límite al incremento del saldo vivo de la deuda del Estado. Aquí aparece una anomalía que conviene señalar en el examen. Como España funciona con los PGE de 2023 prorrogados (1.5.3), el límite vigente en 2024, 2025 y 2026 es el que fijó el artículo 46 de la Ley 31/2022, de PGE para 2023: un incremento del saldo vivo de la deuda del Estado de 96.021,98 millones de euros. En la práctica, el límite no ha sido restrictivo, porque la emisión neta de 2024-2026 se ha situado muy por debajo. Pero ilustra que la autorización parlamentaria del endeudamiento lleva cuatro ejercicios fijada con una ley pensada para 2023.

El cuarto nivel es la normativa de desarrollo. Cada año, una orden ministerial dispone la creación de deuda del Estado y fija los instrumentos y procedimientos autorizados. Para 2026 es la Orden ECM/2/2026, de 9 de enero, que faculta a la Secretaría General del Tesoro y Financiación Internacional para emitir durante 2026 y enero de 2027. Autoriza letras de hasta 24 meses, bonos de 2 a 5 años, obligaciones de más de 5 años y sus variantes indexadas. Contempla como procedimientos la subasta, la sindicación, las operaciones con creadores de mercado y otras técnicas que se consideren adecuadas. A la orden se suman las resoluciones de la Secretaría General, que publican el calendario de subastas, regulan las condiciones de los creadores de mercado y presentan la Estrategia de Financiación anual.

c) Una corrección institucional

[DOC] El documento atribuye la gestión del endeudamiento a la Dirección General del Tesoro y Política Financiera, una de las direcciones generales del Ministerio de Asuntos Económicos y Transformación Digital.

Corrección. Desde la reestructuración ministerial de noviembre de 2023, las competencias de economía corresponden al Ministerio de Economía, Comercio y Empresa, cuya estructura orgánica desarrolla el Real Decreto 410/2024. La gestión de la deuda corresponde a la Secretaría General del Tesoro y Financiación Internacional, dependiente de la Secretaría de Estado de Economía y Apoyo a la Empresa. La Secretaría General se creó a finales de 2011 como Secretaría General del Tesoro y Política Financiera y adoptó su denominación actual en 2018. Dentro de ella, la Dirección General del Tesoro y Política Financiera es la unidad que gestiona la deuda y la tesorería del Estado, junto con la Dirección General de Financiación Internacional. El documento no se equivoca al nombrar la dirección general, sino al omitir el órgano superior y al citar un ministerio que ya no existe.

Un detalle reciente: el Real Decreto 644/2026, de 28 de julio, permite que la titularidad de la Dirección General del Tesoro y Política Financiera recaiga en empleados públicos que no pertenezcan al subgrupo A1 de funcionarios de carrera. Lo justifica por la elevada especialización técnica y proyección estratégica del puesto, que incluye la gestión de la deuda, la regulación financiera y la prevención del blanqueo. Es un movimiento hacia el perfil técnico de las agencias de deuda que se discute en el epígrafe siguiente.

\subsubsection{Organización: ¿ministerio o agencia de deuda?}

a) Lo que dice el documento

[DOC] En una nota, el documento señala que en otros países, como Francia, Alemania o el Reino Unido, la gestión del endeudamiento corresponde a agencias independientes que gozan, al menos en teoría, de mayor autonomía. Añade que en España el Tesoro depende formalmente del ministerio, pero opera con gran autonomía y guiado en todo momento por criterios financieros.

b) El debate en la literatura

La nota del documento apunta a un debate que la literatura sobre gestión de la deuda soberana ha tratado con detalle. Desde finales de los ochenta, varios países trasladaron la gestión de la deuda a agencias especializadas. Irlanda creó la National Treasury Management Agency en 1990. El Reino Unido creó la Debt Management Office en 1998, tras conceder la independencia al Banco de Inglaterra, que hasta entonces gestionaba la deuda. Francia creó la Agence France Trésor en 2001, como servicio con competencia nacional dentro del Ministerio de Economía. Alemania creó la Finanzagentur en 2000-2001, como sociedad pública de responsabilidad limitada. Portugal tiene el IGCP, y Suecia la Riksgälden, con orígenes en el siglo XVIII. En cambio, Italia y España mantienen la gestión dentro de una unidad del ministerio.

Los argumentos a favor de la agencia son cuatro. Primero, la especialización técnica y la capacidad de contratar perfiles de mercado con retribuciones competitivas. Segundo, la separación entre las decisiones de gestión de la deuda y las presiones políticas de corto plazo, como la tentación de emitir a corto plazo para abaratar el coste inmediato a costa de más riesgo. Tercero, un mandato explícito y medible que facilita la rendición de cuentas, en la línea de la independencia de los bancos centrales. Cuarto, la separación nítida respecto del banco central, que evita conflictos entre política de deuda y política monetaria. Los argumentos a favor de la gestión ministerial son la mejor coordinación con la política presupuestaria y con la gestión de la tesorería, la responsabilidad política directa ante el Parlamento y el menor coste administrativo.

La conclusión de la literatura comparada (Currie, Dethier y Togo, 2003, para el Banco Mundial; las Directrices para la gestión de la deuda pública del FMI y el Banco Mundial, revisadas en 2014; y los trabajos de la OCDE) es que no hay un modelo institucional superior. Lo decisivo es que existan un mandato claro, objetivos explícitos de coste y riesgo, una separación operativa respecto de la política monetaria, transparencia y capacidad técnica. El Tesoro español es un caso habitualmente citado de unidad ministerial con elevada autonomía técnica y reputación en los mercados. Así lo reflejan los informes de la OCDE sobre gestión de la deuda, la valoración de las agencias de calificación y la ausencia de episodios de gestión oportunista. El documento tiene razón en este punto.

Hay, sin embargo, un argumento específico del caso español que el documento no menciona. Desde 2018, con la separación entre el Ministerio de Hacienda (presupuesto, ingresos y financiación territorial) y el de Economía (Tesoro), la gestión de la deuda y la gestión presupuestaria dependen de ministerios distintos. Esto aproxima de hecho el modelo español al de agencia en su separación de la política presupuestaria, aunque sin su estatuto jurídico. Esa separación exige coordinación en tres frentes: las previsiones de ingresos y gastos de caja, que determinan las necesidades; la financiación de los mecanismos autonómicos, gestionados por Hacienda pero financiados por el Tesoro; y los préstamos a la Seguridad Social. No he encontrado evaluaciones públicas de esa coordinación.

c) La relación con el banco central: un principio en revisión

El principio de separación entre gestión de la deuda y política monetaria es uno de los pilares de las directrices internacionales. El gestor de la deuda decide los plazos y volúmenes con criterios de coste y riesgo, y el banco central fija los tipos con criterios de estabilidad de precios. La expansión cuantitativa de 2009-2022 ha cuestionado ese principio en la práctica. Cuando el banco central compra deuda a largo plazo para reducir la prima por plazo, y el Tesoro alarga la vida media para aprovechar los tipos bajos, las dos políticas actúan en direcciones opuestas sobre la duración de la deuda en manos del sector privado. Greenwood, Hanson, Rudolph y Summers (2014) documentaron ese conflicto en Estados Unidos: el Tesoro alargaba la vida media mientras la Reserva Federal la acortaba en manos privadas con sus compras. Blommestein y Turner (2012), para la OCDE, advirtieron del riesgo de que la política de deuda se subordine a la monetaria, o al revés, en contextos de dominancia fiscal.

En la zona del euro el problema tiene una dimensión adicional, porque hay un solo banco central y veinte gestores de deuda. El BCE compró deuda de cada país según su clave de capital con el APP y el PEPP, y desde 2022 cuenta con el TPI (2.3.3). El Tesoro español alargó la vida media mientras el Eurosistema absorbía más de la mitad de su emisión neta, lo que el documento recoge y analizaré en 3.5.\footnote{La interacción entre la política monetaria no convencional y la gestión de la deuda (efecto de las compras de activos sobre la prima por plazo y el canal de reequilibrio de carteras): Temas 3A-36/37.
}

\subsubsection{Objetivos de la política de financiación y su fundamento teórico}

a) El objetivo canónico: coste mínimo con un riesgo prudente

El documento no formula expresamente el objetivo del Tesoro. La formulación de referencia internacional es la de las Directrices para la gestión de la deuda pública del FMI y el Banco Mundial (2014): asegurar que las necesidades de financiación del Estado y sus obligaciones de pago se cubran al menor coste posible a medio y largo plazo, con un grado de riesgo prudente. El Tesoro español la hace suya en sus documentos de estrategia: cubrir las necesidades de financiación del Estado al menor coste posible a medio y largo plazo, con un nivel de riesgo prudente, y contribuir al buen funcionamiento del mercado de deuda.

La formulación contiene tres elementos que conviene explicar. El primero es el horizonte. El coste relevante no es el del año, sino el de medio y largo plazo. Emitir todo a muy corto plazo abarata normalmente el coste inmediato, porque la curva de tipos suele tener pendiente positiva, pero expone al Estado a refinanciar toda su deuda a los tipos que haya en el futuro.

El segundo elemento es el riesgo, en varias dimensiones. El riesgo de refinanciación es la posibilidad de no poder renovar los vencimientos, o de hacerlo en condiciones muy desfavorables. El riesgo de tipo de interés es la sensibilidad del coste de la cartera a las variaciones de los tipos. El riesgo de tipo de cambio, hoy prácticamente inexistente en España porque toda la deuda está en euros, fue relevante en los noventa (1.2.4). El riesgo de inflación afecta a los bonos indexados. El riesgo de liquidez es la dificultad de colocar deuda en un mercado poco profundo. Y el riesgo operativo abarca los fallos de procedimiento.

El tercer elemento es la disyuntiva entre coste y riesgo. Alargar la vida media reduce el riesgo de refinanciación y el de tipo de interés. Pero normalmente eleva el coste esperado, porque los inversores exigen una prima por plazo por inmovilizar su dinero más tiempo. El gestor elige un punto de esa frontera en función de su aversión al riesgo, y ese punto es el núcleo de su estrategia. El FMI y el Banco Mundial han formalizado este análisis en su marco de Estrategia de Gestión de la Deuda a Medio Plazo (MTDS, 2009), que compara carteras alternativas por su coste esperado y su coste en escenarios adversos.

b) El fundamento teórico: la deuda como instrumento de cobertura fiscal

Detrás del objetivo canónico hay una literatura teórica que conviene conocer, porque explica por qué los gestores prefieren la deuda larga y por qué emiten bonos indexados. La expongo solo en lo que ayuda a entender la estrategia española y remito su desarrollo al tema correspondiente.\foonote{La teoría de la gestión óptima de la deuda (nivelación impositiva, deuda contingente, estructura de plazos e indexación) y la de las crisis de deuda autocumplidas se desarrollan en el Tema 3A-39. Recordatorio: con impuestos distorsionantes de coste convexo, el gobierno minimiza la pérdida de eficiencia manteniendo constante el tipo impositivo esperado, lo que convierte a la deuda en el amortiguador de las perturbaciones presupuestarias (Barro, 1979).
}

El punto de partida es la nivelación impositiva de Barro (1979). Si los impuestos generan distorsiones crecientes con el tipo, el gobierno debe mantener tipos estables en el tiempo y absorber las perturbaciones del presupuesto con la deuda. De ahí se deriva una pregunta de gestión de la deuda: qué estructura de deuda (plazos, indexación) protege mejor al presupuesto frente a las perturbaciones.

Lucas y Stokey (1983) mostraron que, idealmente, el gobierno emitiría deuda contingente, cuyo rendimiento bajaría cuando el presupuesto sufre una perturbación adversa y subiría en caso contrario. Como esa deuda no existe, la literatura ha buscado aproximaciones. Angeletos (2002) y Buera y Nicolini (2004) mostraron que una combinación adecuada de plazos puede replicar esa contingencia, porque el valor de la deuda larga cae cuando suben los tipos. Barro (2003) concluyó que, en ausencia de deuda contingente, la mejor aproximación es la deuda larga indexada a la inflación. Este es el fundamento del programa de bonos indexados del Tesoro (2014), cuyo atractivo se basa, como dice el documento, en que los ingresos públicos están fuertemente ligados a la inflación. El documento añade una advertencia pertinente: la cobertura se debilita cuando el gasto también se indexa, como ha ocurrido con las pensiones desde 2021.

Una segunda línea destaca la credibilidad. Calvo y Guidotti (1990) y Missale y Blanchard (1994) mostraron que un gobierno con mucha deuda nominal a largo plazo tiene incentivos a generar inflación para reducir su valor real. Por eso, en contextos de baja credibilidad antiinflacionista, la deuda corta o indexada es una señal de compromiso. En la zona del euro, sin política monetaria nacional, ese argumento pierde fuerza para los Estados miembros: España no puede inflar su deuda, porque no controla la política monetaria.

Una tercera línea, la más relevante para la experiencia española de 2010-2012, es la de las crisis de refinanciación autocumplidas. Cole y Kehoe (2000) mostraron que un país con mucha deuda a corto plazo es vulnerable a una crisis en la que los inversores, temiendo que no pueda refinanciarse, se niegan a hacerlo, y provocan así el impago que temían. Alargar la vida media reduce la proporción de deuda que hay que refinanciar cada año y, con ella, la vulnerabilidad a este tipo de crisis. Arellano y Ramanarayanan (2012) formalizaron la disyuntiva: la deuda corta da mejores incentivos al gobierno, porque le disciplina, y es más barata; la larga le protege mejor frente a los pánicos. La elección óptima depende del riesgo de crisis. Esta literatura es el fundamento teórico de la estrategia seguida por el Tesoro español desde 2012: alargar la vida media hasta los 8 años.

c) Los objetivos complementarios y sus tensiones

Al objetivo principal se añaden otros que el Tesoro formula expresamente en sus estrategias y que conviene presentar con sus tensiones internas.

El primero es la liquidez y el buen funcionamiento del mercado. El Tesoro concentra la emisión en pocas referencias de gran volumen, que el documento sitúa en torno a 25.000 millones para bonos y obligaciones, para garantizar su negociabilidad en el mercado secundario. Una deuda más líquida es más atractiva y, por tanto, más barata. Además, la curva de deuda pública es la referencia para la deuda privada.

El segundo es la transparencia y la previsibilidad. El Tesoro publica en enero su estrategia anual, con el calendario de subastas y los rangos objetivo de colocación, y comunica con antelación sus cambios. La lógica es que la previsibilidad reduce la incertidumbre de los inversores y, con ella, la prima que exigen. Es la aplicación a la gestión de la deuda del principio de no sorprender, que en la literatura se asocia a la experiencia de las agencias británica y sueca.

El tercero es la diversificación de la base inversora. Una base amplia y variada (bancos nacionales, aseguradoras, fondos de pensiones, bancos centrales extranjeros, fondos soberanos, minoristas) reduce el riesgo de que la retirada de un grupo de inversores desestabilice la financiación. Esta fue la lección de 2011-2012, cuando la salida de los no residentes dejó a la banca nacional como principal comprador (1.4.2d). Los programas de bonos indexados (2014) y de bonos verdes (2021) responden en parte a este objetivo, como señala el documento.

El cuarto, más reciente, es la sostenibilidad en sentido ambiental. El programa de bonos verdes vincula la emisión a gastos elegibles del presupuesto.

Estos objetivos no siempre son compatibles. La liquidez exige pocas referencias grandes, mientras que la diversificación empuja a crear instrumentos nuevos, como los indexados y los verdes, que fragmentan la oferta y reducen la liquidez de cada segmento. La previsibilidad exige atenerse al calendario, mientras que la minimización del coste invitaría a aprovechar las ventanas de mercado favorables. Y el alargamiento de la vida media reduce el riesgo, pero, con una curva positiva, eleva el coste esperado. La estrategia de cada año es, en esencia, la resolución de estas tensiones. Lo veremos aplicado en 3.3, 3.4 y 3.5.

\subsubsection{Indicadores de coste y riesgo de la cartera española}

Los objetivos anteriores se traducen en un conjunto de indicadores que el Tesoro publica y que permiten evaluar su gestión. Los presento aquí de forma breve, como marco de referencia, porque su evolución se analiza en los subapartados siguientes.

Los indicadores de coste son dos. El coste medio de emisión mide el tipo al que se ha emitido la deuda en el año. Según el documento, en 2021 fue negativo por primera vez en la historia del Tesoro (−0,04%). Con la subida de tipos pasó a superar el 3% en 2023 y a situarse en el 3,16% en 2024, según CaixaBank Research a partir de la Estrategia del Tesoro para 2025. El coste medio de la deuda en circulación mide el tipo medio del conjunto de la cartera. Según el documento, alcanzó su mínimo histórico, del 1,64%, en 2021. Desde entonces sube de forma gradual, a medida que la deuda antigua y barata se refinancia a tipos más altos, y se situaría en torno al 2,2-2,5% en 2025. La diferencia entre ambos indicadores es la clave de la dinámica de la carga de intereses. Mientras el coste de emisión supere al coste medio de la cartera, este seguirá subiendo aunque los tipos de mercado se estabilicen.

Los indicadores de riesgo son varios. La vida media de la deuda del Estado alcanzó por primera vez los 8 años en 2021, según el documento, y se ha mantenido en torno a ese nivel. En 2024 era de unos 8 años, frente a 8,9 en la zona del euro, según el BdE (2.3.2). Los demás indicadores son la proporción de deuda a corto plazo y la proporción de deuda que vence en los próximos doce meses. Según el documento, las letras representan menos del 6% de la deuda del Estado, frente al 62% de 1990. Según el BdE, el 13,9% de la deuda de las AAPP vencía en menos de un año en 2024. Completan el cuadro el perfil de vencimientos, es decir, la concentración de amortizaciones en determinados años, y la composición de la base inversora. El Eurosistema tendría en torno a un 25-26% de la deuda en 2025, según la misma fuente, y los no residentes, en torno al 44%.

A estos indicadores la literatura añade el plazo medio hasta la revisión del tipo (ATR, por sus siglas en inglés), que mide cuánto tarda la cartera en trasladar un cambio de tipos al coste. En España coincide prácticamente con la vida media, porque casi toda la deuda es a tipo fijo, con la excepción de los bonos indexados, cuyo principal se revisa con la inflación. Este dato es el que explica por qué la subida de tipos de 2022-2023 tardó en trasladarse a la carga de intereses (1.5.2a). También explica por qué seguirá trasladándose en 2026-2030, con la amplificación que suponen las subidas del BCE de junio y septiembre de 2026 (1.5.3c).


El documento base mide las necesidades del Tesoro con dos cifras de 2022 que hoy inducen a error. Dice que las necesidades brutas rondan los 250.000 millones anuales, algo más del 20% del PIB. En 2026 la emisión bruta prevista es de unos 285.700 millones, pero con el PIB nominal actual eso ronda el 16% del PIB, no el 20%. Hay además un hecho más importante que el documento no podía prever. La emisión neta del Tesoro lleva tres años estable en 55.000 millones, pero la deuda que el mercado privado tiene que absorber cada año es mucho mayor. El Eurosistema ya no reinvierte sus vencimientos, así que la oferta neta de deuda que llega al sector privado ha pasado de ser negativa en 2015-2021 a rondar probablemente los 90.000-100.000 millones anuales.

Si en el examen analizas las necesidades del Tesoro solo con la emisión neta y la bruta, omites la variable que más ha cambiado desde 2022. Por eso este subapartado distingue tres magnitudes (necesidades netas, brutas y oferta neta al sector privado) y explica qué determina cada una.

\subsection{Necesidades de financiación del Tesoro: evolución y determinantes}

Introducción

[DOC] El documento abre este apartado con una constatación: en los últimos quince años las necesidades de financiación del Tesoro han aumentado mucho. En 2007 la necesidad bruta fue de unos 50.000 millones y la emisión neta, negativa en casi 10.000 millones. Desde entonces la emisión neta ha sido positiva sin interrupción, y superó los 100.000 millones en 2009 y en 2020. Las necesidades brutas se han disparado, por el aumento de la emisión neta y de los vencimientos, hasta unos 250.000 millones anuales.

Este subapartado desarrolla y actualiza esa constatación en cuatro epígrafes. El primero precisa los conceptos: qué son las necesidades netas y brutas y qué relación guardan con el déficit (3.2.1). El segundo analiza sus determinantes: por qué la emisión neta del Tesoro no coincide con el déficit y qué papel desempeñan los préstamos a otras administraciones, los fondos europeos y la tesorería (3.2.2). El tercero recorre su evolución entre 2007 y 2026 (3.2.3). El cuarto introduce el concepto de oferta neta al sector privado y examina las presiones sobre las necesidades en los próximos años (3.2.4).

Las necesidades del Tesoro se miden con tres magnitudes. Las netas equivalen a la variación del saldo de deuda del Estado y cubren su déficit de caja, los préstamos a las comunidades y a la Seguridad Social, las operaciones financieras y la variación de la tesorería, menos la financiación europea. Las brutas añaden los vencimientos, que dependen del saldo y de la vida media. La oferta neta al sector privado añade a las netas los vencimientos que el Eurosistema ya no reinvierte.

La evolución muestra una explosión en 2008-2013: una emisión neta de más de 100.000 millones en 2009, el recurso a las letras y el inicio de la financiación a las comunidades. Siguió una normalización en 2014-2019, con el BCE absorbiendo más que la emisión neta, y un máximo de 277.000 millones de emisión bruta en 2020, amortiguado por el SURE, el Next Generation EU y el PEPP. Desde 2021, la emisión neta ha bajado hasta estabilizarse en 55.000 millones (2024-2026), mientras la bruta ha subido hasta unos 285.700 millones en 2026, en torno al 16% del PIB, por el aumento de los vencimientos.

El rasgo estructural del caso español es que el Estado actúa como prestamista de las demás administraciones, sobre todo de la Seguridad Social, uno de los principales destinos de la emisión neta. El cambio más relevante desde 2022 es que la oferta neta al sector privado ha pasado de ser negativa a rondar probablemente los 90.000-100.000 millones anuales. Esa deuda la están absorbiendo sobre todo inversores no residentes (47,8% de la cartera en 2025). Las presiones futuras (vencimientos a tipos más altos, fin de la financiación europea, pensiones, financiación autonómica y nuevos gastos) anticipan necesidades elevadas y más caras.

\subsubsection{Conceptos: necesidades netas, necesidades brutas y su medición}

a) Del déficit a las necesidades netas

Como vimos en 1.1.4, el Tesoro no financia el déficit de las AAPP en contabilidad nacional, sino las necesidades de endeudamiento del Estado, una magnitud de caja referida solo al Estado. La emisión neta, es decir, el aumento del saldo vivo de la deuda del Estado en el año, cubre la suma de cuatro componentes.

El primero es el déficit de caja del Estado, que difiere del déficit del Estado en contabilidad nacional por los ajustes de devengo y por la periodificación de los pagos. El segundo son los préstamos netos del Estado a otras administraciones, principalmente a las comunidades autónomas, a través del Fondo de Financiación a CCAA, y a la Seguridad Social. El tercero son las demás operaciones financieras del Estado: aportaciones a fondos y organismos, préstamos a otros países o instituciones, adquisición de participaciones y, en sentido contrario, los ingresos por la amortización de préstamos concedidos o la venta de activos. El cuarto es la variación del saldo de tesorería del Estado, su colchón de liquidez en el BdE y en el mercado.

A esa suma hay que restar la financiación que el Estado obtiene fuera del mercado, que en los últimos años ha procedido sobre todo de la Unión: los préstamos del instrumento SURE y los préstamos del Mecanismo de Recuperación. Las subvenciones del Mecanismo de Recuperación no son financiación, sino ingresos del Estado, pero su efecto de caja es el mismo: reducen lo que el Tesoro necesita emitir.

b) De las necesidades netas a las brutas

Las necesidades brutas añaden a las netas los vencimientos de la deuda emitida en años anteriores, que hay que refinanciar. En la estrategia de 2026, por ejemplo, la emisión neta prevista es de 55.000 millones: 50.000 a medio y largo plazo y 5.000 en letras. La emisión bruta de bonos y obligaciones asciende a 176.935 millones, resultado de sumar a la neta unos 126.935 millones de amortizaciones. La emisión bruta de letras ronda los 108.700 millones, porque cada año se renueva todo el saldo de letras. La emisión bruta total ronda por tanto los 285.700 millones. (OJO: la nota oficial del Gobierno da 285.677 millones y algunas fuentes de prensa 285.693. La diferencia está en el importe de letras y es irrelevante para el análisis, pero conviene citar la oficial.)

Los vencimientos dependen de dos variables: el saldo vivo y su vida media. Con una vida media de unos 8 años, cada año vence en promedio alrededor de una octava parte de la deuda a medio y largo plazo, lo que concuerda con el en torno al 13% del saldo que el Tesoro refinancia anualmente. A ello se suman las letras, que vencen todas en menos de un año. Por eso las necesidades brutas crecen de forma casi automática con el saldo de deuda, aunque la emisión neta sea constante. Este fenómeno explica que la emisión bruta haya pasado de 232.570 millones en 2022 a unos 285.700 en 2026 con una emisión neta decreciente.

c) La medición en porcentaje del PIB y los umbrales de riesgo

[DOC] El documento cifra las necesidades brutas en algo más del 20% del PIB. Actualización. La emisión bruta del Estado prevista para 2026 equivale a en torno al 16% del PIB nominal. La proporción ha bajado porque el PIB nominal ha crecido más deprisa que los vencimientos. Pero conviene distinguir esa cifra de las necesidades brutas de financiación de las AAPP, que incluyen además el déficit y los vencimientos de las comunidades y de las entidades locales que se financian en el mercado. La Comisión las situaba cerca del 20% del PIB en su Debt Sustainability Monitor de 2024, entre las más altas de la Unión junto a Italia, Francia y Bélgica (2.3.2).

La cifra importa por su relación con el riesgo de refinanciación. El marco de análisis de sostenibilidad del FMI para países con acceso a los mercados utiliza como referencia para las economías avanzadas unas necesidades brutas del 15% del PIB. Por encima de ese umbral, la vulnerabilidad a un cierre de los mercados aumenta de forma significativa. España supera el umbral en términos de AAPP y lo roza con el Estado solo. Es uno de los motivos por los que la Comisión y el FMI insisten en mantener la vida media larga, porque cada año de vida media adicional reduce los vencimientos anuales y, con ellos, las necesidades brutas.\footnote{Recordatorio: necesidades brutas de financiación (GFN) = déficit + amortizaciones de la deuda a medio y largo plazo + saldo de la deuda a corto plazo. El FMI utiliza, en su marco de análisis de sostenibilidad para países con acceso a los mercados, umbrales de alerta del 15% del PIB para las economías avanzadas. Desarrollo de los indicadores de sostenibilidad y de riesgo de refinanciación en el Tema 3A-39.
}

\subsubsection{Determinantes: por qué el Tesoro emite más de lo que dice el déficit}

a) El Estado como prestamista de las demás administraciones

El rasgo más característico de las necesidades del Tesoro español es que financia no solo el déficit del Estado, sino una parte del de las comunidades autónomas y de la Seguridad Social. Como vimos en 1.1.1 y en el Tema 25, desde 2012 la mayoría de las comunidades se financian a través de los mecanismos estatales (FFPP, FLA y, desde 2015, el Fondo de Financiación a CCAA), en lugar de acudir al mercado. El Estado les presta con cargo a la emisión del Tesoro. Entre 2012 y 2024, esos mecanismos movilizaron 126.895 millones. La Seguridad Social, por su parte, cubre su déficit desde 2017 con préstamos del Estado, además de con transferencias. En 2025 su deuda, casi toda con el Estado, aumentó en unos 10.000 millones, y la de las comunidades en unos 5.700 millones (1.1.4).

Esto tiene tres consecuencias para la política de financiación. La primera, contable: los préstamos del Estado a otras administraciones no alteran la deuda consolidada de las AAPP, porque la deuda del prestatario con el Estado se elimina al consolidar, pero sí aumentan la deuda del Estado y, por tanto, lo que el Tesoro emite. Por eso la deuda del Estado es casi el 90% del total de las AAPP, como recoge el documento. La segunda, financiera: las comunidades y la Seguridad Social se financian al coste del Tesoro, inferior al que obtendrían en el mercado, lo que reduce el coste agregado. La tercera, de incentivos, y es la más debatida. Los mecanismos debilitan la restricción presupuestaria de las comunidades, que ya no afrontan el juicio del mercado, y generan el riesgo moral que analizamos en el Tema 25 con la literatura de la restricción presupuestaria blanda (Kornai, Rodden).\foonote{Mecanismos de financiación autonómica (FFPP, FLA, Facilidad Financiera y FFCCAA), restricción presupuestaria blanda y condonación: Tema 4B-25 (ya redactado). Préstamos del Estado a la Seguridad Social: Tema 4B-24 (ya redactado).
}

La condonación de 83.252 millones de deuda autonómica en tramitación (1.5.3d) no aumentará la emisión del Tesoro, porque la mayor parte de la deuda condonada es deuda con el propio Estado a través del FFCCAA. Lo que hará es convertir un activo del Estado (el derecho de cobro) en nada. El efecto sobre las necesidades será diferido: el Estado dejará de recibir las amortizaciones que las comunidades habrían devuelto en el futuro, y eso aumentará marginalmente sus necesidades netas en los próximos años. No he encontrado una estimación pública de ese efecto sobre la programación del Tesoro.

b) Los fondos europeos: financiación fuera del mercado

[DOC] El documento destaca que el diseño de la respuesta europea a la pandemia contuvo el impacto sobre las necesidades de emisión de España. Sin los préstamos del SURE (21.326 millones en 2020-2021) y las transferencias del programa Next Generation EU (31.036 millones en 2021-2022), las emisiones habrían sido mayores. Añade que esas transferencias se complementarían en 2023-2026 con préstamos reembolsables.

La actualización confirma el mecanismo. Según la Estrategia de 2026, España movilizó unos 16.000 millones de préstamos europeos en 2025 y ha solicitado 6.500 millones para 2026, que complementan la emisión del Tesoro. En conjunto, las subvenciones y los préstamos del Mecanismo de Recuperación han reducido las necesidades de mercado del Tesoro en varias decenas de miles de millones en 2021-2026.

Dos observaciones analíticas completan el dato. La primera es que los préstamos europeos son deuda del Estado español a todos los efectos. Reducen la emisión en el mercado, pero no la deuda. Su ventaja es el coste, cercano al de la Unión, que se financia a tipos inferiores a los españoles, y los plazos largos. Por eso el Gobierno solicitó solo una parte de los préstamos disponibles (1.5.2d): cuando el diferencial entre España y la Unión se estrechó, el ahorro dejó de compensar la condicionalidad y la gestión asociadas. La segunda es que la financiación europea desaparece como amortiguador a partir de 2027. Los préstamos del SURE y del Mecanismo de Recuperación empezarán a amortizarse, y las subvenciones dejarán de llegar. La Estrategia de 2026 incluye además una amortización de unos 3.643 millones del préstamo del MEDE de 2012, que termina de devolverse en 2027 (1.4.2d).

c) La tesorería: el colchón de liquidez

El cuarto determinante suele omitirse, pero en algunos años ha sido cuantitativamente relevante: la variación del saldo de tesorería del Estado. El Tesoro mantiene un colchón de liquidez para cubrir los desfases entre cobros y pagos a lo largo del año, afrontar los grandes vencimientos y no depender de una sola subasta en un momento de tensión. En 2020 lo aumentó sustancialmente como precaución, lo que elevó la emisión neta por encima de lo que exigía el déficit. En años posteriores lo ha reducido, lo que ha permitido emitir menos.

Las reglas del BCE sobre la remuneración de los depósitos de las administraciones en los bancos centrales influyen en esa gestión. Desde diciembre de 2024 están remunerados, como máximo, al €STR menos 20 puntos básicos, lo que incentiva a los Tesoros a colocar sus excedentes en el mercado en lugar de mantenerlos en el banco central. Por eso el Tesoro ha desarrollado sus operaciones de colocación de liquidez (subastas de depósitos y adquisiciones temporales), que se estudian en 3.4. Como vimos en 1.1.3, estas variaciones de tesorería aparecen en la deuda como ajustes déficit-deuda: aumentan la deuda bruta sin aumentar el déficit.

d) Una aproximación a la descomposición de 2026

Con estos elementos puede ensayarse una descomposición de la emisión neta de 2026, que ilustra la bisagra entre la primera y la tercera parte del tema. Esta estimación de orden de magnitud es mía, a partir de las cifras disponibles, y no la ha publicado el Tesoro. De los 55.000 millones de emisión neta, en torno a la mitad correspondería al déficit de caja del Estado propiamente dicho. Una parte importante, del orden de 10.000 a 15.000 millones, correspondería a los préstamos a la Seguridad Social. Otra, menor y variable, a la financiación neta de las comunidades a través del FFCCAA, una vez descontadas sus amortizaciones. El resto correspondería a operaciones financieras. Los préstamos europeos (6.500 millones) y la amortización del MEDE (3.643 millones) actúan en sentido contrario el uno del otro. Lo relevante para el examen no es la cifra exacta, que queda pendiente de verificar con la IGAE y el Tesoro, sino la estructura: la Seguridad Social se ha convertido en uno de los principales destinos de la emisión neta del Tesoro. Es la manifestación financiera del problema de las pensiones analizado en el Tema 24.

\subsubsection{Evolución, 2007-2026}

a) 2007-2013: la explosión de las necesidades

[DOC] En 2007 la necesidad bruta fue de unos 50.000 millones, con una emisión neta negativa de casi 10.000 millones: el Estado tenía superávit y amortizaba deuda. Es la consecuencia financiera de los superávits de 2005-2007 analizados en 1.3. Con la Gran Recesión, la emisión neta superó los 100.000 millones en 2009.

[DOC] En nota, el documento recoge la respuesta del Tesoro. En 2009, con un acceso difícil a los mercados, aumentó las emisiones a corto plazo, más baratas y más fáciles de colocar en condiciones de tensión, y cubrió casi el 30% de sus necesidades netas con letras, que llegaron a representar el 18% de la deuda del Estado. Ese mismo año emitió un bono a tipo variable (Euribor a 12 meses menos 10 puntos básicos) para diversificar la base inversora, una modalidad que hoy no utiliza.

Entre 2010 y 2013 las necesidades se mantuvieron muy elevadas por tres motivos. Los déficits siguieron por encima del 9% del PIB. Desde 2012 el Estado empezó a prestar a las comunidades a través de los mecanismos de liquidez, que sumaron varias decenas de miles de millones al año. Y el Tesoro tuvo que refinanciar vencimientos crecientes en el peor momento de la crisis soberana. Hubo, además, una financiación que no pasó por el Tesoro: los 41.333 millones del MEDE para la recapitalización bancaria se canalizaron a través del FROB, en forma de bonos del propio MEDE, y computaron como deuda de las AAPP (1.4.2d). En 2012 la demanda se sostuvo gracias a la banca nacional, que absorbió deuda con la liquidez de las subastas a tres años del BCE, cuando los no residentes redujeron su participación al 33% (1.4.2d).

b) 2014-2019: la normalización con el BCE como comprador

[DOC] Según el documento, el aumento de las necesidades coincidió con una caída sin precedentes del coste de emisión. La explica por el cambio del patrón de crecimiento hacia uno más equilibrado, que mejoró la calificación y redujo la prima de riesgo, y por los estímulos del BCE. En nota añade que, entre 2015 y 2022, las compras netas de deuda española del BCE con el APP y el PEPP superaron los 500.000 millones, más de la tercera parte de la deuda de las AAPP a finales de 2022. Descontada la deuda en manos del BCE, la deuda de las AAPP a finales de 2022 equivalía a la de principios de 2013.

En 2015-2019 la emisión neta se redujo de forma gradual con el déficit, mientras la bruta se mantenía elevada por los vencimientos. Las compras del programa PSPP superaron en varios años la emisión neta del Tesoro, lo que significa que la deuda en manos del sector privado disminuyó. Es la clave para entender el coste de emisión negativo de 2021 que señala el documento, y la analizaré en 3.5. No he podido reconstruir la serie anual de emisión neta de 2014-2019 con fuentes accesibles desde aquí, y queda pendiente.

c) 2020: el máximo histórico

[DOC] En enero de 2020, antes de la pandemia, el Tesoro esperaba emitir en bruto menos de 200.000 millones por primera vez en casi diez años. La caída de los ingresos y el aumento del gasto obligaron a emitir 277.000 millones, la cifra más alta de su historia hasta entonces, con una emisión neta superior a los 100.000 millones. El Tesoro respondió con el mismo instrumento que en 2009, aumentando las letras al principio de la crisis, pero con un contexto muy distinto: el PEPP absorbió una parte sustancial de la emisión neta de 2020-2021 y el coste de emisión siguió cayendo. Ese mismo año se aumentó también el colchón de tesorería.

d) 2021-2026: la emisión neta baja y la bruta sube

Desde 2021 la emisión neta se ha reducido de forma escalonada: 70.063 millones en 2022, unos 65.000 en 2023, 55.000 previstos en 2024 y 55.000 en 2025 (el plan inicial de 60.000 se recortó en septiembre de 2025 por la mejora de los ingresos) y 55.000 en 2026. En cambio, la emisión bruta ha crecido: 232.570 millones en 2022, unos 252.000 en 2023, unos 257.600 previstos en 2024, unos 274.200 en 2025 y unos 285.700 en 2026. La divergencia se explica por los vencimientos. El saldo de deuda crece cada año en la cuantía de la emisión neta, y además vencen las emisiones de 2016-2021, que fueron muy abultadas.

En la presentación de la Estrategia de 2026, el Tesoro destacó que había mantenido la vida media en torno a los 8 años por quinto año consecutivo (7,93 años a finales de 2025). El coste medio de la deuda en circulación era del 2,31% en 2025, 10 puntos básicos más que en 2024. El coste medio de la nueva emisión, del 2,70% en 2025, era 74 puntos básicos inferior al de 2023 (3,44%). Los no residentes tenían el 47,8% de la deuda, el segundo registro más alto de la serie. La brecha entre el coste de emisión (2,70%) y el de la cartera (2,31%) implica que el coste medio seguirá subiendo gradualmente, más aún con las subidas del BCE de 2026, como anticipé en 3.1.4.

\subsubsection{Oferta neta al sector privado y las presiones de los próximos años}

a) Un concepto necesario: la oferta neta al sector privado

La emisión neta mide cuánto aumenta el saldo de deuda del Estado, pero no cuánta deuda nueva tiene que absorber el mercado privado. Para ello hay que tener en cuenta lo que hace el Eurosistema. Cuando el BCE compra deuda en el mercado secundario por encima de la emisión neta, como ocurrió en 2015-2021, la oferta neta al sector privado es negativa: los inversores privados reducen sus tenencias. Cuando el BCE deja de reinvertir los vencimientos de sus carteras, la deuda que vence en su poder tiene que ser refinanciada por el Tesoro y colocada en el mercado. La oferta neta al sector privado se aproxima entonces a la emisión neta más los vencimientos no reinvertidos del Eurosistema.

El Eurosistema tenía en torno al 25-26% de la deuda española en 2025, según CaixaBank Research a partir de datos del Tesoro. Desde que dejó de reinvertir los vencimientos del APP (julio de 2023) y del PEPP (finales de 2024), sus carteras se reducen cada año por el importe de lo que vence. Si ese importe ronda los 35.000-45.000 millones anuales para la deuda española, la oferta neta al sector privado se situaría en torno a 90.000-100.000 millones anuales en 2025-2026, frente a cifras negativas en 2015-2021. Mi estimación es de orden de magnitud y está pendiente de verificar con los datos de tenencias del BdE.

Este cambio explica dos datos que el Tesoro destaca: que los inversores internacionales se han consolidado como los principales compradores de bonos y obligaciones del Estado, en sustitución del BCE, y que la proporción de no residentes ha subido al 47,8%. La emisión sindicada a diez años de enero de 2026, de 15.000 millones, recibió una demanda récord de 145.000 millones. Los no residentes tomaron el 90%, y los bancos centrales e instituciones oficiales fueron el primer grupo inversor (26,5%). La demanda ha respondido. Pero el precio de esa respuesta es una base inversora más sensible a las condiciones financieras globales y a la comparación con otros emisores, lo que remite a la fragilidad que se vio en 2011-2012. Lo analizo en 3.5.

b) Las presiones sobre las necesidades en los próximos años

Mirando hacia delante, [DOC] el documento anticipaba que las necesidades brutas se mantendrían muy elevadas por la necesidad de refinanciar vencimientos crecientes y por el fin de los préstamos del SURE y de los préstamos y transferencias del Next Generation EU. La previsión se ha cumplido y puede precisarse con cinco factores.

El primero son los vencimientos, que seguirán creciendo con el saldo de deuda y con la llegada a vencimiento de las grandes emisiones de la etapa de tipos cero, que además se refinanciarán a tipos más altos. El segundo es el fin de la financiación europea: las subvenciones del Mecanismo de Recuperación terminan en 2026, y desde 2027-2028 empezarán a amortizarse los préstamos del SURE y del Mecanismo. El tercero son los préstamos a la Seguridad Social, que tenderán a crecer con el gasto en pensiones mientras no se corrija su desequilibrio, que ahora depende de la cláusula de salvaguarda del Tema 24. El cuarto es la financiación autonómica: la condonación y la eventual reforma del modelo de financiación (Tema 25) alterarán los flujos entre el Estado y las comunidades. El quinto son los nuevos gastos, en particular las medidas frente a la crisis de Oriente Medio y el compromiso de defensa, que en ausencia de presupuestos aprobados se canalizan por créditos extraordinarios y aumentan el déficit de caja del Estado.

En sentido contrario, el crecimiento nominal y la buena evolución de la recaudación han permitido mantener la emisión neta estable. La AIReF prevé un déficit de las AAPP del 2,6% en 2026 (1.5.3d). Si se cumple, la emisión neta podría mantenerse en niveles parecidos. El riesgo está en la combinación de un déficit que no baja y unos tipos que suben, que llevaría el coste de la deuda a ser el principal motor de las necesidades futuras.

El documento base presenta la diversificación de instrumentos (bonos indexados, bonos verdes, plazos de hasta 50 años) como uno de los rasgos de la política del Tesoro. Los datos muestran otra cosa: la financiación española está muy concentrada en una sola familia de productos. Los bonos y obligaciones nominales a tipo fijo superan el 90% de la deuda del Estado. Los programas de diversificación son pequeños: en torno al 6% en indexados y unos 22.000 millones en verdes, poco más del 1,5% de la cartera. La diversificación relevante de los últimos años se ha producido en la base inversora, no en los instrumentos: no residentes, bancos centrales extranjeros y, en 2023, un aumento súbito de la inversión minorista en letras.

Esto no es un defecto. Es una decisión deliberada a favor de la liquidez, y así debe explicarse en el examen. Presentar los programas de indexados y verdes como pilares de la financiación sobredimensiona su papel. Además, el documento menciona, pero no desarrolla, que el argumento teórico a favor de los indexados se ha debilitado en España desde que las pensiones se revalorizan con el IPC.

\subsection{Instrumentos de financiación del Tesoro}

Introducción

[DOC] El documento organiza este apartado en tres bloques: las letras del Tesoro, los bonos y obligaciones del Estado (incluidos los indexados y los verdes) y los otros instrumentos (deuda en divisas y préstamos). Parte de una idea de diseño que conviene conservar como hilo conductor: la gama de instrumentos debe atender las preferencias de los inversores, pero evitando la segmentación del mercado que produciría un número excesivo de instrumentos, porque esa segmentación reduciría la liquidez de la deuda y, con ella, su atractivo.

Este subapartado desarrolla esa idea en cinco epígrafes. El primero reconstruye cómo se construyó la gama y con qué principios (3.3.1). Los tres siguientes analizan los instrumentos principales con sus debates: las letras (3.3.2), los bonos y obligaciones nominales (3.3.3) y los programas de diversificación, es decir, los indexados y los verdes (3.3.4). El último trata los instrumentos residuales y los préstamos (3.3.5). Los procedimientos de emisión (subastas, sindicaciones, creadores de mercado) se estudian en 3.4.

La gama de instrumentos del Tesoro se construyó desde 1982 sustituyendo los instrumentos con ventajas fiscales o demanda obligatoria (deuda desgravable, pagarés) por instrumentos de mercado: bonos (1982), obligaciones (1983) y letras (1987). Responde a un principio central, la concentración en pocas referencias grandes y líquidas, aplicado mediante la agregación de emisiones, la curva completa de referencias entre 3 y 50 años y la segregación en strips desde 1997.

El resultado es una cartera dominada por los bonos y obligaciones nominales a tipo fijo, más del 90%, con una vida media de unos 8 años, frente a los menos de 2 de 1990. Las letras (en torno al 7%) cumplen funciones de gestión de tesorería, de amortiguador en las crisis y de activo monetario. En 2023 atrajeron una inversión minorista sin precedentes, de más de 28.000 millones. La lección de 1990, cuando representaban el 62% de la deuda, explica la prudencia con que se usan hoy.

Los programas de diversificación son deliberadamente pequeños. Los indexados, en torno al 6%, tienen un fundamento teórico sólido, pero su cobertura se ha debilitado con la indexación de las pensiones y resultaron caros en 2022-2023. Los verdes, unos 22.000 millones tras la emisión de septiembre de 2026, aportan transparencia y un ahorro modesto, con un debate abierto sobre su adicionalidad. La deuda en divisas ha desaparecido con el euro. Los préstamos oficiales europeos (MEDE, SURE y Mecanismo de Recuperación) han amortiguado las necesidades de mercado en las crisis, pero dejarán de hacerlo desde 2027. La diversificación relevante de los últimos años se ha producido, en suma, en la base inversora y no en los instrumentos.



\subsubsection{La construcción de la gama: liquidez frente a segmentación}

a) El origen: dos instrumentos en 1982

[DOC] Según el documento, cuando en 1982 comenzó el giro hacia la financiación de mercado, el Tesoro disponía solo de dos instrumentos. El primero era la deuda desgravable del Estado, un instrumento a medio plazo dirigido a los contribuyentes del IRPF, cuya demanda se basaba sobre todo en su ventaja fiscal. El segundo eran los pagarés del Tesoro, activos a corto plazo creados por el BdE como instrumento de regulación de la liquidez que, al no tener restricciones de acceso, servían también para financiar al Estado. A finales de 1982 se crearon los bonos del Estado (de 2 a 5 años) y en 1983 las obligaciones del Estado, a plazos más largos, sin ventajas fiscales. En 1987 se crearon las letras del Tesoro, que sustituyeron a los pagarés. Desde entonces la gama se ha ampliado para atender las preferencias de los inversores.

Esta secuencia refleja la transición analizada en 1.2.2 y 3.1.1. Los primeros instrumentos dependían de ventajas fiscales (deuda desgravable, opacidad fiscal de los pagarés) o de demanda obligatoria (coeficientes). Los instrumentos que aparecen desde 1982 se basan en su rentabilidad de mercado. El paso de unos a otros es el paso de una financiación privilegiada a una financiación de mercado.

b) Los principios de diseño

La gama actual responde a cuatro principios que la literatura sobre gestión de la deuda considera buenas prácticas y que el Tesoro formula en sus estrategias.

El primero es la concentración en pocos instrumentos estandarizados. Un mercado con pocas referencias grandes es más líquido que uno con muchas pequeñas, y la liquidez reduce la prima que exigen los inversores. Por eso el Tesoro aplica el llamado sistema de agregación de emisiones: reabre sucesivamente cada referencia hasta que alcanza un volumen que garantiza su liquidez, que el documento sitúa en torno a 25.000 millones para bonos y obligaciones. Lo empezó a aplicar a finales de los ochenta.

El segundo es la curva completa de referencias (benchmarks). [DOC] El Tesoro mantiene referencias a 3, 5, 10, 15, 20, 30 y 50 años. Cuando emite una nueva referencia a un plazo, esta se convierte en el nuevo benchmark, y la anterior pasa a ser una referencia off-the-run. Una curva completa y líquida no solo abarata la financiación del Tesoro. Sirve también de referencia para valorar la deuda de las comunidades, de las empresas y de los bancos españoles, lo que convierte la política de emisión en un bien público para el sistema financiero.

El tercero es la adaptación a la demanda sin fragmentar la oferta. [DOC] El mejor ejemplo, según el documento, es la segregación: todos los bonos y obligaciones son segregables desde 1997. Los creadores de mercado pueden separar cada cupón y el principal en bonos de cupón cero (strips), y reconstituir después el bono original. Los strips atienden a inversores que necesitan flujos únicos en fechas concretas, como aseguradoras y fondos de pensiones que cubren pasivos, o que buscan una ventaja fiscal: sus rendimientos no están sujetos a retención en el Impuesto sobre Sociedades. Y lo hacen sin que el Tesoro tenga que emitir instrumentos nuevos. El documento vincula su creación a la proximidad de la UEM, cuando la deuda española tenía que competir con la de los socios europeos.

El cuarto es la diversificación de la base inversora con programas específicos, los indexados (2014) y los verdes (2021). Estos programas entran en tensión con el primer principio, como vimos en 3.1.3c, y por eso su tamaño es limitado (3.3.4).

c) La evolución de la composición

[DOC] La composición de la cartera ha cambiado radicalmente. En 1990 las letras llegaron a representar el 62% de la deuda en circulación y los bonos y obligaciones el 30%. Hoy los bonos y obligaciones superan el 90% y las letras representan en torno al 6%. Los demás instrumentos (deuda en divisas y préstamos), que a finales de los noventa llegaron a superar el 15%, son hoy menos del 4%. Con ello, y con la emisión a plazos cada vez más largos, la vida media de la deuda del Estado pasó de menos de 2 años en 1990 a superar los 8 en 2021. Según el documento, las agencias de calificación citaron el alargamiento de la vida media entre los argumentos de las subidas de calificación de 2014-2019.

El peso de las letras ha subido ligeramente desde 2022. Con una emisión bruta de letras prevista de unos 108.700 millones en 2026 y un saldo vivo de deuda del Estado de en torno a 1,4 billones, las letras se sitúan hoy cerca del 7-8% de la cartera. La vida media, 7,93 años a finales de 2025 (3.2.3), se ha mantenido cerca del máximo alcanzado en 2021.

\subsubsection{Letras del Tesoro: la flexibilidad a corto plazo}

a) Diseño

[DOC] Según el documento, las letras son valores de renta fija a corto plazo, emitidos al descuento y sin cupón, cuya rentabilidad es implícita: la diferencia entre el precio de compra y el valor nominal. En los años de tipos negativos se emitieron a premio, con rentabilidad negativa. Se introdujeron en 1987 a doce meses. Cuando en 1990 se restringió el acceso del Tesoro al crédito del BdE, se crearon letras a 3 y 6 meses para gestionar los desfases transitorios de tesorería. Hoy se emiten a 3, 6, 9 y 12 meses. Para maximizar su liquidez, cada mes se emite una nueva letra a 12 meses, por un importe inicial elevado, que se reabre cuando su vida residual es de 9, 6 y 3 meses. En la práctica solo hay letras a doce meses y en todo momento circulan doce referencias. La inversión mínima es de 1.000 euros.

La Orden ECM/2/2026 autoriza letras de hasta 24 meses. El Tesoro no ha emitido letras de más de 12 meses en los últimos años, pero el margen legal le da flexibilidad para una situación de tensión.

b) Funciones

Las letras cumplen tres funciones que justifican mantener una proporción de deuda a corto plazo, aunque aumente el riesgo de refinanciación. La primera es la gestión de tesorería: permiten ajustar la emisión a los desfases de caja a lo largo del año sin alterar el programa de bonos y obligaciones, anunciado con antelación. La segunda es servir de amortiguador en las crisis. Como recoge el documento para 2009, y como ocurrió también en 2020, cuando el acceso al mercado se complica el Tesoro puede aumentar las letras, que se colocan más fácilmente porque su precio es menos volátil. La tercera es alimentar el mercado monetario con un activo sin riesgo a corto plazo, útil para fondos monetarios, tesorerías de empresas y bancos.

Una literatura reciente añade un argumento de coste. Greenwood, Hanson y Stein (2015) sostienen que la deuda pública a corto plazo tiene un rendimiento de conveniencia: los inversores la valoran como un sustituto del dinero y aceptan por ella una rentabilidad inferior a la que justificaría su riesgo. El Estado tiene una ventaja comparativa para ofrecer ese activo. Por tanto, desde el punto de vista del coste, le convendría emitir más deuda corta, hasta el punto en que el ahorro marginal iguale al coste marginal del mayor riesgo de refinanciación. Es un contrapeso teórico al argumento de Cole y Kehoe a favor de la deuda larga (3.1.3b). La proporción actual de letras, en torno al 7%, refleja el peso que el Tesoro español da al riesgo de refinanciación tras la experiencia de 2010-2012.

c) La lección histórica: el riesgo de refinanciación de 1990

[DOC] Las letras demostraron desde su creación una gran capacidad de financiación, gracias a los elevados tipos a los que se emitieron (15,5% anual al principio). En 1990 llegaron a suponer el 62% de la deuda, lo que exponía al Tesoro a un elevado riesgo de refinanciación: la necesidad de renovar en menos de un año la mayor parte de la deuda, en condiciones desconocidas. Para corregirlo se promovió la participación de inversores institucionales, que prefieren plazos más largos. En 1990 se introdujeron incentivos fiscales para los no residentes, que elevaron su proporción de bonos y obligaciones del 10% a más del 30% en apenas un año. También se fomentó la demanda minorista con campañas de publicidad y los FONDTESORO, fondos de inversión especializados en deuda del Tesoro con comisiones limitadas por convenio. Las letras bajaron por debajo del 30% de la deuda en 1996.

Este episodio, que se sitúa entre los analizados en 1.2.4, explica la vulnerabilidad de la deuda española en 1992-1995. Con una cartera tan corta, la subida de los tipos para defender la peseta en el SME se trasladó casi de inmediato al coste de la deuda. Es un ejemplo del riesgo de tipo de interés y de refinanciación que la literatura describe.

d) 2023: el aumento súbito de la inversión minorista

Un fenómeno reciente que el documento no recoge merece atención. Con la subida de los tipos del BCE, la rentabilidad de las letras pasó de negativa en 2022 a superar el 3,5-4% en 2023, mientras los depósitos bancarios apenas se remuneraban. Los ahorradores españoles acudieron en masa a las subastas, sobre todo mediante peticiones no competitivas a través de la Cuenta Directa del BdE y de las entidades financieras. Según datos del BdE recogidos por la EFPA, la inversión de los particulares en letras superó los 28.000 millones en 2023, frente a unos 99 millones en 2022. La proporción de letras en manos de hogares e instituciones sin fines de lucro pasó del 0,1% en septiembre de 2022 a cerca del 30% en septiembre de 2023.

El episodio tiene tres lecturas. Para el Tesoro supuso una diversificación de la base inversora en un momento en que el BCE dejaba de comprar. Para el sistema financiero fue una presión competitiva que obligó a los bancos a mejorar la remuneración de los depósitos, lo que confirma el efecto de transmisión de la política monetaria. Para la gestión del riesgo planteó una pregunta: una base minorista concentrada en el corto plazo es sensible a la rentabilidad relativa y puede retirarse con rapidez cuando los tipos bajan, como empezó a ocurrir en 2024-2025 con las bajadas del BCE. No he encontrado datos del saldo minorista en 2026, ahora que los tipos vuelven a subir.

\subsubsection{Los bonos y obligaciones del Estado: el núcleo de la financiación}

a) Diseño

[DOC] Según el documento, los bonos y obligaciones son valores de renta fija a más de dos años que solo se diferencian en el plazo: de 2 a 5 años los bonos y más de 5 años las obligaciones, hasta 50 años en la actualidad. Pagan normalmente un cupón anual fijo, aunque el Tesoro ha llegado a emitirlos con cupón cero y, en la etapa de tipos negativos, al descuento con rentabilidad negativa. La obligación a diez años es la referencia para calcular la prima de riesgo frente a Alemania. Se emiten habitualmente por subasta, salvo el primer tramo de las nuevas obligaciones a plazos largos, que suele colocarse por sindicación. De forma puntual se han hecho colocaciones privadas. La inversión mínima es de 1.000 euros.

b) La emisión sobre la par y bajo la par: una decisión con consecuencias contables

[DOC] El documento dedica una nota extensa a la emisión de reaperturas por encima o por debajo de la par. Ya se analizó en 1.1.1b, así que aquí solo añado su dimensión de gestión.

Cuando los tipos caen, como en 2015-2021, una reapertura de una referencia con cupón alto se emite sobre la par. El Tesoro ingresa más que el nominal, y la diferencia, la prima de emisión, financia el déficit de caja sin aumentar la deuda PDE, que se valora por el nominal. A cambio, el Estado paga cupones altos hasta el vencimiento. Desde 2022 ocurre lo contrario: muchas reaperturas se emiten bajo la par, y el Tesoro necesita emitir más nominal para obtener el mismo efectivo.

La elección entre reabrir una referencia antigua (con cupón muy distinto del tipo de mercado) o crear una nueva (con cupón próximo al de mercado) no es neutral para la deuda publicada. Aunque en valor presente el coste sea equivalente, la reapertura sobre la par reduce la deuda PDE del momento y eleva los intereses futuros en contabilidad nacional. La literatura sobre contabilidad creativa ha señalado este tipo de decisiones como un margen de maniobra de los gestores (von Hagen y Wolff, 2006, en 2.1.3d). No hay evidencia de que el Tesoro español lo haya utilizado con ese fin: su criterio declarado es la liquidez de las referencias. Pero conviene conocer el argumento.

c) Los plazos muy largos: el debate sobre la deuda a 50 años

[DOC] El documento señala que el Tesoro emite obligaciones con vencimiento de hasta 50 años. La primera se emitió en 2016, con posteriores emisiones en la etapa de tipos mínimos. La justificación fue aprovechar unos tipos históricamente bajos para asegurar financiación barata durante medio siglo y atender la demanda de aseguradoras y fondos de pensiones, que necesitan activos muy largos para cubrir sus pasivos.

El debate sobre estas emisiones es doble. A favor está la lógica del seguro: el Estado compra protección frente a subidas futuras de tipos, como la de 2022-2026. En contra está el coste esperado: con una curva de pendiente positiva, la deuda a 50 años es más cara en promedio que la renovación sucesiva de deuda más corta. Además, es poco líquida. La evaluación ex post depende del escenario. Las emisiones a 50 años de 2016-2021 se hicieron a tipos muy inferiores a los actuales y han resultado baratas en comparación con la refinanciación a los tipos de 2023-2026. Pero esa evaluación es sensible al periodo elegido, y en todo caso su volumen es pequeño en relación con la cartera.

d) Las cláusulas de acción colectiva

Una característica jurídica que el documento no menciona afecta a todos los bonos y obligaciones emitidos desde 2013. En aplicación del artículo 12.3 del Tratado del MEDE, todos los valores de deuda pública de la zona del euro a más de un año emitidos desde el 1 de enero de 2013 incluyen cláusulas de acción colectiva (CAC) estandarizadas. Permiten que una mayoría cualificada de tenedores acuerde una reestructuración que vincula a todos, lo que evita que una minoría bloquee el acuerdo. La reforma del Tratado del MEDE de 2021 prevé sustituirlas desde 2022 por cláusulas de agregación única (single-limb), en las que la mayoría se calcula para todas las emisiones a la vez. Pero la reforma no ha entrado en vigor, porque Italia no la había ratificado. Las CAC tienen poca relevancia práctica para España, pero muestran que la deuda soberana de la zona del euro no se considera jurídicamente libre de riesgo de reestructuración, una consecuencia de la experiencia griega de 2012.\foonote{Reestructuraciones de deuda soberana, cláusulas de acción colectiva y la experiencia griega de 2012: Tema 3A-38. Reforma del Tratado del MEDE: Tema 3B-45.
}

\subsubsection{Los programas de diversificación: bonos indexados y bonos verdes}

a) Los bonos indexados a la inflación

[DOC] Según el documento, el programa lanzado en 2014 permite al Tesoro emitir bonos y obligaciones cuya rentabilidad está indexada al IPC armonizado de la zona del euro sin tabaco, publicado por Eurostat. Si la inflación acumulada fuera negativa, se devolvería al menos el nominal. El programa ha diversificado la base inversora, con acceso a inversores interesados en protegerse de la inflación. El Tesoro está en una posición idónea para emitirlos, porque sus ingresos están ligados a la inflación. Pero esa posición se ha debilitado al volver a indexar a la inflación el gasto en pensiones. El saldo vivo equivale a algo más del 6% de la deuda del Estado.

El fundamento teórico es el que se vio en 3.1.3b. Barro (2003) concluyó que, en ausencia de deuda contingente, la deuda larga indexada es la que mejor protege al presupuesto, y Campbell y Shiller (1996) defendieron su emisión por el ahorro de la prima de riesgo de inflación: los inversores exigen a la deuda nominal una compensación por la incertidumbre inflacionista, que el Estado se ahorra si asume él ese riesgo.

La evaluación del programa español tiene, sin embargo, cuatro matices. El primero es el que señala el documento, y es el más relevante para España. La cobertura funciona si los ingresos públicos crecen con la inflación más que los gastos. Desde la Ley 21/2021, que revaloriza las pensiones con el IPC, una parte muy importante del gasto también está indexada, y la cobertura natural se reduce: el Estado soporta el riesgo de inflación por los dos lados de su balance.

El segundo es el coste ex post de 2022-2023. Con una inflación de la zona del euro de más del 8% en 2022, el principal de los indexados se revalorizó en la misma cuantía, con un coste superior al de la deuda nominal emitida en su momento. No es un argumento contra el programa, porque se trata de un seguro y su coste es la contrapartida de la protección que ofrece a los inversores. Pero sí es un argumento a favor de mantenerlo pequeño mientras el gasto esté indexado. El Reino Unido, cuya deuda indexada supera el 20% del total, sufrió un aumento muy fuerte de su carga de intereses en 2022-2023 por este motivo.

El tercero es el índice de referencia. El Tesoro indexa a la inflación de la zona del euro, no a la española, porque es el índice con un mercado líquido de inversores y de coberturas. Pero los ingresos españoles dependen de la inflación española. Si la española se desvía de la media, como en 2022 (en algunos meses por encima y en otros por debajo), la cobertura se debilita.

El cuarto es la liquidez. El segmento indexado es menos líquido que el nominal, y los inversores exigen una prima de liquidez que compensa en parte el ahorro de la prima de inflación. La literatura sobre el caso británico y el estadounidense ha encontrado que, en algunos periodos, la prima de liquidez ha superado a la de inflación y ha hecho los indexados más caros que los nominales. El Tesoro ha seguido emitiendo indexados en los últimos años, incluida una nueva referencia a unos doce años. Su saldo sigue en torno al 6% de la cartera, lo que indica que los mantiene como instrumento de diversificación y no como apuesta de coste.

b) Los bonos verdes soberanos

[DOC] Según el documento, el programa de bonos verdes se lanzó en 2021 con un doble objetivo: diversificar la base inversora y reforzar el compromiso de España con la transición ecológica. Los fondos obtenidos deben destinarse a gastos elegibles: mitigación y adaptación al cambio climático, uso sostenible y protección de los recursos hídricos y marinos, economía circular, prevención de la contaminación y protección de la biodiversidad. Esa vinculación limita el volumen que se puede emitir. La elevada demanda, unida a la escasez relativa, ha permitido obtener un ahorro frente a un bono convencional (el greenium), que en la emisión inaugural fue de 2 puntos básicos.

La emisión inaugural, de septiembre de 2021, fue una obligación a 20 años con vencimiento en 2042 por 5.000 millones, reabierta varias veces después. A principios de 2026 el saldo de bonos verdes rondaba los 18.000 millones. El 9 de septiembre de 2026 el Tesoro colocó por sindicación una nueva referencia verde a 20 años (vencimiento en 2047) por 4.000 millones, a una rentabilidad en torno al 4,2%, con una demanda de 68.500 millones, diecisiete veces la oferta. El saldo verde se acerca así a los 22.000 millones, en torno al 1,5% de la deuda del Estado.

El programa suscita tres debates. El primero es la adicionalidad. El dinero es fungible: el Estado habría hecho los gastos verdes elegibles con o sin bono verde, y la emisión no altera la composición del presupuesto. Los defensores del programa responden que su valor está en la transparencia. Los informes anuales de asignación y de impacto obligan a identificar y medir los gastos verdes, y el programa envía una señal de compromiso a los inversores y a la sociedad. Los críticos lo ven como una etiqueta de marketing sin efecto real.

El segundo es la magnitud del greenium. La evidencia empírica sobre bonos verdes soberanos (Doronzo, Siracusa y Antonelli, 2021, para Italia; los estudios del BCE y del BPI) encuentra ahorros positivos pero pequeños, de pocos puntos básicos, y variables en el tiempo. Una parte de la diferencia puede reflejar la menor liquidez de los bonos verdes y no una preferencia de los inversores. El ahorro sobre el volumen total es, por tanto, modesto.

El tercero es la competencia europea. La Comisión, que financia una parte del Next Generation EU con bonos verdes, se ha convertido en el mayor emisor verde del mundo, lo que reduce la escasez relativa que sostenía el greenium de los emisores soberanos nacionales. El diferencial de solo 7 puntos básicos de la emisión de septiembre de 2026 respecto a la referencia verde de 2043 sugiere que el mercado sigue valorando la curva verde española como un segmento propio.

c) Un instrumento que el Tesoro no emite: la deuda vinculada al PIB

Para cerrar el epígrafe, conviene mencionar un instrumento que la teoría de la gestión óptima de la deuda recomienda y que ningún gran emisor utiliza. Siguiendo la lógica de Lucas y Stokey (1983), la deuda cuyo pago se vincula al crecimiento del PIB nominal protegería al presupuesto frente a las recesiones: se pagaría menos cuando la economía cae y más cuando crece. Shiller la propuso hace décadas, y Borensztein y Mauro (2004) la defendieron. En 2016-2017, el Banco de Inglaterra y un grupo de trabajo del G20 elaboraron un modelo de contrato (el London Term Sheet). Los obstáculos que han impedido su adopción son varios: el riesgo de revisiones y manipulación del PIB, la falta de liquidez de un instrumento nuevo, la prima de novedad que exigirían los inversores y la ausencia de un primer emisor de referencia. La experiencia más próxima son las cláusulas de deuda resistente a catástrofes climáticas, que suspenden los pagos tras un desastre natural y que han adoptado algunos países y el Reino Unido en sus préstamos a países en desarrollo. Tras la DANA de 2024, la idea tendría cierta lógica para España, pero no figura en la agenda del Tesoro.\footnote{La deuda contingente al estado de la economía como solución teórica a la nivelación impositiva (Lucas y Stokey, 1983) y la aproximación mediante la estructura de plazos e indexación: Tema 3A-39.
}

\subsubsection{Otros instrumentos: deuda en divisas y préstamos}

a) La deuda en divisas

[DOC] Según el documento, la deuda en divisas tuvo un papel destacado para atraer inversores internacionales en los ochenta y noventa, y llegó a superar el 10% de la deuda a finales de los noventa. Con el euro dejó de ser necesaria para acceder a los no residentes. La última emisión en divisas fue en 2014, y quedan emisiones residuales en yenes, dólares y libras de menos del 0,1% de la deuda. El Tesoro solía cubrir el riesgo de tipo de cambio con permutas financieras (swaps).

La historia de la deuda en divisas española ilustra el riesgo cambiario que analizamos en 1.2.4. Las devaluaciones de 1992-1995 elevaron el valor en pesetas de la deuda en divisas no cubierta y contribuyeron al salto de la ratio de 1993. Con el euro, el riesgo cambiario ha desaparecido para la deuda del Estado, pero la renuncia a la moneda propia tiene la contrapartida analizada en el Bloque II: la deuda está denominada en una moneda que el país no controla.

b) Los préstamos

[DOC] Según el documento, España ha recibido préstamos que han permitido al Tesoro reducir sus emisiones: el del MEDE en 2012 para la recapitalización bancaria (41.333 millones), los del SURE en 2020-2021 (21.326 millones), varios del Banco Europeo de Inversiones y algunos préstamos Schuldschein con inversores privados. Estos últimos son préstamos sujetos a la legislación alemana, adaptados a las necesidades de cada inversor, y se utilizaron por última vez en 2014. Todos se obtuvieron a tipos inferiores a los que el Tesoro pagaba en cada momento.

Actualización y matiz. A estos préstamos se suman ahora los del Mecanismo de Recuperación y Resiliencia. España ha solicitado unos 21.500 millones en total: unos 16.000 millones movilizados en 2025 y 6.500 millones solicitados para 2026 (3.2.2b). Los préstamos del MEDE y del SURE no fueron préstamos al Tesoro en sentido estricto. El del MEDE se desembolsó al FROB en bonos del propio MEDE, y el del SURE se financió con emisiones de la Unión garantizadas por los Estados miembros. Pero ambos computan como deuda de las AAPP españolas y reducen lo que el Tesoro necesita emitir en el mercado.

La observación analítica relevante es que la financiación oficial europea ha actuado en 2012 y en 2020-2026 como amortiguador de las necesidades de mercado. En 2012 lo hizo con condicionalidad financiera, en el contexto de una crisis de acceso. En 2020-2026 lo hizo sin estigma, en el marco de una respuesta común. El cambio refleja la evolución institucional de la Unión que se analizó en 1.5.1c y 2.2.4. Pero desaparece a partir de 2027, lo que devolverá al Tesoro a una dependencia completa del mercado (3.2.4b).


El documento base afirma que la subasta española permite al Tesoro colocar los valores con un menor coste que el que resultaría de utilizar el sistema de subasta holandesa. Esa afirmación es un error de razonamiento que un tribunal con formación en economía puede señalar. Compara los formatos manteniendo constantes las pujas, cuando la teoría de subastas enseña que los inversores cambian su forma de pujar según el formato. Con pujas estratégicas no hay una ordenación general de los formatos por ingreso, y la evidencia empírica es mixta. La justificación sostenible del formato español no es que sea más barato, sino que combina un incentivo a pujar alto con una protección frente a la maldición del ganador, y que funciona bien en un mercado con creadores de mercado. He construido 3.4 alrededor de ese debate, con la teoría en nota al pie y remisión al tema 3A-14, como acordamos.


\subsection{Procedimientos de emisión y el funcionamiento del mercado}

Introducción

[DOC] El documento explica que el paso a la financiación de mercado desde principios de los ochenta exigía procedimientos que asegurasen el abastecimiento regular del mercado y la financiación al menor coste posible. En 1983 se adoptó la subasta competitiva para las obligaciones del Estado, y después se extendió a los bonos y a las letras. Sigue siendo el principal procedimiento. A la subasta se añaden la sindicación, las colocaciones privadas, los canjes y los préstamos, y todo el sistema se apoya en los creadores de mercado.

Este subapartado analiza los procedimientos en cinco epígrafes. El primero trata el diseño de la subasta española y el debate teórico sobre los formatos (3.4.1). El segundo, el calendario y los mecanismos de transparencia que acompañan a las subastas (3.4.2). El tercero, las sindicaciones y su disyuntiva entre coste y colocación (3.4.3). El cuarto, los procedimientos residuales: colocaciones privadas, canjes y recompras (3.4.4). El quinto, los creadores de mercado y el mercado secundario (3.4.5).

### Síntesis de 3.4

La subasta sigue siendo el procedimiento central del Tesoro, con más del 80% de la emisión bruta. El formato español es híbrido: las pujas por encima del precio medio ponderado se adjudican a ese precio, y las demás, al precio ofrecido. El documento lo presenta como más barato que la subasta holandesa, pero esa conclusión no es válida. Con pujas estratégicas no hay una ordenación general de los formatos (Vickrey, 1961; Ausubel y Cramton, 2002), y la evidencia empírica es mixta. Su justificación es que corrige parcialmente la rebaja de las pujas altas y la reducción estratégica de la demanda. Las peticiones no competitivas han ganado importancia con la entrada de los minoristas desde 2023.

La previsibilidad (calendario anual, 48 subastas en 2026, rangos objetivo desde 1995) aplica a la gestión de la deuda la lógica de las reglas frente a la discrecionalidad. Las sindicaciones, algo más del 10% de la emisión en años normales y casi el 20% en 2020, permiten lanzar nuevas referencias largas con volumen y distribución controlada, a cambio de una prima de nueva emisión y de una dependencia del sindicato. Las de 2026 han registrado demandas récord (145.000 millones en enero) y una base inversora internacional (90% no residentes).

Los canjes, las recompras y las operaciones de tesorería gestionan la estructura de la cartera y la liquidez. El sistema de creadores de mercado remunera de forma implícita sus compromisos con privilegios como la segunda vuelta, que equivale a una opción gratuita. Afronta además el reto de absorber una oferta neta creciente en un entorno de restricciones regulatorias sobre el balance de los bancos, ahora que el BCE se retira.



\subsubsection{La subasta: el formato español y el debate sobre el diseño}

a) El funcionamiento

[DOC] Según el documento, la subasta es el procedimiento más utilizado por el Tesoro: con ella obtiene más del 80% de la emisión bruta anual. El sistema español combina los dos formatos que utilizan otros Tesoros. En la subasta holandesa o de precio único, todas las peticiones aceptadas se adjudican al precio marginal, es decir, el mínimo aceptado. En la subasta convencional o de precios múltiples, cada petición aceptada se adjudica al precio que ofreció.

En el sistema español, el Tesoro fija primero el precio mínimo aceptado y calcula después el precio medio ponderado de las peticiones aceptadas. Las peticiones a un precio superior al medio ponderado se adjudican a ese precio medio, y las demás, al precio que ofrecieron. El Tesoro decide el precio mínimo en la propia subasta a la vista de las peticiones recibidas, del importe que desea emitir (anunciado antes como un rango orientativo), del número de referencias subastadas y del saldo en circulación de cada una. En las subastas de letras se ofrecen siempre dos referencias, y en las de bonos y obligaciones, normalmente tres o cuatro.

Los inversores pueden presentar también peticiones no competitivas, en las que solo indican la cantidad, que se adjudican al precio medio ponderado. Salvo excepciones, no pueden superar los 5 millones de euros, y con carácter general no se adjudican si ese precio implica una rentabilidad negativa. Son las que utilizan habitualmente los inversores minoristas.

Una precisión sobre el papel del BdE. El documento dice que interviene, pero únicamente proveyendo la infraestructura tecnológica. Es más preciso decir que actúa como agente financiero del Estado: recibe las peticiones, gestiona la resolución técnica de la subasta, liquida las operaciones y publica los resultados. Las decisiones de fondo (el importe y el precio mínimo) corresponden al Tesoro.

b) El debate teórico: ¿qué formato minimiza el coste?

Aquí está el error del documento. Su argumento es el siguiente: con el formato español, las pujas altas se adjudican al precio medio ponderado, que es superior al marginal. Por tanto, el Tesoro ingresa más que con la subasta holandesa, en la que todo se adjudicaría al marginal. El razonamiento sería correcto si los inversores pujasen igual en ambos formatos. Pero no lo hacen.

La teoría de subastas, desde Vickrey (1961), muestra que el comportamiento de los pujadores depende del formato.\footnote{En una subasta de unidades múltiples con valoraciones privadas independientes y demanda unitaria, el teorema de equivalencia de ingresos (Vickrey, 1961; Myerson, 1981) implica que los formatos de precio único y de precios múltiples generan el mismo ingreso esperado. Con demanda de varias unidades y valoraciones con componente común, como en la deuda pública, la equivalencia se rompe. En la subasta de precios múltiples, los pujadores rebajan sus pujas para no pagar de más y para protegerse de la maldición del ganador. En la de precio único, pujan cerca de su valoración en las unidades inframarginales, pero reducen su demanda en las marginales para bajar el precio que pagan por todas (Ausubel y Cramton, 2002). El efecto neto sobre el ingreso es ambiguo y depende de la distribución de valoraciones, del número de pujadores y de su información. El formato español combina los dos: el pujador paga el mínimo entre su puja y el precio medio ponderado. ([\authorfont{Ver Tema 3.A.14}])
}

En una subasta de precios múltiples, cada inversor paga lo que puja. Tiene incentivo a pujar por debajo de su valoración real (bid shading) para no pagar de más. Y teme la maldición del ganador: si gana, es probablemente porque ha sobrevalorado el título respecto a los demás. Ambos efectos reducen las pujas y el ingreso del emisor.

En una subasta de precio único, el inversor paga el precio marginal, determinado por las pujas ajenas, lo que le anima a pujar más cerca de su valoración. Friedman (1960) defendió este formato por esa razón. Pero Ausubel y Cramton (2002) mostraron que, cuando cada inversor demanda varias unidades, surge otro incentivo: reducir la demanda en las unidades marginales para bajar el precio que pagará por todas las demás. Con ese comportamiento, el formato de precio único puede generar menos ingresos.

En los contextos de unidades múltiples y valoraciones interdependientes, que son los de la deuda pública, no se cumple el teorema de equivalencia de ingresos. No hay un formato que domine teóricamente a los demás en ingreso esperado.

La evidencia empírica es igualmente mixta. El experimento del Tesoro de Estados Unidos, que entre 1992 y 1998 utilizó el formato de precio único en parte de sus subastas y lo generalizó en 1998, encontró ahorros modestos y difíciles de atribuir con seguridad al formato (Malvey y Archibald, 1998). Umlauf (1993), con el caso de México, encontró ingresos mayores con el precio único. Hortaçsu y McAdams (2010), con datos de Turquía, estimaron diferencias pequeñas entre formatos. Hortaçsu, Kastl y Zhang (2018) muestran para Estados Unidos que el comportamiento estratégico de los creadores de mercado, que tienen información sobre el flujo de órdenes de los clientes, influye en el resultado de forma importante.

Sobre el formato español existe una literatura específica, aunque reducida. Álvarez y Mazón (2007) lo compararon teóricamente con el de precios múltiples y encontraron que puede reducir el bid shading sin eliminar la protección frente a pujas excesivamente altas. Abbink, Brandts y Pezanis-Christou (2006), en un experimento de laboratorio, compararon los tres formatos y encontraron que el español no genera sistemáticamente ingresos menores que los otros dos, pero tampoco mayores en todas las condiciones.

La conclusión que se puede defender en el examen es que el formato español no es más barato por construcción. Su justificación es que corrige parcialmente los dos problemas de los formatos puros: reduce el incentivo a rebajar las pujas altas, porque las pujas por encima de la media no se pagan a su precio, y limita el problema de la reducción de la demanda, porque las pujas por debajo de la media sí se pagan a su precio. A ello se suma su antigüedad y el conocimiento que tienen de él los inversores, lo que reduce la incertidumbre.

c) Las peticiones no competitivas y los minoristas

Las peticiones no competitivas tienen un papel que se ha vuelto relevante desde 2023, con el aumento súbito de la inversión minorista en letras (3.3.2d). Permiten a los pequeños ahorradores participar sin necesidad de estimar el precio, y les garantizan el precio medio, que no es el peor de la subasta. Para el Tesoro, sus importes son relativamente previsibles y constituyen una demanda estable, aunque sensible a la rentabilidad relativa. La regla que excluye su adjudicación con rentabilidades negativas protegía a los minoristas en la etapa de tipos negativos. Desde 2025 el Tesoro ofrece además a los particulares una plataforma propia para invertir directamente en letras, bonos y obligaciones, con guías prácticas, como parte de su estrategia de diversificación.

\subsubsection{El calendario, los rangos y la transparencia}

a) El calendario anual

[DOC] Según el documento, para asegurar la estabilidad de las emisiones y el acceso continuado a la financiación, cada año se publica en el BOE el calendario de subastas. Ese acceso es esencial desde que en 1994 se prohibió la financiación del BdE. Cada mes hay dos subastas de letras, una a 6 y 12 meses y otra a 3 y 9 meses, y dos de bonos y obligaciones. El Tesoro anuncia el viernes anterior las referencias concretas que subastará y, desde 1995, el lunes de la semana de la subasta anuncia el rango objetivo de colocación. El documento explica que esta última práctica se introdujo para evitar las grandes oscilaciones que había hasta entonces en los volúmenes emitidos. La Estrategia de 2026 prevé 48 subastas ordinarias.

b) Por qué importa la previsibilidad

El calendario y los rangos son la aplicación del objetivo de previsibilidad y transparencia (3.1.3c), y su fundamento es doble. Por un lado, reducen la incertidumbre de los inversores sobre la oferta futura, lo que disminuye la prima que exigen por esa incertidumbre. Por otro, reducen el margen de discrecionalidad del Tesoro. Un emisor que decidiera cada semana cuánto emite según las condiciones del momento podría aprovechar las ventanas favorables, pero los inversores anticiparían ese comportamiento: esperarían que el Tesoro emita más cuando los precios le son favorables y menos cuando no. Eso aumentaría la volatilidad y la prima exigida. Es una aplicación a la gestión de la deuda de la idea de reglas frente a discrecionalidad que Kydland y Prescott (1977) formularon para la política monetaria.

El rango objetivo tiene, sin embargo, una contrapartida. Al comprometerse con un intervalo, el Tesoro renuncia a ajustar mucho el importe cuando la demanda es débil. En las subastas de 2011-2012, cuando la demanda caía, colocar en la parte baja del rango a tipos altos era preferible a no colocar, porque no hacerlo habría enviado una señal de pérdida de acceso. No he encontrado un análisis del Tesoro sobre cuántas subastas se han resuelto por debajo del rango.

\subsubsection{Las sindicaciones: colocación frente a coste}

a) Funcionamiento y ventajas

[DOC] Según el documento, las sindicaciones son el segundo procedimiento en importancia y se usan, con carácter general, para emitir el primer tramo de las nuevas obligaciones. En una sindicación, el Tesoro contrata a un grupo de bancos (el sindicato), fija el precio de emisión y los inversores indican la cantidad que desean a ese precio. Sus ventajas son tres. La primera es que permiten emitir de una vez un volumen mayor que una subasta, lo que da liquidez al nuevo título desde el inicio: en plena pandemia el Tesoro captó 15.000 millones de una nueva obligación en una sola operación. La segunda es que dan al Tesoro el control de la distribución. Como la demanda suele superar mucho la oferta, el Tesoro puede asignar los títulos por categorías de inversores y dar un porcentaje mayor a los que previsiblemente los mantendrán más tiempo, como bancos centrales e instituciones oficiales, o a los de regiones con menor inversión acumulada en deuda española. La tercera es la notoriedad.

Según el documento, su inconveniente principal es un mayor coste. El Tesoro no llega a conocer el precio marginal al que los inversores estarían dispuestos a comprar, como sí ocurre en la subasta, y debe pagar comisiones al sindicato. Además, depende del sindicato, que tiene más información de primera mano sobre los inversores finales. Normalmente el Tesoro obtiene algo más del 10% de su emisión bruta anual mediante sindicaciones, y en 2020 llegó a casi el 20%.

b) La actualidad: demandas récord y una base inversora internacional

Las sindicaciones recientes muestran un apetito por la deuda española muy superior al de hace una década. En enero de 2026, el Tesoro emitió 15.000 millones de una nueva obligación a diez años con una demanda de 145.000 millones, un récord histórico que supera los 139.000 millones de 2025. El cupón era del 3,30%, y la rentabilidad del 3,323% quedaba 6 puntos básicos por encima de la referencia a diez años vigente. Los no residentes tomaron el 90% de la emisión, y los bancos centrales e instituciones oficiales fueron el primer grupo inversor (26,5%). En mayo de 2026 colocó otra obligación a diez años por 13.000 millones, con una demanda diez veces superior. En septiembre de 2026 colocó 4.000 millones de un nuevo bono verde a veinte años, con una demanda de 68.500 millones (3.3.4b).

c) El debate: ¿cuánto cuesta la sindicación?

La literatura analiza el coste de la sindicación en dos componentes. El primero es la comisión pagada al sindicato, que es explícita y pequeña en relación con el importe. El segundo, más importante, es la prima de nueva emisión, es decir, la diferencia entre el precio de emisión y el precio al que cotiza el título después en el mercado secundario. El diferencial de 6 puntos básicos de la emisión de enero de 2026 sobre la curva es un ejemplo de esa prima: el Tesoro paga algo más para asegurar una colocación grande y bien distribuida. La cuestión es si ese coste se compensa con los beneficios de una base inversora más estable y de la liquidez inmediata del título.

Los análisis de la OCDE y la práctica de los gestores de deuda europeos sugieren que la respuesta es afirmativa para las nuevas referencias largas. En ellas la subasta es arriesgada, porque no hay precio de referencia en el mercado secundario y el volumen necesario para garantizar la liquidez es grande. Por eso los Tesoros europeos, incluido el español, reservan la sindicación para el lanzamiento de nuevas referencias a plazos largos, para los programas especiales (verdes e indexados) y para situaciones de tensión. Reabren después esas referencias por subasta. La experiencia de 2020 muestra que la sindicación es también un instrumento de gestión de crisis, porque permite asegurar grandes volúmenes en un solo día cuando la volatilidad hace arriesgada la subasta.

Una crítica que se puede formular es la de la asimetría de información que señala el documento. Los bancos del sindicato conocen la demanda de sus clientes antes que el Tesoro, y pueden fijar el precio de forma favorable para ellos o para sus clientes. El Tesoro la mitiga con varios mecanismos: rota a los bancos del sindicato según su desempeño como creadores de mercado, publica los criterios de asignación y compara el precio con el de la curva secundaria.

\subsubsection{Otros procedimientos: colocaciones privadas, canjes y recompras}

[DOC] Según el documento, el Tesoro utiliza otros procedimientos de importancia residual. Las colocaciones privadas se realizan a propuesta de los inversores y se ejecutan si contribuyen a diversificar la base inversora, reducen la carga de intereses y encajan en la estrategia. Los canjes se utilizan de forma excepcional para aumentar la liquidez y suavizar el perfil de vencimientos: retiran referencias poco líquidas o próximas al vencimiento a cambio de otras. El objetivo de los préstamos es reducir la carga de intereses y alargar la vida media.

Los canjes y las recompras (en las que el Tesoro compra en el mercado referencias propias antes de su vencimiento) son instrumentos de gestión de pasivos, no de financiación. Su objetivo no es captar recursos, sino modificar la estructura de la cartera. Sirven para suavizar los picos de vencimientos, que concentran el riesgo de refinanciación en fechas concretas, y para retirar referencias poco líquidas que fragmentan el mercado. Su uso ha sido limitado en España y suele concentrarse en los años previos a grandes amortizaciones. En un contexto de vencimientos crecientes como el actual (3.2.4b), podrían ganar importancia para suavizar el perfil de amortizaciones, algo que otros Tesoros europeos hacen con más frecuencia.

A estos procedimientos se añaden las operaciones de gestión de tesorería. Desde que el BCE limitó la remuneración de los depósitos de las administraciones (3.2.2c), el Tesoro coloca sus excedentes de liquidez en el mercado mediante subastas de depósitos y adquisiciones temporales de activos con las entidades. Eso le permite remunerarlos sin mantenerlos en el banco central. Estas operaciones no aumentan la deuda, pero forman parte de la gestión integrada de la deuda y la tesorería que las directrices internacionales recomiendan.

\subsubsection{Los creadores de mercado y el mercado secundario}

a) El sistema de creadores de mercado

[DOC] Según el documento, los creadores de mercado son entidades financieras que se comprometen con el Tesoro a favorecer la liquidez del mercado secundario y a colaborar en la difusión de la deuda del Estado. Hay dos grupos, uno de letras y otro de bonos y obligaciones, aunque casi todas las entidades pertenecen a ambos: en 2022 eran 19 y 18, respectivamente. Para acceder a la condición deben cumplir requisitos técnicos y humanos y demostrar compromiso con el mercado español. Adquieren obligaciones de participación en las subastas y de provisión de liquidez en el secundario. A cambio, tienen condiciones más favorables en las subastas (solo ellos pueden presentar peticiones en los treinta minutos previos al cierre), forman parte de los sindicatos, acceden a las segundas vueltas y pueden segregar y reconstituir valores. El Tesoro evalúa mensualmente su desempeño con criterios objetivos y públicos, y los resultados determinan el acceso a las segundas vueltas y la selección de los sindicatos. En 2022, los mejores creadores de bonos y obligaciones fueron Deutsche Bank, BBVA y JP Morgan, y los de letras, Crédit Agricole, Banco Santander y BBVA.

b) La segunda vuelta: una opción gratuita

[DOC] El documento explica con precisión la segunda vuelta, introducida en 1991. Permite a los creadores de mercado solicitar, hasta un máximo del 24% del nominal que adquirieron en la subasta, títulos adicionales al precio medio ponderado redondeado. El porcentaje de cada creador depende de su participación y de su desempeño. Como se celebra dos o tres días después de la subasta, funciona como una opción de compra gratuita: el creador la ejercerá si los tipos han bajado y el precio ha subido en ese intervalo.

El análisis económico de este privilegio es interesante. Para el Tesoro tiene un coste esperado, el valor de la opción que regala, que es mayor cuanto más volátil es el mercado. Ese coste es la remuneración implícita de los compromisos de los creadores de mercado (participar en las subastas, aunque sean difíciles, y dar precios en el secundario). Es una forma de pagar por un servicio sin hacerlo de forma explícita. La literatura sobre sistemas de creadores de mercado (Arnone e Iden, 2003, para el FMI) discute si es preferible remunerar esos servicios con privilegios de este tipo o con comisiones explícitas. La mayoría de los Tesoros europeos utiliza combinaciones de privilegios parecidas a la española.

c) El mercado secundario

La deuda del Estado se negocia en el Mercado de Deuda Pública en Anotaciones y se registra en Iberclear. La mayor parte de la negociación entre profesionales se realiza en plataformas electrónicas (MTS España, BrokerTec, y plataformas con clientes como Tradeweb o Bloomberg) y en operaciones bilaterales. A ello se suma un mercado de repos muy activo, que permite financiar las posiciones en deuda pública y hace de ella un colateral esencial para el sistema financiero. La liquidez del secundario, medida por los diferenciales entre precios de compra y de venta y por el volumen negociado, es uno de los objetivos del Tesoro. Se ha mantenido elevada en los últimos años, incluso durante la retirada del BCE.

d) El debate sobre el futuro del modelo

El modelo de creadores de mercado afronta dos retos que la literatura reciente ha señalado. El primero es regulatorio. Las reformas bancarias posteriores a 2008 (la ratio de apalancamiento de Basilea III y los requisitos de capital y liquidez) han encarecido que los bancos mantengan grandes carteras de deuda pública y actúen como intermediarios. Duffie (2020) mostró cómo esta restricción contribuyó a las tensiones del mercado de deuda estadounidense en marzo de 2020, cuando los intermediarios no pudieron absorber las ventas masivas. En la zona del euro, el problema se ha amortiguado por la presencia del BCE. Pero con la retirada del Eurosistema y el aumento de la oferta neta al sector privado (3.2.4a), la capacidad de los creadores de mercado para absorber la emisión vuelve a ser una cuestión relevante.

El segundo es estructural. La electronificación y la entrada de nuevos intermediarios no bancarios, como los fondos de alta frecuencia y los creadores de mercado electrónicos, han cambiado la forma en que se provee la liquidez. Algunos Tesoros europeos han revisado sus sistemas de creadores de mercado para adaptarlos, y el español se ha mantenido estable, con ajustes en los criterios de evaluación. No he podido verificar el número actual de creadores de mercado ni la última revisión de su regulación.



La tesis central del documento base sobre este periodo solo es cierta a medias. Según el documento, el Tesoro anticipó el endurecimiento monetario alargando la vida media de la deuda hasta los 8 años para anclar los bajos costes. Eso es cierto para la deuda del Tesoro. Pero si se consolida el balance del Estado con el del BdE, el cuadro cambia. Entre 2015 y 2022 el Eurosistema compró más de 500.000 millones de deuda española y los financió con reservas bancarias remuneradas al tipo de la facilidad de depósito, que es un tipo a un día. Para el sector público en su conjunto, una parte muy relevante de la deuda larga se sustituyó por pasivos a un día.

Cuando el BCE subió los tipos, ese coste llegó por otra vía. El BdE perdió 8.389 millones en sus operaciones de política monetaria en 2024 y no transfirió beneficios al Tesoro en 2023 ni en 2024. Si en el examen presentas el alargamiento de la vida media como una protección completa frente a la subida de tipos, un tribunal puede objetarte que omites el balance del banco central. Por eso he dedicado a esta cuestión un epígrafe propio.

\subsection{El BCE y la estrategia de financiación del Tesoro en 2026}

Introducción

[DOC] El documento cierra el bloque con la estrategia de financiación en la actualidad y un recuadro sobre los programas de compra de activos del BCE. Deja una nota interna que advierte de la necesidad de actualizarlo cada año con la estrategia que el Tesoro publica en enero. Su versión describe la estrategia de 2023. Tras el fin de las compras netas y la primera subida de tipos en más de diez años, en 2022, el Tesoro preveía estabilizar la vida media, cubrir toda la financiación neta (70.000 millones) con instrumentos a medio y largo plazo y reducir las letras en 5.000 millones.

Este último subapartado actualiza esa estrategia a 2026 y la sitúa en su contexto decisivo, la relación con el BCE, en cinco epígrafes. El primero analiza la etapa del BCE como comprador, entre 2015 y 2022 (3.5.1). El segundo examina esa etapa desde el balance consolidado del sector público, que es la perspectiva que el documento omite (3.5.2). El tercero analiza la retirada del BCE y la vuelta de las subidas de tipos, entre 2022 y 2026 (3.5.3). El cuarto presenta la estrategia de 2026 (3.5.4). El quinto ofrece una valoración (3.5.5). La teoría de la política monetaria no convencional se remite a su tema.\footnote{Política monetaria no convencional (expansión cuantitativa, canales de transmisión, efectos sobre la prima por plazo) y el TPI: Temas 3A-36/37.
}

Entre 2015 y 2022, el BCE compró más de 500.000 millones de deuda española con el APP y el PEPP, más de un tercio de la deuda de las AAPP. Eso permitió al Tesoro financiarse a coste negativo en 2021 y alargar la vida media por encima de 8 años, incluso renunciando al ahorro de las letras durante la pandemia.

Desde el balance consolidado del sector público, la expansión cuantitativa acortó la duración efectiva de la deuda. Una parte se financió con reservas remuneradas a un día, y la subida de tipos trasladó ese coste al BdE: pérdidas de 8.389 millones en política monetaria en 2024 y ninguna transferencia al Tesoro en 2023-2024. Es un coste fiscal que no aparece en la cartera del Tesoro y que ha alimentado el debate sobre la remuneración de las reservas.

La retirada del BCE, con el fin de las reinversiones del APP en julio de 2023 y del PEPP en enero de 2025, ha multiplicado la oferta neta de deuda que absorbe el sector privado. La han cubierto sobre todo los no residentes (47,8%) sin aumento de la prima, que está por debajo de 60 puntos básicos, gracias a los fundamentales, al TPI y a la gradualidad. Las subidas del BCE de junio y septiembre de 2026 (hasta el 2,50%) marcan el fin del dinero barato.

La estrategia de 2026 es de continuidad: 55.000 millones netos, unos 285.700 millones brutos, una vida media de unos 8 años, 48 subastas, sindicaciones con demandas récord y bonos verdes. Afronta riesgos de tipos, fiscales, de contagio y de concentración de la base inversora. La valoración final es que la gestión del Tesoro es una buena práctica internacional, pero gestiona el coste y el riesgo de la deuda, no su volumen, que depende de una consolidación fiscal que, como han mostrado los bloques anteriores, España sigue sin completar.

\subsubsection{2015-2022: el BCE como comprador}

a) Los programas

[DOC] Según el documento, a finales de 2014 las medidas convencionales no bastaban para devolver la inflación a su objetivo, y el BCE anunció un paquete de medidas no convencionales. Incluía las operaciones de financiación a largo plazo con objetivo específico (TLTRO) y el programa de compra de activos (APP), que abarcaba deuda pública (PSPP), bonos corporativos (CSPP), titulizaciones (ABSPP) y bonos garantizados (CBPP). En marzo de 2020 se añadió el programa de compras de emergencia frente a la pandemia (PEPP). En ambos casos las compras de deuda pública se reparten según la clave de capital de cada banco central nacional, con cierta flexibilidad temporal. Las compras netas del APP alcanzaron su máximo en 2016-2017, con 80.000 millones mensuales, y las del PEPP llegaron a unos 120.000 millones en junio de 2020. En nota a 3.2, el documento cifra las compras netas de deuda española entre 2015 y 2022 en más de 500.000 millones, más de la tercera parte de la deuda de las AAPP a finales de 2022. Descontada esa deuda, la deuda de las AAPP en manos de otros inversores era a finales de 2022 equivalente a la de principios de 2013.

Dos precisiones de diseño son relevantes para el Tesoro. La primera es que el PSPP limitaba las tenencias del Eurosistema al 33% de cada emisión y de cada emisor. El límite se estableció para no bloquear las cláusulas de acción colectiva (3.3.3d) y para evitar que el BCE se convirtiera en el acreedor dominante de un Estado. La segunda es que el PEPP flexibilizó esos límites y la clave de capital, lo que permitió comprar más deuda de los países más afectados en 2020. Esa flexibilidad fue decisiva para contener las primas en la primavera de 2020 (1.5.1c).

b) Los efectos sobre el Tesoro

[DOC] Según el documento, estos programas redujeron de forma significativa el coste de emisión, porque absorbieron la mayor parte de la emisión neta desde 2015. En 2021 el coste medio de emisión fue negativo por primera vez en la historia del Tesoro (−0,04%), y el coste de la deuda en circulación alcanzó un mínimo histórico del 1,64%. El Tesoro aprovechó esas condiciones para alargar la vida media de la deuda, que en 2021 superó por primera vez los 8 años. Según el documento, mantuvo la estrategia de reducir las letras incluso durante la pandemia, renunciando al ahorro inmediato de la deuda corta para asegurar costes bajos a largo plazo, a diferencia de muchos países de su entorno.

Esta última decisión merece destacarse, porque ilustra la disyuntiva entre coste y riesgo de 3.1.3a. En 2020-2021 la curva era muy plana y los tipos a corto eran negativos. Emitir letras habría sido más barato a corto plazo. El Tesoro eligió comprar seguro frente a una subida de tipos futura, alargando la vida media. Ex post, con la subida de 2022-2023, la decisión resultó acertada para la cartera del Tesoro.

c) El debate sobre la compra de deuda pública por el BCE

Los programas del BCE generaron un debate jurídico y económico que conviene conocer en lo que afecta al Tesoro. En el plano jurídico, el TJUE avaló el OMT en la sentencia Gauweiler (2015) y el PSPP en la sentencia Weiss (2018). Consideró que no vulneraban la prohibición de financiación monetaria del artículo 123 del TFUE, siempre que existieran garantías como los límites por emisión, la compra solo en el mercado secundario y un periodo de espera tras la emisión. En mayo de 2020, el Tribunal Constitucional alemán consideró la sentencia Weiss ultra vires y exigió al BCE justificar la proporcionalidad del programa. Fue un episodio de tensión institucional sin consecuencias prácticas duraderas.

En el plano económico, el debate gira en torno a la dominancia fiscal. Los críticos sostienen que, al absorber buena parte de la emisión neta de los Estados durante años, el BCE redujo la disciplina de mercado y facilitó que los gobiernos retrasaran la consolidación. España, que no corrigió su déficit estructural en 2015-2019 (1.4.3), sería un ejemplo. Los defensores responden que los programas se diseñaron con un objetivo de estabilidad de precios, en un contexto de inflación persistentemente baja. Añaden que las compras según la clave de capital no favorecían a ningún país, y que su efecto sobre las primas fue un efecto secundario de la reducción del riesgo de redenominación, no una financiación selectiva.\footnote{Dominancia fiscal y la interacción entre política monetaria y fiscal (Sargent y Wallace, 1981; teoría fiscal del nivel de precios): Tema 3A-39. Sentencias Gauweiler (C-62/14), Weiss (C-493/17) y la del Tribunal Constitucional alemán de 5 de mayo de 2020: Temas 3A-36/37 y 3B-45.
}

\subsubsection{La perspectiva del balance consolidado: el coste que no se ve en el Tesoro}

a) El argumento

Este epígrafe desarrolla la cuestión que el documento omite. Cuando el Eurosistema compra deuda del Estado, en la práctica sustituye en manos del sector privado un bono del Tesoro (largo y a tipo fijo) por una reserva bancaria en el banco central (a la vista y remunerada al tipo de la facilidad de depósito). Si se consolida el balance del Estado con el de su banco central nacional, que es su propiedad y le transfiere sus beneficios, el efecto neto de la expansión cuantitativa es un acortamiento de la duración de la deuda del sector público consolidado. Una parte de la deuda larga que el Tesoro emitió pasa a estar financiada con pasivos a un día.

Greenwood, Hanson, Rudolph y Summers (2014) documentaron este efecto en Estados Unidos: mientras el Tesoro alargaba la vida media, la Reserva Federal la acortaba en manos privadas. En la zona del euro el razonamiento se aplica al BdE, que tenía en su balance la mayor parte de la deuda española comprada por el Eurosistema, porque las compras de deuda soberana se hacían sobre todo a través de los bancos centrales nacionales y sin compartir riesgos.

b) La materialización: las pérdidas del BdE

La consecuencia se hizo visible con la subida de tipos de 2022-2023. Los bancos centrales del Eurosistema tenían en el activo bonos comprados a tipos muy bajos y en el pasivo reservas bancarias que empezaron a remunerarse al 4% en 2023. El margen se volvió negativo. El BdE registró pérdidas en sus operaciones de política monetaria, que en 2024 ascendieron a 8.389 millones. Pudo absorberlas con las provisiones acumuladas en años anteriores y cerró 2023 y 2024 con un resultado nulo, sin transferencia al Tesoro. En 2025, gracias a las bajadas de la facilidad de depósito entre junio de 2024 y junio de 2025, volvió a beneficios (234 millones, íntegramente transferidos al Tesoro), tras liberar 541 millones de provisiones y con unas pérdidas por política monetaria reducidas a 2.071 millones. En los años previos a 2023, el BdE transfería al Tesoro beneficios del orden de miles de millones anuales.

En términos fiscales, la subida de tipos de 2022-2024 no solo elevó el coste de la nueva emisión del Tesoro. Eliminó además durante dos años un ingreso público relevante, los beneficios del banco central. Esa pérdida de ingreso es la contrapartida del acortamiento de la deuda consolidada que produjo la expansión cuantitativa. Con las subidas de junio y septiembre de 2026, el margen del BdE volverá a estrecharse. Es probable que sus beneficios de 2026 vuelvan a reducirse, aunque el efecto depende de la evolución de su balance, que se reduce a medida que vencen los bonos.

c) El debate sobre la remuneración de las reservas

Esta cuestión ha dado lugar a un debate de política con implicaciones fiscales directas. De Grauwe y Ji (2023) calcularon que, con unas reservas de 4,3 billones en el Eurosistema, remunerarlas al tipo de la facilidad de depósito suponía transferir a los bancos unos 129.000 millones anuales. Propusieron un sistema de dos niveles, con una parte de las reservas sujeta a un coeficiente obligatorio no remunerado, que habría reducido esa transferencia a unos 32.000 millones anuales con un tipo del 3%. El BCE adoptó una medida en esa línea, aunque limitada. En julio de 2023 fijó en el 0% la remuneración de las reservas mínimas obligatorias, cuyo volumen es pequeño, pero no amplió el coeficiente. Los críticos de la propuesta advierten de que un coeficiente amplio no remunerado actuaría como un impuesto sobre los bancos de la zona del euro, podría desplazar depósitos hacia fuera del sistema bancario y dificultaría la transmisión de la política monetaria.

Para este tema, la conclusión es que la estrategia del Tesoro y la política del BCE no pueden evaluarse por separado. El alargamiento de la vida media protegió la cartera del Tesoro, pero la expansión cuantitativa trasladó al banco central una parte del riesgo de tipo de interés que el Tesoro había cubierto, y ese riesgo se materializó en la cuenta de resultados pública por otra vía.\footnote{Sobre el reparto de riesgos en el Eurosistema (las compras de deuda soberana del PSPP se hicieron en un 90% sin mutualización de riesgos, a través de los bancos centrales nacionales) y el tratamiento de las pérdidas de los bancos centrales: Temas 3A-36/37.
}

\subsubsection{2022-2026: la retirada del BCE y la vuelta de las subidas de tipos}

a) La cronología y una corrección

[DOC] El documento señala que el BCE dejó de hacer compras netas en 2022 y que desde marzo de 2023 reinvierte solo parcialmente los vencimientos, con lo que empezó la reducción de su balance (endurecimiento cuantitativo o QT).

Precisión. La afirmación es correcta para el APP pero incompleta, y el calendario conviene tenerlo exacto. Las compras netas del PEPP terminaron en marzo de 2022 y las del APP en junio de 2022. En el APP, la reinversión se redujo en 15.000 millones mensuales entre marzo y junio de 2023 y se suprimió por completo desde julio de 2023. En el PEPP, los vencimientos se reinvirtieron íntegramente hasta junio de 2024, se redujeron en unos 7.500 millones mensuales en el segundo semestre de 2024 y dejaron de reinvertirse desde enero de 2025. Desde entonces el Eurosistema no reinvierte nada, y sus carteras se reducen al ritmo de los vencimientos. El BCE lo reiteró en su decisión de septiembre de 2026.

En cuanto a los tipos, el BCE subió la facilidad de depósito del −0,50% al 4% entre julio de 2022 y septiembre de 2023. La bajó después, desde junio de 2024 hasta el 2% en junio de 2025. Tras la crisis de Oriente Medio, la subió al 2,25% en junio de 2026 y al 2,50% en septiembre de 2026. En marzo de 2024 el BCE revisó su marco operativo, redujo a 15 puntos básicos el diferencial entre el tipo de las operaciones principales de financiación y la facilidad de depósito, y confirmó que esta seguirá siendo el tipo que orienta la política monetaria. Por eso, tras la subida de septiembre de 2026, las operaciones principales se remuneran al 2,65%.

b) Los efectos sobre la financiación española

La retirada del BCE ha tenido tres efectos sobre el Tesoro, que se han analizado en 3.2 y 3.3 y aquí se reúnen. El primero es el aumento de la oferta neta al sector privado, que ha pasado de ser negativa a rondar probablemente los 90.000-100.000 millones anuales (3.2.4a). El segundo es el cambio en la base inversora. Los no residentes han pasado a ser los principales compradores de bonos y obligaciones y tenían el 47,8% de la deuda en 2025, y los minoristas volvieron a las letras en 2023. [DOC] El documento ya lo anticipaba: la deuda del Estado en manos de personas físicas tocó mínimo en abril de 2022, con 898 millones (el 0,08% del total), y empezó a aumentar después. A finales de 2022 los minoristas españoles tenían apenas el 0,26% de la deuda del Estado, frente al 3,8% de media de la Unión, el 7,9% de Italia y el 13,2% de Portugal. El documento añadía que, como esa demanda se concentra en el corto plazo, el Tesoro podría plantearse aumentar ligeramente la emisión de letras. El tercero es el coste. El coste medio de emisión subió al 3,44% en 2023, bajó al 3,16% en 2024 y al 2,70% en 2025 con las bajadas del BCE. El de la cartera subió del 1,64% de 2021 al 2,31% de 2025.

Lo que no ha ocurrido es un aumento de la prima de riesgo. Contra lo que temían muchos analistas en 2022, la retirada del BCE ha coincidido con una caída de la prima española por debajo de 60 puntos básicos (2.3.3a). Hay tres explicaciones complementarias. La primera es la mejora de los fundamentales relativos de España. La segunda es la existencia del TPI, que actúa como red de seguridad aunque no se haya activado nunca. La tercera es la gradualidad de la retirada, que el BCE ha descrito como un proceso en segundo plano. El riesgo es que esa calma dependa de que el TPI siga siendo creíble y de que España cumpla las condiciones de elegibilidad, entre ellas el respeto del marco fiscal europeo. Es precisamente lo que la AIReF pone en duda para 2026 (1.5.3d).

\subsubsection{La estrategia de financiación de 2026}

a) Las cifras

La Estrategia de Financiación de 2026, presentada en diciembre de 2025, prevé una emisión neta de 55.000 millones, igual que en 2025: 50.000 en bonos y obligaciones y 5.000 en letras. La emisión bruta será de unos 285.677 millones, un 4% más que en 2025: 176.935 millones de bonos y obligaciones, para cubrir unas amortizaciones de unos 126.935 millones, y unos 108.742 millones de letras. El programa se ejecuta mediante 48 subastas ordinarias y varias sindicaciones, para nuevas referencias y para el programa verde. Se completa con 6.500 millones de préstamos del Mecanismo de Recuperación, y en 2026 se amortizan unos 3.643 millones del préstamo del MEDE.

El Tesoro se fija como objetivo preservar la vida media de la cartera, que fue de 7,93 años en 2025, en torno a su máximo histórico por quinto año consecutivo, y limitar los riesgos de refinanciación. Mantiene los bonos verdes como elemento estructural del programa y parte de una base inversora en la que los internacionales se han consolidado como principales compradores. Destaca las mejoras de calificación y la bajada de la prima de riesgo a mínimos de 2009.

b) La ejecución en 2026

La ejecución hasta septiembre de 2026 ha confirmado la solidez de la demanda. En enero, la sindicación a diez años de 15.000 millones recibió una demanda récord de 145.000 millones. En mayo hubo otra sindicación a diez años de 13.000 millones, y en septiembre una nueva referencia verde a veinte años de 4.000 millones, con una demanda de 68.500 millones. Pero también muestra el cambio del entorno. La obligación a diez años de enero se emitió con un cupón del 3,30%, y la verde a veinte años de septiembre a una rentabilidad cercana al 4,2%. Ambas están muy por encima del 2,31% de coste medio de la cartera, lo que garantiza que el coste medio seguirá subiendo en 2026 y 2027.

c) La lógica de las decisiones

La estrategia de 2026 puede leerse a la luz de los objetivos de 3.1.3. La continuidad (misma emisión neta, misma vida media, mismo calendario) responde al objetivo de previsibilidad. En un año de shocks (la guerra en Oriente Medio, las subidas del BCE, la ausencia de presupuestos), el Tesoro evita añadir incertidumbre sobre la oferta.

El mantenimiento de la vida media en 8 años, con tipos subiendo, responde a la disyuntiva entre coste y riesgo. Acortar la deuda abarataría la emisión inmediata si la curva mantiene su pendiente, pero aumentaría la exposición a nuevas subidas de tipos y el volumen a refinanciar cada año, ya en torno al 16% del PIB. Con un BCE que sube tipos por un choque de oferta de duración incierta, el Tesoro opta por el seguro frente al ahorro, como en 2020-2021.

El ligero aumento neto de las letras (5.000 millones), tras años de reducción, puede interpretarse de dos formas compatibles. Por un lado, recoge la demanda minorista y de fondos monetarios a corto plazo que el documento anticipaba. Por otro, da un pequeño margen de flexibilidad de tesorería en un año con necesidades inciertas por las medidas de emergencia.

d) Los riesgos del escenario

La estrategia de 2026 afronta riesgos que conviene exponer en el examen como cierre del tema. El primero son los tipos de interés. Si la inflación energética se prolonga y el BCE sigue subiendo, el coste de la nueva emisión aumentará, y con él la carga de intereses proyectada. A ello se suma el contexto global: la rentabilidad del bono estadounidense a diez años se acercaba al 5% en septiembre de 2026. El segundo es el riesgo fiscal: la AIReF prevé un crecimiento del gasto neto del 6,4% frente al 3,5% comprometido (1.5.3d), y España lleva tres años con los PGE de 2023 prorrogados. Un incumplimiento prolongado del marco europeo podría afectar a la percepción de los mercados y, en el límite, a la elegibilidad para el TPI. El tercero es el contagio. La inestabilidad fiscal y política de Francia, cuya prima frente a España está en máximos, puede trasladarse al conjunto de la zona del euro en un episodio de aversión al riesgo. El cuarto es la concentración de la base inversora en no residentes (47,8%), más sensibles a las condiciones financieras globales. El quinto es el fin de la financiación europea desde 2027 (3.2.4b).

En sentido contrario, hay factores favorables. Las calificaciones siguen mejorando: Scope subió a España a A+ el 25 de septiembre de 2026. El crecimiento está por encima de la media, el saldo primario cerca del equilibrio y la vida media es larga. Todo ello da al Tesoro un margen que no tenía en 2011-2012.


\subsection{Una valoración}

La política de financiación del Tesoro español es, con criterios internacionales, un ejemplo de buena práctica. Tiene un marco jurídico claro, un mandato de coste y riesgo, procedimientos previsibles y transparentes, una cartera larga, instrumentos estandarizados y líquidos y una base inversora diversificada. Su gestión en 2020-2021 (alargar la vida media en lugar de aprovechar los tipos cortos negativos) resultó acertada con la subida de 2022-2023.

Pero la gestión de la deuda tiene un límite que conviene subrayar al cerrar el tema. El Tesoro gestiona la composición de la deuda, su coste y su riesgo, no su volumen, que depende de la política fiscal. La mejor estrategia de financiación no puede compensar un déficit estructural que no se corrige. Como mostró el Bloque I, la reducción de la deuda en 2021-2025 se debió al crecimiento nominal y no a superávits primarios. Como mostró el Bloque II, España no ha generado superávits primarios en las fases favorables del ciclo, a diferencia de Portugal. La calidad de la gestión del Tesoro reduce el coste y el riesgo de esa deuda, pero no sustituye a la consolidación fiscal.

[DOC] En su conclusión, el documento plantea este límite cuando advierte que el éxito de la política de financiación dependerá de la capacidad de la economía para crecer y de acometer una nueva consolidación fiscal, porque niveles muy elevados de deuda son fuente de vulnerabilidad y pueden afectar al crecimiento. En nota cita el debate sobre el umbral del 90% de Reinhart y Rogoff (2010) y su crítica por Herndon, Ash y Pollin (2014), que cuestionan el umbral y sugieren que la causalidad puede ir del bajo crecimiento a la deuda alta.\footnote{Reinhart y Rogoff (2010), la crítica de Herndon, Ash y Pollin (2014) y la literatura posterior sobre la relación entre deuda y crecimiento (causalidad inversa, heterogeneidad entre países, ausencia de umbral universal): Tema 3A-38.
} Es la idea con la que conviene cerrar el tema, y la retomaré en la conclusión general.

















































\clearpage
\section{Saldo Presupuestario: }
\puntosclave{
    Aunque el déficit público no aparece en España hasta la Transición. la situación de equilibrio presupuestario previa a ese momento ocultaba el atraso del sistema tributario español y sólo era posible por el reducido tamaño del sector público, consecuencia a su vez del escaso desarrollo del estado de bienestar.
    
    La creación del estado de bienestar y de las autonomías llevó aparejado un importante aumento del gasto público de carácter estructural que no fue acompañado de medidas equivalentes por el lado de los ingresos. iniciándose un periodo de déficits recurrentes. Esos déficits han sido estructuralmente negativos de manera casi sistemática y. además, procíclicos, algo especialmente indeseable.

    En 1996 comienza una etapa muy favorable para el saneamiento de las finanzas públicas que permitió una mejora continuada del saldo presupuestario, que llegó a presentar los tres únicos superávits de la democracia entre 2005 y 2007.
    
    La profunda recesión en la que entró España en 2008 provocó un gran deterioro de las cuentas públicas, llegando a registrarse déficits anuales superiores al 10\% del PIB, que provocaron que la deuda pública alcanzase en 2014 el 105\% del PIB.

    En 2005 comenzó un lento proceso de consolidación presupuestaria que se vio frenado en seco con el estallido de la pandemia del covid-19, que ha situado nuestro coeficiente de deuda pública como el 7° más alto de las economías avanzadas y el 4° más alto de la zona euro, con un 113,1\% del PIB en 2022.
    
    El problema estructural de las cuentas públicas españolas, por tanto. no lo es solo en términos absolutos, sino que también queda de manifiesto cuando comparamos los datos con los de los países de nuestro entorno.
}


El presente tema trata sobre el saldo presupuestario3 y la deuda de las AAPP en España, por lo que vamos a comenzar definiendo esos dos conceptos:\footnote{Para un mayor detalle, puede consultarse el tema 4B 1.

Cabe destacar que. como veremos en la última parte del tema. lo relevante a efectos de la política de financiación del Tesoro Público no será el saldo presupuestario de las AAPP. sino la necesidad de endeudamiento del Estado
}

• El saldo presupuestario es la diferencia entre los ingresos y gastos totales de las AAPP durante el ejercicio presupuestario, es decir, la capacidad o necesidad de financiación de las AAPP. Siguiendo el SEC-201 O, el saldo presupuestario corresponde con el saldo de la cuenta de capital del sector institucional AAPP, y es un indicador esencial para realizar comparaciones internacionales y para valorar el cumplimiento de las reglas fiscales de la Unión Europea.

• La deuda, por su parte, se suele definir como el saldo en circulación en un momento determinado de los pasivos contraídos para financiar aquellos gastos que no pueden ser financiados con ingresos corrientes y de capital. Existen diferentes conceptos de deuda, en función de los instrumentos incluidos y del método de valoración utilizado, siendo el más relevante a efectos de este tema el definido en el Protocolo de Déficit Excesivo (PDE). Dicho criterio valora la deuda considerada utilizando su valor nominal o facial, evitando así las grandes variaciones en los niveles de deuda que podrían producirse si se utilizara el valor de mercado.

Respecto a la delimitación de las AAPP, podemos hacerla siguiendo las indicaciones del PDE,según el cual estas se componen de:

1. La Administración Central, constituida por aquellas unidades institucionales con competencias en todo el territorio nacional. Incluye el Estado (Casa del Rey, Gobierno, Cortes Generales, Tribunal Constitucional, etc.), los organismos de la administración central y las empresas que se clasifican como administración central.

2. Las comunidades autónomas (CCAA), que incluyen los órganos de gobierno de las CCAA y los organismos autónomos de carácter administrativo y similares. Esto último incluye las universidades dependientes de cada CCAA y las empresas que se clasifican como CCAA.

3. Las corporaciones locales, constituidas por los ayuntamientos, diputaciones y cabildos, sus mancomunidades y agrupaciones, las ciudades autónomas (Ce uta y Melilla) y los organismos autónomos de carácter administrativo y similares dependientes.

4. Las administraciones de la Seguridad Social, que incluyen las unidades institucionales de ámbito territorial diverso que llevan a cabo funciones relacionadas con la provisión de prestaciones sociales.

Además, con carácter general, se consideran parte de las AAPP aquellas unidades institucionales controladas por las AAPP y que son productores no de mercado, definidos por el SEC 201 O como aquellas unidades institucionales que no cubren más del 50% de sus costes de producción con los ingresos procedentes de la venta de sus productos.

Una vez definidos estos conceptos, pasamos a analizar la evolución reciente del saldo presupuestario y de la deuda de las AAPP en España.


\subsection{Desarrollo histórico: }
Como veíamos en la introducción, los años previos a la Transición, España fue capaz de mantener el equilibrio presupuestario e incluso pequeños superávits, algo que se veía facilitado por el reducido nivel de gasto público, que rondaba el 25% del PlB (20 puntos porcentuales inferior a la media de los países comunitarios). El reducido tamaño del sector público, consecuencia del escaso desarrollo del estado de bienestar, permitía disimular el atraso del sistema tributa1io español, que contaba con una escasa capacidad recaudatoria y era ineficiente y escasamente redistributivo.

Además, el recurso al sefioreaje era frecuente, por lo que el stock de deuda pública en España en 1975 se situaba por debajo del 10\% del PlB, frente al casi 50\% de media en los países comunitarios.

En definitiva, bajo la aparente solidez de las cuentas públicas de 1975 se escondía un sistema tributario poco eficiente y un gasto público insuficiente.

i. Desde la Transición hasta la adhesión a la CEE (1975-1985)

La Constitución de 1978 consagra el estado de bienestar y de las autonomías. lo que llevó aparejado un importante aumento del gasto público de carácter estructural. El aumento del gasto público en este periodo, no obstante, tuvo además un importante componente coyuntural derivado de la crisis económica que atravesaba Éspafia que, en el contexto internacional de las crisis del petróleo, hizo necesaria una reconversión industrial que absorbió grandes cantidades de recursos públicos.

Para responder a estos aumentos del gasto y profundizar en la creación del estado de bienestar, en los Pactos de la Moncloa ( 1977) se acordó llevar a cabo una profunda reforma del sistema fiscal, por la que se revisaron el IRPF y el Impuesto sobre Sociedades (dotándoles de la forma básica que hoy conocemos), se creó el Impuesto sobre el Patrimonio y se eliminaron algunas figuras arcaicas.\footnote{Se derogaron el impuesto de producto y el impuesto --general sobre la renta. sustituyéndolos por un sistema de retenciones en la fuente sobre algunos rendimientos -trabajo y capital. básicamente- y por un impuesto sobre la renta de las personas físicas (1 RPF) de tarifa progresiva.
}

Pese a ello, las reformas por el lado de los ingresos tardaron mucho más en hacerse sentir que las reformas por el lado del gasto (que en 1985 superaba ya el 40% del PIB), iniciándose una etapa de omnipresencia del déficit público. Ese mismo año. el déficit alcanzó el 6,9% del PIB y la deuda el 4 5%.

ii. 1986-1990: reducción del déficit público

La adhesión de España a la CEE en 1986 supuso un cambio muy importante para el país. Entre otras cosas, la mejor coyuntura económica permitió conseguir crecimientos anuales del entorno del 5%, lo que redujo a su vez la presión al alza sobre el gasto causada por la actuación de los estabilizadores automáticos. Además, la legislación comunitaria hizo necesario reducir o eliminar muchas de las subvenciones puestas en marcha para hacer frente al proceso de reconversión industrial.

Por el lado de los ingresos, comienzan a dar sus frutos las reformas tributarias de los Pactos de la Moncloa, y a esos ingresos se le añaden desde 1 9 86 los proporcionados por el IV A, un impuesto de obligatoria creación para la integración en la CEE. 

En estas circunstancias, se puede hablar de un cierto proceso de consolidación fiscal, situándose el déficit público en 1990 por debajo del 3\% del PIB. La ratio de deuda pública sobre PIB, por su parte, se consiguió estabilizar en los niveles de 1985.

iii. 1990-1995: incremento sustancial del gasto y del déficit público

Desde 1990 se ralentiza el ritmo de crecimiento de la economía española, lo que desembocó en la crisis económica para España de 1992-93. Esa crisis se produjo, en parte, por las debilidades estructurales en el diseño del Sistema Monetario Europeo (SME), al que España se había incorporado en 19 89. De hecho, fueron varias las economías europeas que sufrieron una crisis en esos años, y algunas de ellas (Reino Unido'. Italia, Finlandia, Suecia y Noruega) abandonaron el SME, mientras que otras, como España, tuvieron que devaluar sus monedas con respecto al ECU.\footnote{En el caso de la peseta, se produjeron hasta tres devaluaciones con respecto al ECU. la última en marzo de 1993. En agosto de ese año, las presiones persistentes sobre el tipo de cambio obligaron a ampliar la banda de fluctuación del tipo de cambio de las ocho divisas que seguían conformando el SME con respecto al ECU desde el ±6\% hasta el ± 15\%.
}

Esos años coinciden además con un periodo de fuerte inversión en infraestructuras, destacando la construcción de la primera línea de A VE entre Madrid y Sevilla; y de gasto asociado a la organización de los JJ.OO. de Barcelona y de la Expo de Sevilla. A ello se le añadió el aumento del gasto en prestaciones por desempleo y del coste del servicio de la deuda, consecuencia tanto de la mayor deuda, que llegó a alcanzar el 65\% del PIB en 1995, como de los mayores tipos de interés exigidos por los inversores ante el aumento del déficit público, que llegó a superar el 7\% del PIB en 1993.\footnote{La reforma del mercado laboral de 1984 incentivó la contratación temporal lo que provocó que, cuando llegó la crisis, el aumento del desempleo fuera muy acusado. Así, el desempleo (cuyas prestaciones. además, se habían incrementado tras la huelga general de 1988) alcanzó el 24,5% en 1994.
} Como resultado, el gasto público llegó a suponer más del 45\% del PIB en 1993. 

Entre los requisitos definidos por el Tratado de Maastricht de 1992 para entrar a formar parte de la UEM, no obstante, se incluyó situar el déficit público por debajo del 3\% del PIB, lo que hizo necesario iniciar un proceso de consolidación fiscal que se vio favorecido desde 1995 por la mejora de la coyuntura económica.

iv. 1996-2007: proceso de consolidación fiscal

Entre 1996 y la crisis de 2007 tuvo lugar una etapa muy favorable para el saneamiento de las finanzas públicas que permitió una mejora continuada del saldo presupuestario, que llegó a presentar los tres únicos superávits de la democracia entre 2005 y 2007. Fue un periodo de crecimiento, en el que:
(i) la tasa de variación del PIB no bajó nunca del 2%, lo que favoreció el aumento de la base imponible de los principales impue_stos y, con ello, de la recaudación;
(ii) el desempleo llegó a situarse por debajo del 8% en 2007, reduciendo el gasto en prestaciones por desempleo;
(iii) se puso en marcha el Programa de Modernización del Sector Público Empresarial del Estado, por el que se privatizaron hasta 36 empresas en 4 años, que permitieron recaudar 30.000 millones de euros. un importe equivalente a más del 6% del PIB de 1996.

Desde el punto de vista estructural, se llevaron a cabo reformas del sistema de pensiones, del sistema de financiación autonómico y del sistema de prestaciones por desempleo, encaminadas todas ellas a reducir el nivel de gasto público. También se pusieron en marcha algunas medidas para reducir el gasto público a corto plazo, como la congelación de los salarios de los funcionarios en 1997.

Todo ello permitió a España cumplir en 1998 con los criterios de convergencia, registrando un déficit público del 2,7\% del PIB y un nivel de deuda pública que, aunque ligeramente superior al límite del 60\%, se aceptó dada la tendencia a la baja que venía mostrando.

La entrada en la UEM coincidió en España con un periodo de fuerte expansión económica que se extendió hasta el año 2007. La buena coyuntura, unida a la reducción de los tipos de interés que siguió a los primeros años de pertenencia a la UEM, permitió a España reducir significativamente sus niveles de deuda y de déficit públicos.\footnote{Estos buenos datos no deben hacer olvidar los importantes desequilibrios que se fueron originando durante ese periodo, entre los que destacan el enorme crecimiento del endeudamiento privado, el mantenimiento de un diferencial de inflación positivo respecto a nuestros socios de la zona euro (con la consiguiente pérdida de competitividad precio) y el creciente déficit por cuenta corriente
} Concretamente, España mantuvo desde 1999 hasta 2004 déficits públicos de menos del 1\% del PIB y. entre 2005 y 2007, posiciones superavitarias, que pennitieron situar los niveles de deuda pública por debajo del 36% del PIB en 2007.

Esta evolución significó también que España, a diferencia de países como Francia o Alemania, pudo cumplir inicialmente sin excesivos problemas con el Pacto de Estabilidad y Crecimiento (PEC), aprobado en el Consejo Europeo de Ámsterdam de 1997 con el objetivo de garantizar una cierta disciplina fiscal una vez constituida la zona euro.\footnote{Fue. de hecho, la no imposición de sanciones a Francia y Alemania por su incumplimiento del PEC en 2003 lo que llevó a su refonna en 2005 para fle:-.ibilizar su aplicación. La no imposición de sanciones a Francia y Alemania por parte del Ecofin en 2003 provocó incluso una demanda de la Comisión contra el Consejo de la Unión. si bien el Tribunal de Justicia de la Unión Europea dictaminó que el Consejo de la Unión tenía derecho a no adoptar las recomendaciones de la Comisión (que. en este caso. pasaban por la imposición de sanciones a Francia y Alemania).
} El PEC llevaría a la aprobación en España en 2001 de la Ley General de Estabilidad Presupuestaria, que insertaba las exigencias del Pacto en la legislación española y que sería el marco regulador de la política fiscal española hasta el año 2012, cuando se aprobó la Ley Orgánica de Estabilidad Presupuestaria y Sostenibilidad Financiera (LOEPSF) (previamente, la LGEP había sido modificada en 2006, tras la revisión del PEC en 2005).

1 1 . 1 1 1 . Desde la crisis financiera internacional hasta la crisis de la covid- 1 9 (2008-
2019)

En el verano de 2007 se produce el estallido en EE. UU. de la crisis financiera internacional, que se contagia rápidamente a la economía europea. Así, España entró en 2008 en una profunda recesión que afectó al saldo de las cuentas públicas, pasándose de un superávit del 1,9\% del PIB en 2007 a un déficit del 4, 1 % en 2008 y de más del 11 % en 2009. Hasta 2012, además, el déficit no bajó del 9,5% del PIB. Como resultado, la ratio de deuda sobre PIB aumentó en casi 70 puntos porcentuales, desde el 36% del PIB en 2007 hasta el 1 05% del PIB en 2014.

Este gran deterioro de las cuentas públicas se explica por varios factores, entre los que destacan (i) la acción de los estabilizadores automáticos, que llevaron al aumento del gasto público (la tasa de paro, por ejemplo, llegó a superar el 25\%) y a la disminución de los ingresos, provocada no solo por la reducción del PIB, sino también por la caída de la presión fiscal (concretamente, la recaudación como porcentaje del PIB cayó más de 6 puntos porcentuales, una caída que no se dio en ningún otro país de nuestro entorno); (ii) la adopción de políticas de demanda expansivas para sostener la demanda interna en los inicios de la crisis (p.ej. Plan E); (iii) las ayudas al sector bancario (incluida su recapitalización en 2012), muy afectado por el estallido de la burbuja inmobiliaria; y, finalmente, .(iv) la crisis de deuda soberana9, que disparó el coste de financiar el déficit público.\footnote{El origen de la crisis de deuda soberana de la zona euro se encuentra la pérdida de confianza de los inversores en la solvencia de algunos de sus Estados miembros y, en particular, de los países periféricos. Esa pérdida de confianza se explica por varios factores: • La revisión al alza de los datos de déficit de Grecia por parte del nuevo gobierno a finales de 2009, que provocó que los inversores empezaran a dudar de la capacidad del país para hacer frente a sus deudas. • . El Acuerdo de Deauville de octubre de 201 O, por el que Merkel y Sarkozy acordaron proponer, entre otras cosas, que los fu turos programas de asistencia financiera del MEDE requirieran la imposición de quitas entre los acreedores privados. Hasta ese momento. la posición oficial había sido que la deuda griega era sostenible, pero la declaración franco-alemana hizo explícita la posibilidad de que los acreedores privados tuvieran que aceptar quitas de su deuda soberana. • La posibilidad de una ruptura de la zona euro y la vuelta a las monedas nacionales, con el consiguiente riesgo de redenominación de la deuda, que hubiera perjudicado especialmente a los tenedores de deuda de países cuyas monedas se hubieran eventualmente devaluado (como, por ej emplo, la española).

En este contexto, se produjeron ataques especulativos masivos contra la deuda soberana de los países periféricos de la zona euro, lo que encareció de fonna notable la financiación de esos países y puso en peligro el gran supuesto implícito del que partía el proceso de integración monetaria europeo: su irreversibilidad. Solo el famoso whatever it takes'· de Mario Draghi a finales de julio de 2012 y el anuncio de la creación del programa OMT, al eliminar por completo la posible ruptura de la zona euro y el consiguiente riesgo de redenominación. permitió que tuviera lugar una bajada de las primas de riesgo.
}

El deterioro de las cuentas públicas provocó la apertura de un procedimiento de déficit excesivo contra España en el año 2009, cambiándose radicalmente de estrategia presupuestaria en 2010.\footnote{Ese procedimiento, que en un primer momento se esperaba que tenninara en 2012, no se cerraría finalmente hasta junio de 2019, cuando el Consejo de la Unión Europea certificó que, en 2018, España había reducido su déficit hasta un nivel inferior al 3\%.
} Así, se adoptan una serie de medidas de consolidación fiscal encaminadas a situar el déficit por debajo del 3\% lo antes posible, si bien dicho objetivo no se conseguiría hasta 2018. Entre las medidas adoptadas desde el 2010, destacan:

• La reducción del gasto no financiero del Estado, con medidas como el recorte y congelación posterior del sueldo de los trabajadores públicos, la reducción de las prestaciones por desempleo, la reducción de la inversión pública o el descenso del gasto en educación, sanidad y ayuda oficial al desarrollo, entre otras cosas.\footnote{El límite máximo de gasto no financiero del presupuesto del Estado pasó de 182.439 millones de euros en 201 O a 117.353 millones de euros en 2012.
}

• Medidas de aumento de los ingresos, tanto de carácter extraordinario (amnistía fiscal en 2012 para incentivar la regularización tributaria de los contribuyentes del IRPF, del IS y del IRNR;\footnote{Los ingresos públicos, que habían alcanzado los 308.559 millones de euros en 2008, cayeron hasta 272.363 millones de euros en 2011, iniciando a partir de 2012 una lenta recuperación. El nivel de ingresos públicos de 2008. no obstante, no volvería a superarse hasta 2021.
} privatización parcial de AENA en 2015) como de aumento de la presión impositiva (establecimiento de un recargo de solidaridad en el I RPF entre 2012 y 2014 e incremento del tipo de gravamen aplicable a la base imponible del ahorro; restablecimiento del impuesto sobre el patrimonio; establecimiento de nuevos límites a la aplicación de deducciones en el impuesto sobre sociedades; incremento de los tipos de
gravamen general y reducido del IV A y de algunos impuestos especiales, entre otrap
cosas).

A nivel europeo, la crisis de deuda soberana dejó en evidencia los fallos de diseño de la UEM y. tras una fase inicial en la que la falta de consenso entre los países de la zona euro impidió dar una respuesta coordinada a la crisis, desde mayo de 201 O se empiezan a poner en marcha medidas coordinadas (\textcolor{red}{\textbf{Ver anexo 3.}}). Al margen de las medidas de carácter coyuntural tomadas para evitar el colapso del sistema, se decidió renovar en profundidad el entramado institucional y la gobemanza económica de la UEM. A este respecto, se avanzó en la coordinación de las políticas económicas -Semestre Europeo, Pacto por el Euro Plus-. se reformó el marco fiscal de la UEM -PEC reforzado, Tratado de Estabilidad, Coordinación y Gobernanza-, y se creó un mecanismo para tratar de evitar que surgieran desequilibrios macroeconómicos como los que habían llevado a lacrisis de deuda soberana -Procedimiento de Desequilibrios Macroeconómicos-.

En España, la reforma del marco fiscal europeo se tradujo en la revisión del artículo 135 de la Constitución, que se aprobó en septiembre de 2011. Por la misma se elevó a rango constitucional la obligación de que todas las AAPP adecúen sus actuaciones al principio de estabilidad presupuestaria. En 2012, además, se aprobó la Ley Orgánica de Estabilidad Presupuestaria y Sostenibilidad Financiera (LOEPSF), que desarrolla la reforma constitucional y reemplaza las leyes de estabilidad vigentes hasta ese momento, que se habían mostrado insuficientes para disciplinar las finanzas públicas.\footnote{En 2014. además. comienza a funcionar la Autoridad I ndependiente de Responsabilidad Fiscal (AIReF), una institución pública encargada de velar por la sostenibilidad de las finanzas públicas españolas y por el cumplimiento efectivo por las AAPP del principio de estabilidad presupuestaria previsto en el artículo 135 de la Constitución. Aunque es un organismo adscrito, a efectos organizativos y presupuestarios, al Ministerio de Hacienda y Función Pública, ejerce sus funciones con plena independencia respecto de las AAPP. Tiene encomendadas la evaluación continua del ciclo presupuestario (con la finalidad de detectar de fomia temprana las posibles desviaciones en los objetivos perseguidos). del endeudamiento público y de las previsiones macroeconómicas que se incorporen a los proyectos de presupuestos y escenarios a medio plazo. Su presidente es nombrado por un periodo de seis años. no prorrogables. por el Consejo de Ministros. y debe ser aprobado por mayoría absoluta del Congreso de los Diputados o, en su caso, mayoría simple del Senado.
} 

Para el cumplimiento de los objetivos de estabilidad presupuestaria y sostenibilidad financiera, la LOEPSF se sirve de tres reglas fiscales:\footnote{Se ha de resaltar que, como en la mayoría de los Estados miembro. la norma española no coincide exactamente con la europea, pese a ser compatible jurídicamente con el marco europeo aprobado entonces. Así, la regla de gasto de la normativa nacional difiere de la de la Unión tanto en la definición de la variable sobre la que se aplica. corno en el ritmo de ajuste estructural anual requerido. En cuanto a lo primero, mientras que en la Unión la regla se aplica a todo el gasto salvo el de intereses de la deuda y el correspondiente al desempleo cíclico, en España se aplica a todo el gasto salvo el de intereses de la deuda. desempleo cíclico y pensiones. En cuanto a lo segundo. mientras que según la LOEPSF España debería haber realizado un ajuste estructural anual de 0,8 puntos porcentuales entre 2013 y 2020, los requerimientos anuales de ajuste estructural europeo en ese mismo periodo oscilaron entre I ,1 y -0,4 puntos porcentuales. Para más detalle, véase Crespo. Ramos, Rodríguez (2020).
}
• Equilibrio o superávit estructural: el artículo 1 1 de la LOEPSF establece que ninguna AAPP podrá incurrir en déficit estructural, entendido como el déficit ajustado del ciclo, neto de medidas excepcionales y temporales. Se establecen, no obstante, ciertas excepciones en caso de catástrofe natural, recesión económica grave o situaciones de emergencia extraordinaria que escapen al control de las AAPP, así como la posibilidad de alcanzar un déficit estructural del 0,4\% del PlB cuando el motivo sea la puesta en marcha de reformas estructurales con efectos presupuestarios a largo plazo.
• Límite a la ratio de deuda del 60\% del PIB: el objetivo de esta regla es incorporar la estabilidad en un contexto intertemporal, para lo que se establece un límite de deuda del conjunto de las AAPP del 60\% del PIB16, equivalente al límite previsto en la normativa europea.\footnote{A la Administración Central le corresponde un 44 % del PIB: al conjunto de CCAA. un 13 % (este límite se aplica también a cada una de ellas con respecto a su PIB regional); y a las corporaciones locales. un 3%. En el caso del déficit estructural: 0,26% del PIB para la Administración Central y O, 14% para las CCAA.
}
• Regla de gasto: la variación del gasto computable de las AAPP17 no podrá superar la tasa de crecimiento del PIB a medio plazo de la economía española18.\footnote{El gasto computable no incluye los intereses de la deuda, el gasto no discrecional en prestaciones por desempleo, el gasto en pensiones, y otros elementos que escapan al control de la autoridad fiscal. Además, si se aprueban cambios normativos que incrementen (reduzcan) permanentemente la recaudación, el límite de gasto deberá aumentar (disminuir) en una cuantía equivalente (i.e. si se aprueban medidas que aumenten la recaudación, el gasto podrá crecer por encima de la tasa de crecimiento del PIB a medio plazo).

Se define como la media de las estimaciones (según la metodología utilizada por la Comisión) del crecimiento potencial del PIB de I O años (los últimos 5 años, el año corriente y las proyecciones para los siguientes 4 años), y su objetivo es evitar generar un comportamiento procíclico del gasto público.
} 

Estas reglas, no obstante, se encuentran suspendidas desde el año 2020, y cabe prever que sean modificadas en línea con la revisión de las reglas fiscales europeas (\textcolor{red}{\textbf{Ver anexo 4}}).

En paralelo a lo anterior, España solicitó, en junio de 2012, asistencia financiera del Mecanismo Europeo de Estabilidad (MEDE) para culminar el proceso de reestructuración y recapitalización de su sistema bancario.\footnote{E l Mecanismo Europeo d e Estabilidad (MEDE) fue e l mecanismo d e rescate de carácter permanente que sustituyó a los dos instrumentos financieros europeos de carácter temporal creados en un primer momento para brindar asistencia financiera a los países que lo necesitaron: el Mecanismo Europeo de Estabilización Financiera (MEEF) y el Fondo Europeo de Estabilidad Financiera (FEEF).
} Esta asistencia financiera sería aprobada tras la firma del Memorandum of Understanding entre España y la Unión Europea en julio de 2012, y se concretó en una línea de crédito del MEDE de un máximo de 100.000 millones de euros, de los que España utilizó 4 1.333 millones22.\footnote{La ayuda quedó sometida a tres tipos de condicionantes:
• En el caso de los bancos directamente afectados, se estableció que los planes de recapitalización que implicasen el uso de fondos públicos tendrían que demostrar la viabilidad a largo plazo del banco sin mantener ayudas de Estado. Para ello. se aprobaron planes de reestructuración en los que se contemplaba una reducción significativa de negocios no rentables, la reestructuración de sucursales y plantilla, y la reducción de la dependencia del suministro de liquidez del BCE.
• Se estableció también una condicionalidad horizontal. a aplicar sobre todo el sistema financiero. Esta implicaba la reforma de la legislación sobre provisiones para préstamos morosos. el establecimiento de un control trimestral de la liquidez de las entidades de crédito receptoras de ayudas de Estado o con un déficit de capital en los test de estrés, la reforma de las cajas de ahorro, una mayor independencia operativa del BdE, reformas en materia de protección del consumidor, el fortalecimiento de la intermediación financiera no bancaria, y la revisión de la gobernanza del FROB y del Fondo de Garantía de Depósitos, tomando en cuenta los conflictos de interés existentes.
• Aunque finalmente no se incluyeron nuevos condicionantes no financieros, sí que se incluyeron recordatorios de compromisos ya adquiridos por el Gobierno español, como el cumplimiento de los objetivos de déficit, el establecimiento de una institución presupuestaria independiente, o la puesta en marcha de diversas reformas estructurales.

El hecho de que las necesidades finales del programa de reestructuración y recapitalización fueran mucho menores de las previstas devolvió parte de la confianza perdida a la deuda española. Además, España acordó con el MEDE devolver de forma anticipada 17.612 millones de euros entre 2014 y 2018, aprovechando la mejora de las condiciones de financiación y la posibilidad de sustituir la deuda con el MEDE por nueva deuda emitida a tipos de interés. inferiores. Los 23.721 millones de euros restantes se terminarán de devolver en 2027.
} La solicitud de asistencia financiera se produjo en el momento más delicado para la deuda pública española, cuya prima de riesgo llegó a alcanzar los 638 puntos básicos en julio de 2012, situando la rentabilidad del bono a 1 O años por encima del 7,5\%.\footnote{A la delicada situación de la deuda pública española contribuyeron también las sucesivas reducciones de la calificación crediticia de la deuda española por parte de las principales agencias de rating. Moody's. por ejemplo. rebajó la calificación crediticia de la deuda española en nueve escalones (pasando de Aaa a Baa3) en un plazo de menos de dos años, situándola únicamente un escalón por encima del grado especulativo en junio de 2012 (ver gráfico 6).
}


En paralelo a la refonna de la gobemanza económica de la zona euro y a la reestructuración de su sistema financiero. España llevó a cabo diversas reformas estructurales -refonna de las pensiones (20 1 1 ). reforma laboral (2012), Ley de Garantía de Unidad de Mercado (2013), reforma del Sector Eléctrico (2013) , etc.-, lo que contribuyó a que, pese a ser uno de los países cuya recesión fue más profunda, la economía española volviera a la senda de crecimiento desde el año 20 I 3, convirtiéndose poco después en la primera gran economía occidental en términos de crecimiento.

Todas las medidas anteriores empezaron a dejar sentir sus efectos sobre el saldo presupuestario español desde el año 201 5, cuando el déficit público empieza a reducirse hasta situarse por debajo del 3\% en 201 8, tras más de I O años de déficits superiores a esa magnitud. El coeficiente de deuda pública sobre PIB, por su parte, inicia también ese año una senda ligeramente descendente que permitió su reducción desde el 105\% en 2014 hasta el 98\% en 2019. El estallido de la pandemia del covid-J 9, sin embargo, frenó por completo esta buena dinámica.


JI.IV. Situación actual

En 2020, el déficit público español volvió a situarse por encima del 10\% del PIB y la deuda pública superó el 120\% del PIB. Ello se explica, por el lado del denominador, por la caída sin precedentes del PIB, de más del 11\% y, por el lado del numerador, por la actuación de los estabilizadores automáticos y por las medidas discrecionales puestas en marcha para hacer frente al impacto económico y social de la pandemia.\footnote{Según la actualización del programa de estabilidad 2021-2024, el total de la ayuda movilizada mediante medidas discrecionales en el periodo 2020-2021 ascendió a cerca de 2 3 1 .000 millones de euros, el 20,6% del PIB de 2020. Del total, más de 73.000 millones de euros (6,5% del PIB) fueron ayudas directas y 1 5 8.000 millones de euros ( 14, 1 % del PIB), financiación movilizada a través de garantías, avales y moratorias
} Entre estas últimas medidas cabe destacar:
• La flexibilización del sistema de Expedientes de Regulación Temporal de Empleo (ERTEs), por la que en un primer momento se exoneró totalmente a las empresas del pago de los salarios y de las cotizaciones sociales de los trabajadores acogidos a este régimen a cambio de mantener su vínculo laboral con el trabajador en cuestión. Este sistema flexibilizado se mantuvo, con distintas modificaciones, hasta marzo de 2022, y benefició a 6,6 millones de empleados.
• Las medidas de apoyo a la solvencia empresarial, entre las que hay que distinguir aquellas con impacto directo en gasto de las que no lo tuvieron. Entre las primeras, destacan:
o La creación de dos fondos para recapitalizar empresas afectadas por la pandemia: el primero, gestionado por la SEPI y dirigido a apoyar la solvencia de empresas estratégicas, se dotó con 10.000 millones de euros de los que se tenninaron utilizando 3.2 55,8 millones; el segundo, gestionado por COFIDES y dirigido a apoyar la solvencia de empresas medianas, se dotó con 1.000 millones de euros de los que se tenninaron usando 779 millones.
o La creación de un fondo de ayudas directas dirigido a empresas y autónomos de los sectores identificados como más vulnerables que hubieran visto caer sus ingresos al menos un 30\% respecto de 2019. Este fondo se dotó con 7.000 millones de euros, de los que se terminaron utilizando algo más de 6.000 millones, y fue gestionado por las CCAA.
o Entre las medidas· sin impacto inmediato en gasto, cabe destacar la puesta en funcionamiento de distintas líneas de avales por un valor total de 143.000 millones de euros, la mayoría gestionados por el ICO.\footnote{Posteriormente, se desarrolló un marco de reestructuración de los créditos avalados públicamente, dirigido a aliviar la carga linanciera de las empresas avaladas por el ICO mediante la reestructuración de sus deudas bancarias.

Según el informe del ICO de junio de 2022, se han desplegado avales por importe de I 07. 187 millones de euros que han permitido movilizar 140.737 millones de euros en linanciación hacia el tejido productivo en cerca de 1,2 millones de operaciones, de las que más del 98% han sido suscritas por pymes y autónomos.
} Aunque estos avales solo implicarán un gasto en caso de que las empresas deudoras no devuelvan los préstamos avalados, se dotó una provisión de pérdida esperada con impacto en el saldo de déficit del ejercicio fiscal de 2020.

• Las moratorias de créditos y el aplazamiento de algunas obligaciones tributarias, dirigidas a mejorar la posición liquidez de los agentes económicos a través del aplazamiento de los pagos a más corto plazo.
• El incremento del gasto social y, en especial, en sanidad. Aquí ha de incluirse también la aprobación del Ingreso Mínimo Vital, un instrumento de apoyo que busca prevenir el riesgo de pobreza y exclusión social de las personas sin recursos suficientes para la cobertura de sus necesidades más básicas. 

Por lo que se refiere a la situación de la deuda pública, el importante incremento que experimentó en 2020 ha intensificado ef debate sobre su sostenibilidad.\footnote{Desde una perspectiva general, lo primero que cabe señalar es que la sostenibilidad de la deuda pública de un país depende de multitud de factores, entre los que destacan la calidad institucional y democrática y la estabilidad política (efecto positivo), el grado de apertura externa de la economía (efecto positivo). la vida media de la deuda (efecto positivo), el porcentaje de deuda en manos de inversores extranjeros (efecto ambiguo), el envejecimiento demográlico (efecto ambiguo), el grado de concentración de la estructura productiva (efecto negativo), la pertenencia a una unión monetaria (efecto negativo) y, por supuesto. el contexto macroeconómico del país y. en particular, su tasa de crecimiento real, tipo de interés real y saldo primario. Así, nos encontramos con que Japón, por ejemplo, ha sido capaz de mantener una deuda superior al 200% del PIB durante muchos años, mientras que Ucrania tuvo que reestructurar su deuda en 2015. cuando está no superaba el 80% de su PIB.
} Si repasamos las previsiones contenidas en el último informe del observatorio de deuda de la AIReF, estas sitúan el coeficiente de deuda sobre PIB de España ligeramente por debajo del 110\% en 2024. Dicha reducción, no obstante, vendrá sustentada principalmente por el crecimiento del PIB nominal, con una contribución muy notable del deflactor del PIB. Así, una vez que finalice el impulso del crecimiento, las proyecciones de la AIReF dibujan una dinámica desfavorable de la ratio de deuda bajo un escenario a políticas constantes. La AIReF también alerta de que el aumento de los tipos de interés podría agravar esta dinámica, en la medida en que unos tipos de interés más elevados generarán una mayor carga financiera.\footnote{A corto plazo, no obstante, el efecto del aumento de la carga financiera sobre la ratio de deuda sobre PIB será más que compensado por el aumento del PIB nominal provocado por los elevados niveles de inflación.
}

De hecho, el rápido aumento de los tipos de interés desde 2022 ha puesto en entredicho el supuesto que empezaba a aceptarse antes del estallido de la pandemia, según el cual se habían producido una serie de cambios estructurales en las economías avanzadas que provocarían la persistencia de tipos de interés reales reducidos o incluso negativos, que harían sostenibles niveles de deuda pública más elevados que en el pasado.\footnote{Véase, entre otros, Galesi, Nuño, Thomas (2017).

El estallido de la pandemia ya había dejado en evidencia la fragilidad de este supuesto. Así, cuando el 12 de marzo de 2020 (punto I en el gráfico) la presidenta del BCE, Christine Lagarde, alirmó que el BCE no tenía entre su mandato el cierre de los diferenciales entre los tipos de interés de la deuda de los países de la UEM (we are no/ here to e/ose spreads'), la rentabilidad del bono español a I O años aumentó casi 100 puntos básicos en apenas una semana. Poco después, el 1 8 de marzo y el 22 de abril de 2020 (puntos 2 y 3 en el gráfico). el BCE anunciaría el PEPP y la flexibilización de su política de colateral, lo que permitió que la rentabilidad del bono español a I O años se redujera hasta situarse, en julio de 2020. por debajo de su nivel a principios de ese año. Estos eventos pusieron de manifiesto, una vez más, la absoluta dependencia de los tipos de interés de la deuda (y, con ello, de su sostenibilidad) de las decisiones de política monetaria.
} Por ahora, el aumento de los tipos de interés ha venido provocado por un crecimiento aún mayor de los precios, de forma que en las economías avanzadas los tipos de interés reales presentan todavía niveles negativos, pero no puede asumirse que esa situación se mantendrá indefinidamente.\footnote{A cierre de 2022, por ejemplo, la rentabilidad del bono español a J O años era del 3,66%, mientras que la tasa de variación anual del IPC se situó en el 5,7%,
}

En el caso de España, que como veremos en un momento es una de las economías avanzadas con un nivel más alto de deuda pública, parecería necesario acometer un proceso de consolidación fiscal que permita reducir dichos niveles y, con ello, la vulnerabilidad asociada a los mismos. Al mismo tiempo, no obstante, resulta complicado poner en marcha ese prqceso, al ser España uno de los pocos países que, a finales de 2022, no había recuperado todavía los niveles de PIB previos a la pandemia (de hecho, las previsiones del FMI sitúan a España como una de las economías avanzadas que más tiempo tardará en recuperar el nivel de PIB previo a la pandemia). En efecto, como se analiza en detalle en los temas 3A3 8 y 3A39, una consolidación fiscal demasiado agresiva puede tener un efecto indeseado sobre la ratio de deuda sobre PIB, pues el efecto positivo de la consolidación sobre el volumen de deuda (efecto numerador) puede verse superado por el efecto negativo de la misma sobre el PIB (efecto denominador).

11.V. Com paración internacional


Como veíamos en la introducción, fue con las crisis del petróleo de los 70 cuando el déficit público se convierte en un rasgo estructural de las economías avanzadas. En esa década, no obstante, el déficit medio en los países de la OCDE siguió siendo bajo, de entre el 1 % y el 2% del PIB, si bien por encima del prácticamente equilibrio presupuestario medio español.

En la década de los 80, el déficit medio de la OCDE alcanzó el 3%, mientras que en España se situó algo por encima del 4%. Sin embargo, fue a principios de los 90 cuando, en un contexto de crisis mundial, tiene lugar un aumento generalizado del déficit público, que fue de nuevo especialmente acusado en España. Así, el déficit medio de la OCDE en 1993 fue del 4,9%, mientras que en España se llegó al 7,3%.

A partir de entonces se inicia en los países de la OCDE un intenso proceso de consolidación fiscal, que permitió a países como EE. UU. mantener superávits en 1999 y 2000. Japón, por su parte, atravesó su particular década perdida, y los superávits de principios de los 90 se transformaron en déficits que se han mantenido hasta la actualidad, situando su deuda pública por encima del 250\% del PIB. En Europa, los 1 1 países que integraron el primer grupo de la UEM pasaron de presentar un déficit medio del 5, 7% en 1993 a uno por debajo del 2% a finales de la década de los 90. Con la llegada de la moneda única, la evolución del déficit fue muy dispar en los distintos países de la zona euro. España o Irlanda continuaron con el proceso de consolidación fiscal, mientras que países como Francia, Alemania, Italia o Portugal mostraron grandes dificultades para cumplir con el límite del 3% de déficit previsto en el PEC.

Tras el estallido de la crisis financiera internacional, prácticamente todos los países de la OCDE incurrieron en déficit o vieron agravado el que ya presentaban, y el déficit medio pasó de poco más del 1 % en 2007 a más del 8% en 2009, con grandes diferencias entre unos países y otros.

Centrándonos en eventos más recientes, en la crisis del covid-19 encontramos diferencias importantes en la evolución de los saldos presupuestarios de los distintos grupos de países:

• En el caso de las economías avanzadas, el déficit promedio pasó de menos del 3% del PIB en 2019 a casi el 1 1 % en 2020, habiéndose reducido rápidamente para situarse, dos años después, en niveles próximos, aunque todavía algo superiores, a los de 2019.

• En el caso de las economías emergentes (excluida China), el déficit promedio pasó de casi el 4% en 2019 a casi el 8% en 2020, habiendo vuelto a niveles cercanos al 4% en 2022. El caso de China es particular, pues su déficit en 2019 era ya cercano al 6%, creció hasta casi el 10% en 2020 y, aunque volvió al 6% en 2021, en 2022 volvió a alcanzar el 9% del PIB. Esto se explica por la particular respuesta sanitaria de China a la pandemia, que ha hecho necesario que el estímulo fiscal para evitar un frenazo excesivo de la actividad económica se extienda durante más tiempo.

• Finalmente, las economías de bajo ingreso tenían un déficit cercano al 4% del PIB en 2019, que creció hasta casi el 5% del PIB en 2020 y que se ha mantenido en esos niveles desde entonces.

Se observan, además, grandes diferencias dentro del grupo de las economías avanzadas. Si analizamos los datos de España, observamos que en 2019 tenía el 4° déficit público más elevado de entre las economías avanzadas, superado únicamente por EE. UU., Australia e Israel, e igualado por Francia. En 2020 fue la 8ª economía avanzada con mayor nivel de déficit público. y en 2022 seguía ocupando una posición parecida.

A la hora de hacer comparaciones internacionales, no obstante, resulta más interesante recurrir a los datos de deuda pública que a los del saldo presupuestario, dada la mayor variabilidad de este último indicador. De acuerdo con los datos del FMJ, el endeudamiento público mundial aumentó en 1 5 puntos porcentuales en 2020, el mayor aumento en un año desde la 11 Guerra Mundial. Como resultado, el coeficiente de deuda pública alcanzó un nivel sin precedentes, del 99% del PIB. Eso supone el 40% de la deuda mundial total, la mayor proporción desde mediados de los años 70.
Pero tanto los niveles absolutos, como los cambios experimentados recientemente en los niveles de deuda pública, varían enormemente entre países. De hecho. más del 90% del aumento de la deuda en 2020 tuvo lugar en las economías avanzadas y en China, que se pudieron beneficiar en mayor medida de la actuación de sus bancos centrales y de la reducción de los tipos de interés gracias al mayor desarrollo de sus mercados financieros. Así, pueden identificarse dos fenómenosa nivel general:
• En las economías avanzadas, la deuda pública aumentó en 19 puntos porcentuales de PIB, un aumento similar al ocurrido en los dos años más duros de la crisis financiera internacional (2008 y 2009). Como resultado, la deuda pública de estos países se situó en 2020 en el 124% del PIB, una cifra que se ve muy impactada por el dato de EE. UU., cuya deuda pública llegó a suponer casi el 1 3 5% del PJB.
• En las economías emergentes (excluida China) y de bajo ingreso, en cambio, el aumento nominal de la deuda fue muy reducido, dadas las restricciones financieras existentes. Concretamente, el aumento de la deuda pública fue de alrededor de 7 puntos porcentuales de PIB, y se explica en su mayor parte por la reducción del PIB, y no tanto por el aumento de la deuda. A pesar de ello, en 2020 la deuda pública alcanzó un máximo histórico en las economías emergentes, y en las economías de bajo ingreso, niveles no observados desde comienzos de siglo, cuando muchos de esos países se beneficiaron de iniciativas de alivio de la deuda.

Dentro de las economías avanzadas, entre las que se encuentra España, el país que experimentó un mayor aumento de su deuda pública en 2020 fue Canadá, con un crecimiento de más de 30 puntos porcentuales de PIB. Taiwán, por el contrario, vio cómo su deuda pública incluso se reducía en 2020, aunque en apenas O, 1 puntos porcentuales de PIB. España fue la 7ª de las economías consideradas por el FMI como avanzadas con mayor crecimiento de su deuda pública en 2020, con un aumento de más de 21 puntos porcentuales de PIB.

No obstante, resulta también interesante analizar la evolución d�� la deuda pública en los años posteriores a la pandemia pues, en algunos casos, el aumento del coeficiente de deuda pública en 2020 se debió en buena medida a la reducción del PIB. En esos casos, el crecimiento del P I B en los años posteriores a 2020 ha sido en general más rápido, y también ha sido más rápida la reducción del coeficiente de deuda. Así, si comparamos el coeficiente de deuda de 2019 con el de 2022, encontramos que España es la 5ª economía avanzada que ha experimentado un mayor crecimiento de su deuda pública, superada únicamente por San Marino, Japón. Nueva Zelanda y Malta. Concretamente, el crecimiento de nuestro coeficiente de deuda pública entre esos años fue  de 1 5 puntos porcentuales.

En términos absolutos, nuestro coeficiente de deuda pública es el 7º más alto de las economías avanzadas, con un 113,1\% del PIB en 2022. Se trata, además, del 4° nivel más alto de la zona euro, situándose 20 puntos porcentuales por encima del promedio.

El problema estructural de las cuentas públicas españolas, por tanto, no lo es solo en términos absolutos, como veíamos en el apartado anterior, sino que también queda de manifiesto cuando comparamos los datos con los de los países de nuestro entorno.






\clearpage
\section{}
\puntosclave{
    La deuda del Estado supone casi el 90\% de la deuda total de las AAPP en España. y su gestión es competencia de la Dirección General del Tesoro y Política Financiera.
    
    Desde la crisis financiera internacional se ha producido un gran aumento de las necesidades de financiación del Tesoro. Concretamente, las necesidades de financiación brutas del Tesoro han pasado de 50.000 millones de euros en 2007 a alrededor de 250.000 millones de euros anuales en la actualidad. una cifra equivalente a algo más del 20\% del PIB.
    
    El crecimiento de las necesidades de financiación ha obligado al Tesoro a estandarizar los procedimientos de emisión y a recurrir a instrumentos que permitan emitir mayores volúmenes de deuda de forma continuada.
    
    El aumento de las necesidades de financiación del Tesoro ha coincidido en el tiempo con los estímulos monetarios extraordinarios puestos en marcha por el BCE, lo que ha permitido reducir el coste de emisión ( en 2021, el Tesoro llegó a tener, por primera vez en su historia. un coste medio de emisión de su deuda negativo) al tiempo que se aumentaba la vida media de la deuda en circulación (que en 2021 alcanzó por primera vez los 8 años) para reducir el riesgo de refinanciación.

    El brusco cambio de tendencia de la política monetaria en 2022 para hacer frente a los elevados niveles de inflación, obligará al Tesoro a replantear la estrategia seguida estos últimos años. Cualquier cambio en la estrategia de financiación del Tesoro se realizará, en cualquier caso, de forma muy gradual.
    
    Las necesidades de financiación brutas del Tesoro se mantendrán en niveles muy elevados los próximos años, ante la necesidad de refinanciar unos vencimit!ntos crecientes y el final de los préstamos del fondo SURE y de los préstamos y transferencias que está recibiendo España en el marco del programa Next Generation EU (NGEU).
}


Una vez hemos visto que el déficit público se ha convertido en un rasgo estructural de la economía española, adquiere una gran importancia el estudio de su forma de financiación. Su relevancia deriva (i) de los efectos económicos que genera y (ii) de que la aplicación de una política de financiación óptima ayuda a reducir las necesidades de financiación. Esto último es especialmente importante en la actualidad, dadas las grandes necesidades de financiación de las AAPP españolas.

En España, la gestión del endeudamiento del Estado (cuya deuda supone casi el 90\% de la deuda total de las AAPP) es competencia de la Dirección General del Tesoro y Política Financiera, una de las direcciones generales del Ministerio de Asuntos Económicos y Transformación Digital.\footnote{En otros países (p.ej. Francia. Alemania o Reino Unido), la gestión del endeudamiento del Estado es realizada por agencias independientes que gozan. al menos en teoría. de una mayor autonomía en su actuación. En España, aunque el Tesoro depende formalmente del Ministerio de Asuntos Económicos y Transformación Digital, opera con gran autonomía en la gestión del endeudamiento del Estado, que está guiada en todo momento por criterios financieros.
} En la última sección del terna, vamos a estudiar la política de financiación del Tesoro en la actualidad, para lo que veremos la evolución reciente de sus necesidades de financiación y los instrumentos y procedimientos de emisión de los que se ha ido dotando a lo largo del tiempo. Terminaremos repasando brevemente la estrategia de financiación del Tesoro en la actualidad y tratando de anticipar algunos de los cambios que pueden producirse en el futuro próximo.


111.1. Evoluc ión reciente de las necesidades de financiación del Tesoro

En los últimos 15 años se ha producido un gran aumento de las necesidades de financiación del Tesoro. Así, en el año 2007 el Tesoro tuvo una necesidad de financiación bruta de 50.000 millones  de euros, con una emisión neta (es decir, aminorada del volumen de deuda amortizada durante el ejercicio) negativa de casi 10.000 millones de euros. Desde entonces, la emisión neta ha sido positiva de forma ininterrumpida, llegando a superar los 1 00.000 millones de euros en 2009 y en 2020.

Las necesidades de financiación brutas, por su parte, se han disparado como consecuencia del aumento tanto de la emisión neta como de los vencimientos anuales (consecuencia, a su vez, del crecimiento continuado de la deuda pública en circulación), hasta situarse estos últimos años alrededor de los 250.000 millones de euros anuales, una cifra equivalente a algo más del 20% del PIB. Y ello a pesar de que el diseño de la respuesta europea a la pandemia ha permitido contener en cierta medida su impacto sobre las necesidades de emisión de las economías más afectadas y, en particular, de España, cuyas emisiones habrían sido mayores de no ser por los préstamos del fondo SURE ( que ascendieron a 2 1 .326 millones de euros entre 2020 y 202 I ) y por las transferencias recibidas en el marco del programa NGEU (31.036 millones entre 2021 y 2022), que se verán complementadas en el periodo 2023-2026 por préstamos de carácter reembolsable.

En enero de 2020, antes del estallido de la pandemia del covid-19, el Tesoro esperaba realizar una emisión bruta anual inferior a los 200.000 millones de euros por primera vez en casi 1 O años, pero la reducción de los ingresos públicos y el aumento del gasto hizo necesario emitir 277.000 millones de euros, la cifra más alta de la historia del Tesoro.

El aumento de las necesidades de financiación del Tesoro ha coincidido en el tiempo con una caída sin precedentes de sus costes de emisión, favorecida tanto por el cambio del patrón de crecimiento de la economía española hacia uno más equilibrado y sostenible en el tiempo, que ha provocado la mejora de la calificación crediticia de la deuda española y la reducción de su primade riesgo; como por los estímulos monetarios extraordinarios puestos en marcha por el BCE.\footnote{Entre 20 1 5 y 2022, las compras netas de deuda pública española por parte del BCE mediante el APP y el PEPP superaron los 500.000 millones de euros, un importe equivalente a más de la tercera parte de la deuda en circulación de las AAPP a finales de 2022. Si restamos 'ta deuda en manos del BCE, la deuda de las AAPP a finales de 2022 era equivalente a la que había a principios de 20 1 3 (ver gráfico 8).
}

Como veremos a continuación, el crecimiento de las necesidades de financiación ha traído consigo algunos cambios tanto en los instrumentos a utilizar como en los procedimientos de emisión:

• En lo que se refiere a los instrumentos, se ha primado la emisión de letras del Tesoro y de bonos y obligaciones del Estado, en detrimento de otros instrumentos como las emisiones en divisa extranjera. Además, se han ampliado el número de referencias benchmark y los plazos de vencimiento (llegando a emitirse obligaciones der Estado con vencimiento a 50 años).\footnote{Las referencias benchmark son aquellas referencias utilizadas para construir la curva de rentabilidades del Tesoro, además de ser las referencias emitidas más frecuentemente en las subastas. Concretamente. son las últimas referencias emitidas por el Tesoro a los plazos de 3, 5, 1 O, 15. 20, 30 y 50 años. Cuando se emite una nueva referencia a un determinado plazo, esta se convierte en el nuevo benchmark a ese plazo, y la anterior referencia benchmark pasa a recibir el nombre de o/f-1he-run.
} También se ha creado el programa de bonos y obligaciones indexados a la inflación europea y el programa de bonos verdes, tratando de ampliar y de diversificar la base inversora.

• En lo que se refiere a los procedimientos de emisión, las subastas y las sindicaciones han pasado a ser la modalidad prácticamente exclusiva de emisión, teniendo otros procedimientos, como las colocaciones privadas o los canjes, una importancia residual. Lo que se ha buscado, en general, es estandarizar los procedimientos de emisión y recurrir a instrumentos que permitan emitir mayores volúmenes de deuda de forma continuada. Veamos, en cualquier caso, todos los instrumentos y procedimientos de emisión a disposición del Tesoro.


111.11. Instrumentos


Corno veíamos en la introducción, hasta 1982 la financiación del déficit público en España recayó, en buena medida, en el BdE. Ese año, más del 80% del déficit del Estado se financió mediante el recurso al crédito del BdE. Sin embargo, el fuerte crecimiento del déficit público desde finales de los años 70 comenzó a complicar la instrumentación de una política monetaria enfocada en la reducción de la inflación, que en 198 1 superaba el 14%. Para hacer frente a ese problema, en 1982 se comienza a cambiar la política de financiación del Estado para orientarla hacia los mecanismos de mercado. Ello obliga al Tesoro Público a desarrollar una política de deuda pública, entendida como el manejo de una serie de instrumentos que le permitieran financiar el déficit público sin recurrir a la monetización y al menor coste posible. En aquel momento, el Tesoro disponía únicamente de dos instrumentos, la deuda desgravable del Estado y los pagarés del Tesoro:
• La deuda desgravable del Estado era un instrumento a medio plazo dirigido a los sujetos pasivos del IRPF, cuya demanda estaba motivada, sobre todo, por razones fiscales.
• Los pagarés del Tesoro, por su parte, eran activos a corto plazo (6 meses en un primer momento, y 18 meses más tarde) que se habían creado como instrumento de regulación de la liquidez por el BdE, pero al no existir restricciones a su adquisición por parte de los inversores particulares, desempeñaban también el papel de instrumento definanciación del Estado.

Se hacía necesario, por tanto, aumentar la gama de instrumentos disponibles con activos a medio y largo plazo que se adaptasen a las necesidades de los inversores, pero evitando la segmentación del mercado que podría producir la introducción de un número de instrumentos demasiado elevado, pues esa segmentación habría provocado una menor liquidez de la deuda y, en consecuencia, un menor atractivo de la misma para los inversores.\footnote{La creación de la gama de instrumentos que vamos a ver era una condición necesaria, pero no suficiente, para la financiación ortodoxa del déficit público. Así, había que contar también con demandantes de esos valores, algo dificil de conseguir en un primer momento por el escaso desarrollo de los mercados de valores en España y la práctica inexistencia de inversores institucionales. Por este motivo. en 1984 se estableció el coeficiente de inversión obligatoria en pagarés del Tesoro, que obligaba a los bancos y a las cajas de ahorros a mantener en estos activos un 12\% de sus pasivos computables. Como resultado, durante los primeros años de esta etapa el sistema bancario sustituyó al BdE como el principal financiador del déficit público. El coeficiente de inversión obligatoria en pagarés del Tesoro se fue reduciendo desde 1989, y se eliminó completamente en 1992.
}

Así, a finales de 1982 se crean los bonos del Estado (con vencimientos a 2 a 5 años) y, en 1983, las obligaciones del Estado, con plazos de vencimiento más largos. Estos activos financieros, que  como veremos constituyen el principal instrumento de financiación del Tesoro en la actualidad, no ofrecían ya ningún tipo de ventaja fiscal.

En 1987 se crean las letras del Tesoro como activo a corto plazo sustitutivo de los pagarés del Tesoro y, desde entonces, la gama de instrumentos del Tesoro se ha ido ampliando para atender las preferencias y necesidades de los inversores en deuda pública. Veamos en algo más de detalle la gama actual de instrumentos del Tesoro.

i. Letras del Tesoro

Las letras del Tesoro son valores de renta fija a corto plazo. Cuando se introdujeron en 1987 tenían un vencimiento a 12 meses, pero al restringirse legalmente en 1990 el acceso del Tesoro al crédito del BdE (en lo que sería un paso previo a su prohibición absoluta en 1994, de acuerdo con las exigencias del Tratado de Maastricht) se hizo necesario introducir letras del Tesoro con vencimientos a 3 y 6 meses como herramienta para la gestión de los desfases transitorios de tesorería a lo largo del año. Aunque en algún momento se llegaron a emitir letras con vencimiento a 18 meses, en la actualidad se emiten únicamente letras con vencimiento a 3. 6, 9 y 12 meses. 

Las letras del Tesoro no ofrecen un interés periódico en forma de cupón, por lo que su rentabilidad es implícita y se calcula a partir de la diferencia entre el precio de compra y su valor nominal. Tradicionalmente, la emisión de las letras se hacía al descuento (es decir, a un precio inferior a su valor nominal), pero en los últimos años ha sido habitual que el Tesoro emitiera letras a premio (i.e. con rentabilidad negativa). Las letras se emiten mediante subasta, en las que las peticiones han de ser múltiplos de 1.000 euros (siendo, por tanto, el importe mínimo de cada petición de 1.000 euros).

La liquidez de las letras del Tesoro es uno de sus principales atractivos, y para conseguirlo el Tesoro minimiza el número de referencias en circulación en cada momento. Para ello, cada mes se emite una nueva letra a 12 meses, que se reabre mediante subasta cuando su vida residual es de 9, 6 y 3 meses.\footnote{Formalmente, por tanto, no hay letras del Tesoro a 3. 6 y 9 meses. únicamente hay letras a 12 meses de las que se emite más volumen cuando su vida residual es de 9. 6 y 3 meses.
} La emisión inicial de cada letra a 12 meses se hace por un mayor importe, garantizando así su liquidez desde el primer momento. Gracias a este sistema, existen en todo momento únicamente 12 referencias de letras del Tesoro en circulación.

Las letras del Tesoro demostraron, desde su creación, una gran capacidad de financiación, gracias a los elevados tipos de interés a los que se emitieron en un primer momento (del 15,5\% anual). Así, en 1 990 las letras llegaron a suponer el 62\% de la deuda en circulación, dejando al Tesoro expuesto a un elevado riesgo de refinanciación.\footnote{El riesgo de refinanciación de las letras del Tesoro deriva de la necesidad de emitir nueva deuda en un plazo muy corto (como máximo. 12 meses) para poder hacer frente a su amortización en su fecha de vencimiento. Esa nueva deuda tendrá que emitirse en un momento para el que no se conoce cuáles serán las condiciones de emisión. de ahí que se considere un riesgo. Este riesgo se produce únicamente en caso de tener que refinanciar (de ahí su nombre) las letras. algo que ocurrirá siempre que la emisión neta del Tesoro no sea negativa.
} Para hacer frente a esta situación, se trató depromover la participación en el mercado de inversores institucionales, que por su propia naturaleza tienden a invertir en instrumentos con vencimientos más largos. Así, en 1990 se establecieron incentivos fiscales a la inversión en deuda pública de inversores no residentes, lo que permitió que el porcentaje de bonos y obligaciones en su poder pasara del 10\% a más del 30\% en apenas un año.\footnote{Además, también se trató de incentivar la demanda minorista en valores del Estado. para lo que. entre otras cosas, se empezaron a poner en marcha campañas de publicidad y se crearon los FONDTESORO, una modalidad especial de fondos de inversión que invierten gran parte de su patrimonio en valores del Tesoro. Los FONDTESORO perrniten a los inversores minoristas invertir en valores del Tesoro sin necesidad de mantener directamente dichos valores, cediendo por tanto la gestión de la cartera a una sociedad gestora a cambio del pago de unas comisiones muy competitivas. al encontrarse limitadas por el convenio de colaboración firrnado entre las sociedades gestoras y el Tesoro.
} Así, el porcentaje de las letras sobre el total de la deuda en circulación cayó por debajo del 30\% en 1996, y ha seguido reduciéndose de forma continuada desde entonces (con la excepción del periodo entre 2009 y 2013, cuando creció ligeramente), hasta representar menos del 6\% en la actualidad.\footnote{En 2009, el dificil acceso a los mercados obligó al Tesoro a modificar su estrategia de financiación, aumentando sus emisiones a corto plazo, que tienen un menor coste de emisión y que. especialmente en condiciones de financiación complicadas, son más fáciles de colocar (dada su menor volatilidad). Así, ese año el Tesoro cubrió casi el 30% de sus necesidades de financiación netas mediante el aumento del saldo en circulación de letras del Tesoro. Corno resultado, las letras del Tesoro llegaron a representar ese año el 18% de la deuda del Estado en circulación.
}


ii. Bonos y obligaciones del Estado (linkers, bonos verdes)

Los bonos y obligaciones del Estado son valores de renta fija emitidos por el Tesoro a un plazo a superior a dos años. Los bonos y las obligaciones del Estado son iguales en todas sus características menos en el plazo, que en el caso de los bonos oscila entre 2 y 5 años, mientras que en el de las obligaciones es superior a 5 años.\footnote{La obligación a 10 años es la que suele utilizarse corno referencia al comparar la rentabilidad de la deuda de distintos países. De hecho, cuando se hace referencia a la prima de riesgo, esta suele estar calculada corno el diferencial entre la rentabilidad de la obligación española a 1 0 años y la rentabilidad de la obligación alemana a 1 0 años.
} Concretamente, el Tesoro emite en la actualidad obligaciones con vencimiento hasta 50 años. A diferencia de las letras, los bonos y las obligaciones del Estado son títulos que ofrecen, con carácter general, un interés periódico en forma de cupón.\footnote{Habitualmente, el importe del cupón queda fijado en el momento de la emisión. En 2009, no obstante, el Tesoro decidió emitir un bono a tipo de interés variable, cuyo cupón, pagadero trimestralmente, se determinaba corno el Euribor a 1 2 meses menos I O puntos básicos. Esa decisión se tornó con el objetivo de diversificar la base inversora. ante el repentino aumento de las necesidades de financiación del Tesoro. Actualmente el Tesoro no realiza emisiones a tipo de interés variable.
} En los últimos años, no obstante, el Tesoro ha llegado a emitir bonos y obligaciones del Estado con cupón cero y al descuento, ofreciendo por tanto una rentabilidad negativa.

Habitualmente, estos valores se emiten mediante subasta, aunque la emisión del primer tramo de las obligaciones se realiza mediante sindicación, como veremos más adelante. De manera puntual, también se han emitido bonos y obligaciones del Estado mediante colocación privada a un inversor o grupo de inversores. Una vez se ha emitido un nuevo bono u obligación del Estado, se emiten sucesivos tramos del mismo instrumento hasta que su volumen en circulación alcanza un nivel que permita garantizar su liquidez en el mercado secundario.\footnote{Este sistema, conocido corno sistema de agregación de emisiones. es el mismo que se aplica a las letras del Tesoro, y empezó a utilizarse desde finares de los 80 para aumentar la liquidez de los valores del Tesoro y hacerlos así más atractivos para los inversores.
} Dicho nivel se sitúa, actualmente, en torno a 25.000 millones de euros. En cuanto a la adquisición de . estos instrumentos, igual que ocurría con las letras, debe hacerse en cantidades múltiplos de 1.000 euros.

La emisión inicial de los bonos y obligaciones del Estado se hace a valores próximos a la par, pero sus reaperturas podrán emitirse a la par, bajo la par o sobre la par, en función de la evolución de los tipos de interés\footnote{Si la TIR es mayor que el cupón se emitirá bajo la par. y si la TIR es menor que el cupón se emitirá sobre la par. Cuando los tipos presentan una tendencia descendente ( corno ocurrió, en general, entre 2 0 1 5 y 202 1 ) es habitual la emisión de bonos y obligaciones sobre la par (ya que se emiten nuevos tramos de bonos y obligaciones con cupones muy por encima de la TIR). Ello tiene gran relevancia desde el punto de vista de la contabilidad nacional, pues corno veíamos al principio del terna, el PDE valora la deuda utilizando su valor nominal o facial. Esto significa que, si se emite un bono u obligación sobre la par, a efectos del PDE sólo computa corno deuda el valor nominal (que será inferior a la cantidad ingresada por el Tesoro). La diferencia entre la cantidad ingresada por el Tesoro y el valor nominal de la deuda se conoce corno prima de emisión, y constituye una forma de financiar el déficit de caja sin aumentar el valor de la deuda PDE. Como aspecto negativo, la emisión de bonos y obligaciones con cupones altos obliga a pagar esos cupones elevados hasta el vencimiento del instrumento (lo cual afecta al déficit de los años siguientes). Actualmente. y corno consecuencia de la rápida subida de los tipos de interés ocurrida desde 2022, nos encontrarnos con la situación opuesta, es decir, con la emisión de muchos bonos y obligaciones bajo la par. Cuando esto ocurre, el Tesoro ingresa una cantidad inferior al valor nominal de la deuda, haciendo necesaria la emisión de un mayor volumen de deuda nominal para la financiación del déficit de caja.
}

En la última década. además, el Tesoro ha lanzado dos programas de bonos y obligaciones que merecen una mención especial: el programa de bonos y obligaciones del Estado indexados a la inflación europea y el programa de bonos verdes soberanos del Reino de España.

• El primero, lanzado en 2014, permite al Tesoro emitir bonos y obligaciones del Estado cuya rentabilidad está indexada al Índice de Precios al Consumo armonizado ex-tabaco para la zona euro publicado por Eurostat. La rentabilidad de esos bonos. también conocidos como linkers, depende por tanto no sólo de su cupón y rentabilidad implícita, sino también de la evolución del IPCA ex-tabaco entre la fecha de emisión y la fecha de vencimiento.\footnote{S i la inflación entre l a fecha de emisión y l a fecha de vencimiento fuera negativa, s e devolvería u n importe equivalente al valor nominal.
} El desarrollo de este programa ha permitido al Tesoro diversificar su base inversora, accediendo a inversores nacionales e internacionales que tienen un interés particular en activos que eviten la pérdida de valor provocada por la inflación. Además, en la medida en que los ingresos públicos están fuertemente ligados a la inflación. el Tesoro se encuentra en una posición idónea para emitir este tipo de instrumentos.\footnote{Dicha posición. no obstante. se ha visto perjudicada en los últimos años, en los que se ha decidido volver a indexar a la inflación una parte muy importante del gasto público. como es el gasto en pensiones.
} Actualmente. el saldo en circulación de los bonos y obligaciones del Estado ligados a la inflación europea equivale a algo más del 6\% de la deuda del Estado en circulación.

• El programa de bonos verdes soberanos del Reino de España, por su parte. se lanzó en el año 2021 con un doble objetivo: diversificar la base inversora del Tesoro y reforzar el compromiso de España con la transición ecológica. Para ello, el Tesoro se compromete a emplear los fondos obtenidos de la emisión de estos instrumentos a proyectos destinados a la mitigación y adaptación al cambio climático, el uso sostenible y la protección de los recursos hídricos y marítimos, la transición hacia una economía circular, la prevención y control de l a contaminación, y la protección y recuperación de la biodiversidad y los ecosistemas. La necesidad de emplear los fondos obtenidos de la emisión de bonos verdes a esos proyectos limita de manera muy importante la cantidad de este tipo de instrumentos que puede ser emitida por el Tesoro. Finalmente, cabe señalar que la elevada demanda de bonos verdes por parte de los inversores (cada vez más concienciados de la importancia de estas cuestiones), unida a su escasez relativa, ha permitido al Tesoro obtener un ahorro financiero respecto a la emisión de un bono convencional a plazo equivalente.\footnote{Este ahorro financiero obtenido en algunas emisiones verdes frente a las convencionales se conoce popularmente como greenium
} En la emisión inaugural del programa, por ejemplo, dicho ahorro ascendió a 2 puntos básicos.

El peso de los bonos y obligaciones del Estado en el total de la deuda en circulación ha experimentado una evolución contraria al de las letras, pasando del 30% en 1 990 a más del 90% en la actualidad. Este hecho, unido a la emisión de obligaciones con vencimientos cada vez más largos (de hasta 50 años en la última década), ha permitido aumentar considerablemente la vida media de la deuda del Estado en circulación, que ha pasado de menos de 2 años en 1 990 a superar los 8 en 202 1 .

El alargamiento de la vida media de la deuda contribuye a aumentar su sostenibilidad, al reducir el riesgo de refinanciación y la emisión bruta a corto plazo. De hecho, este fue uno de los argumentos empleados por las agencias de calificación crediticia para justificar las subidas de rating realizadas entre 2014 y 2019.\footnote{Otros fueron el mayor crecimiento, la consecución de un superávit continuado por cuenta corriente o el proceso de consolidación fiscal iniciado en 201 5.
}

\begin{specialblock}{\textbf{Strips}}
    Todos los bonos y obligaciones del Estado emitidos por el Tesoro en la actualidad son ··segregables'. Esta característica permite a los creadores de mercado separar cada bono u obligación en en valores (denominados strips).\footnote{Los creadores de mercado también pueden realizar la operación inversa a la segregación. es decir. la reconstitución del valor original a partir de los bonos cupón cero procedentes de su segregación.
    } Concretamente. podrán crear un strip por cada uno de los pagos que dé derecho a recibir el valor en cuestión, siendo cada strip un bono cupón cero cuya fecha de vencimiento y valor de reembolso coincide con la del cupón/principal del valor segregado.\footnote{Por ejemplo. de un bono a 5 años se pueden obtener 6 strips, uno por cada uno de los cupones que da derecho a recibir y otro por el principal
    } Los strips presentan ciertas ventajas que los hacen especialmente atractivos para determinados inversores, de forma que la posibilidad de segregación permite al Tesoro cubrir esa demanda sin necesidad de aumentar la gama de instrumentos disponibles, evitando la excesiva segmentación del mercado y la consecuente menor liquidez de la deuda.\footnote{Una de esas ventajas es fiscal, en la medida en que el beneficio de los strips no está sujeto a retención en el Impuesto sobre Sociedades. retención que si se produce en el caso de los cupones de los bonos y obligaciones no segregados
    } El Tesoro comenzó a emitir bonos y obligaciones segregables en 1997. cuando la proximidad de la integración en la UEM hacía necesario empezar a ofrecer instrumentos con un atractivo que permitiera a la deuda española competir con la de nuestros socios europeos.
\end{specialblock}


iii. Otros instrumentos (deuda en divisas, préstamos)

Como señalábamos al comienzo de esta sección, el crecimiento de las necesidades de financiación del Tesoro ha llevado a primar la emisión de letras del Tesoro y de bonos y obligaciones del  Estado sobre otros instrumentos, que en la actualidad suponen menos del 4% de la deuda del Estado en circulación, frente al más del 15% que llegaron a suponer a finales de los 90. Entre estos otros instrumentos cabe destacar la deuda en divisas y los préstamos otorgados por el MEDE y por el fondo SURE de la Unión Europea:

• En cuanto a la deuda en divisas, esta tuvo un papel destacado en la atracción de inversores internacionales en la década de los 80 y 90, llegando a suponer a finales de los 90 más  del 10% de la deuda total en circulación. Sin embargo, desde la adopción del euro la emisión en divisas alternativas se ha hecho cada vez menos necesaria para acceder a los inversores internacionales. Así, la última emisión de deuda en divisas tuvo I ugar en 2014, quedando emisiones residuales en yenes, dólares y libras esterlinas equivalentes a menos del O, 1 % de la deuda total en circulación. Habitualmente, el Tesoro recurría a swaps de divisas para cubrirse frente al riesgo de tipo de cambio de estas operaciones. 

• Préstamos: en los últimos años, España ha recibido varios préstamos que han permitido al Tesoro reducir sus emisiones de otros instrumentos. Los más importantes han sido el acordado con el Mecanismo Europeo de Estabilidad (MEDE) en 2012 para la recapitalización del sistema financiero, que ascendió a 41.333 millones de euros y, más recientemente, los recibidos a través del fondo SURE de la U E en 2020 y 2021, que han ascendido a 21.326 millones de euros. A ellos hay que añadir varios préstamos de menor importe recibidos del Banco Europeo de Inversiones; y algunos préstamos Schuldschein concertados con inversores privados para adaptarlos a sus necesidades específicas. Estos últimos préstamos, sometidos a la legislación alemana, fueron utilizados por última vez por el Tesoro en 2014. Todos estos préstamos se han recibido a tipos de interés inferiores a aquellos a los que se podía financiar el Tesoro en cada momento, permitiendo así obtener un ahorro financiero.

111.111. Procedimien tos de emisión

Además de crear una gama adecuada de instrumentos y una base inversora lo más diversificada posible, el paso a la financiación que hoy llamaríamos ortodoxa (es decir, sin recurrir a la monetización) del déficit público que se llevó a cabo en España desde principios de los 80 hacía necesario establecer una serie de procedimientos de emisión de los valores que asegurasen tanto el abastecimiento regular del mercado como la consecución de la financiación al menor coste posible para el Tesoro. Así, en 1983 se adopta como método de emisión de las obligaciones del Estado la subasta competitiva. Este método de emisión se extendería más tarde a los bonos del Estado y a las letras del Tesoro, y sigue siendo actualmente el principal procedimiento de emisión del Tesoro. Veamos en más detalle las peculiaridades del sistema de subasta adoptado por el Tesoro, así como los otros métodos que utiliza para la emisión de los valores.

i. Subastas

Se trata del método de emisión más utilizado por el Tesoro, permitiéndole obtener más del 80% de la emisión bruta anual. El BdE interviene, pero únicamente proveyendo la infraestructura tecnológica.

El sistema de subasta utilizado por el Tesoro es una combinación de los dos sistemas de subasta utilizados por otros Tesoros, que son la subasta holandesa o de precio único y la súbasta convencional o de precios múltiples. En la subasta holandesa se adjudican todas las peticiones aceptadas al precio mínimo aceptado o precio marginal, mientras que en la subasta convencional cada petición aceptada se adjudica al precio solicitado.

En el sistema de subasta utilizado por el Tesoro se determina el precio mínimo aceptado y, una vez determinado, se calcula el precio medio ponderado de todas las peticiones aceptadas.\footnote{Ese precio mínimo lo determina el Tesoro en la propia subasta en función de las peticiones recibidas, de la cantidad total a emitir en la subasta (para lo que se publica un rango de emisión objetivo orientativo con carácter previo), del número de instrumentos a subastar (que serán siempre dos en el caso de las subastas de letras del Tesoro, pero que pueden ser más -habitualmente, tres o cuatro- en el caso de las subastas de bonos y obligaciones) y del saldo en circulación de cada instrumento subastado.
} Las peticiones realizadas a un precio superior al precio medio ponderado se adjudican a ese precio, y las demás se adjudican al precio solicitado.

La subasta española. que es como se conoce el sistema de subasta utilizado por el Tesoro, incentiva a los inversores a introducir en la subasta pujas a precios elevados para garantizarse obtener los valores (al igual que ocurre en la subasta holandesa), pero a diferencia de la subasta holandesa, en la que todas las peticiones aceptadas se adjudican al precio mínimo aceptado, en la subasta española las peticiones a precios superiores al medio ponderado se adjudican a ese precio (que es un precio superior al precio mínimo aceptado y, por ende, un.precio que permite al Tesoro colocar los valores con un menor coste que el que resultaría de utilizar el sistema de subasta holandesa). Evidentemente, las pujas introducidas a precios elevados aumentarán el precio medio ponderado. pero lo harán. en general, de forma moderada. y por tanto con poco impacto para los inversores. También se ofrece a los inversores la posibilidad de introducir peticiones no competitivas (i.e. peticiones en las que no se establece un precio. únicamente la cantidad solicitada), que se adjudican al precio medio ponderado resultante de la subasta (salvo si ese precio implica una rentabilidad negativa, en cuyo caso, con carácter general. las peticiones no competitivas no se adjudican). Este tipo de petición es la utilizada habitualmente por los inversores minoristas y, salvo en determinados casos. el importe de cada petición no competitiva no puede superar los 5 millones de euros.

\begin{specialblock}{\textbf{Calendario de emisiones}}
    Para asegurar la estabilidad de las emisiones a lo largo del año y garantizar un acceso continuado del Tesoro a la financiación (algo esencial. dada la prohibición absoluta de recurrir a la financiación del BdE desde 1994 ), cada año se publica en el Boletín Oficial del Estado el calendario anual de subastas. En el caso de las letras, se celebran dos subastas mensuales, la primera de letras a 6 y 12 meses, y la segunda, de letras a 3 y 9 meses. En el caso de las subastas de bonos y obligaciones, también se celebran dos cada mes, y el Tesoro anuncia el viernes anterior a cada subasta los bonos y obligaciones concretos a emitir.
    
    Además, el lunes de aquellas semanas en las que hay subasta de letras o de bonos y obligaciones, el Tesoro anuncia el rango objetivo de colocación de la subasta. Esto último se empezó a hacer en 1995, con el objetivo de evitar las grandes oscilaciones que había hasta ese momento en los volúmenes emitidos en cada subasta.
\end{specialblock}


ii. Sindicaciones
Las sindicaciones bancarias son el segundo método de emisión más importante para el Tesoro, y se utilizan, con carácter general, para realizar la emisión del primer tramo de las obligaciones del Estado. Ello se explica por algunas de las ventajas que ofrece este mecanismo de emisión respecto de la subasta, entre las que se incluyen:

• La posibilidad de emitir una mayor cantidad de valores que en una subasta ordinaria, lo que favorece la liquidez del nuevo instrumento desde su lanzamiento. Por poner un ejemplo, en plena pandemia del covid-19, el Tesoro utilizó este procedimiento de emisión para captar 15.000 millones de euros de una nueva obligación del Estado en una sola operación.

• El control de la distribución de la emisión: a diferencia de las subastas, en las que el procedimiento es ciego y la distribución de la emisión depende de las pujas realizadas por los inversores, en las sindicaciones el emisor ( en este caso, el Tesoro) define el precio de emisión y los inversores indican la cantidad de valores que estarían dispuestos a adquirir a ese precio. Habitualmente, la demanda supera ampliamente la cantidad que el Tesoro desea emitir, lo que permite al Tesoro determinar la cantidad final a asignar a cada inversor y, con ello, garantizar la diversificación de la base inversora inicial en ese título.\footnote{Para garantizar que esto último se hace de forma objetiva, se clasifica a los inversores por categoría (bancos centrales. instituciones oficiales, aseguradoras, fondos de pensiones, gestoras de fondos, bancos, hedge fimds, mesas de trading) y se definen unos porcentajes de asignación sobre las peticiones realizadas para cada categoría de inversor. El porcentaje de asignación es mayor para aquellas categorías de inversores más deseables para el Tesoro, viniendo determinada esa deseabilidad por la estabilidad de la inversión del inversor. Así. se adjudica un mayor porcentaje de sus peticiones a aquellos inversores que cabe esperar que mantengan los valores en su canera durante periodos de tiempo más largos, como los bancos centrales o las instituciones oficiales. Otro criterio a tener en cuenta para asignar la emisión es el país o región de origen del inversor, favoreciéndose a aquellos inversores de regiones con menor inversión acumulada en deuda del Estado
}

• La publicidad o notoriedad que ofrecen al Tesoro, al favorecer su presencia mediática.

Su inconveniente principal es su mayor coste, pues no se llega a conocer el precio marginal al que el inversor estaría dispuesto a adquirir los valores, como sí ocurre en las subastas. Además, es necesario pagar comisiones al sindicato bancario, algo que no ocurre en las subastas. Finalmente, en las sindicaciones hay una especial dependencia del Tesoro del sindicato bancario, que tiene un mayor acceso a información de primera mano de los inversores finales que el Tesoro.

Normalmente, el Tesoro obtiene algo más del I 0% de la emisión bruta anual mediante sindicaciones. En 2020, no obstante, el fuerte crecimiento de las necesidades de financiación obligó al Tesoro a aumentar el recurso a este procedimiento de emisión, que sirvió para cubrir casi el 20% del programa de financiación de ese año.


iii. Otros procedimientos de emisión: colocaciones privadas, canjes
Además de las subastas y sindicaciones, existen otros procedimientos de emisión con una importancia residual en la estrategia de financiación del Tesoro. Es el caso (i) de las colocaciones privadas, por las que se emite un valor de forma directa a un inversor o grupo de inversores; (ii) de los canjes de deuda, por los que se retira total o parcialmente del mercado una determinada referencia entregando a cambio a los inversores interesados otras referencias distintas; y (iii) de los préstamos, que veíamos al analizar los instrumentos a disposición del Tesoro.

Las colocaciones privadas se realizan a propuesta de los propios inversores, y se ejecutan en la medida en que contribuyan a la diversificación de la base inversora del Tesoro, permitan reducir la carga de intereses de la deuda del Estado y encajen en las líneas estratégicas establecidas por el Tesoro.

Los canjes de deuda, por su parte, se utilizan de forma excepcional con el objetivo de aumentar la liquidez del mercado y de suavizar el perfil de vencimientos de la deuda del Estado. Así, las referencias emitidas mediante este procedimiento sustituyen habitualmente otras referencias poco líquidas y con un vencimiento próximo en el tiempo.

Finalmente, el objetivo principal de l_os préstamos es la reducción de la carga de intereses de la deuda del Estado y el alargamiento de su vida media, que en el caso de los préstamos se sitúa habitualmente por encima de la de la deuda del Estado en circula_ción.

\begin{specialblock}{\textbf{Mercado:} Creadores de mercado de deuda público España}
    Los creadores de mercado son un grupo de entidades financieras que adquieren con el Tesoro el compromiso de favorecer la liquidez del mercado secundario de deuda pública y de colaborar en la difusión exterior e interior de la deuda del Estado.
    
    Existen dos grupos de creadores de mercado, el grupo de creadores de mercado de letras del Tesoro y el grupo de creadores de mercado de bonos y obligaciones del Estado, si bien la práctica totalidad de creadores pertenece a ambos grupos. Así, actualmente hay 19 creadores de mercado de letras del Tesoro y 18 creadores de mercado de bonos y obligaciones del Estado. El acceso a la condición de creador de mercado exige cumplir con ciertos requisitos en cuanto a los medios técnicos y humanos disponibles y demostrar con su actividad el compromiso con el mercado español de deuda pública. Además. los creadores de mercado adquieren una serie de compromisos en cuanto a participación en las subastas y provisión de liquidez en el mercado secundario. A cambio, pueden participar en las subastas en condiciones más favorables (siendo los únicos que pueden introducir peticiones en los treinta minutos previos a la subasta), ser parte de los sindicatos bancarios contratados por el Tesoro para llevar a cabo emisiones sindicadas (a cambio de lo que reciben. como veíamos antes, comisiones), acceder a las segundas vueltas de las subastas o segregar y reconstituir valores de deuda del Estado.\footnote{Las segundas vueltas de las subastas, introducidas por el Tesoro en 1 99 1 , penniten a los creadores de mercado solicitar hasta un máximo del 24% del importe nominal que adquirieron en la subasta al precio medio ponderado redondeado resultante de esta. El porcentaje que puede solicitar cada creador depende de su participación en la subasta en cuestión y de su desempeño como creador de mercado. Dado que la segunda vuelta se celebra entre 2 y 3 días después de la subasta, en función de la evolución del precio de las letras/bonos y obligaciones. los creadores de mercado pueden tener una oportunidad de arbitraje. Así. la segundad vuelta es equivalente a una opción .. catr' gratuita para los creadores de mercado. que ejercerán en caso de que los tipos caigan entre la subasta y la segunda vuelta
    }
    
    El Tesoro elabora mensualmente una evaluación del desempeño de los creadores de mercado en base a criterios objetivos y públicos. Los resultados de esa evaluación se utilizan para determinar el acceso de los creadores a las segundas vueltas de las subastas o para seleccionar a las entidades que componen los sindicatos bancarios contratados por el Tesoro para sus emisiones sindicadas, entre otras cosas. En 2022, los creadores de mercado de bonos y obligaciones del Estado con mejor desempeño fueron Deutsche Bank, BB V A y JP Morgan. En el caso de los creadores de mercado de letras del Tesoro, los que tuvieron un mejor desempeño fueron Crédit Agricole, Banco Santander y BBV A.
\end{specialblock}


111.IV. La estrategia de financiación del Tesoro en la actualidad

\textcolor{red}{\textbf{NOTA.-} Para actualizar este apartado: Consultar estrategia anual de financiación del Tesoro. Se publica anualmente los primeros días del mes de enero y está disponible en la página web del Tesoro.}

El año 2022 supuso el fin a una larga etapa de estímulos monetarios por parte del BCE que, para hacer frente a los elevados niveles de inflación, terminó con las compras netas de deuda bajo sus dos programas de compra de activos y subió sus tipos de interés por primera vez en más de 10 años.\footnote{El Asset Purchase Programme o APP y el programa de compras de emergencia frente a la pandemia o PEPP. por sus siglas en inglés.
} Ese cambio de tendencia de la política monetaria se produjo tras un año, el 2021, en el que el Tesoro había tenido, por primera vez en su historia, un coste medio de emisión de su deuda negativo, del -0,04%. El coste de toda la deuda en circulación, por su parte, alcanzó también en 2021 un mínimo histórico del 1,64%.

El cambio de tendencia de la política monetaria había sido anticipado por el Tesoro, que llevaba años preparando su cartera de deuda pública mediante el alargamiento de la vida media de la deuda en circulación (que en 2021 alcanzó por primera vez los 8 años) para anclar los bajos costes de financiación, y medíante la diversificación de su base inversora.

Pese a ello, el nuevo escenario obligará al Tesoro a replantear la estrategia seguida estos últimos años, encaminada precisamente a prepararse ante una eventual subida de los tipos de interés. Así, la estrategia anual de financiación del Tesoro de 2023 prevé la estabilización de la vida media de la deuda en circulación en los niveles actuales. Para conseguirlo, el Tesoro espera obtener toda la financiación neta (q ue ascenderá, en principio, a 70.000 millones de euros) a través de la emisión de instrumentos a medio y largo plazo, realizando una emisión neta de letras del Tesoro negativa por importe de 5.000 millones de euros (de forma que el peso de las letras del Tesoro sobre el total de la deuda del Estado en circulación seguirá cayendo en 2023, hasta alrededor del 5\% del total).\footnote{El Tesoro mantuvo su estrategia de reducir la emisión neta de letras incluso durante la pandemia. renunciando al ahorro a corto plazo que ello habría supuesto a cambio de anclar los bajos coste de financiación que había en ese momento a medio y largo plazo. Muchos países de nuestro entorno, en cambio, decidieron aumentar sus emisiones a corto plazo para hacer frente a las necesidades de financiación extraordinarias provocadas por la pandemia.
}

Para terminar, cabe mencionar que recientemente se ha producido un fenómeno desconocido para el Tesoro en los últimos años, como es el au mento de la den:ianda de deuda por parte de los inversores minoristas. La deuda del Estado en manos de personas fisicas alcanzó un mínimo de 898 millones de euros (un 0,0 8\% del total) en abril de 2022 y comenzó a aumentar desde entonces, de manera especialmente intensa en los primeros meses de 2023.\footnote{A pesar de ello, los inversores minoristas españoles tenían a finales de 2022 apenas el 0,26% de la deuda del Estado en circulación, por debajo del 3.8% de media de la Unión y muy lejos del 7,9% de Italia o del 13,2% de Portugal.
} La demanda minorista, no obstante, se concentra en los plazos más cortos (y, en particular, en las letras del Tesoro), lo que podría llevar al Tesoro a plantearse aumentar ligeramente su emisión.


\begin{specialblock}{\textbf{Programas BCE:} Programas compra de activos}
    A finales de 2014, ante la incapacidad de las medidas convencionales de política monetaria para situar los niveles de inflación en su objetivo, el BCE anunció un paquete de medidas no convencionales entre las que se incluyó un programa de compra de activos, el A sset Purchase Programme o APP.\footnote{Entre las medidas anunciadas se incluyeron también las TL TRO o largeted longer-term rejinan cing operalions que, mediante el ofrecimiento a los bancos de financiación a largo plazo en condiciones muy favorables, perseguían mejorar el flujo de crédito bancario al sector privado no financiero.
    } Los programas de compra de activos también se denominan medidas de expansión cuantitativa (o Quantitative Easing o QE, por sus siglas en inglés) por implicar una expansión notable del tamaño del balance del banco central.

    A través del APP, el BCE ha adquirido títulos de deuda pública (a través del subprograma PSPP), bonos corporativos (CSPP), bonos de titulización de activos (ABSPP) y bonos garantizados (CBPP). A este primer programa se le sumó, en marzo de 2020, el programa de compras de emergencia frente a la pandemia (PEPP). En ambos casos, las compras de títulos de deuda pública se realizan en proporción a la clave de capital de cada banco central nacional en el Eurosistema, si bien se permite una cierta flexibilidad que hace que sean posibles desviaciones temporales.
    
    El ritmo de compra de activos bajo estos programas ha ido variando a lo largo del tiempo. En el caso del APP, las compras netas máximas se produjeron entre 2016 y 2017, cuando llegaron a alcanzar los 80.000 millones de euros mensuales. En el caso del PEPP, las compras netas mensuales llegaron a alcanzar los 120.000 millones en junio de 2020. Actualmente, el BCE no. realiza compras netas de activos bajo ninguno de estos programas, y desde marzo de 2023 reinvierte solo parcialmente los vencimientos de la deuda adquirida en el pasado, habiendo comenzado por tanto la reducción del tamaño de su balance (algo conocido como quantitative tightening o QT, por sus siglas en inglés). 

    Estos programas han permitido reducir significativamente el coste de emisión de la deuda del Estado, al absorber la mayor parte de la emisión neta realizada desde 20 15. En ese sentido, resulta interesante analizar la evolución del stock de deuda nominal del Estado detrayendo la deuda adquirida por el BCE bajo estos dos programas.
\end{specialblock}



































\clearpage
\section{Conclusión} 

La Transición española es uno de los acontecimientos más relevantes de nuestra historia reciente y simboliza el cierre de un periodo histórico caracterizado por la turbulencia política y el aislamiento internacional. Sin embargo, para hacerla posible los principales partidos políticos, asociaciones empresariales y sindicatos acordaron priorizar las cuestiones políticas y sociales sobre las presupuestarias. 

Como resultado, el saldo presupuestario español empieza a mostrar en ese momento déficits recurrentes, una situación que, con un breve paréntesis entre 2005 y 2007, se ha mantenido hasta nuestros días. Ello ha provocado un gran crecimiento de nuestra deuda pública, que llegó a superar en 2020 el 1 20% del PIB, un nivel no visto desde el siglo XIX.

Lo anterior ha convertido la política de financiación del Tesoro en una cuestión clave, y así seguirá siendo en los próximos años, dada la necesidad de refinanciar unos vencimientos crecientes y el final de los préstamos del fondo SURE y de los préstamos y transferencias del programa NGEU.

En la última década, la actuación del Tesoro se ha visto favorecida por el cambio del patrón de crecimiento de la economía española hacia uno más equilibrado y sostenible en el tiempo, que ha aumentado el atractivo de la deuda española para los inversores; y por la política de estímulos monetarios del BCE, que ha contribuido a reducir el coste de emisión a niveles mínimos históricos. Como resultado, la carga de intereses como porcentaje de los ingresos públicos se mantenía en 2022 en niveles similares a los de los años previos a la crisis financiera internacional ( de hecho, la carga de intereses en valor absoluto fue menor en 2021 que en 2011, cuando la deuda pública representaba menos del 70% del PIB).

Lo anterior ha sido posible incluso aumentando la vida media de la deuda en circulación a niveles nunca vistos, lo que ha permitido anclar los bajos costes de financiación y reducir el riesgo de refinanciación. A pesar de ello, la necesidad de financiación bruta anual del Tesoro estos últimos años es equivalente a algo más del 20% del PIB, y la rápida subida de los tipos de interés provocada por el brusco cambio de tendencia de la política monetaria en 2022 se dejará sentir pronto sobre la carga de intereses de la deuda.

A largo plazo, por tanto, el éxito de la política de financiación del Tesoro dependerá en gran medida de la capacidad de la economía para crecer de forma sostenible y para acometer un nuevo proceso de consolidación fiscal, pues si bien es evidente que en el contexto actual resultan sostenibles niveles de deuda pública superiores a los previstos en la normativa europea y en la Ley Orgánica de Estabilidad Presupuestaria y Sostenibilidad Financiera, también lo es que niveles muy elevados de deuda pública suponen una importante fuente de vulnerabilidad y pueden tener un efecto negativo sobre e l crecimiento.\footnote{La teoría económica discute si existe un umbral a partir del cual la deuda pública se convierte en un hándicap para el crecimiento económico a largo plazo. Así, mientras autores como Reinhart y RogolT (201 O) situaron dicho umbral en el 90%, Herndon, Ash y Pollin (2014) y otros muchos autores discuten esa conclusión, sugiriendo que la relación de causalidad podría ser inversa, es decir. que serían precisamente los bajos niveles de crecimiento los que provocarían los elevados niveles de deuda pública en muchas economías.
}


\subsection{Recapitulación}


\subsection{Extensiones y relación con otras partes del Temario}



\subsection{Opinión}

\subsubsection{Idea final (Salida o cierra)}

\clearpage
\pagenumbering{gobble}
\MyBibliography




- AIReF (2025). Informe sobre la ejecución presupuestaria, deuda pública y regla de gasto 2025.
- AIReF (2026). Informe de seguimiento 2026 del Plan Fiscal y Estructural de Medio Plazo 2025-2028.
- BdE (2026). La deuda de las Administraciones Públicas se situó en el 100,7% del PIB al cierre de 2025. Nota de prensa, 31 de marzo de 2026.
- Blanchard, O. (2019). Public Debt and Low Interest Rates. American Economic Review, 109(4).
- Fedea (2025). Estimación del gasto futuro en intereses de la deuda pública. Apuntes 2025/28.
- Ministerio de Hacienda (2026). España cierra 2025 con un déficit público del 2,2% del PIB. Nota de prensa, 31 de marzo de 2026.
- Reglamento (CE) 479/2009 del Consejo, relativo a la aplicación del Protocolo sobre el procedimiento aplicable en caso de déficit excesivo.
- Reglamento (Unión) 549/2013 (SEC 2010).

Pendiente de verificar:
- la estimación más reciente del saldo estructural de España (Comisión y AIReF, 2025 y 2026);
- el gasto en intereses del conjunto de las AAPP en 2025 en porcentaje del PIB;
- las estimaciones revisadas del saldo estructural de 2007;
- el tratamiento de la DANA en la evaluación europea;
- la fecha de reclasificación del FADE y su saldo vivo;
- el tema del tercer ejercicio donde se desarrollan el producto potencial y la brecha de producción.

Fuentes de actualidad consultadas:
- [Infobae (31 de marzo de 2026): déficit de 2025 del 2,18%, o del 2,39% con la DANA](https://www.infobae.com/america/agencias/2026/03/31/el-deficit-publico-cierra-2025-en-el-218-del-pib-con-36780-millones-y-mejora-el-objetivo-del-gobierno/)
- [Infobae (31 de marzo de 2026): la deuda pública cerró 2025 en el 100,7% del PIB, con desglose por subsectores](https://www.infobae.com/espana/agencias/2026/03/31/la-deuda-publica-cerro-2025-en-el-1007-del-pib-el-nivel-mas-bajo-desde-la-pandemia/)
- [El Diario de Madrid (junio de 2026): el Estado gastó en 2025 34.900 millones en intereses y 7.684 en inversión](https://www.eldiariodemadrid.es/articulo/economia/espana-gasto-45-euros-intereses-deuda-cada-euro-invertido-2025/20260617080650134214.html)
- [El Español (noviembre de 2025): más de 40.000 millones al año en intereses](https://www.elespanol.com/invertia/economia/20251106/espana-pagara-millones-intereses-deuda-ano-equivalente-recaudacion-sociedades/1003744001419_0.html)
- [La Crónica del Henares (septiembre de 2026): intereses hasta julio de 2026](https://www.cronicadelhenares.com/2026/09/deficit-publico-espana-2026-intereses-deuda-comunidades.html)
- [El Español (23 de marzo de 2021): la Sareb sumará 35.000 millones a la deuda pública](https://www.elespanol.com/invertia/empresas/banca/20210323/sareb-sumara-millones-publica-gobierno-puerta-extender/568194202_0.html)
- [El Independiente (23 de marzo de 2021): Europa obliga a incluir a la Sareb en la deuda](https://www.elindependiente.com/economia/2021/03/23/europa-obliga-al-gobierno-a-que-incluya-a-la-sareb-en-el-computo-de-deuda-publica-que-sube-al-120-en-2020/)
- [Fedea: estimación del gasto futuro en intereses de la deuda pública](https://documentos.fedea.net/pubs/ap/2025/ap2025-28.pdf)
- [Forbes: la estrategia del Tesoro para 2026](https://forbes.es/economia/842646/la-emision-neta-del-tesoro-se-mantendra-en-55-000-millones-en-2026-pero-la-bruta-subira-un-42/)
- [AIReF: seguimiento 2026 del Plan Fiscal y Estructural de Medio Plazo](https://www.airef.es/es/centro-documental/informes/informe-de-seguimiento-2026-del-plan-fiscal-y-estructural-de-medio-plazo-2025-2028/)
Referencias nuevas
- Alesina, A. y Drazen, A. (1991), Why Are Stabilizations Delayed?, American Economic Review.
- Bentolila, S. y Dolado, J. J. (1994), Labour Flexibility and Wages: Lessons from Spain, Economic Policy.
- Comín, F. (1996), Historia de la Hacienda pública, II: España (1808-1995), Crítica.
- Comín, F. y Díaz Fuentes, D. (2005), Sector público administrativo y Estado del bienestar, en Carreras, A. y Tafunell, X. (coords.), Estadísticas históricas de España, Fundación BBVA.
- Comín, F. y Díaz Fuentes, D. (2023), Las transformaciones fiscales de la Hacienda y las Administraciones públicas españolas, 1898-2023, ICE, Revista de Economía, 933.
- Escario, R., Sabaté, M. y Gadea, M. D. (2011), La sombra monetaria del déficit en la España de la peseta, Investigaciones de Historia Económica.
- Esteve, V. y Prats, M. A. (2024), Evolución histórica de las variables macroeconómicas fiscales de la economía española, 1964-2022, ALdE (blog académico).
- Fuentes Quintana, E. (1990), Las reformas tributarias en España. Teoría, historia y propuestas, Crítica.
- Galí, J. y Perotti, R. (2003), Fiscal Policy and Monetary Integration in Europe, Economic Policy.
- Giavazzi, F. y Pagano, M. (1988), The Advantage of Tying One's Hands: EMS Discipline and Central Bank Credibility, European Economic Review.
- Lago Peñas, S. (2024), La dinámica de la deuda pública en España: presente, pasado y futuro, Funcas, Documentos de trabajo.
- Prats Albentosa, M. A. y Rocamora Martí, A. M. (2016), Análisis de la sostenibilidad de la deuda pública en España, Revista de Ciencias Sociales, 22(2).
- Reinhart, C. y Sbrancia, M. B. (2015), The Liquidation of Government Debt, Economic Policy.
- Roubini, N. y Sachs, J. (1989), Political and Economic Determinants of Budget Deficits in the Industrial Democracies, European Economic Review.
- Sargent, T. y Wallace, N. (1981), Some Unpleasant Monetarist Arithmetic, Federal Reserve Bank of Minneapolis Quarterly Review.
- Sudrià, C. (2014), Las crisis bancarias en España: una perspectiva histórica, Estudios de Economía Aplicada.
- von Hagen, J. (1992), Budgeting Procedures and Fiscal Performance in the European Communities, Economic Papers, 96, Comisión.

Pendiente de verificar
- La serie homogénea del déficit AAPP para 1985, 1989, 1990, 1993 y 1995 (AMECO, SEC-2010 retropolado). En particular, si el máximo del periodo es 1993 o 1995 y si el mínimo de la fase 1986-1990 es 1989.
- La presión fiscal OCDE de España en 1975, 1985 y 1990 y la media OCDE: las cifras ~18%, ~27% y ~31-32% están tomadas de memoria.
- La media de deuda de la CEE en 1975, para sustituir el casi 50% del documento.
- El déficit primario de 1993: el 0,5% de Prats y Rocamora frente a mi estimación de ~2%.
- La descomposición de la variación de la deuda en 1993 (bola de nieve, primario y ajustes). Mi cifra es una estimación aritmética.
- Los términos exactos de la conversión del saldo deudor del Tesoro con el BdE en 1994 (norma, importe y plazo). La Ley 13/1994 no los recoge en su articulado.
- La vida media de la deuda del Estado a principios de los noventa.
- La fecha de supresión definitiva del coeficiente de inversión obligatoria (Ley 13/1985 y su desmantelamiento, probablemente a principios de los noventa).
- Las cifras del documento para el gasto público: >40% en 1985 y >45% en 1993.

Fuentes de actualidad consultadas
- [Ley 13/1994, de Autonomía del BdE (BOE)](https://www.boe.es/buscar/doc.php?id=BOE-A-1994-12553)
- [Escario, Sabaté y Gadea (2011), La sombra monetaria del déficit en la España de la peseta](https://recyt.fecyt.es/index.php/IHE/article/download/70317/42540/219443)
- [Sudrià (2014), Las crisis bancarias en España: una perspectiva histórica](https://www.redalyc.org/pdf/301/30130732001.pdf)
- [Prats Albentosa y Rocamora (2016), Análisis de la sostenibilidad de la deuda pública en España](https://www.redalyc.org/journal/280/28049145002/html/)
- [Lago Peñas (2024), Funcas: La dinámica de la deuda pública en España](https://www.funcas.es/documentos_trabajo/la-dinamica-de-la-deuda-publica-en-espana-presente-pasado-y-futuro/)
- [Esteve y Prats (2024), ALdE: Evolución histórica de las variables fiscales 1964-2022](https://alde.es/blog/evolucion-historica-de-las-variables-macroeconomicas-fiscales-de-la-economia-espanola/)
- [Comín y Díaz Fuentes (2023), ICE 933](https://www.revistasice.com/index.php/ICE/article/view/7685)
- [FMI, base histórica de deuda pública (DataMapper, GG_DEBT_GDP)](https://www.imf.org/external/datamapper/api/v1/GG_DEBT_GDP/ESP)
- [Ley 13/1985, de coeficientes de inversión (BOE)](https://www.boe.es/buscar/doc.php?id=BOE-A-1985-9680)
Referencias nuevas
- Alberola, E., Estrada, Á. y Santabárbara, D. (2014), Growth and imbalances in Spain: a reassessment of the output gap, SERIEs.
- Alesina, A. y Ardagna, S. (2010), Large Changes in Fiscal Policy: Taxes versus Spending, Tax Policy and the Economy.
- Alesina, A. y Perotti, R. (1995), Fiscal Expansions and Adjustments in OECD Countries, Economic Policy.
- De Castro, F., Estrada, Á., Hernández de Cos, P. y Martí, F. (2008), Una aproximación al componente transitorio del saldo público en España, Boletín Económico del BdE, junio.
- De Grauwe, P. (2011), The Governance of a Fragile Eurozone, CEPS Working Document.
- Estrada, Á., Jimeno, J. F. y Malo de Molina, J. L. (2009), La economía española en la UEM: los diez primeros años, Documentos Ocasionales del BdE, 0901.
- Fernández-Villaverde, J., Garicano, L. y Santos, T. (2013), Political Credit Cycles: The Case of the Eurozone, Journal of Economic Perspectives.
- Giavazzi, F. y Pagano, M. (1990), Can Severe Fiscal Contractions Be Expansionary? Tales of Two Small European Countries, NBER Macroeconomics Annual.
- Guajardo, J., Leigh, D. y Pescatori, A. (2014), Expansionary Austerity? International Evidence, Journal of the European Economic Association.
- Instituto Monetario Europeo (1998), Informe de Convergencia, marzo.
- Martí, F. y Pérez, J. J. (2015), Spanish Public Finances through the Financial Crisis, Fiscal Studies.
- Perotti, R. (2011), The 'Austerity Myth': Gain without Pain?, NBER Working Paper 17571.

Pendiente de verificar
- Los importes y el número de empresas del Programa de Modernización del Sector Público Empresarial (30.000 M€, 36 empresas, 1996-2000/2001), en el balance de la SEPI.
- La serie de intereses de las AAPP sobre el PIB (máximo 1995-1996 en ~5,3% y 2007 en ~1,6%) con AMECO o IGAE.
- El saldo primario de 1996 y 2007 y mi descomposición aproximada de la caída de la deuda (dos tercios de superávit primario y un tercio de bola de nieve).
- La vida media de la deuda del Estado en 1995 y 2007, y la proporción de no residentes en 2007-2008.
- Los importes de la asunción de deuda de RENFE (2004) y RTVE (2006-2007) y su impacto en el saldo del año.
- Los umbrales exactos de crecimiento de la reforma de la LGEP de 2006.
- El déficit de 1999 con la serie SEC-2010.
- El saldo estructural ex post de 2007 en la última vintage de AMECO.

Fuentes de actualidad consultadas
- [IME (1998), Informe de Convergencia (versión española)](https://www.ecb.europa.eu/pub/pdf/conrep/cr1998es.pdf)
- [Ministerio de Economía y Hacienda (2008), presentación del cierre de 2007](https://www.hacienda.gob.es/GabineteMinistro/presentaciones/21-02-08%20presentacion%20cierre%202007.pdf)
- [De Castro, Estrada, Hernández de Cos y Martí (2008), Boletín Económico del BdE](https://www.bde.es/f/webbde/SES/Secciones/Publicaciones/InformesBoletinesRevistas/BoletinEconomico/08/Jun/Fich/art4.pdf)
- [Martí y Pérez (2015/2016), Documento de Trabajo del BdE 1620](https://www.bde.es/f/webbde/SES/Secciones/Publicaciones/PublicacionesSeriadas/DocumentosTrabajo/16/Fich/dt1620e.pdf)
- [Alberola, Estrada y Santabárbara (2014), SERIEs](https://link.springer.com/article/10.1007/s13209-014-0112-z)
- [Ministerio de Inclusión: el Fondo de Reserva alcanzó 57.223 M€ en 2008](https://www.inclusion.gob.es/w/el-fondo-de-reserva-de-la-seguridad-social-alcanzo-57.223-17-millones-de-euros-en-2008)
- [Conde-Ruiz, Lecciones del pasado: el fondo de reserva que pudo ser y no fue (Nada es Gratis)](https://nadaesgratis.es/j-ignacio-conde-ruiz/lecciones-del-pasado-el-fondo-de-reserva-que-pudo-ser-y-no-fue)
Referencias nuevas
- Alesina, A., Favero, C. y Giavazzi, F. (2019), Austerity: When It Works and When It Doesn't, Princeton University Press.
- BdE (2015), La evolución de la deuda pública en España en 2014, Boletín Económico, julio.
- BdE (2020), La evolución de la deuda pública en España en 2019, Notas Económicas, Boletín Económico 3/2020.
- Blanchard, O. y Leigh, D. (2013), Growth Forecast Errors and Fiscal Multipliers, American Economic Review: Papers & Proceedings.
- De Grauwe, P. y Ji, Y. (2013), Self-fulfilling Crises in the Eurozone: An Empirical Test, Journal of International Money and Finance.
- Guajardo, J., Leigh, D. y Pescatori, A. (2014), Expansionary Austerity? International Evidence, Journal of the European Economic Association.
- Tribunal de Cuentas (2025), actualización del Informe de fiscalización sobre el proceso de reestructuración bancaria (coste a 1 de enero de 2022).

Pendiente de verificar
- La serie vigente de déficit y deuda 2008-2019 en Eurostat, en especial el efecto exacto de la Sareb sobre 2012-2014. Uso 2012: −11,5%, 2014: 105,1% y 2019: 98,2%.
- El impacto en el déficit de las ayudas bancarias de 2012 (~3,7 puntos) con la serie actual.
- La recaudación final de la regularización fiscal de 2012.
- El saldo pendiente del préstamo del MEDE en 2026 y si hubo nuevas amortizaciones anticipadas después de 2018.
- El techo de gasto de 2011 y 2012, para cuantificar cuánto de la caída se debe al cambio de modelo autonómico.
- El saldo estructural de 2019 en la última estimación de la Comisión y de la AIReF.
- La presión fiscal de 2007 y 2009 en la serie SEC-2010.

Fuentes de actualidad consultadas
- [BdE, La evolución de la deuda pública en España en 2014](https://www.bde.es/f/webbde/SES/Secciones/Publicaciones/InformesBoletinesRevistas/BoletinEconomico/15/Jul/Fich/be1507-art4.pdf)
- [BdE, La evolución de la deuda pública en España en 2019](https://www.bde.es/f/webbde/SES/Secciones/Publicaciones/InformesBoletinesRevistas/NotasEconomicas/20/T3/descargar/Fich/be2003-ne05.pdf)
- [Tribunal de Cuentas: actualización del coste de la reestructuración bancaria (71.833 M€)](https://www.tcu.es/es/sala-de-prensa/noticias/El-TCu-actualiza-hasta-los-71.833-millones-de-euros-el-coste-de-la-reestructuracion-bancaria/)
- [Hoy Aragón: detalle del informe del Tribunal de Cuentas (marzo de 2025)](https://www.hoyaragon.es/articulo/noticias-espana/tribunal-cuentas-rescate-bancario/20250316192541090616.html)
- [countryeconomy.com: serie de déficit de España 2007-2025 (datos Eurostat)](https://countryeconomy.com/deficit/spain)
- [Martí y Pérez, Documento de Trabajo del BdE 1620](https://www.bde.es/f/webbde/SES/Secciones/Publicaciones/PublicacionesSeriadas/DocumentosTrabajo/16/Fich/dt1620e.pdf)
Referencias nuevas
- AIReF (2026), Informe sobre la ejecución presupuestaria, deuda pública y regla de gasto 2026, julio.
- BdE (2025), La evolución de la deuda pública en España en 2024, Boletín Económico.
- Banco Central Europeo (2026), Decisiones de política monetaria, 11 de junio y 10 de septiembre.
- Blanchard, O. (2019), Public Debt and Low Interest Rates, American Economic Review.
- Consejo de la Unión (2025), Recomendación de 21 de enero de 2025 por la que se aprueba el plan fiscal-estructural nacional a medio plazo de España, DOUE.
- FMI (2021), Fiscal Monitor: Database of Country Fiscal Measures in Response to the COVID-19 Pandemic.
- Galesi, A., Nuño, G. y Thomas, C. (2017), El tipo de interés natural: concepto, determinantes e implicaciones para la política monetaria, Boletín Económico del BdE.
- Real Decreto-ley 25/2026, de 29 de septiembre, en el marco del Plan Integral de Respuesta a la Crisis en Oriente Medio, BOE.
- Reglamento (Unión) 2024/1263 del Parlamento Europeo y del Consejo, relativo a la coordinación eficaz de las políticas económicas y a la supervisión presupuestaria multilateral.

Pendiente de verificar
- La serie vigente de déficit de 2020-2024 tras la revisión del PIB de 2024; uso 2020 en torno al 10%, 2024 en el 3,2% con DANA y 2,8% sin ella.
- El número máximo de trabajadores en ERTE (abril de 2020) frente al acumulado de 6,6 millones del documento.
- La pérdida efectiva de los avales ICO a 2026.
- El coste total de las medidas frente a la crisis energética de 2022-2023 (~45.000 M€).
- La presión fiscal de 2024 y 2025 en porcentaje del PIB.
- El volumen total del Plan de Recuperación tras la adenda, los préstamos efectivamente recibidos y la solicitud del último desembolso (plazo del 30 de septiembre de 2026).
- El saldo estructural de 2025 y 2026 según la Comisión y la AIReF, y la valoración del cumplimiento de la senda de gasto neto en 2025.
- La fecha en que España recuperó el nivel de PIB real previo a la pandemia (probablemente 2023).
- El coste final del Plan de Respuesta a la crisis de Oriente Medio y su tratamiento como medida excepcional.
- La tramitación del proyecto de ley de condonación de deuda autonómica en los próximos meses.
- El peso máximo del Eurosistema en la deuda española y su evolución desde 2023.

Fuentes de actualidad consultadas
- [BdE, La evolución de la deuda pública en España en 2024](https://www.bde.es/f/webbe/SES/Secciones/Publicaciones/PublicacionesSeriadas/DocumentosOcasionales/25/Fich/be2503-art04.pdf)
- [BCE, decisiones de política monetaria del 10 de septiembre de 2026](https://www.ecb.europa.eu/press/pr/date/2026/html/ecb.mp260910~314e508016.es.html)
- [BdE, el BCE sube los tipos 25 puntos básicos en junio de 2026](https://www.bde.es/wbe/en/noticias-eventos/actualidad-bce/decisiones-politica-monetaria/bce-tipos-junio26.html)
- [La Moncloa, Plan de respuesta a la guerra en Oriente Próximo (31 de agosto de 2026)](https://www.lamoncloa.gob.es/serviciosdeprensa/notasprensa/presidencia/paginas/2026/plan-respuesta-guerra-oriente.aspx)
- [BOE, Real Decreto-ley 25/2026, de 29 de septiembre](https://www.boe.es/diario_boe/txt.php?id=BOE-A-2026-20265)
- [Qué.es, el IPC de agosto de 2026 sube al 4,3%](https://www.que.es/2026/09/15/ipc-agosto-2026-maximo-tres-anos/)
- [The Objective, la AIReF reclama un ajuste del 0,6% del PIB (15 de julio de 2026)](https://theobjective.com/economia/2026-07-15/airef-gobierno-ajuste-pib-normas-fiscales/)
- [BOE, Recomendación del Consejo por la que se aprueba el plan fiscal-estructural de España](https://www.boe.es/buscar/doc.php?id=DOUE-Z-2025-70011)
- [Infobae, España encara el final de los fondos europeos con 25.862 millones pendientes (31 de agosto de 2026)](https://www.infobae.com/espana/2026/08/31/espana-encara-el-final-de-los-fondos-europeos-con-25862-millones-aun-pendientes-tras-haberse-embolsado-78000-millones/)
- [Ministerio de Hacienda, España recibe 6.234 millones del sexto desembolso](https://www.hacienda.gob.es/es-es/prensa/noticias/paginas/2026/20260811-np-pago-sexto-desembolso-prtr.aspx)
- [Infobae, el Congreso reactiva la ley de condonación de deuda autonómica (29 de septiembre de 2026)](https://www.infobae.com/america/agencias/2026/09/29/el-congreso-reactiva-la-ley-para-condonar-deuda-a-las-comunidades-autonomas-tras-casi-un-ano-guardada-en-un-cajon/)
- [El Economista, la recaudación fiscal de 2025 supera los 325.000 millones](https://www.eleconomista.es/economia/noticias/13850918/03/26/la-recaudacion-fiscal-en-2025-hace-historia-al-superar-los-325000-millones-y-subir-un-104.html)
- [SEPG, Presupuestos Generales del Estado prorrogados para 2026](https://www.sepg.pap.hacienda.gob.es/sitios/sepg/es-ES/Presupuestos/PGE/PGE2025Prorroga/Paginas/PGE2025Prorroga.aspx)
Referencias nuevas
- Abbas, S. A., Belhocine, N., El-Ganainy, A. y Horton, M. (2010), A Historical Public Debt Database, IMF Working Paper, 10/245.
- Alesina, A. y Tabellini, G. (1990), A Positive Theory of Fiscal Deficits and Government Debt, Review of Economic Studies.
- Baldwin, R. y Giavazzi, F. (eds.) (2015), The Eurozone Crisis: A Consensus View of the Causes and a Few Possible Solutions, CEPR Press (VoxEU).
- Buchanan, J. y Wagner, R. (1977), Democracy in Deficit: The Political Legacy of Lord Keynes, Academic Press.
- Buiter, W., Corsetti, G. y Roubini, N. (1993), Excessive Deficits: Sense and Nonsense in the Treaty of Maastricht, Economic Policy.
- Grilli, V., Masciandaro, D. y Tabellini, G. (1991), Political and Monetary Institutions and Public Financial Policies in the Industrial Countries, Economic Policy.
- Koen, V. y van den Noord, P. (2005), Fiscal Gimmickry in Europe: One-off Measures and Creative Accounting, OECD Economics Department Working Papers, 417.
- Lane, P. R. (2012), The European Sovereign Debt Crisis, Journal of Economic Perspectives.
- Milesi-Ferretti, G. M. (2003), Good, Bad or Ugly? On the Effects of Fiscal Rules with Creative Accounting, Journal of Public Economics.
- Persson, T. y Svensson, L. (1989), Why a Stubborn Conservative Would Run a Deficit: Policy with Time-Inconsistent Preferences, Quarterly Journal of Economics.
- Tanzi, V. y Schuknecht, L. (2000), Public Spending in the 20th Century: A Global Perspective, Cambridge University Press.
- von Hagen, J. y Wolff, G. (2006), What Do Deficits Tell Us about Debt? Empirical Evidence on Creative Accounting with Fiscal Rules in the EU, Journal of Banking & Finance.

Pendiente de verificar
- El déficit medio de la OCDE en los setenta, en los ochenta y en 1993 (1-2%, ~3% y 4,9% según el documento), en OECD Economic Outlook.
- La media ponderada de deuda de la CEE de los nueve en 1975 (mi estimación: 35-40%).
- El déficit de Suecia en 1993 y el de Canadá en 1992-1993.
- La cuantía de la eurotasa italiana y del pago de France Télécom (1997).
- El alcance exacto de la revisión de las estadísticas griegas de 2004 (años afectados y déficit revisado).
- El déficit por cuenta corriente de España en 2007 (~9,6% del PIB) y la deuda privada.
- La deuda neta de Japón en la última edición del Fiscal Monitor.
- La diferencia cuantitativa entre la deuda de Estados Unidos según el GFSM y según el criterio de Maastricht.

Fuentes de actualidad consultadas
- [FMI, DataMapper: deuda de las AAPP (% del PIB), base histórica](https://www.imf.org/external/datamapper/api/v1/GG_DEBT_GDP/ITA/BEL/DEU/FRA/GBR/USA/JPN/IRL/GRC/PRT/SWE/CAN/ESP)
- [Abbas et al. (2010), IMF WP 10/245](https://www.imf.org/external/pubs/ft/wp/2010/wp10245.pdf)
- [Grilli, Masciandaro y Tabellini (1991), Economic Policy](https://academic.oup.com/economicpolicy/article-abstract/6/13/341/2392405)
- [von Hagen y Wolff (2006), Journal of Banking & Finance (EconPapers)](https://econpapers.repec.org/RePEc:eee:jbfina:v:30:y:2006:i:12:p:3259-3279)
- [IME (1998), Informe de Convergencia](https://www.ecb.europa.eu/pub/pdf/conrep/cr1998es.pdf)
- [Martí y Pérez, Documento de Trabajo del BdE 1620 (deuda de la zona del euro en 2007)](https://www.bde.es/f/webbde/SES/Secciones/Publicaciones/PublicacionesSeriadas/DocumentosTrabajo/16/Fich/dt1620e.pdf)
- Dolls, M., Fuest, C. y Peichl, A. (2012), Automatic Stabilizers and Economic Crisis: US vs. Europe, Journal of Public Economics.
- Eurostat (2011), Euro area and EU27 government deficit at 6.2% and 6.6% of GDP respectively, Euro indicators, 21 de octubre.
- Fatás, A. y Summers, L. (2018), The Permanent Effects of Fiscal Consolidations, Journal of International Economics.
- FMI (2021), Database of Country Fiscal Measures in Response to the COVID-19 Pandemic, octubre.
- House, C., Proebsting, C. y Tesar, L. (2020), Austerity in the Aftermath of the Great Recession, Journal of Monetary Economics.
- Laeven, L. y Valencia, F. (2018), Systemic Banking Crises Revisited, IMF Working Paper, 18/206.
- Sgaravatti, G., Tagliapietra, S., Trasi, C. y Zachmann, G. (2023), National Fiscal Policy Responses to the Energy Crisis, Bruegel Datasets.

Pendiente de verificar
- Las caídas del PIB de 2009 (España −3,8%, Alemania −5,7%, Italia −5,5%, Reino Unido −4,5%) en las series actuales.
- La consolidación estructural primaria acumulada en 2010-2013 por país (Grecia, Portugal, Irlanda y España).
- El coste fiscal bruto de las crisis bancarias de Irlanda y España según Laeven y Valencia (2018).
- Los saldos primarios de Portugal, Grecia y Chipre en 2022-2023, y el primer superávit total portugués (2019).
- El ranking de déficits de las economías avanzadas en 2019 con los datos actuales del FMI.
- Las cifras de apoyo directo y de liquidez de la base de medidas del FMI (octubre de 2021) para España y sus pares.
- La cifra de apoyo energético de España en la base de Bruegel.
- La deuda irlandesa sobre la GNI en 2019 y 2023.
- La media de deuda de la zona del euro en 2014 y 2023.

Fuentes de actualidad consultadas
- [FMI, DataMapper: deuda de las AAPP (% del PIB), periferia europea](https://www.imf.org/external/datamapper/api/v1/GG_DEBT_GDP/PRT/GRC/IRL/ESP/NLD/BEL/CYP)
- [FMI, DataMapper: deuda de las AAPP (% del PIB), grandes economías](https://www.imf.org/external/datamapper/api/v1/GG_DEBT_GDP/ESP/ITA/FRA/DEU/PRT/GRC/IRL/GBR/USA/JPN/NLD/BEL)
- [Eurostat (21 de octubre de 2011), notificación de déficit y deuda 2007-2010](https://ec.europa.eu/eurostat/documents/2995521/5037534/2-21102011-AP-EN.PDF/3eb0463f-c019-44de-ac16-fc360b4ef4f9?version=1.0)
- [FMI, Database of Country Fiscal Measures (octubre de 2021)](https://www.imf.org/-/media/files/topics/covid/fm-database/october-2021/oct-2021-country-fiscal-measures-database-publication.pdf)
- [Rappler/Bruegel: el gasto europeo en la crisis energética se acerca a 800.000 millones](https://www.rappler.com/business/european-union-national-fiscal-policy-responses-energy-crisis-bruegel-february-2023/)
- [Bruegel, base de datos de respuestas fiscales a la crisis energética](https://www.bruegel.org/dataset/national-policies-shield-consumers-rising-energy-prices)
- [El Blog Salmón: el coste del rescate al sector financiero en la Unión](https://www.elblogsalmon.com/economia/este-es-el-coste-del-rescate-al-sector-financiero-en-la-ue)
Referencias nuevas
- Comisión (2024), Debt Sustainability Monitor 2023, Institutional Paper, marzo de 2024.
- Comisión (2026), Debt Sustainability Monitor 2025, 12 de febrero de 2026.
- Eurostat (2026), Euro area government deficit at 2.9% and EU at 3.1% of GDP, Euro indicators, 22 de abril de 2026.
- Eurostat (2026), Government debt at 87.8% of GDP in euro area, Euro indicators, 22 de abril de 2026.
- Eurostat (2026), Government Finance Statistics, Statistics Explained.
- Comité de personas expertas (2022), Libro Blanco sobre la Reforma Tributaria, Instituto de Estudios Fiscales.

Pendiente de verificar
- La clasificación de riesgo de España en la ficha país del Debt Sustainability Monitor 2025 (corto, medio y largo plazo) y sus indicadores S1 y S2.
- El saldo primario de 2025 de Italia, Grecia, Portugal y Chipre.
- El gasto y los ingresos públicos de España en 2025 en porcentaje del PIB.
- La posición de inversión internacional neta de España en 2025 y la deuda privada consolidada.
- La prima de riesgo española y la rentabilidad del bono a diez años a finales de septiembre de 2026, y el diferencial franco-español.
- Las calificaciones de Francia e Italia a septiembre de 2026 y las revisiones de S&P y Fitch del 11 de septiembre de 2026 sobre España.
- La deuda y el déficit de Estados Unidos, Japón y el Reino Unido en el Fiscal Monitor de abril de 2026.
- El déficit alemán de 2025 y su trayectoria tras la reforma del freno de la deuda.

Fuentes de actualidad consultadas
- [Eurostat, déficit de la zona del euro del 2,9% y de la Unión del 3,1% en 2025 (22 de abril de 2026)](https://ec.europa.eu/eurostat/web/products-euro-indicators/w/2-22042026-ap)
- [Eurostat, deuda del 87,8% del PIB en la zona del euro (22 de abril de 2026)](https://ec.europa.eu/eurostat/web/products-euro-indicators/w/2-22042026-bp)
- [Eurostat, Government finance statistics (Statistics Explained)](https://ec.europa.eu/eurostat/statistics-explained/index.php?title=Government_finance_statistics)
- [Comisión, Debt Sustainability Monitor 2025](https://economy-finance.ec.europa.eu/publications/debt-sustainability-monitor-2025_en)
- [Vozpópuli, Bruselas ve alto riesgo en la deuda pública de España a medio plazo](https://www.vozpopuli.com/economia/macroeconomia/bruselas-alto-riesgo-deuda-publica-espana.html)
- [Forbes España, el rating de España tras las subidas de 2025](https://forbes.es/economia/850407/el-rating-de-espana-buscara-el-sobresaliente-en-2026-tras-las-subidas-cosechadas-en-2025/)
- [La Crónica del Henares, Scope sube a España a A+ (septiembre de 2026)](https://www.cronicadelhenares.com/2026/09/scope-rating-comunidad-madrid-a-plus-techo-espana.html)
- [El Economista, el diferencial de la deuda francesa con la española en máximos históricos (septiembre de 2026)](https://www.eleconomista.es/mercados-cotizaciones/noticias/14011436/09/26/el-diferencial-de-la-deuda-francesa-con-la-espanola-se-situa-en-maximos-historicos.html)
- [El Economista, Francia paga la mayor prima de su historia frente a España (julio de 2026)](https://www.eleconomista.es/mercados-cotizaciones/noticias/13999958/07/26/francia-paga-la-mayor-prima-de-riesgo-de-su-historia-frente-a-espana-por-su-debilidad-fiscal-con-el-bono-a-un-paso-del-4.html)
- [El Boletín, España recorta su prima por debajo de 60 puntos básicos (septiembre de 2026)](https://www.elboletin.com/italia-lidera-la-caida-de-las-primas-de-riesgo-en-la-eurozona-en-septiembre-y-espana-recorta-por-debajo-de-60-pb/)
- [BdE, La evolución de la deuda pública en España en 2024](https://www.bde.es/f/webbe/SES/Secciones/Publicaciones/PublicacionesSeriadas/DocumentosOcasionales/25/Fich/be2503-art04.pdf)
Referencias nuevas
- Angeletos, G.-M. (2002), Fiscal Policy with Noncontingent Debt and the Optimal Maturity Structure, Quarterly Journal of Economics.
- Arellano, C. y Ramanarayanan, A. (2012), Default and the Maturity Structure in Sovereign Bonds, Journal of Political Economy.
- Barro, R. (1979), On the Determination of the Public Debt, Journal of Political Economy.
- Barro, R. (2003), Optimal Management of Indexed and Nominal Debt, Annals of Economics and Finance.
- Blommestein, H. y Turner, P. (2012), Interactions between Sovereign Debt Management and Monetary Policy under Fiscal Dominance and Financial Instability, OECD Working Papers on Sovereign Borrowing and Public Debt Management, 3.
- Buera, F. y Nicolini, J. P. (2004), Optimal Maturity of Government Debt without State Contingent Bonds, Journal of Monetary Economics.
- Calvo, G. y Guidotti, P. (1990), Indexation and Maturity of Government Bonds: An Exploratory Model, en Dornbusch y Draghi (eds.), Public Debt Management: Theory and History, Cambridge University Press.
- Cole, H. y Kehoe, T. (2000), Self-Fulfilling Debt Crises, Review of Economic Studies.
- Currie, E., Dethier, J.-J. y Togo, E. (2003), Institutional Arrangements for Public Debt Management, World Bank Policy Research Working Paper, 3021.
- FMI y Banco Mundial (2014), Revised Guidelines for Public Debt Management.
- FMI y Banco Mundial (2009), Developing a Medium-Term Debt Management Strategy (MTDS): Guidance Note for Country Authorities.
- Greenwood, R., Hanson, S., Rudolph, J. y Summers, L. (2014), Government Debt Management at the Zero Lower Bound, Hutchins Center Working Paper, 5, Brookings.
- Lucas, R. y Stokey, N. (1983), Optimal Fiscal and Monetary Policy in an Economy without Capital, Journal of Monetary Economics.
- Missale, A. y Blanchard, O. (1994), The Debt Burden and Debt Maturity, American Economic Review.
- Orden ECM/2/2026, de 9 de enero, por la que se dispone la creación de Deuda del Estado durante el año 2026 y enero de 2027, BOE.
- Real Decreto 410/2024, de 23 de abril, estructura orgánica básica del Ministerio de Economía, Comercio y Empresa, BOE.
- Real Decreto 644/2026, de 28 de julio, BOE.

Pendiente de verificar
- El mecanismo exacto del requisito de 1984 en pagarés del Tesoro (coeficiente de inversión o tramo del coeficiente de caja) y su porcentaje.
- La fecha y la norma de la restricción del crédito del BdE al Tesoro en 1990.
- La formulación literal del objetivo de la política de financiación en la Estrategia del Tesoro para 2026.
- El coste medio de la deuda en circulación y la vida media a finales de 2025, en las estadísticas del Tesoro (no accesibles desde aquí).
- La proporción de deuda en manos del Eurosistema y de no residentes a mediados de 2026.
- La emisión bruta de 2026: 285.677 millones según la nota oficial, frente a los 285.693 que di en 1.1.4 (corregir allí).
- Evaluaciones de la OCDE sobre la gestión de la deuda española.

Fuentes de actualidad consultadas
- [BOE, Orden ECM/2/2026, de creación de Deuda del Estado para 2026](https://www.boe.es/eli/es/o/2026/01/09/ecm2)
- [BOE, Real Decreto 410/2024, estructura del Ministerio de Economía, Comercio y Empresa](https://www.boe.es/buscar/act.php?id=BOE-A-2024-8194)
- [BOE, Real Decreto 644/2026, de 28 de julio](https://www.boe.es/diario_boe/txt.php?id=BOE-A-2026-16441)
- [Wikipedia, Secretaría General del Tesoro](https://es.wikipedia.org/wiki/Secretar%C3%ADa_General_del_Tesoro)
- [Plan de Recuperación, el Tesoro presenta la Estrategia de Financiación para 2026](https://planderecuperacion.gob.es/noticias/tesoro-publico-presenta-estrategia-financiacion-2026-prtr)
- [CaixaBank Research, la Estrategia del Tesoro 2025](https://www.caixabankresearch.com/en/economics-markets/public-sector/2025-treasury-strategy-context-reduction-spains-public-deficit)
- [BdE, nota sobre la remuneración de los depósitos de las administraciones en el BCE (2024)](https://www.bde.es/f/webbe/GAP/Secciones/SalaPrensa/ComunicadosBCE/NotasInformativasBCE/24/presbce2024-42.pdf)
- [BdE, La evolución de la deuda pública en España en 2024](https://www.bde.es/f/webbe/SES/Secciones/Publicaciones/PublicacionesSeriadas/DocumentosOcasionales/25/Fich/be2503-art04.pdf)

Referencias nuevas
- CaixaBank Research (2025), La Estrategia del Tesoro 2025 en un contexto de reducción del déficit público.
- FMI (2022), Staff Guidance Note on the Sovereign Risk and Debt Sustainability Framework for Market Access Countries.
- La Moncloa (2023), El Tesoro prevé una emisión neta de 70.000 millones este año y cierra 2022 con un sólido acceso al mercado, nota de prensa, 12 de enero.
- La Moncloa (2026), El Tesoro emite 15.000 millones de un nuevo bono a 10 años, con la mayor demanda de su historia, nota de prensa, 20 de enero.
- Ministerio de Economía, Comercio y Empresa (2025), Las mejoras de rating y la bajada de la prima de riesgo a mínimos de 2009 refuerzan el buen acceso a financiación del Tesoro (Estrategia de Financiación 2026), nota de prensa, diciembre.

Pendiente de verificar
- La serie anual de emisión neta y bruta del Tesoro entre 2008 y 2021, en las estadísticas del Tesoro (no accesibles desde aquí).
- La emisión neta y bruta realizada en 2023, 2024 y 2025 (he usado cifras previstas o de prensa).
- La descomposición de la emisión neta de 2026 por destinos: déficit del Estado, préstamos a la Seguridad Social, FFCCAA y operaciones financieras.
- Los vencimientos anuales de la cartera del Eurosistema en deuda española y la oferta neta al sector privado en 2025-2026.
- El umbral exacto de necesidades brutas del marco de sostenibilidad del FMI para las economías avanzadas.
- El calendario de amortización de los préstamos del SURE y del Mecanismo de Recuperación.
- El importe de la emisión bruta de 2026: 285.677 millones (oficial) frente a 285.693 (algunas fuentes de prensa).
- La evolución del colchón de tesorería del Estado en 2020-2025.

Fuentes de actualidad consultadas
- [elDiario.es, el Tesoro mantendrá en 2026 una estrategia similar, con emisión neta de 55.000 millones](https://www.eldiario.es/economia/tesoro-mantendra-2026-estrategia-similar-emision-neta-55-000-millones_1_12882825.html)
- [Público, el Tesoro prevé emitir 55.000 millones de deuda neta en 2026](https://www.publico.es/economia/tesoro-preve-emitir-55-000-millones-2026-deuda-neta-2026-ano.html)
- [Plan de Recuperación, el Tesoro presenta la Estrategia de Financiación 2026](https://planderecuperacion.gob.es/noticias/tesoro-publico-presenta-estrategia-financiacion-2026-prtr)
- [La Moncloa, cierre del programa de 2022 y estrategia de 2023](https://www.lamoncloa.gob.es/serviciosdeprensa/notasprensa/asuntos-economicos/Paginas/2023/120123-tesoro.aspx)
- [El Economista, emisión neta de 55.000 millones en 2024](https://www.eleconomista.es/economia/noticias/12615361/01/24/cuerpo-anuncia-una-emision-neta-del-tesoro-55000-millones-en-2024-10000-millones-menos.html)
- [La Moncloa, el Tesoro recorta la emisión neta de 2025 en 5.000 millones (septiembre de 2025)](https://www.lamoncloa.gob.es/serviciosdeprensa/notasprensa/economia-comercio-empresa/Paginas/2025/300925-emision-deuda.aspx)
- [La Moncloa, emisión sindicada a 10 años de enero de 2026](https://www.lamoncloa.gob.es/serviciosdeprensa/notasprensa/economia-comercio-empresa/Paginas/2026/200126-bono-tesoro.aspx)
- [CaixaBank Research, la Estrategia del Tesoro 2025](https://www.caixabankresearch.com/en/economics-markets/public-sector/2025-treasury-strategy-context-reduction-spains-public-deficit)
- [BdE, nota del BCE sobre la remuneración de los depósitos de las administraciones](https://www.bde.es/f/webbe/GAP/Secciones/SalaPrensa/ComunicadosBCE/NotasInformativasBCE/24/presbce2024-42.pdf)

Referencias nuevas
- Borensztein, E. y Mauro, P. (2004), The Case for GDP-Indexed Bonds, Economic Policy.
- Campbell, J. y Shiller, R. (1996), A Scorecard for Indexed Government Debt, NBER Macroeconomics Annual.
- Doronzo, R., Siracusa, V. y Antonelli, S. (2021), Green Bonds: The Sovereign Issuers' Perspective, Banca d'Italia, Markets, Infrastructures, Payment Systems, 3.
- EFPA España (2024), 2023, el año en que los inversores particulares se fijaron en las Letras del Tesoro.
- Greenwood, R., Hanson, S. y Stein, J. (2015), A Comparative-Advantage Approach to Government Debt Maturity, Journal of Finance.
- Tratado constitutivo del Mecanismo Europeo de Estabilidad (2012), artículo 12.3.

Pendiente de verificar
- El peso actual de las letras en la deuda del Estado (mi estimación: 7-8%) y el saldo vivo total de la deuda del Estado a mediados de 2026.
- El saldo y el peso actuales de los bonos indexados y la fecha y el importe de la nueva referencia indexada a unos doce años.
- El saldo de bonos verdes tras la emisión de septiembre de 2026 y el último informe de asignación e impacto.
- Las fechas e importes de las emisiones a 50 años (2016 y posteriores) y sus rentabilidades.
- El saldo de inversión minorista en letras en 2024-2026.
- El estado de ratificación de la reforma del Tratado del MEDE a septiembre de 2026.
- La proporción de deuda indexada del Reino Unido y su coste en 2022-2023.

Fuentes de actualidad consultadas
- [Merca2, el Tesoro coloca 4.000 millones en bonos verdes ante una demanda récord (9 de septiembre de 2026)](https://www.merca2.es/2026/09/09/bonos-verdes-espana-tesoro-emision-2450986/)
- [Infobae, reapertura del bono verde (mayo de 2025)](https://www.infobae.com/espana/agencias/2025/05/22/el-tesoro-vende-1646-millones-en-una-nueva-reapertura-del-bono-verde-y-reduce-su-interes/)
- [La Moncloa, primera reapertura del bono verde en 2023](https://www.lamoncloa.gob.es/serviciosdeprensa/notasprensa/asuntos-economicos/Paginas/2023/130423-bono-verde-soberano-espanol.aspx)
- [EFPA, 2023, el año en que los inversores particulares se fijaron en las Letras del Tesoro](https://www.asesoresfinancierosefpa.es/actualidad/2023-el-ano-en-que-los-inversores-particulares-se-fijaron-en-las-letras-del-tesoro/)
- [Bankinter, primera emisión de un bono indexado a la inflación (2014)](https://www.bankinter.com/blog/lo-ultimo/primera-emision-de-un-bono-indexado-a-la-inflacion-deuda-bajo-presion)
- [BOE, Orden ECM/2/2026](https://www.boe.es/eli/es/o/2026/01/09/ecm2)
- [Público, Estrategia del Tesoro 2026 (letras y bonos verdes)](https://www.publico.es/economia/tesoro-preve-emitir-55-000-millones-2026-deuda-neta-2026-ano.html)
- [Tratado del MEDE (texto en español)](https://www.esm.europa.eu/sites/default/files/migration_files/20150203_-_esm_treaty_-_es.pdf)

Referencias nuevas
- Abbink, K., Brandts, J. y Pezanis-Christou, P. (2006), Auctions for Government Securities: A Laboratory Comparison of Uniform, Discriminatory and Spanish Designs, Journal of Economic Behavior & Organization.
- Álvarez, F. y Mazón, C. (2007), Comparing the Spanish and the Discriminatory Auction Formats: A Discrete Model with Private Information, European Journal of Operational Research.
- Arnone, M. e Iden, G. (2003), Primary Dealers in Government Securities: Policy Issues and Selected Countries' Experience, IMF Working Paper, 03/45.
- Ausubel, L. y Cramton, P. (2002), Demand Reduction and Inefficiency in Multi-Unit Auctions, University of Maryland Working Paper.
- Duffie, D. (2020), Still the World's Safe Haven? Redesigning the U.S. Treasury Market after the COVID-19 Crisis, Hutchins Center Working Paper, 62, Brookings.
- Friedman, M. (1960), A Program for Monetary Stability, Fordham University Press.
- Hortaçsu, A. y McAdams, D. (2010), Mechanism Choice and Strategic Bidding in Divisible Good Auctions: An Empirical Analysis of the Turkish Treasury Auction Market, Journal of Political Economy.
- Hortaçsu, A., Kastl, J. y Zhang, A. (2018), Bid Shading and Bidder Surplus in the US Treasury Auction System, American Economic Review.
- Kydland, F. y Prescott, E. (1977), Rules Rather than Discretion: The Inconsistency of Optimal Plans, Journal of Political Economy.
- Malvey, P. y Archibald, C. (1998), Uniform-Price Auctions: Update of the Treasury Experience, U.S. Department of the Treasury.
- Umlauf, S. (1993), An Empirical Study of the Mexican Treasury Bill Auction, Journal of Financial Economics.
- Vickrey, W. (1961), Counterspeculation, Auctions, and Competitive Sealed Tenders, Journal of Finance.

Pendiente de verificar
- La proporción de emisión bruta obtenida por subasta y por sindicación en 2024-2025.
- El número actual de creadores de mercado de letras y de bonos y obligaciones, y la resolución vigente que regula su actuación.
- Los resultados de Álvarez y Mazón (2007) y de Abbink et al. (2006) sobre el formato español, para citarlos con precisión.
- Las operaciones de canje o recompra del Tesoro en 2023-2026.
- Las operaciones de colocación de tesorería (subastas de depósitos y adquisiciones temporales) y su volumen.
- La nueva plataforma del Tesoro para inversores particulares (fecha y funcionamiento).
- El número de subastas en 2026 (48 ordinarias según la Estrategia) y su reparto por tipo.

Fuentes de actualidad consultadas
- [La Moncloa, emisión sindicada a 10 años de enero de 2026 (demanda de 145.000 millones)](https://www.lamoncloa.gob.es/serviciosdeprensa/notasprensa/economia-comercio-empresa/Paginas/2026/200126-bono-tesoro.aspx)
- [La Moncloa, emisión sindicada a 10 años de mayo de 2026](https://www.lamoncloa.gob.es/serviciosdeprensa/notasprensa/economia-comercio-empresa/paginas/2026/270526-emision-bonos-tesoro-publico.aspx)
- [Merca2, sindicación verde de septiembre de 2026](https://www.merca2.es/2026/09/09/bonos-verdes-espana-tesoro-emision-2450986/)
- [Forbes España, emisión sindicada a 20 años](https://forbes.es/ultima-hora/1013933/el-tesoro-emite-4-000-millones-en-una-emision-sindicada-a-20-anos-con-demanda-inicial-de-55-000-millones/)
- [elDiario.es, Estrategia del Tesoro 2026 (48 subastas ordinarias)](https://www.eldiario.es/economia/tesoro-mantendra-2026-estrategia-similar-emision-neta-55-000-millones_1_12882825.html)
- [Tesoro Público, Guía práctica de compra de deuda pública por inversores particulares 2025](https://www.tesoropublico.gob.es/sites/default/files/Gu%C3%ADa_Pr%C3%A1ctica_TP_2025_Interactivo.pdf)
- [BOE, Orden ECM/2/2026](https://www.boe.es/eli/es/o/2026/01/09/ecm2)

Referencias nuevas
- BdE (2026), Cuentas anuales del BdE 2025.
- Banco Central Europeo (2023), decisión de 27 de julio de 2023 sobre la remuneración de las reservas mínimas.
- Banco Central Europeo (2024), Changes to the operational framework for implementing monetary policy, 13 de marzo.
- Banco Central Europeo (2026), Decisiones de política monetaria, 10 de septiembre.
- De Grauwe, P. y Ji, Y. (2023), Monetary Policies with Fewer Subsidies for Banks: A Two-Tier System of Minimum Reserve Requirements, VoxEU/CEPR, marzo.
- Herndon, T., Ash, M. y Pollin, R. (2014), Does High Public Debt Consistently Stifle Economic Growth? A Critique of Reinhart and Rogoff, Cambridge Journal of Economics.
- Reinhart, C. y Rogoff, K. (2010), Growth in a Time of Debt, American Economic Review: Papers & Proceedings.
- Sentencia del TJUE de 16 de junio de 2015, Gauweiler, C-62/14.
- Sentencia del TJUE de 11 de diciembre de 2018, Weiss, C-493/17.

Pendiente de verificar
- Las pérdidas del BdE por operaciones de política monetaria en 2023 y las transferencias al Tesoro en 2019-2022.
- El volumen de deuda española en el balance del Eurosistema a mediados de 2026 y el ritmo anual de vencimientos no reinvertidos.
- El calendario exacto de reducción de las reinversiones del PEPP en el segundo semestre de 2024.
- El coste medio de emisión en lo que va de 2026 y la previsión del Tesoro para el año.
- Las emisiones realizadas por el Tesoro en 2026 hasta septiembre, frente al programa anunciado.
- La rentabilidad del bono estadounidense a diez años en septiembre de 2026.
- Si el Tesoro ha revisado la emisión neta de 2026 tras las medidas de emergencia (en 2025 la revisó en septiembre).

Fuentes de actualidad consultadas
- [BCE, decisiones de política monetaria del 10 de septiembre de 2026](https://www.ecb.europa.eu/press/pr/date/2026/html/ecb.mp260910~314e508016.es.html)
- [BdE, el BCE sube los tipos 25 puntos básicos en junio de 2026](https://www.bde.es/wbe/en/noticias-eventos/actualidad-bce/decisiones-politica-monetaria/bce-tipos-junio26.html)
- [Bolsamanía, el BdE gana 234 millones en 2025 tras las pérdidas de 2023 y 2024](https://www.bolsamania.com/noticias/mercados/economiafinanzas--banco-de-espana-gana-234-millones-en-2025-tras-liberar-541-millones-en-provisiones--22209834.html)
- [BdE, Cuentas Anuales 2025](https://www.bde.es/f/webbe/SES/Secciones/Publicaciones/PublicacionesAnuales/Cuentas_anuales_del_BdE/2025/fich/BE_CuentasAnuales_2025.pdf)
- [De Grauwe y Ji (2023), VoxEU/CEPR](https://cepr.org/voxeu/columns/monetary-policies-fewer-subsidies-banks-two-tier-system-minimum-reserve-requirements)
- [Plan de Recuperación, el Tesoro presenta la Estrategia de Financiación 2026](https://planderecuperacion.gob.es/noticias/tesoro-publico-presenta-estrategia-financiacion-2026-prtr)
- [elDiario.es, detalles de la Estrategia de 2026](https://www.eldiario.es/economia/tesoro-mantendra-2026-estrategia-similar-emision-neta-55-000-millones_1_12882825.html)
- [La Moncloa, emisión sindicada a 10 años de enero de 2026](https://www.lamoncloa.gob.es/serviciosdeprensa/notasprensa/economia-comercio-empresa/Paginas/2026/200126-bono-tesoro.aspx)
- [Merca2, emisión verde a 20 años de septiembre de 2026](https://www.merca2.es/2026/09/09/bonos-verdes-espana-tesoro-emision-2450986/)
- [La Crónica del Henares, Scope sube a España a A+ (septiembre de 2026)](https://www.cronicadelhenares.com/2026/09/scope-rating-comunidad-madrid-a-plus-techo-espana.html)




% ANEXOS
\clearpage






\end{document}